% !TEX program = lualatex
% ---------------------------------------------------------------------------------------------
%  Readback note 41: The filling data: Bessel subgroups, links of frames, Siegel's lemma and strong approximation
%
%  Build:  python3 tools/rb.py note <slug>
%
%  One unit of a Lean artifact, rewritten as ordinary mathematics. The prose is hand-written. Lean names,
%  line ranges and counts come from data.tex; quoted Lean statements are copied from the artifact by the build.
%
%  Rules of the text: standard terminology; every symbol is explained where it first appears and listed again
%  in the notation at the end; no letter has two meanings; a proof is complete or says what it takes for granted;
%  Lean names appear only in the LEAN lines and in the concordance.
%
%  The signature rule (the build enforces it): a function, relation or predicate that the artifact defines is
%  written only through its notation macro, which the build makes from the spec (`notation`) with one slot for
%  each explicit argument of the Lean definition. No argument is fixed and dropped, also not "for brevity" and
%  not for one section. Every argument and every hypothesis of a stated declaration stands inside its
%  statement, or in a Setting that the statement cites by number; only implicit and instance arguments (the
%  ambient objects) may stand in the conventions.
% ---------------------------------------------------------------------------------------------
\documentclass[11pt,a4paper]{article}
\ifdefined\pdfvariable\pdfvariable suppressoptionalinfo \numexpr 1+2+4+8\relax\fi
\usepackage[a4paper,hmargin=35mm,top=22mm,bottom=27mm,footskip=12mm,nohead]{geometry}

\usepackage{fontspec}
\directlua{
  luaotfload.add_fallback("gfallback", {"STIXTwoMath-Regular.otf:mode=node;", "DejaVuSans.ttf:mode=node;"})
  luaotfload.add_fallback("gmonofallback", {"DejaVuSansMono.ttf:mode=node;", "STIXTwoMath-Regular.otf:mode=node;"})
}
\setmainfont{NewCM10}[Extension=.otf,UprightFont=*-Book,ItalicFont=*-BookItalic,
  BoldFont=*-Bold,BoldItalicFont=*-BoldItalic,SmallCapsFont=*-Book,SmallCapsFeatures={Letters=SmallCaps},
  RawFeature={fallback=gfallback}]
\setsansfont{NewCMSans10}[Extension=.otf,UprightFont=*-Book,ItalicFont=*-BookOblique,
  BoldFont=*-Bold,BoldItalicFont=*-BoldOblique,RawFeature={fallback=gfallback}]
\setmonofont{JetBrainsMono}[Extension=.otf,UprightFont=*-Regular,ItalicFont=*-Italic,
  BoldFont=*-SemiBold,BoldItalicFont=*-SemiBoldItalic,Scale=MatchLowercase,
  RawFeature={-calt;-liga;fallback=gmonofallback}]
\usepackage{amsmath}
\usepackage{amsthm}
\usepackage[math-style=TeX]{unicode-math}
\setmathfont{NewCMMath-Book.otf}
\emergencystretch=2.5em\relax

\usepackage{xcolor}
\usepackage{tikz}
\usetikzlibrary{calc,arrows.meta,positioning,backgrounds}
\usepackage[protrusion=true,expansion=true]{microtype}
\usepackage{enumitem}
\usepackage{needspace}
\usepackage{booktabs}
\usepackage{array}
\usepackage{longtable}
\usepackage{fancyhdr}
\usepackage{lastpage}
\usepackage{multicol}
\usepackage[hidelinks,pdfusetitle]{hyperref}

\definecolor{ink}{HTML}{0F172A}
\definecolor{inksoft}{HTML}{475569}
\definecolor{inkfaint}{HTML}{94A3B8}
\definecolor{rulegrey}{HTML}{CBD5E1}
\definecolor{arith}{HTML}{C2410C}
\definecolor{est}{HTML}{7C3AED}
\definecolor{ana}{HTML}{0891B2}
\definecolor{ok}{HTML}{047857}
\definecolor{leancol}{HTML}{3730A3}

\title{The filling data: Bessel subgroups, links of frames, Siegel's lemma and strong approximation}
\author{grh-compact}
\date{}

% ------------------------------------------------------------------------------------- macros
\newcommand{\leanfont}{\ttfamily\fontsize{8.6}{10.5}\selectfont\color{leancol}}
\newcommand{\lean}[1]{\csname LN@#1\endcsname}
\newcommand{\fillin}[1]{{\color{arith}\itshape[#1]}}  % a place of the template that is not written yet; the build reports every one that is left
\newcommand{\leanquote}[1]{\csname LQ@#1\endcsname}
% generated by rb.py note; do not edit
\newcommand{\ConcDecls}{232}
\newcommand{\Pin}{72fead10}
\newcommand{\PinText}{that of commit {\leanfont 72fead10}}
\newcommand{\ArtifactName}{grh-compact}
\expandafter\newcommand\csname LN@openCellConst\endcsname{\href{https://github.com/lsernau/grh-lean/blob/72fead102f663d5c7a407964a13dc0cf9d977207/grh-compact/src/New/Filling/BesselSystem.lean\#L18-L18}{{\leanfont Shape.\allowbreak{}openCellConst}}}
\expandafter\newcommand\csname LN@openCellConstpos\endcsname{\href{https://github.com/lsernau/grh-lean/blob/72fead102f663d5c7a407964a13dc0cf9d977207/grh-compact/src/New/Filling/BesselSystem.lean\#L20-L20}{{\leanfont Shape.\allowbreak{}openCellConst\_pos}}}
\expandafter\newcommand\csname LN@bandConst\endcsname{\href{https://github.com/lsernau/grh-lean/blob/72fead102f663d5c7a407964a13dc0cf9d977207/grh-compact/src/New/Filling/BesselSystem.lean\#L24-L24}{{\leanfont Shape.\allowbreak{}bandConst}}}
\expandafter\newcommand\csname LN@besselSize\endcsname{\href{https://github.com/lsernau/grh-lean/blob/72fead102f663d5c7a407964a13dc0cf9d977207/grh-compact/src/New/Filling/BesselSystem.lean\#L29-L29}{{\leanfont besselSize}}}
\expandafter\newcommand\csname LN@InCut\endcsname{\href{https://github.com/lsernau/grh-lean/blob/72fead102f663d5c7a407964a13dc0cf9d977207/grh-compact/src/New/Filling/BesselSystem.lean\#L33-L33}{{\leanfont InCut}}}
\expandafter\newcommand\csname LN@frameBox\endcsname{\href{https://github.com/lsernau/grh-lean/blob/72fead102f663d5c7a407964a13dc0cf9d977207/grh-compact/src/New/Filling/BesselSystem.lean\#L36-L36}{{\leanfont frameBox}}}
\expandafter\newcommand\csname LN@gramDiscReal\endcsname{\href{https://github.com/lsernau/grh-lean/blob/72fead102f663d5c7a407964a13dc0cf9d977207/grh-compact/src/New/Filling/BesselSystem.lean\#L40-L40}{{\leanfont gramDiscReal}}}
\expandafter\newcommand\csname LN@cellMass\endcsname{\href{https://github.com/lsernau/grh-lean/blob/72fead102f663d5c7a407964a13dc0cf9d977207/grh-compact/src/New/Filling/BesselSystem.lean\#L44-L45}{{\leanfont cellMass}}}
\expandafter\newcommand\csname LN@besselCorrelation\endcsname{\href{https://github.com/lsernau/grh-lean/blob/72fead102f663d5c7a407964a13dc0cf9d977207/grh-compact/src/New/Filling/BesselSystem.lean\#L49-L50}{{\leanfont besselCorrelation}}}
\expandafter\newcommand\csname LN@SchurLe\endcsname{\href{https://github.com/lsernau/grh-lean/blob/72fead102f663d5c7a407964a13dc0cf9d977207/grh-compact/src/New/Filling/BesselSystem.lean\#L54-L55}{{\leanfont SchurLe}}}
\expandafter\newcommand\csname LN@HasSchurBound\endcsname{\href{https://github.com/lsernau/grh-lean/blob/72fead102f663d5c7a407964a13dc0cf9d977207/grh-compact/src/New/Filling/BesselSystem.lean\#L59-L62}{{\leanfont HasSchurBound}}}
\expandafter\newcommand\csname LN@normCorrelation\endcsname{\href{https://github.com/lsernau/grh-lean/blob/72fead102f663d5c7a407964a13dc0cf9d977207/grh-compact/src/New/Filling/BesselSystem.lean\#L67-L69}{{\leanfont normCorrelation}}}
\expandafter\newcommand\csname LN@HasBandLimits\endcsname{\href{https://github.com/lsernau/grh-lean/blob/72fead102f663d5c7a407964a13dc0cf9d977207/grh-compact/src/New/Filling/BesselSystem.lean\#L74-L77}{{\leanfont HasBandLimits}}}
\expandafter\newcommand\csname LN@Side\endcsname{\href{https://github.com/lsernau/grh-lean/blob/72fead102f663d5c7a407964a13dc0cf9d977207/grh-compact/src/New/Filling/BesselSystem.lean\#L80-L80}{{\leanfont Side}}}
\expandafter\newcommand\csname LN@IsCutPair\endcsname{\href{https://github.com/lsernau/grh-lean/blob/72fead102f663d5c7a407964a13dc0cf9d977207/grh-compact/src/New/Filling/BesselSystem.lean\#L85-L96}{{\leanfont IsCutPair}}}
\expandafter\newcommand\csname LN@BesselSystem\endcsname{\href{https://github.com/lsernau/grh-lean/blob/72fead102f663d5c7a407964a13dc0cf9d977207/grh-compact/src/New/Filling/BesselSystem.lean\#L102-L109}{{\leanfont BesselSystem}}}
\expandafter\newcommand\csname LN@reductionapplycongr\endcsname{\href{https://github.com/lsernau/grh-lean/blob/72fead102f663d5c7a407964a13dc0cf9d977207/grh-compact/src/New/Filling/CrudeBoundDefs.lean\#L16-L17}{{\leanfont Filling.\allowbreak{}reduction\_apply\_congr}}}
\expandafter\newcommand\csname LN@existsreduceAboveGeq\endcsname{\href{https://github.com/lsernau/grh-lean/blob/72fead102f663d5c7a407964a13dc0cf9d977207/grh-compact/src/New/Filling/CrudeBoundDefs.lean\#L19-L23}{{\leanfont Filling.\allowbreak{}exists\_reduceAboveG\_eq}}}
\expandafter\newcommand\csname LN@vectorFixer\endcsname{\href{https://github.com/lsernau/grh-lean/blob/72fead102f663d5c7a407964a13dc0cf9d977207/grh-compact/src/New/Filling/CrudeBoundDefs.lean\#L27-L29}{{\leanfont Filling.\allowbreak{}vectorFixer}}}
\expandafter\newcommand\csname LN@FixesNormAtLeast\endcsname{\href{https://github.com/lsernau/grh-lean/blob/72fead102f663d5c7a407964a13dc0cf9d977207/grh-compact/src/New/Filling/CrudeBoundDefs.lean\#L33-L34}{{\leanfont Filling.\allowbreak{}FixesNormAtLeast}}}
\expandafter\newcommand\csname LN@fixingSet\endcsname{\href{https://github.com/lsernau/grh-lean/blob/72fead102f663d5c7a407964a13dc0cf9d977207/grh-compact/src/New/Filling/CrudeBoundDefs.lean\#L38-L39}{{\leanfont Filling.\allowbreak{}fixingSet}}}
\expandafter\newcommand\csname LN@zFormPair\endcsname{\href{https://github.com/lsernau/grh-lean/blob/72fead102f663d5c7a407964a13dc0cf9d977207/grh-compact/src/New/Filling/CrudeBoundDefs.lean\#L45-L45}{{\leanfont Filling.\allowbreak{}zFormPair}}}
\expandafter\newcommand\csname LN@commonLink\endcsname{\href{https://github.com/lsernau/grh-lean/blob/72fead102f663d5c7a407964a13dc0cf9d977207/grh-compact/src/New/Filling/CrudeBoundDefs.lean\#L51-L52}{{\leanfont commonLink}}}
\expandafter\newcommand\csname LN@HasCommonFlagVector\endcsname{\href{https://github.com/lsernau/grh-lean/blob/72fead102f663d5c7a407964a13dc0cf9d977207/grh-compact/src/New/Filling/CrudeBoundDefs.lean\#L56-L57}{{\leanfont HasCommonFlagVector}}}
\expandafter\newcommand\csname LN@IsNormalForm\endcsname{\href{https://github.com/lsernau/grh-lean/blob/72fead102f663d5c7a407964a13dc0cf9d977207/grh-compact/src/New/Filling/CrudeBoundDefs.lean\#L61-L62}{{\leanfont IsNormalForm}}}
\expandafter\newcommand\csname LN@ratGram\endcsname{\href{https://github.com/lsernau/grh-lean/blob/72fead102f663d5c7a407964a13dc0cf9d977207/grh-compact/src/New/Filling/CrudeBoundDefs.lean\#L68-L68}{{\leanfont Filling.\allowbreak{}ratGram}}}
\expandafter\newcommand\csname LN@linkForm\endcsname{\href{https://github.com/lsernau/grh-lean/blob/72fead102f663d5c7a407964a13dc0cf9d977207/grh-compact/src/New/Filling/CrudeBoundDefs.lean\#L72-L72}{{\leanfont Filling.\allowbreak{}linkForm}}}
\expandafter\newcommand\csname LN@classCoord\endcsname{\href{https://github.com/lsernau/grh-lean/blob/72fead102f663d5c7a407964a13dc0cf9d977207/grh-compact/src/New/Filling/CrudeBoundDefs.lean\#L76-L76}{{\leanfont Filling.\allowbreak{}classCoord}}}
\expandafter\newcommand\csname LN@classCoordX\endcsname{\href{https://github.com/lsernau/grh-lean/blob/72fead102f663d5c7a407964a13dc0cf9d977207/grh-compact/src/New/Filling/CrudeBoundDefs.lean\#L80-L80}{{\leanfont Filling.\allowbreak{}classCoordX₀}}}
\expandafter\newcommand\csname LN@countWeight\endcsname{\href{https://github.com/lsernau/grh-lean/blob/72fead102f663d5c7a407964a13dc0cf9d977207/grh-compact/src/New/Filling/CrudeBoundDefs.lean\#L85-L86}{{\leanfont Filling.\allowbreak{}countWeight}}}
\expandafter\newcommand\csname LN@HasLevel\endcsname{\href{https://github.com/lsernau/grh-lean/blob/72fead102f663d5c7a407964a13dc0cf9d977207/grh-compact/src/New/Filling/CrudeBoundDefs.lean\#L92-L93}{{\leanfont Filling.\allowbreak{}HasLevel}}}
\expandafter\newcommand\csname LN@CongruentTerms\endcsname{\href{https://github.com/lsernau/grh-lean/blob/72fead102f663d5c7a407964a13dc0cf9d977207/grh-compact/src/New/Filling/CrudeBoundDefs.lean\#L97-L98}{{\leanfont Filling.\allowbreak{}CongruentTerms}}}
\expandafter\newcommand\csname LN@termDiscNorm\endcsname{\href{https://github.com/lsernau/grh-lean/blob/72fead102f663d5c7a407964a13dc0cf9d977207/grh-compact/src/New/Filling/CrudeBoundDefs.lean\#L102-L102}{{\leanfont Filling.\allowbreak{}termDiscNorm}}}
\expandafter\newcommand\csname LN@boxedPoints\endcsname{\href{https://github.com/lsernau/grh-lean/blob/72fead102f663d5c7a407964a13dc0cf9d977207/grh-compact/src/New/Filling/CrudeBoundDefs.lean\#L106-L107}{{\leanfont Filling.\allowbreak{}boxedPoints}}}
\expandafter\newcommand\csname LN@IsNear\endcsname{\href{https://github.com/lsernau/grh-lean/blob/72fead102f663d5c7a407964a13dc0cf9d977207/grh-compact/src/New/Filling/CrudeBoundDefs.lean\#L111-L111}{{\leanfont Filling.\allowbreak{}IsNear}}}
\expandafter\newcommand\csname LN@CharsAgree\endcsname{\href{https://github.com/lsernau/grh-lean/blob/72fead102f663d5c7a407964a13dc0cf9d977207/grh-compact/src/New/Filling/CrudeBoundDefs.lean\#L115-L116}{{\leanfont Filling.\allowbreak{}CharsAgree}}}
\expandafter\newcommand\csname LN@latMiddleGram\endcsname{\href{https://github.com/lsernau/grh-lean/blob/72fead102f663d5c7a407964a13dc0cf9d977207/grh-compact/src/New/Filling/EichlerFlag.lean\#L17-L19}{{\leanfont latMiddleGram}}}
\expandafter\newcommand\csname LN@middleBaseChange\endcsname{\href{https://github.com/lsernau/grh-lean/blob/72fead102f663d5c7a407964a13dc0cf9d977207/grh-compact/src/New/Filling/EichlerFlag.lean\#L29-L29}{{\leanfont middleBaseChange}}}
\expandafter\newcommand\csname LN@besselLatticeBasis\endcsname{\href{https://github.com/lsernau/grh-lean/blob/72fead102f663d5c7a407964a13dc0cf9d977207/grh-compact/src/New/Filling/EichlerFlag.lean\#L33-L37}{{\leanfont besselLatticeBasis}}}
\expandafter\newcommand\csname LN@frameDisc\endcsname{\href{https://github.com/lsernau/grh-lean/blob/72fead102f663d5c7a407964a13dc0cf9d977207/grh-compact/src/New/Filling/EichlerFlag.lean\#L43-L43}{{\leanfont frameDisc}}}
\expandafter\newcommand\csname LN@grambesselLatticeBasis\endcsname{\href{https://github.com/lsernau/grh-lean/blob/72fead102f663d5c7a407964a13dc0cf9d977207/grh-compact/src/New/Filling/EichlerFlag.lean\#L45-L48}{{\leanfont gram\_besselLatticeBasis}}}
\expandafter\newcommand\csname LN@framePair\endcsname{\href{https://github.com/lsernau/grh-lean/blob/72fead102f663d5c7a407964a13dc0cf9d977207/grh-compact/src/New/Filling/EichlerFlag.lean\#L52-L52}{{\leanfont framePair}}}
\expandafter\newcommand\csname LN@isotropicFramePair\endcsname{\href{https://github.com/lsernau/grh-lean/blob/72fead102f663d5c7a407964a13dc0cf9d977207/grh-compact/src/New/Filling/EichlerFlag.lean\#L55-L55}{{\leanfont isotropicFramePair}}}
\expandafter\newcommand\csname LN@linkMatrix\endcsname{\href{https://github.com/lsernau/grh-lean/blob/72fead102f663d5c7a407964a13dc0cf9d977207/grh-compact/src/New/Filling/EichlerFlag.lean\#L59-L59}{{\leanfont linkMatrix}}}
\expandafter\newcommand\csname LN@link\endcsname{\href{https://github.com/lsernau/grh-lean/blob/72fead102f663d5c7a407964a13dc0cf9d977207/grh-compact/src/New/Filling/EichlerFlag.lean\#L63-L63}{{\leanfont link}}}
\expandafter\newcommand\csname LN@dampedPower\endcsname{\href{https://github.com/lsernau/grh-lean/blob/72fead102f663d5c7a407964a13dc0cf9d977207/grh-compact/src/New/Filling/EichlerFlag.lean\#L68-L68}{{\leanfont dampedPower}}}
\expandafter\newcommand\csname LN@headWeight\endcsname{\href{https://github.com/lsernau/grh-lean/blob/72fead102f663d5c7a407964a13dc0cf9d977207/grh-compact/src/New/Filling/EichlerFlag.lean\#L72-L72}{{\leanfont headWeight}}}
\expandafter\newcommand\csname LN@SieveSystem\endcsname{\href{https://github.com/lsernau/grh-lean/blob/72fead102f663d5c7a407964a13dc0cf9d977207/grh-compact/src/New/Filling/EichlerFlag.lean\#L81-L91}{{\leanfont SieveSystem}}}
\expandafter\newcommand\csname LN@pi\endcsname{\href{https://github.com/lsernau/grh-lean/blob/72fead102f663d5c7a407964a13dc0cf9d977207/grh-compact/src/New/Filling/EichlerFlag.lean\#L96-L112}{{\leanfont SieveSystem.\allowbreak{}pi}}}
\expandafter\newcommand\csname LN@stepsum\endcsname{\href{https://github.com/lsernau/grh-lean/blob/72fead102f663d5c7a407964a13dc0cf9d977207/grh-compact/src/New/Filling/EichlerFlag.lean\#L116-L120}{{\leanfont SieveSystem.\allowbreak{}step\_sum}}}
\expandafter\newcommand\csname LN@nonnegandlowerSemicontinuous\endcsname{\href{https://github.com/lsernau/grh-lean/blob/72fead102f663d5c7a407964a13dc0cf9d977207/grh-compact/src/New/Filling/EichlerFlag.lean\#L125-L139}{{\leanfont SieveSystem.\allowbreak{}nonneg\_and\_lowerSemicontinuous}}}
\expandafter\newcommand\csname LN@finHeadWeight\endcsname{\href{https://github.com/lsernau/grh-lean/blob/72fead102f663d5c7a407964a13dc0cf9d977207/grh-compact/src/New/Filling/EichlerFlag.lean\#L145-L146}{{\leanfont finHeadWeight}}}
\expandafter\newcommand\csname LN@isUnimodularlatMiddleGram\endcsname{\href{https://github.com/lsernau/grh-lean/blob/72fead102f663d5c7a407964a13dc0cf9d977207/grh-compact/src/New/Filling/EichlerFlag.lean\#L148-L150}{{\leanfont isUnimodular\_latMiddleGram}}}
\expandafter\newcommand\csname LN@frameImageonemulVec\endcsname{\href{https://github.com/lsernau/grh-lean/blob/72fead102f663d5c7a407964a13dc0cf9d977207/grh-compact/src/New/Filling/EichlerFlag.lean\#L152-L153}{{\leanfont frameImage\_one\_mulVec}}}
\expandafter\newcommand\csname LN@frameImagemulVec\endcsname{\href{https://github.com/lsernau/grh-lean/blob/72fead102f663d5c7a407964a13dc0cf9d977207/grh-compact/src/New/Filling/EichlerFlag.lean\#L155-L156}{{\leanfont frameImage\_mulVec}}}
\expandafter\newcommand\csname LN@existsstarvecMuleq\endcsname{\href{https://github.com/lsernau/grh-lean/blob/72fead102f663d5c7a407964a13dc0cf9d977207/grh-compact/src/New/Filling/EichlerFlag.lean\#L158-L159}{{\leanfont exists\_star\_vecMul\_eq}}}
\expandafter\newcommand\csname LN@existsmembesselUeqeichlerMap\endcsname{\href{https://github.com/lsernau/grh-lean/blob/72fead102f663d5c7a407964a13dc0cf9d977207/grh-compact/src/New/Filling/EichlerFlag.lean\#L165-L175}{{\leanfont exists\_mem\_besselU\_eq\_eichlerMap}}}
\expandafter\newcommand\csname LN@finAdeleMat\endcsname{\href{https://github.com/lsernau/grh-lean/blob/72fead102f663d5c7a407964a13dc0cf9d977207/grh-compact/src/New/Filling/GramLiouville.lean\#L18-L18}{{\leanfont Filling.\allowbreak{}finAdeleMat}}}
\expandafter\newcommand\csname LN@finAdeleMatconjTranspose\endcsname{\href{https://github.com/lsernau/grh-lean/blob/72fead102f663d5c7a407964a13dc0cf9d977207/grh-compact/src/New/Filling/GramLiouville.lean\#L20-L20}{{\leanfont Filling.\allowbreak{}finAdeleMat\_conjTranspose}}}
\expandafter\newcommand\csname LN@mapconjTranspose\endcsname{\href{https://github.com/lsernau/grh-lean/blob/72fead102f663d5c7a407964a13dc0cf9d977207/grh-compact/src/New/Filling/GramLiouville.lean\#L22-L22}{{\leanfont Filling.\allowbreak{}map\_conjTranspose}}}
\expandafter\newcommand\csname LN@finAdeleMatsmul\endcsname{\href{https://github.com/lsernau/grh-lean/blob/72fead102f663d5c7a407964a13dc0cf9d977207/grh-compact/src/New/Filling/GramLiouville.lean\#L25-L26}{{\leanfont Filling.\allowbreak{}finAdeleMat\_smul}}}
\expandafter\newcommand\csname LN@muleqzeroofrangeleker\endcsname{\href{https://github.com/lsernau/grh-lean/blob/72fead102f663d5c7a407964a13dc0cf9d977207/grh-compact/src/New/Filling/GramLiouville.lean\#L28-L29}{{\leanfont Filling.\allowbreak{}mul\_eq\_zero\_of\_range\_le\_ker}}}
\expandafter\newcommand\csname LN@existsbasisextendker\endcsname{\href{https://github.com/lsernau/grh-lean/blob/72fead102f663d5c7a407964a13dc0cf9d977207/grh-compact/src/New/Filling/GramLiouville.lean\#L33-L40}{{\leanfont Filling.\allowbreak{}exists\_basis\_extend\_ker}}}
\expandafter\newcommand\csname LN@sublatticeBasis\endcsname{\href{https://github.com/lsernau/grh-lean/blob/72fead102f663d5c7a407964a13dc0cf9d977207/grh-compact/src/New/Filling/GramLiouville.lean\#L49-L54}{{\leanfont Filling.\allowbreak{}sublatticeBasis}}}
\expandafter\newcommand\csname LN@sublatticeDisc\endcsname{\href{https://github.com/lsernau/grh-lean/blob/72fead102f663d5c7a407964a13dc0cf9d977207/grh-compact/src/New/Filling/GramLiouville.lean\#L58-L59}{{\leanfont Filling.\allowbreak{}sublatticeDisc}}}
\expandafter\newcommand\csname LN@sublatticeDiscmulindexsqle\endcsname{\href{https://github.com/lsernau/grh-lean/blob/72fead102f663d5c7a407964a13dc0cf9d977207/grh-compact/src/New/Filling/GramLiouville.lean\#L64-L79}{{\leanfont Filling.\allowbreak{}sublatticeDisc\_mul\_index\_sq\_le}}}
\expandafter\newcommand\csname LN@continuousbesselFn\endcsname{\href{https://github.com/lsernau/grh-lean/blob/72fead102f663d5c7a407964a13dc0cf9d977207/grh-compact/src/New/Filling/LatticeChart.lean\#L14-L14}{{\leanfont Filling.\allowbreak{}continuous\_besselFn}}}
\expandafter\newcommand\csname LN@continuousbesselKernel\endcsname{\href{https://github.com/lsernau/grh-lean/blob/72fead102f663d5c7a407964a13dc0cf9d977207/grh-compact/src/New/Filling/LatticeChart.lean\#L16-L17}{{\leanfont Filling.\allowbreak{}continuous\_besselKernel}}}
\expandafter\newcommand\csname LN@enormbesselKernel\endcsname{\href{https://github.com/lsernau/grh-lean/blob/72fead102f663d5c7a407964a13dc0cf9d977207/grh-compact/src/New/Filling/LatticeChart.lean\#L19-L22}{{\leanfont Filling.\allowbreak{}enorm\_besselKernel}}}
\expandafter\newcommand\csname LN@enormamplitudele\endcsname{\href{https://github.com/lsernau/grh-lean/blob/72fead102f663d5c7a407964a13dc0cf9d977207/grh-compact/src/New/Filling/LatticeChart.lean\#L24-L25}{{\leanfont Filling.\allowbreak{}enorm\_amplitude\_le}}}
\expandafter\newcommand\csname LN@isIntMatrixmiddleBaseChangemul\endcsname{\href{https://github.com/lsernau/grh-lean/blob/72fead102f663d5c7a407964a13dc0cf9d977207/grh-compact/src/New/Filling/LatticeChart.lean\#L34-L38}{{\leanfont Filling.\allowbreak{}isIntMatrix\_middleBaseChange\_mul}}}
\expandafter\newcommand\csname LN@forallinclHFinsuboneiff\endcsname{\href{https://github.com/lsernau/grh-lean/blob/72fead102f663d5c7a407964a13dc0cf9d977207/grh-compact/src/New/Filling/LatticeChart.lean\#L40-L43}{{\leanfont Filling.\allowbreak{}forall\_inclHFin\_sub\_one\_iff}}}
\expandafter\newcommand\csname LN@isIntMatrixscaleAtsubone\endcsname{\href{https://github.com/lsernau/grh-lean/blob/72fead102f663d5c7a407964a13dc0cf9d977207/grh-compact/src/New/Filling/LatticeChart.lean\#L49-L55}{{\leanfont Filling.\allowbreak{}isIntMatrix\_scaleAt\_sub\_one}}}
\expandafter\newcommand\csname LN@inclHFinmemlatStabilizer\endcsname{\href{https://github.com/lsernau/grh-lean/blob/72fead102f663d5c7a407964a13dc0cf9d977207/grh-compact/src/New/Filling/LatticeChart.lean\#L57-L65}{{\leanfont Filling.\allowbreak{}inclHFin\_mem\_latStabilizer}}}
\expandafter\newcommand\csname LN@reductioninclHFineqoneiff\endcsname{\href{https://github.com/lsernau/grh-lean/blob/72fead102f663d5c7a407964a13dc0cf9d977207/grh-compact/src/New/Filling/LatticeChart.lean\#L67-L72}{{\leanfont Filling.\allowbreak{}reduction\_inclHFin\_eq\_one\_iff}}}
\expandafter\newcommand\csname LN@classSumeqsum\endcsname{\href{https://github.com/lsernau/grh-lean/blob/72fead102f663d5c7a407964a13dc0cf9d977207/grh-compact/src/New/Filling/LatticeChart.lean\#L74-L75}{{\leanfont Filling.\allowbreak{}classSum\_eq\_sum}}}
\expandafter\newcommand\csname LN@cardresFieldeqlevelPrime\endcsname{\href{https://github.com/lsernau/grh-lean/blob/72fead102f663d5c7a407964a13dc0cf9d977207/grh-compact/src/New/Filling/LatticeChart.lean\#L77-L78}{{\leanfont Filling.\allowbreak{}card\_resField\_eq\_levelPrime}}}
\expandafter\newcommand\csname LN@mapsplitGramlatMiddle\endcsname{\href{https://github.com/lsernau/grh-lean/blob/72fead102f663d5c7a407964a13dc0cf9d977207/grh-compact/src/New/Filling/LevelBridge.lean\#L29-L30}{{\leanfont Filling.\allowbreak{}map\_splitGram\_latMiddle}}}
\expandafter\newcommand\csname LN@mapmemunitriangularOn\endcsname{\href{https://github.com/lsernau/grh-lean/blob/72fead102f663d5c7a407964a13dc0cf9d977207/grh-compact/src/New/Filling/LevelBridge.lean\#L32-L33}{{\leanfont Membership.\allowbreak{}mem.\allowbreak{}map\_mem\_unitriangularOn}}}
\expandafter\newcommand\csname LN@toLatBasismembesselU\endcsname{\href{https://github.com/lsernau/grh-lean/blob/72fead102f663d5c7a407964a13dc0cf9d977207/grh-compact/src/New/Filling/LevelBridge.lean\#L39-L41}{{\leanfont Filling.\allowbreak{}toLatBasis\_mem\_besselU}}}
\expandafter\newcommand\csname LN@ratToFinmembesselU\endcsname{\href{https://github.com/lsernau/grh-lean/blob/72fead102f663d5c7a407964a13dc0cf9d977207/grh-compact/src/New/Filling/LevelBridge.lean\#L43-L44}{{\leanfont Filling.\allowbreak{}ratToFin\_mem\_besselU}}}
\expandafter\newcommand\csname LN@inclHFinmembesselH\endcsname{\href{https://github.com/lsernau/grh-lean/blob/72fead102f663d5c7a407964a13dc0cf9d977207/grh-compact/src/New/Filling/LevelBridge.lean\#L46-L46}{{\leanfont Filling.\allowbreak{}inclHFin\_mem\_besselH}}}
\expandafter\newcommand\csname LN@rangeinclHFin\endcsname{\href{https://github.com/lsernau/grh-lean/blob/72fead102f663d5c7a407964a13dc0cf9d977207/grh-compact/src/New/Filling/LevelBridge.lean\#L48-L50}{{\leanfont Filling.\allowbreak{}range\_inclHFin}}}
\expandafter\newcommand\csname LN@ratToFintoBesselH\endcsname{\href{https://github.com/lsernau/grh-lean/blob/72fead102f663d5c7a407964a13dc0cf9d977207/grh-compact/src/New/Filling/LevelBridge.lean\#L52-L53}{{\leanfont Filling.\allowbreak{}ratToFin\_toBesselH}}}
\expandafter\newcommand\csname LN@classBasis\endcsname{\href{https://github.com/lsernau/grh-lean/blob/72fead102f663d5c7a407964a13dc0cf9d977207/grh-compact/src/New/Filling/LevelBridge.lean\#L57-L57}{{\leanfont Filling.\allowbreak{}classBasis}}}
\expandafter\newcommand\csname LN@coetermPosition\endcsname{\href{https://github.com/lsernau/grh-lean/blob/72fead102f663d5c7a407964a13dc0cf9d977207/grh-compact/src/New/Filling/LevelBridge.lean\#L59-L60}{{\leanfont Filling.\allowbreak{}coe\_termPosition}}}
\expandafter\newcommand\csname LN@termPositionmembesselU\endcsname{\href{https://github.com/lsernau/grh-lean/blob/72fead102f663d5c7a407964a13dc0cf9d977207/grh-compact/src/New/Filling/LevelBridge.lean\#L62-L64}{{\leanfont Membership.\allowbreak{}mem.\allowbreak{}termPosition\_mem\_besselU}}}
\expandafter\newcommand\csname LN@mul\endcsname{\href{https://github.com/lsernau/grh-lean/blob/72fead102f663d5c7a407964a13dc0cf9d977207/grh-compact/src/New/Filling/LevelBridge.lean\#L66-L67}{{\leanfont Filling.\allowbreak{}termsOver.\allowbreak{}mul}}}
\expandafter\newcommand\csname LN@inva\endcsname{\href{https://github.com/lsernau/grh-lean/blob/72fead102f663d5c7a407964a13dc0cf9d977207/grh-compact/src/New/Filling/LevelBridge.lean\#L69-L69}{{\leanfont Filling.\allowbreak{}termsOver.\allowbreak{}inv}}}
\expandafter\newcommand\csname LN@simpleRootSumadelicToFin\endcsname{\href{https://github.com/lsernau/grh-lean/blob/72fead102f663d5c7a407964a13dc0cf9d977207/grh-compact/src/New/Filling/LevelBridge.lean\#L71-L75}{{\leanfont Filling.\allowbreak{}simpleRootSum\_adelicToFin}}}
\expandafter\newcommand\csname LN@simpleRootSumratToAdelic\endcsname{\href{https://github.com/lsernau/grh-lean/blob/72fead102f663d5c7a407964a13dc0cf9d977207/grh-compact/src/New/Filling/LevelBridge.lean\#L77-L78}{{\leanfont Filling.\allowbreak{}simpleRootSum\_ratToAdelic}}}
\expandafter\newcommand\csname LN@normselfDualChar\endcsname{\href{https://github.com/lsernau/grh-lean/blob/72fead102f663d5c7a407964a13dc0cf9d977207/grh-compact/src/New/Filling/LevelBridge.lean\#L80-L80}{{\leanfont Filling.\allowbreak{}norm\_selfDualChar}}}
\expandafter\newcommand\csname LN@ratToInfmemunitriangularOn\endcsname{\href{https://github.com/lsernau/grh-lean/blob/72fead102f663d5c7a407964a13dc0cf9d977207/grh-compact/src/New/Filling/LevelBridge.lean\#L82-L83}{{\leanfont Membership.\allowbreak{}mem.\allowbreak{}ratToInf\_mem\_unitriangularOn}}}
\expandafter\newcommand\csname LN@besselFnadelicToInfmul\endcsname{\href{https://github.com/lsernau/grh-lean/blob/72fead102f663d5c7a407964a13dc0cf9d977207/grh-compact/src/New/Filling/LevelBridge.lean\#L88-L96}{{\leanfont Filling.\allowbreak{}besselFn\_adelicToInf\_mul}}}
\expandafter\newcommand\csname LN@besselFnratToInfmul\endcsname{\href{https://github.com/lsernau/grh-lean/blob/72fead102f663d5c7a407964a13dc0cf9d977207/grh-compact/src/New/Filling/LevelBridge.lean\#L98-L100}{{\leanfont Filling.\allowbreak{}besselFn\_ratToInf\_mul}}}
\expandafter\newcommand\csname LN@isUnimodularlatMiddle\endcsname{\href{https://github.com/lsernau/grh-lean/blob/72fead102f663d5c7a407964a13dc0cf9d977207/grh-compact/src/New/Filling/LevelBridge.lean\#L108-L109}{{\leanfont Filling.\allowbreak{}isUnimodular\_latMiddle}}}
\expandafter\newcommand\csname LN@transposesplitGram\endcsname{\href{https://github.com/lsernau/grh-lean/blob/72fead102f663d5c7a407964a13dc0cf9d977207/grh-compact/src/New/Filling/LevelBridge.lean\#L114-L115}{{\leanfont Filling.\allowbreak{}transpose\_splitGram}}}
\expandafter\newcommand\csname LN@idOrInvAdjoint\endcsname{\href{https://github.com/lsernau/grh-lean/blob/72fead102f663d5c7a407964a13dc0cf9d977207/grh-compact/src/New/Filling/LevelBridge.lean\#L119-L120}{{\leanfont Filling.\allowbreak{}idOrInvAdjoint}}}
\expandafter\newcommand\csname LN@reduceAboveG\endcsname{\href{https://github.com/lsernau/grh-lean/blob/72fead102f663d5c7a407964a13dc0cf9d977207/grh-compact/src/New/Filling/LevelBridge.lean\#L125-L126}{{\leanfont Filling.\allowbreak{}reduceAboveG}}}
\expandafter\newcommand\csname LN@intGLmapreduceAboveG\endcsname{\href{https://github.com/lsernau/grh-lean/blob/72fead102f663d5c7a407964a13dc0cf9d977207/grh-compact/src/New/Filling/LevelBridge.lean\#L130-L136}{{\leanfont Filling.\allowbreak{}intGL\_map\_reduceAboveG}}}
\expandafter\newcommand\csname LN@gIdeleval\endcsname{\href{https://github.com/lsernau/grh-lean/blob/72fead102f663d5c7a407964a13dc0cf9d977207/grh-compact/src/New/Filling/LevelBridge.lean\#L138-L138}{{\leanfont Filling.\allowbreak{}gIdele\_val}}}
\expandafter\newcommand\csname LN@gIdelemem\endcsname{\href{https://github.com/lsernau/grh-lean/blob/72fead102f663d5c7a407964a13dc0cf9d977207/grh-compact/src/New/Filling/LevelBridge.lean\#L140-L140}{{\leanfont Filling.\allowbreak{}gIdele\_mem}}}
\expandafter\newcommand\csname LN@forallselfDualCharmuleqoneiff\endcsname{\href{https://github.com/lsernau/grh-lean/blob/72fead102f663d5c7a407964a13dc0cf9d977207/grh-compact/src/New/Filling/LevelBridge.lean\#L142-L145}{{\leanfont Filling.\allowbreak{}forall\_selfDualChar\_mul\_eq\_one\_iff}}}
\expandafter\newcommand\csname LN@selfDualChareqoneofmem\endcsname{\href{https://github.com/lsernau/grh-lean/blob/72fead102f663d5c7a407964a13dc0cf9d977207/grh-compact/src/New/Filling/LevelBridge.lean\#L147-L148}{{\leanfont Filling.\allowbreak{}selfDualChar\_eq\_one\_of\_mem}}}
\expandafter\newcommand\csname LN@invgIdelesmulmemiff\endcsname{\href{https://github.com/lsernau/grh-lean/blob/72fead102f663d5c7a407964a13dc0cf9d977207/grh-compact/src/New/Filling/LevelBridge.lean\#L150-L152}{{\leanfont Filling.\allowbreak{}inv\_gIdele\_smul\_mem\_iff}}}
\expandafter\newcommand\csname LN@principalLevelequnitaryLevel\endcsname{\href{https://github.com/lsernau/grh-lean/blob/72fead102f663d5c7a407964a13dc0cf9d977207/grh-compact/src/New/Filling/LevelBridge.lean\#L154-L157}{{\leanfont Filling.\allowbreak{}principalLevel\_eq\_unitaryLevel}}}
\expandafter\newcommand\csname LN@siegellemma\endcsname{\href{https://github.com/lsernau/grh-lean/blob/72fead102f663d5c7a407964a13dc0cf9d977207/grh-compact/src/New/Filling/NearTerms.lean\#L27-L60}{{\leanfont Filling.\allowbreak{}siegel\_lemma}}}
\expandafter\newcommand\csname LQ@siegellemma\endcsname{\leanquoteline{0}{lemma Filling.siegel\_lemma \{E\} [Field E] [CMType E] [StablyUnram E] \{κ\} [Fintype κ] u \{T : Submodule}\leanquoteline{2}{E \_\} \{s\} (hs : finrank E T + s = card κ) : ∃ X : Matrix (Fin s) κ E, Xᵀᴴ.mulVecLin.ker = T ∧}\leanquoteline{2}{(∀ i, X i ∈ ((stdSpace ..).move (1, u)).baseLattice ∧ X i ≠ 0) ∧}\leanquoteline{2}{(∏ i, ‖((stdSpace ..).move (1, u)).archCoords (X i)‖) \textasciicircum{} (2 * numInfPlaces E) ≤}\leanquoteline{2}{(8 * (numInfPlaces E * s) * rootDiscr E) \textasciicircum{} (numInfPlaces E * s) * sublatticeDisc u.1.3 T}}
\expandafter\newcommand\csname LN@tsumsublatticeDiscinvle\endcsname{\href{https://github.com/lsernau/grh-lean/blob/72fead102f663d5c7a407964a13dc0cf9d977207/grh-compact/src/New/Filling/NearTerms.lean\#L65-L83}{{\leanfont Filling.\allowbreak{}tsum\_sublatticeDisc\_inv\_le}}}
\expandafter\newcommand\csname LN@besselUdense\endcsname{\href{https://github.com/lsernau/grh-lean/blob/72fead102f663d5c7a407964a13dc0cf9d977207/grh-compact/src/New/Filling/NearTerms.lean\#L94-L110}{{\leanfont besselU\_dense}}}
\expandafter\newcommand\csname LN@unipotentCharone\endcsname{\href{https://github.com/lsernau/grh-lean/blob/72fead102f663d5c7a407964a13dc0cf9d977207/grh-compact/src/New/Filling/NearTerms.lean\#L120-L121}{{\leanfont Filling.\allowbreak{}unipotentChar\_one}}}
\expandafter\newcommand\csname LN@existsbesselUapprox\endcsname{\href{https://github.com/lsernau/grh-lean/blob/72fead102f663d5c7a407964a13dc0cf9d977207/grh-compact/src/New/Filling/NearTerms.lean\#L125-L131}{{\leanfont Filling.\allowbreak{}exists\_besselU\_approx}}}
\expandafter\newcommand\csname LQ@existsbesselUapprox\endcsname{\leanquoteline{0}{lemma Filling.exists\_besselU\_approx \{x\} (hx : x ∈ besselU) y : ∃ u ∈ besselU,}\leanquoteline{2}{y⁻¹ * (x⁻¹ * S.ratToFin u) * y ∈}\leanquoteline{2}{S.principalLevel}}
\expandafter\newcommand\csname LN@unipotentChareqofmemprincipalLevel\endcsname{\href{https://github.com/lsernau/grh-lean/blob/72fead102f663d5c7a407964a13dc0cf9d977207/grh-compact/src/New/Filling/NearTerms.lean\#L135-L141}{{\leanfont Filling.\allowbreak{}unipotentChar\_eq\_of\_mem\_principalLevel}}}
\expandafter\newcommand\csname LN@eqandmemclassUnipofbesselS\endcsname{\href{https://github.com/lsernau/grh-lean/blob/72fead102f663d5c7a407964a13dc0cf9d977207/grh-compact/src/New/Filling/NearTerms.lean\#L145-L150}{{\leanfont Filling.\allowbreak{}eq\_and\_mem\_classUnip\_of\_besselS}}}
\expandafter\newcommand\csname LQ@eqandmemclassUnipofbesselS\endcsname{\leanquoteline{0}{lemma Filling.eq\_and\_mem\_classUnip\_of\_besselS [ok : S.IsAdmissible] \{c d γ\} (h : γ ∈ besselS) (o :}\leanquoteline{2}{S.termPosition c d γ ∈ S.principalLevel) : c = d ∧ γ ∈ S.classUnip c}}
\expandafter\newcommand\csname LN@existsratapproxunipotentChareq\endcsname{\href{https://github.com/lsernau/grh-lean/blob/72fead102f663d5c7a407964a13dc0cf9d977207/grh-compact/src/New/Filling/NearTerms.lean\#L152-L156}{{\leanfont Filling.\allowbreak{}exists\_rat\_approx\_unipotentChar\_eq}}}
\expandafter\newcommand\csname LN@existsbesselSmemtermsOver\endcsname{\href{https://github.com/lsernau/grh-lean/blob/72fead102f663d5c7a407964a13dc0cf9d977207/grh-compact/src/New/Filling/NearTerms.lean\#L158-L165}{{\leanfont Filling.\allowbreak{}exists\_besselS\_mem\_termsOver}}}
\expandafter\newcommand\csname LN@eqoneofreduceAteqone\endcsname{\href{https://github.com/lsernau/grh-lean/blob/72fead102f663d5c7a407964a13dc0cf9d977207/grh-compact/src/New/Filling/NearTerms.lean\#L170-L176}{{\leanfont Filling.\allowbreak{}eq\_one\_of\_reduceAt\_eq\_one}}}
\expandafter\newcommand\csname LN@frameImagemul\endcsname{\href{https://github.com/lsernau/grh-lean/blob/72fead102f663d5c7a407964a13dc0cf9d977207/grh-compact/src/New/Filling/NearTerms.lean\#L187-L187}{{\leanfont Filling.\allowbreak{}frameImage\_mul}}}
\expandafter\newcommand\csname LN@frameImagemulofmembesselU\endcsname{\href{https://github.com/lsernau/grh-lean/blob/72fead102f663d5c7a407964a13dc0cf9d977207/grh-compact/src/New/Filling/NearTerms.lean\#L189-L190}{{\leanfont Filling.\allowbreak{}frameImage\_mul\_of\_mem\_besselU0}}}
\expandafter\newcommand\csname LN@frameImageinvmuleqiff\endcsname{\href{https://github.com/lsernau/grh-lean/blob/72fead102f663d5c7a407964a13dc0cf9d977207/grh-compact/src/New/Filling/NearTerms.lean\#L192-L193}{{\leanfont Filling.\allowbreak{}frameImage\_inv\_mul\_eq\_iff}}}
\expandafter\newcommand\csname LN@starmullasteqofframeImageeq\endcsname{\href{https://github.com/lsernau/grh-lean/blob/72fead102f663d5c7a407964a13dc0cf9d977207/grh-compact/src/New/Filling/NearTerms.lean\#L201-L203}{{\leanfont Filling.\allowbreak{}star\_mul\_last\_eq\_of\_frameImage\_eq}}}
\expandafter\newcommand\csname LN@frameProductsmulVecofframeImageeq\endcsname{\href{https://github.com/lsernau/grh-lean/blob/72fead102f663d5c7a407964a13dc0cf9d977207/grh-compact/src/New/Filling/NearTerms.lean\#L205-L207}{{\leanfont Filling.\allowbreak{}frameProducts\_mulVec\_of\_frameImage\_eq}}}
\expandafter\newcommand\csname LN@eqzeroofratMiddleformeqzero\endcsname{\href{https://github.com/lsernau/grh-lean/blob/72fead102f663d5c7a407964a13dc0cf9d977207/grh-compact/src/New/Filling/NearTerms.lean\#L209-L210}{{\leanfont Filling.\allowbreak{}eq\_zero\_of\_ratMiddle\_form\_eq\_zero}}}
\expandafter\newcommand\csname LN@normsqmiddleparteq\endcsname{\href{https://github.com/lsernau/grh-lean/blob/72fead102f663d5c7a407964a13dc0cf9d977207/grh-compact/src/New/Filling/NearTerms.lean\#L215-L223}{{\leanfont Filling.\allowbreak{}norm\_sq\_middle\_part\_eq}}}
\expandafter\newcommand\csname LN@maptermDisceqnorm\endcsname{\href{https://github.com/lsernau/grh-lean/blob/72fead102f663d5c7a407964a13dc0cf9d977207/grh-compact/src/New/Filling/NearTerms.lean\#L225-L228}{{\leanfont Filling.\allowbreak{}map\_termDisc\_eq\_norm}}}
\expandafter\newcommand\csname LN@hasLevellast\endcsname{\href{https://github.com/lsernau/grh-lean/blob/72fead102f663d5c7a407964a13dc0cf9d977207/grh-compact/src/New/Filling/NearTerms.lean\#L230-L231}{{\leanfont Filling.\allowbreak{}hasLevel\_last}}}
\expandafter\newcommand\csname LN@gramDiscRealrealFrameGramratToInf\endcsname{\href{https://github.com/lsernau/grh-lean/blob/72fead102f663d5c7a407964a13dc0cf9d977207/grh-compact/src/New/Filling/NearTerms.lean\#L233-L236}{{\leanfont Filling.\allowbreak{}gramDiscReal\_realFrameGram\_ratToInf}}}
\expandafter\newcommand\csname LN@ratToInfmembesselU\endcsname{\href{https://github.com/lsernau/grh-lean/blob/72fead102f663d5c7a407964a13dc0cf9d977207/grh-compact/src/New/Filling/NearTerms.lean\#L238-L239}{{\leanfont Filling.\allowbreak{}ratToInf\_mem\_besselU0}}}
\expandafter\newcommand\csname LN@mulmemboxedPointsiff\endcsname{\href{https://github.com/lsernau/grh-lean/blob/72fead102f663d5c7a407964a13dc0cf9d977207/grh-compact/src/New/Filling/NearTerms.lean\#L241-L243}{{\leanfont Filling.\allowbreak{}mul\_mem\_boxedPoints\_iff}}}
\expandafter\newcommand\csname LN@ratToInf\endcsname{\href{https://github.com/lsernau/grh-lean/blob/72fead102f663d5c7a407964a13dc0cf9d977207/grh-compact/src/New/Filling/NearTerms.lean\#L245-L250}{{\leanfont IsNormalForm.\allowbreak{}ratToInf}}}
\expandafter\newcommand\csname LN@isIntMatrixtermPosition\endcsname{\href{https://github.com/lsernau/grh-lean/blob/72fead102f663d5c7a407964a13dc0cf9d977207/grh-compact/src/New/Filling/NearTerms.lean\#L252-L253}{{\leanfont Filling.\allowbreak{}termsOver.\allowbreak{}isIntMatrix\_termPosition}}}
\expandafter\newcommand\csname LN@classCoordconjsplitGram\endcsname{\href{https://github.com/lsernau/grh-lean/blob/72fead102f663d5c7a407964a13dc0cf9d977207/grh-compact/src/New/Filling/NearTerms.lean\#L255-L257}{{\leanfont Filling.\allowbreak{}classCoord\_conj\_splitGram}}}
\expandafter\newcommand\csname LN@frameImagetermPosition\endcsname{\href{https://github.com/lsernau/grh-lean/blob/72fead102f663d5c7a407964a13dc0cf9d977207/grh-compact/src/New/Filling/NearTerms.lean\#L259-L260}{{\leanfont Filling.\allowbreak{}frameImage\_termPosition}}}
\expandafter\newcommand\csname LN@frameGramtermPosition\endcsname{\href{https://github.com/lsernau/grh-lean/blob/72fead102f663d5c7a407964a13dc0cf9d977207/grh-compact/src/New/Filling/NearTerms.lean\#L262-L266}{{\leanfont Filling.\allowbreak{}frameGram\_termPosition}}}
\expandafter\newcommand\csname LN@mulofmembesselU\endcsname{\href{https://github.com/lsernau/grh-lean/blob/72fead102f663d5c7a407964a13dc0cf9d977207/grh-compact/src/New/Filling/NearTerms.lean\#L268-L270}{{\leanfont Filling.\allowbreak{}IsNear.\allowbreak{}mul\_of\_mem\_besselU0}}}
\expandafter\newcommand\csname LN@membesselSiffframeImage\endcsname{\href{https://github.com/lsernau/grh-lean/blob/72fead102f663d5c7a407964a13dc0cf9d977207/grh-compact/src/New/Filling/NearTerms.lean\#L278-L279}{{\leanfont Filling.\allowbreak{}mem\_besselS\_iff\_frameImage}}}
\expandafter\newcommand\csname LN@powletailScale\endcsname{\href{https://github.com/lsernau/grh-lean/blob/72fead102f663d5c7a407964a13dc0cf9d977207/grh-compact/src/New/Filling/NearTerms.lean\#L281-L283}{{\leanfont Filling.\allowbreak{}pow\_le\_tailScale}}}
\expandafter\newcommand\csname LN@tailScalepowfour\endcsname{\href{https://github.com/lsernau/grh-lean/blob/72fead102f663d5c7a407964a13dc0cf9d977207/grh-compact/src/New/Filling/NearTerms.lean\#L285-L286}{{\leanfont Filling.\allowbreak{}tailScale\_pow\_four}}}
\expandafter\newcommand\csname LN@ofsizebounds\endcsname{\href{https://github.com/lsernau/grh-lean/blob/72fead102f663d5c7a407964a13dc0cf9d977207/grh-compact/src/New/Filling/NearTerms.lean\#L288-L293}{{\leanfont Filling.\allowbreak{}of\_size\_bounds}}}
\expandafter\newcommand\csname LN@normframeProductsle\endcsname{\href{https://github.com/lsernau/grh-lean/blob/72fead102f663d5c7a407964a13dc0cf9d977207/grh-compact/src/New/Filling/NearTerms.lean\#L295-L300}{{\leanfont Filling.\allowbreak{}IsNear.\allowbreak{}norm\_frameProducts\_le}}}
\expandafter\newcommand\csname LN@normmapleadingMinorle\endcsname{\href{https://github.com/lsernau/grh-lean/blob/72fead102f663d5c7a407964a13dc0cf9d977207/grh-compact/src/New/Filling/NearTerms.lean\#L302-L305}{{\leanfont Filling.\allowbreak{}IsNear.\allowbreak{}norm\_map\_leadingMinor\_le}}}
\expandafter\newcommand\csname LN@eqandmemclassUnip\endcsname{\href{https://github.com/lsernau/grh-lean/blob/72fead102f663d5c7a407964a13dc0cf9d977207/grh-compact/src/New/Filling/NearTerms.lean\#L309-L313}{{\leanfont Filling.\allowbreak{}termsOver.\allowbreak{}eq\_and\_mem\_classUnip₀}}}
\expandafter\newcommand\csname LN@isIntegralframeProducts\endcsname{\href{https://github.com/lsernau/grh-lean/blob/72fead102f663d5c7a407964a13dc0cf9d977207/grh-compact/src/New/Filling/NearTerms.lean\#L315-L317}{{\leanfont Filling.\allowbreak{}isIntegral\_frameProducts}}}
\expandafter\newcommand\csname LN@Filling\endcsname{\href{https://github.com/lsernau/grh-lean/blob/72fead102f663d5c7a407964a13dc0cf9d977207/grh-compact/src/New/Filling/Setup.lean\#L22-L58}{{\leanfont Filling}}}
\expandafter\newcommand\csname LN@g\endcsname{\href{https://github.com/lsernau/grh-lean/blob/72fead102f663d5c7a407964a13dc0cf9d977207/grh-compact/src/New/Filling/Setup.lean\#L70-L70}{{\leanfont Filling.\allowbreak{}g}}}
\expandafter\newcommand\csname LN@ratMiddle\endcsname{\href{https://github.com/lsernau/grh-lean/blob/72fead102f663d5c7a407964a13dc0cf9d977207/grh-compact/src/New/Filling/Setup.lean\#L74-L74}{{\leanfont Filling.\allowbreak{}ratMiddle}}}
\expandafter\newcommand\csname LN@latMiddle\endcsname{\href{https://github.com/lsernau/grh-lean/blob/72fead102f663d5c7a407964a13dc0cf9d977207/grh-compact/src/New/Filling/Setup.lean\#L78-L78}{{\leanfont Filling.\allowbreak{}latMiddle}}}
\expandafter\newcommand\csname LN@G\endcsname{\href{https://github.com/lsernau/grh-lean/blob/72fead102f663d5c7a407964a13dc0cf9d977207/grh-compact/src/New/Filling/Setup.lean\#L82-L82}{{\leanfont Filling.\allowbreak{}G}}}
\expandafter\newcommand\csname LN@GFin\endcsname{\href{https://github.com/lsernau/grh-lean/blob/72fead102f663d5c7a407964a13dc0cf9d977207/grh-compact/src/New/Filling/Setup.lean\#L86-L86}{{\leanfont Filling.\allowbreak{}GFin}}}
\expandafter\newcommand\csname LN@GRes\endcsname{\href{https://github.com/lsernau/grh-lean/blob/72fead102f663d5c7a407964a13dc0cf9d977207/grh-compact/src/New/Filling/Setup.lean\#L91-L91}{{\leanfont Filling.\allowbreak{}GRes}}}
\expandafter\newcommand\csname LN@BesselCoord\endcsname{\href{https://github.com/lsernau/grh-lean/blob/72fead102f663d5c7a407964a13dc0cf9d977207/grh-compact/src/New/Filling/Setup.lean\#L95-L95}{{\leanfont Filling.\allowbreak{}BesselCoord}}}
\expandafter\newcommand\csname LN@Classes\endcsname{\href{https://github.com/lsernau/grh-lean/blob/72fead102f663d5c7a407964a13dc0cf9d977207/grh-compact/src/New/Filling/Setup.lean\#L99-L99}{{\leanfont Filling.\allowbreak{}Classes}}}
\expandafter\newcommand\csname LN@frameProducts\endcsname{\href{https://github.com/lsernau/grh-lean/blob/72fead102f663d5c7a407964a13dc0cf9d977207/grh-compact/src/New/Filling/Setup.lean\#L107-L107}{{\leanfont Filling.\allowbreak{}frameProducts}}}
\expandafter\newcommand\csname LN@latStabilizer\endcsname{\href{https://github.com/lsernau/grh-lean/blob/72fead102f663d5c7a407964a13dc0cf9d977207/grh-compact/src/New/Filling/Setup.lean\#L111-L111}{{\leanfont Filling.\allowbreak{}latStabilizer}}}
\expandafter\newcommand\csname LN@reduction\endcsname{\href{https://github.com/lsernau/grh-lean/blob/72fead102f663d5c7a407964a13dc0cf9d977207/grh-compact/src/New/Filling/Setup.lean\#L115-L116}{{\leanfont Filling.\allowbreak{}reduction}}}
\expandafter\newcommand\csname LN@principalLevel\endcsname{\href{https://github.com/lsernau/grh-lean/blob/72fead102f663d5c7a407964a13dc0cf9d977207/grh-compact/src/New/Filling/Setup.lean\#L120-L120}{{\leanfont Filling.\allowbreak{}principalLevel}}}
\expandafter\newcommand\csname LN@reducedImage\endcsname{\href{https://github.com/lsernau/grh-lean/blob/72fead102f663d5c7a407964a13dc0cf9d977207/grh-compact/src/New/Filling/Setup.lean\#L124-L124}{{\leanfont Filling.\allowbreak{}reducedImage}}}
\expandafter\newcommand\csname LN@kappa\endcsname{\href{https://github.com/lsernau/grh-lean/blob/72fead102f663d5c7a407964a13dc0cf9d977207/grh-compact/src/New/Filling/Setup.lean\#L128-L128}{{\leanfont Filling.\allowbreak{}kappa}}}
\expandafter\newcommand\csname LN@basisScale\endcsname{\href{https://github.com/lsernau/grh-lean/blob/72fead102f663d5c7a407964a13dc0cf9d977207/grh-compact/src/New/Filling/Setup.lean\#L132-L133}{{\leanfont Filling.\allowbreak{}basisScale}}}
\expandafter\newcommand\csname LN@mapsplitGramratMiddle\endcsname{\href{https://github.com/lsernau/grh-lean/blob/72fead102f663d5c7a407964a13dc0cf9d977207/grh-compact/src/New/Filling/Setup.lean\#L138-L139}{{\leanfont Filling.\allowbreak{}map\_splitGram\_ratMiddle}}}
\expandafter\newcommand\csname LN@selfDualChar\endcsname{\href{https://github.com/lsernau/grh-lean/blob/72fead102f663d5c7a407964a13dc0cf9d977207/grh-compact/src/New/Filling/Setup.lean\#L142-L142}{{\leanfont Filling.\allowbreak{}selfDualChar}}}
\expandafter\newcommand\csname LN@varpinezero\endcsname{\href{https://github.com/lsernau/grh-lean/blob/72fead102f663d5c7a407964a13dc0cf9d977207/grh-compact/src/New/Filling/Setup.lean\#L144-L145}{{\leanfont Filling.\allowbreak{}varpi\_ne\_zero}}}
\expandafter\newcommand\csname LN@frameScaleSq\endcsname{\href{https://github.com/lsernau/grh-lean/blob/72fead102f663d5c7a407964a13dc0cf9d977207/grh-compact/src/New/Filling/Setup.lean\#L148-L148}{{\leanfont Filling.\allowbreak{}frameScaleSq}}}
\expandafter\newcommand\csname LN@frameScaleProd\endcsname{\href{https://github.com/lsernau/grh-lean/blob/72fead102f663d5c7a407964a13dc0cf9d977207/grh-compact/src/New/Filling/Setup.lean\#L152-L152}{{\leanfont Filling.\allowbreak{}frameScaleProd}}}
\expandafter\newcommand\csname LN@unipCovolume\endcsname{\href{https://github.com/lsernau/grh-lean/blob/72fead102f663d5c7a407964a13dc0cf9d977207/grh-compact/src/New/Filling/Setup.lean\#L156-L157}{{\leanfont Filling.\allowbreak{}unipCovolume}}}
\expandafter\newcommand\csname LN@siegelConst\endcsname{\href{https://github.com/lsernau/grh-lean/blob/72fead102f663d5c7a407964a13dc0cf9d977207/grh-compact/src/New/Filling/Setup.lean\#L161-L162}{{\leanfont Filling.\allowbreak{}siegelConst}}}
\expandafter\newcommand\csname LN@termDisc\endcsname{\href{https://github.com/lsernau/grh-lean/blob/72fead102f663d5c7a407964a13dc0cf9d977207/grh-compact/src/New/Filling/Setup.lean\#L167-L167}{{\leanfont Filling.\allowbreak{}termDisc}}}
\expandafter\newcommand\csname LN@basisScalebesselFrame\endcsname{\href{https://github.com/lsernau/grh-lean/blob/72fead102f663d5c7a407964a13dc0cf9d977207/grh-compact/src/New/Filling/Setup.lean\#L169-L170}{{\leanfont Filling.\allowbreak{}basisScale\_besselFrame}}}
\expandafter\newcommand\csname LN@coekappa\endcsname{\href{https://github.com/lsernau/grh-lean/blob/72fead102f663d5c7a407964a13dc0cf9d977207/grh-compact/src/New/Filling/Setup.lean\#L172-L173}{{\leanfont Filling.\allowbreak{}coe\_kappa}}}
\expandafter\newcommand\csname LN@kappanezero\endcsname{\href{https://github.com/lsernau/grh-lean/blob/72fead102f663d5c7a407964a13dc0cf9d977207/grh-compact/src/New/Filling/Setup.lean\#L175-L175}{{\leanfont Filling.\allowbreak{}kappa\_ne\_zero}}}
\expandafter\newcommand\csname LN@frameScaleSqpos\endcsname{\href{https://github.com/lsernau/grh-lean/blob/72fead102f663d5c7a407964a13dc0cf9d977207/grh-compact/src/New/Filling/Setup.lean\#L177-L177}{{\leanfont Filling.\allowbreak{}frameScaleSq\_pos}}}
\expandafter\newcommand\csname LN@frameScaleProdpos\endcsname{\href{https://github.com/lsernau/grh-lean/blob/72fead102f663d5c7a407964a13dc0cf9d977207/grh-compact/src/New/Filling/Setup.lean\#L179-L179}{{\leanfont Filling.\allowbreak{}frameScaleProd\_pos}}}
\expandafter\newcommand\csname LN@basisScalenezero\endcsname{\href{https://github.com/lsernau/grh-lean/blob/72fead102f663d5c7a407964a13dc0cf9d977207/grh-compact/src/New/Filling/Setup.lean\#L181-L181}{{\leanfont Filling.\allowbreak{}basisScale\_ne\_zero}}}
\expandafter\newcommand\csname LN@basisScaleGL\endcsname{\href{https://github.com/lsernau/grh-lean/blob/72fead102f663d5c7a407964a13dc0cf9d977207/grh-compact/src/New/Filling/Setup.lean\#L185-L185}{{\leanfont Filling.\allowbreak{}basisScaleGL}}}
\expandafter\newcommand\csname LN@grambasisScaleGL\endcsname{\href{https://github.com/lsernau/grh-lean/blob/72fead102f663d5c7a407964a13dc0cf9d977207/grh-compact/src/New/Filling/Setup.lean\#L187-L188}{{\leanfont Filling.\allowbreak{}gram\_basisScaleGL}}}
\expandafter\newcommand\csname LN@latBasis\endcsname{\href{https://github.com/lsernau/grh-lean/blob/72fead102f663d5c7a407964a13dc0cf9d977207/grh-compact/src/New/Filling/Setup.lean\#L192-L192}{{\leanfont Filling.\allowbreak{}latBasis}}}
\expandafter\newcommand\csname LN@toLatBasis\endcsname{\href{https://github.com/lsernau/grh-lean/blob/72fead102f663d5c7a407964a13dc0cf9d977207/grh-compact/src/New/Filling/Setup.lean\#L196-L197}{{\leanfont Filling.\allowbreak{}toLatBasis}}}
\expandafter\newcommand\csname LN@ratToAdelic\endcsname{\href{https://github.com/lsernau/grh-lean/blob/72fead102f663d5c7a407964a13dc0cf9d977207/grh-compact/src/New/Filling/Setup.lean\#L201-L201}{{\leanfont Filling.\allowbreak{}ratToAdelic}}}
\expandafter\newcommand\csname LN@adelicToInf\endcsname{\href{https://github.com/lsernau/grh-lean/blob/72fead102f663d5c7a407964a13dc0cf9d977207/grh-compact/src/New/Filling/Setup.lean\#L205-L206}{{\leanfont Filling.\allowbreak{}adelicToInf}}}
\expandafter\newcommand\csname LN@adelicToFin\endcsname{\href{https://github.com/lsernau/grh-lean/blob/72fead102f663d5c7a407964a13dc0cf9d977207/grh-compact/src/New/Filling/Setup.lean\#L210-L211}{{\leanfont Filling.\allowbreak{}adelicToFin}}}
\expandafter\newcommand\csname LN@ratToInfa\endcsname{\href{https://github.com/lsernau/grh-lean/blob/72fead102f663d5c7a407964a13dc0cf9d977207/grh-compact/src/New/Filling/Setup.lean\#L215-L215}{{\leanfont Filling.\allowbreak{}ratToInf}}}
\expandafter\newcommand\csname LN@ratToFin\endcsname{\href{https://github.com/lsernau/grh-lean/blob/72fead102f663d5c7a407964a13dc0cf9d977207/grh-compact/src/New/Filling/Setup.lean\#L219-L219}{{\leanfont Filling.\allowbreak{}ratToFin}}}
\expandafter\newcommand\csname LN@inclHFin\endcsname{\href{https://github.com/lsernau/grh-lean/blob/72fead102f663d5c7a407964a13dc0cf9d977207/grh-compact/src/New/Filling/Setup.lean\#L223-L223}{{\leanfont Filling.\allowbreak{}inclHFin}}}
\expandafter\newcommand\csname LN@h\endcsname{\href{https://github.com/lsernau/grh-lean/blob/72fead102f663d5c7a407964a13dc0cf9d977207/grh-compact/src/New/Filling/Setup.lean\#L229-L229}{{\leanfont Filling.\allowbreak{}h}}}
\expandafter\newcommand\csname LN@termPosition\endcsname{\href{https://github.com/lsernau/grh-lean/blob/72fead102f663d5c7a407964a13dc0cf9d977207/grh-compact/src/New/Filling/Setup.lean\#L233-L233}{{\leanfont Filling.\allowbreak{}termPosition}}}
\expandafter\newcommand\csname LN@termsOver\endcsname{\href{https://github.com/lsernau/grh-lean/blob/72fead102f663d5c7a407964a13dc0cf9d977207/grh-compact/src/New/Filling/Setup.lean\#L237-L237}{{\leanfont Filling.\allowbreak{}termsOver}}}
\expandafter\newcommand\csname LN@classUnip\endcsname{\href{https://github.com/lsernau/grh-lean/blob/72fead102f663d5c7a407964a13dc0cf9d977207/grh-compact/src/New/Filling/Setup.lean\#L240-L241}{{\leanfont Filling.\allowbreak{}classUnip}}}
\expandafter\newcommand\csname LN@classUnipa\endcsname{\href{https://github.com/lsernau/grh-lean/blob/72fead102f663d5c7a407964a13dc0cf9d977207/grh-compact/src/New/Filling/Setup.lean\#L244-L244}{{\leanfont Filling.\allowbreak{}classUnip₀}}}
\expandafter\newcommand\csname LN@boxSize\endcsname{\href{https://github.com/lsernau/grh-lean/blob/72fead102f663d5c7a407964a13dc0cf9d977207/grh-compact/src/New/Filling/Setup.lean\#L247-L247}{{\leanfont Filling.\allowbreak{}boxSize}}}
\expandafter\newcommand\csname LN@IsBoxPartition\endcsname{\href{https://github.com/lsernau/grh-lean/blob/72fead102f663d5c7a407964a13dc0cf9d977207/grh-compact/src/New/Filling/Setup.lean\#L252-L255}{{\leanfont Filling.\allowbreak{}IsBoxPartition}}}
\expandafter\newcommand\csname LN@besselFn\endcsname{\href{https://github.com/lsernau/grh-lean/blob/72fead102f663d5c7a407964a13dc0cf9d977207/grh-compact/src/New/Filling/Setup.lean\#L259-L259}{{\leanfont Filling.\allowbreak{}besselFn}}}
\expandafter\newcommand\csname LN@besselKernel\endcsname{\href{https://github.com/lsernau/grh-lean/blob/72fead102f663d5c7a407964a13dc0cf9d977207/grh-compact/src/New/Filling/Setup.lean\#L263-L264}{{\leanfont Filling.\allowbreak{}besselKernel}}}
\expandafter\newcommand\csname LN@columnDim\endcsname{\href{https://github.com/lsernau/grh-lean/blob/72fead102f663d5c7a407964a13dc0cf9d977207/grh-compact/src/New/Filling/Setup.lean\#L268-L268}{{\leanfont Filling.\allowbreak{}columnDim}}}
\expandafter\newcommand\csname LN@besselDegree\endcsname{\href{https://github.com/lsernau/grh-lean/blob/72fead102f663d5c7a407964a13dc0cf9d977207/grh-compact/src/New/Filling/Setup.lean\#L271-L271}{{\leanfont Filling.\allowbreak{}besselDegree}}}
\expandafter\newcommand\csname LN@besselDegreepos\endcsname{\href{https://github.com/lsernau/grh-lean/blob/72fead102f663d5c7a407964a13dc0cf9d977207/grh-compact/src/New/Filling/Setup.lean\#L273-L273}{{\leanfont Filling.\allowbreak{}besselDegree\_pos}}}
\expandafter\newcommand\csname LN@mainEnergy\endcsname{\href{https://github.com/lsernau/grh-lean/blob/72fead102f663d5c7a407964a13dc0cf9d977207/grh-compact/src/New/Filling/Setup.lean\#L276-L276}{{\leanfont Filling.\allowbreak{}mainEnergy}}}
\expandafter\newcommand\csname LN@unipOrder\endcsname{\href{https://github.com/lsernau/grh-lean/blob/72fead102f663d5c7a407964a13dc0cf9d977207/grh-compact/src/New/Filling/Setup.lean\#L280-L280}{{\leanfont Filling.\allowbreak{}unipOrder}}}
\expandafter\newcommand\csname LN@packetConst\endcsname{\href{https://github.com/lsernau/grh-lean/blob/72fead102f663d5c7a407964a13dc0cf9d977207/grh-compact/src/New/Filling/Setup.lean\#L284-L284}{{\leanfont Filling.\allowbreak{}packetConst}}}
\expandafter\newcommand\csname LN@besselEnergy\endcsname{\href{https://github.com/lsernau/grh-lean/blob/72fead102f663d5c7a407964a13dc0cf9d977207/grh-compact/src/New/Filling/Setup.lean\#L289-L289}{{\leanfont Filling.\allowbreak{}besselEnergy}}}
\expandafter\newcommand\csname LN@tailScale\endcsname{\href{https://github.com/lsernau/grh-lean/blob/72fead102f663d5c7a407964a13dc0cf9d977207/grh-compact/src/New/Filling/Setup.lean\#L292-L292}{{\leanfont Filling.\allowbreak{}tailScale}}}
\expandafter\newcommand\csname LN@codimFactor\endcsname{\href{https://github.com/lsernau/grh-lean/blob/72fead102f663d5c7a407964a13dc0cf9d977207/grh-compact/src/New/Filling/Setup.lean\#L296-L296}{{\leanfont Filling.\allowbreak{}codimFactor}}}
\expandafter\newcommand\csname LN@IsAdmissible\endcsname{\href{https://github.com/lsernau/grh-lean/blob/72fead102f663d5c7a407964a13dc0cf9d977207/grh-compact/src/New/Filling/Setup.lean\#L301-L318}{{\leanfont Filling.\allowbreak{}IsAdmissible}}}
\expandafter\newcommand\csname LN@adelicToInfapply\endcsname{\href{https://github.com/lsernau/grh-lean/blob/72fead102f663d5c7a407964a13dc0cf9d977207/grh-compact/src/New/Filling/Setup.lean\#L320-L321}{{\leanfont Filling.\allowbreak{}adelicToInf\_apply}}}
\expandafter\newcommand\csname LN@ratToInfapply\endcsname{\href{https://github.com/lsernau/grh-lean/blob/72fead102f663d5c7a407964a13dc0cf9d977207/grh-compact/src/New/Filling/Setup.lean\#L323-L324}{{\leanfont Filling.\allowbreak{}ratToInf\_apply}}}
\expandafter\newcommand\csname LN@coeratToFin\endcsname{\href{https://github.com/lsernau/grh-lean/blob/72fead102f663d5c7a407964a13dc0cf9d977207/grh-compact/src/New/Filling/Setup.lean\#L326-L326}{{\leanfont Filling.\allowbreak{}coe\_ratToFin}}}
\expandafter\newcommand\csname LN@finiteclassUnipintercompact\endcsname{\href{https://github.com/lsernau/grh-lean/blob/72fead102f663d5c7a407964a13dc0cf9d977207/grh-compact/src/New/Filling/Setup.lean\#L332-L339}{{\leanfont Filling.\allowbreak{}finite\_classUnip\_inter\_compact}}}
\expandafter\newcommand\csname LN@existscutoffclassUnip\endcsname{\href{https://github.com/lsernau/grh-lean/blob/72fead102f663d5c7a407964a13dc0cf9d977207/grh-compact/src/New/Filling/Setup.lean\#L343-L352}{{\leanfont Filling.\allowbreak{}exists\_cutoff\_classUnip}}}
\expandafter\newcommand\csname LN@classCutoff\endcsname{\href{https://github.com/lsernau/grh-lean/blob/72fead102f663d5c7a407964a13dc0cf9d977207/grh-compact/src/New/Filling/Setup.lean\#L356-L357}{{\leanfont Filling.\allowbreak{}classCutoff}}}
\expandafter\newcommand\csname LN@amplitude\endcsname{\href{https://github.com/lsernau/grh-lean/blob/72fead102f663d5c7a407964a13dc0cf9d977207/grh-compact/src/New/Filling/Setup.lean\#L361-L362}{{\leanfont Filling.\allowbreak{}amplitude}}}
\expandafter\newcommand\csname LN@amplitudeAbs\endcsname{\href{https://github.com/lsernau/grh-lean/blob/72fead102f663d5c7a407964a13dc0cf9d977207/grh-compact/src/New/Filling/Setup.lean\#L366-L367}{{\leanfont Filling.\allowbreak{}amplitudeAbs}}}
\expandafter\newcommand\csname LN@amplitudeSum\endcsname{\href{https://github.com/lsernau/grh-lean/blob/72fead102f663d5c7a407964a13dc0cf9d977207/grh-compact/src/New/Filling/Setup.lean\#L370-L370}{{\leanfont Filling.\allowbreak{}amplitudeSum}}}
\expandafter\newcommand\csname LN@abskappaeq\endcsname{\href{https://github.com/lsernau/grh-lean/blob/72fead102f663d5c7a407964a13dc0cf9d977207/grh-compact/src/New/Filling/Setup.lean\#L372-L373}{{\leanfont Filling.\allowbreak{}abs\_kappa\_eq}}}
\expandafter\newcommand\csname LN@realFrameGramratToInf\endcsname{\href{https://github.com/lsernau/grh-lean/blob/72fead102f663d5c7a407964a13dc0cf9d977207/grh-compact/src/New/Filling/Setup.lean\#L375-L377}{{\leanfont Filling.\allowbreak{}realFrameGram\_ratToInf}}}
\expandafter\newcommand\csname LN@detsubmatrixrealFrameGramratToInf\endcsname{\href{https://github.com/lsernau/grh-lean/blob/72fead102f663d5c7a407964a13dc0cf9d977207/grh-compact/src/New/Filling/Setup.lean\#L379-L382}{{\leanfont Filling.\allowbreak{}det\_submatrix\_realFrameGram\_ratToInf}}}
\expandafter\newcommand\csname LN@unipCovolumepos\endcsname{\href{https://github.com/lsernau/grh-lean/blob/72fead102f663d5c7a407964a13dc0cf9d977207/grh-compact/src/New/Filling/Setup.lean\#L384-L384}{{\leanfont Filling.\allowbreak{}unipCovolume\_pos}}}
\expandafter\newcommand\csname LN@gIdele\endcsname{\href{https://github.com/lsernau/grh-lean/blob/72fead102f663d5c7a407964a13dc0cf9d977207/grh-compact/src/New/Filling/Setup.lean\#L387-L387}{{\leanfont Filling.\allowbreak{}gIdele}}}
\expandafter\newcommand\csname LN@unipotentChar\endcsname{\href{https://github.com/lsernau/grh-lean/blob/72fead102f663d5c7a407964a13dc0cf9d977207/grh-compact/src/New/Filling/Setup.lean\#L391-L391}{{\leanfont Filling.\allowbreak{}unipotentChar}}}
\expandafter\newcommand\csname LN@Agrees\endcsname{\href{https://github.com/lsernau/grh-lean/blob/72fead102f663d5c7a407964a13dc0cf9d977207/grh-compact/src/New/Filling/Setup.lean\#L395-L396}{{\leanfont Filling.\allowbreak{}Agrees}}}
\expandafter\newcommand\csname LN@IsStable\endcsname{\href{https://github.com/lsernau/grh-lean/blob/72fead102f663d5c7a407964a13dc0cf9d977207/grh-compact/src/New/Filling/Setup.lean\#L400-L400}{{\leanfont Filling.\allowbreak{}IsStable}}}
\expandafter\newcommand\csname LN@termsOfLink\endcsname{\href{https://github.com/lsernau/grh-lean/blob/72fead102f663d5c7a407964a13dc0cf9d977207/grh-compact/src/New/Filling/Setup.lean\#L404-L404}{{\leanfont Filling.\allowbreak{}termsOfLink}}}
\expandafter\newcommand\csname LN@splitGrammulVec\endcsname{\href{https://github.com/lsernau/grh-lean/blob/72fead102f663d5c7a407964a13dc0cf9d977207/grh-compact/src/New/Filling/StableCells.lean\#L15-L16}{{\leanfont splitGram\_mulVec}}}
\expandafter\newcommand\csname LN@linkMatrixmulVec\endcsname{\href{https://github.com/lsernau/grh-lean/blob/72fead102f663d5c7a407964a13dc0cf9d977207/grh-compact/src/New/Filling/StableCells.lean\#L18-L19}{{\leanfont linkMatrix\_mulVec}}}
\expandafter\newcommand\csname LN@isotropicFramePairmulVeceqzeroiff\endcsname{\href{https://github.com/lsernau/grh-lean/blob/72fead102f663d5c7a407964a13dc0cf9d977207/grh-compact/src/New/Filling/StableCells.lean\#L21-L24}{{\leanfont isotropicFramePair\_mulVec\_eq\_zero\_iff}}}
\expandafter\newcommand\csname LN@memlink\endcsname{\href{https://github.com/lsernau/grh-lean/blob/72fead102f663d5c7a407964a13dc0cf9d977207/grh-compact/src/New/Filling/StableCells.lean\#L28-L32}{{\leanfont mem\_link}}}
\expandafter\newcommand\csname LN@invb\endcsname{\href{https://github.com/lsernau/grh-lean/blob/72fead102f663d5c7a407964a13dc0cf9d977207/grh-compact/src/New/Filling/StableCells.lean\#L34-L36}{{\leanfont IsNormalForm.\allowbreak{}inv}}}
\expandafter\newcommand\csname LN@existsframeImageeqofisNormalForm\endcsname{\href{https://github.com/lsernau/grh-lean/blob/72fead102f663d5c7a407964a13dc0cf9d977207/grh-compact/src/New/Filling/StableCells.lean\#L38-L44}{{\leanfont Filling.\allowbreak{}exists\_frameImage\_eq\_of\_isNormalForm}}}
\expandafter\newcommand\csname LN@besselUmul\endcsname{\href{https://github.com/lsernau/grh-lean/blob/72fead102f663d5c7a407964a13dc0cf9d977207/grh-compact/src/New/Filling/StableCells.lean\#L46-L49}{{\leanfont IsNormalForm.\allowbreak{}besselU\_mul}}}
\expandafter\newcommand\csname LN@linkbesselUmul\endcsname{\href{https://github.com/lsernau/grh-lean/blob/72fead102f663d5c7a407964a13dc0cf9d977207/grh-compact/src/New/Filling/StableCells.lean\#L51-L55}{{\leanfont Filling.\allowbreak{}link\_besselU\_mul}}}
\expandafter\newcommand\csname LN@frameDiscframeGrambesselUmul\endcsname{\href{https://github.com/lsernau/grh-lean/blob/72fead102f663d5c7a407964a13dc0cf9d977207/grh-compact/src/New/Filling/StableCells.lean\#L57-L60}{{\leanfont Filling.\allowbreak{}frameDisc\_frameGram\_besselU\_mul}}}
\expandafter\newcommand\csname LN@isStablehasLevel\endcsname{\href{https://github.com/lsernau/grh-lean/blob/72fead102f663d5c7a407964a13dc0cf9d977207/grh-compact/src/New/Filling/StableCells.lean\#L64-L67}{{\leanfont Filling.\allowbreak{}isStable\_hasLevel}}}
\expandafter\newcommand\csname LN@memlinkifforthogonal\endcsname{\href{https://github.com/lsernau/grh-lean/blob/72fead102f663d5c7a407964a13dc0cf9d977207/grh-compact/src/New/Filling/StableCells.lean\#L71-L79}{{\leanfont mem\_link\_iff\_orthogonal}}}
\expandafter\newcommand\csname LQ@memlinkifforthogonal\endcsname{\leanquoteline{0}{lemma mem\_link\_iff\_orthogonal \{γ : S.G E\} \{p\} : p ∈ link γ ↔ ∀ q ∈ link γ,}\leanquoteline{2}{star q ⬝ᵥ S.linkForm *ᵥ p = 0}}
\expandafter\newcommand\csname LN@frameDiscweylElemmul\endcsname{\href{https://github.com/lsernau/grh-lean/blob/72fead102f663d5c7a407964a13dc0cf9d977207/grh-compact/src/New/Filling/StableCells.lean\#L87-L91}{{\leanfont Filling.\allowbreak{}frameDisc\_weylElem\_mul}}}
\expandafter\newcommand\csname LN@linkbigcelleqiff\endcsname{\href{https://github.com/lsernau/grh-lean/blob/72fead102f663d5c7a407964a13dc0cf9d977207/grh-compact/src/New/Filling/StableCells.lean\#L96-L100}{{\leanfont Filling.\allowbreak{}link\_big\_cell\_eq\_iff}}}
\newcommand{\ConcordanceRows}{\crow{openCellConst}{\href{https://github.com/lsernau/grh-lean/blob/72fead102f663d5c7a407964a13dc0cf9d977207/grh-compact/src/New/Filling/BesselSystem.lean\#L18-L18}{{\leanfont Shape.\allowbreak{}openCellConst}}}{BesselSystem.lean}{18--18}
\crow{openCellConstpos}{\href{https://github.com/lsernau/grh-lean/blob/72fead102f663d5c7a407964a13dc0cf9d977207/grh-compact/src/New/Filling/BesselSystem.lean\#L20-L20}{{\leanfont Shape.\allowbreak{}openCellConst\_pos}}}{BesselSystem.lean}{20--20}
\crow{bandConst}{\href{https://github.com/lsernau/grh-lean/blob/72fead102f663d5c7a407964a13dc0cf9d977207/grh-compact/src/New/Filling/BesselSystem.lean\#L24-L24}{{\leanfont Shape.\allowbreak{}bandConst}}}{BesselSystem.lean}{24--24}
\crow{besselSize}{\href{https://github.com/lsernau/grh-lean/blob/72fead102f663d5c7a407964a13dc0cf9d977207/grh-compact/src/New/Filling/BesselSystem.lean\#L29-L29}{{\leanfont besselSize}}}{BesselSystem.lean}{29--29}
\crow{InCut}{\href{https://github.com/lsernau/grh-lean/blob/72fead102f663d5c7a407964a13dc0cf9d977207/grh-compact/src/New/Filling/BesselSystem.lean\#L33-L33}{{\leanfont InCut}}}{BesselSystem.lean}{33--33}
\crow{frameBox}{\href{https://github.com/lsernau/grh-lean/blob/72fead102f663d5c7a407964a13dc0cf9d977207/grh-compact/src/New/Filling/BesselSystem.lean\#L36-L36}{{\leanfont frameBox}}}{BesselSystem.lean}{36--36}
\crow{gramDiscReal}{\href{https://github.com/lsernau/grh-lean/blob/72fead102f663d5c7a407964a13dc0cf9d977207/grh-compact/src/New/Filling/BesselSystem.lean\#L40-L40}{{\leanfont gramDiscReal}}}{BesselSystem.lean}{40--40}
\crow{cellMass}{\href{https://github.com/lsernau/grh-lean/blob/72fead102f663d5c7a407964a13dc0cf9d977207/grh-compact/src/New/Filling/BesselSystem.lean\#L44-L45}{{\leanfont cellMass}}}{BesselSystem.lean}{44--45}
\crow{besselCorrelation}{\href{https://github.com/lsernau/grh-lean/blob/72fead102f663d5c7a407964a13dc0cf9d977207/grh-compact/src/New/Filling/BesselSystem.lean\#L49-L50}{{\leanfont besselCorrelation}}}{BesselSystem.lean}{49--50}
\crow{SchurLe}{\href{https://github.com/lsernau/grh-lean/blob/72fead102f663d5c7a407964a13dc0cf9d977207/grh-compact/src/New/Filling/BesselSystem.lean\#L54-L55}{{\leanfont SchurLe}}}{BesselSystem.lean}{54--55}
\crow{HasSchurBound}{\href{https://github.com/lsernau/grh-lean/blob/72fead102f663d5c7a407964a13dc0cf9d977207/grh-compact/src/New/Filling/BesselSystem.lean\#L59-L62}{{\leanfont HasSchurBound}}}{BesselSystem.lean}{59--62}
\crow{normCorrelation}{\href{https://github.com/lsernau/grh-lean/blob/72fead102f663d5c7a407964a13dc0cf9d977207/grh-compact/src/New/Filling/BesselSystem.lean\#L67-L69}{{\leanfont normCorrelation}}}{BesselSystem.lean}{67--69}
\crow{HasBandLimits}{\href{https://github.com/lsernau/grh-lean/blob/72fead102f663d5c7a407964a13dc0cf9d977207/grh-compact/src/New/Filling/BesselSystem.lean\#L74-L77}{{\leanfont HasBandLimits}}}{BesselSystem.lean}{74--77}
\crow{Side}{\href{https://github.com/lsernau/grh-lean/blob/72fead102f663d5c7a407964a13dc0cf9d977207/grh-compact/src/New/Filling/BesselSystem.lean\#L80-L80}{{\leanfont Side}}}{BesselSystem.lean}{80--80}
\crow{IsCutPair}{\href{https://github.com/lsernau/grh-lean/blob/72fead102f663d5c7a407964a13dc0cf9d977207/grh-compact/src/New/Filling/BesselSystem.lean\#L85-L96}{{\leanfont IsCutPair}}}{BesselSystem.lean}{85--96}
\crow{BesselSystem}{\href{https://github.com/lsernau/grh-lean/blob/72fead102f663d5c7a407964a13dc0cf9d977207/grh-compact/src/New/Filling/BesselSystem.lean\#L102-L109}{{\leanfont BesselSystem}}}{BesselSystem.lean}{102--109}
\crow{reductionapplycongr}{\href{https://github.com/lsernau/grh-lean/blob/72fead102f663d5c7a407964a13dc0cf9d977207/grh-compact/src/New/Filling/CrudeBoundDefs.lean\#L16-L17}{{\leanfont Filling.\allowbreak{}reduction\_apply\_congr}}}{CrudeBoundDefs.lean}{16--17}
\crow{existsreduceAboveGeq}{\href{https://github.com/lsernau/grh-lean/blob/72fead102f663d5c7a407964a13dc0cf9d977207/grh-compact/src/New/Filling/CrudeBoundDefs.lean\#L19-L23}{{\leanfont Filling.\allowbreak{}exists\_reduceAboveG\_eq}}}{CrudeBoundDefs.lean}{19--23}
\crow{vectorFixer}{\href{https://github.com/lsernau/grh-lean/blob/72fead102f663d5c7a407964a13dc0cf9d977207/grh-compact/src/New/Filling/CrudeBoundDefs.lean\#L27-L29}{{\leanfont Filling.\allowbreak{}vectorFixer}}}{CrudeBoundDefs.lean}{27--29}
\crow{FixesNormAtLeast}{\href{https://github.com/lsernau/grh-lean/blob/72fead102f663d5c7a407964a13dc0cf9d977207/grh-compact/src/New/Filling/CrudeBoundDefs.lean\#L33-L34}{{\leanfont Filling.\allowbreak{}FixesNormAtLeast}}}{CrudeBoundDefs.lean}{33--34}
\crow{fixingSet}{\href{https://github.com/lsernau/grh-lean/blob/72fead102f663d5c7a407964a13dc0cf9d977207/grh-compact/src/New/Filling/CrudeBoundDefs.lean\#L38-L39}{{\leanfont Filling.\allowbreak{}fixingSet}}}{CrudeBoundDefs.lean}{38--39}
\crow{zFormPair}{\href{https://github.com/lsernau/grh-lean/blob/72fead102f663d5c7a407964a13dc0cf9d977207/grh-compact/src/New/Filling/CrudeBoundDefs.lean\#L45-L45}{{\leanfont Filling.\allowbreak{}zFormPair}}}{CrudeBoundDefs.lean}{45--45}
\crow{commonLink}{\href{https://github.com/lsernau/grh-lean/blob/72fead102f663d5c7a407964a13dc0cf9d977207/grh-compact/src/New/Filling/CrudeBoundDefs.lean\#L51-L52}{{\leanfont commonLink}}}{CrudeBoundDefs.lean}{51--52}
\crow{HasCommonFlagVector}{\href{https://github.com/lsernau/grh-lean/blob/72fead102f663d5c7a407964a13dc0cf9d977207/grh-compact/src/New/Filling/CrudeBoundDefs.lean\#L56-L57}{{\leanfont HasCommonFlagVector}}}{CrudeBoundDefs.lean}{56--57}
\crow{IsNormalForm}{\href{https://github.com/lsernau/grh-lean/blob/72fead102f663d5c7a407964a13dc0cf9d977207/grh-compact/src/New/Filling/CrudeBoundDefs.lean\#L61-L62}{{\leanfont IsNormalForm}}}{CrudeBoundDefs.lean}{61--62}
\crow{ratGram}{\href{https://github.com/lsernau/grh-lean/blob/72fead102f663d5c7a407964a13dc0cf9d977207/grh-compact/src/New/Filling/CrudeBoundDefs.lean\#L68-L68}{{\leanfont Filling.\allowbreak{}ratGram}}}{CrudeBoundDefs.lean}{68--68}
\crow{linkForm}{\href{https://github.com/lsernau/grh-lean/blob/72fead102f663d5c7a407964a13dc0cf9d977207/grh-compact/src/New/Filling/CrudeBoundDefs.lean\#L72-L72}{{\leanfont Filling.\allowbreak{}linkForm}}}{CrudeBoundDefs.lean}{72--72}
\crow{classCoord}{\href{https://github.com/lsernau/grh-lean/blob/72fead102f663d5c7a407964a13dc0cf9d977207/grh-compact/src/New/Filling/CrudeBoundDefs.lean\#L76-L76}{{\leanfont Filling.\allowbreak{}classCoord}}}{CrudeBoundDefs.lean}{76--76}
\crow{classCoordX}{\href{https://github.com/lsernau/grh-lean/blob/72fead102f663d5c7a407964a13dc0cf9d977207/grh-compact/src/New/Filling/CrudeBoundDefs.lean\#L80-L80}{{\leanfont Filling.\allowbreak{}classCoordX₀}}}{CrudeBoundDefs.lean}{80--80}
\crow{countWeight}{\href{https://github.com/lsernau/grh-lean/blob/72fead102f663d5c7a407964a13dc0cf9d977207/grh-compact/src/New/Filling/CrudeBoundDefs.lean\#L85-L86}{{\leanfont Filling.\allowbreak{}countWeight}}}{CrudeBoundDefs.lean}{85--86}
\crow{HasLevel}{\href{https://github.com/lsernau/grh-lean/blob/72fead102f663d5c7a407964a13dc0cf9d977207/grh-compact/src/New/Filling/CrudeBoundDefs.lean\#L92-L93}{{\leanfont Filling.\allowbreak{}HasLevel}}}{CrudeBoundDefs.lean}{92--93}
\crow{CongruentTerms}{\href{https://github.com/lsernau/grh-lean/blob/72fead102f663d5c7a407964a13dc0cf9d977207/grh-compact/src/New/Filling/CrudeBoundDefs.lean\#L97-L98}{{\leanfont Filling.\allowbreak{}CongruentTerms}}}{CrudeBoundDefs.lean}{97--98}
\crow{termDiscNorm}{\href{https://github.com/lsernau/grh-lean/blob/72fead102f663d5c7a407964a13dc0cf9d977207/grh-compact/src/New/Filling/CrudeBoundDefs.lean\#L102-L102}{{\leanfont Filling.\allowbreak{}termDiscNorm}}}{CrudeBoundDefs.lean}{102--102}
\crow{boxedPoints}{\href{https://github.com/lsernau/grh-lean/blob/72fead102f663d5c7a407964a13dc0cf9d977207/grh-compact/src/New/Filling/CrudeBoundDefs.lean\#L106-L107}{{\leanfont Filling.\allowbreak{}boxedPoints}}}{CrudeBoundDefs.lean}{106--107}
\crow{IsNear}{\href{https://github.com/lsernau/grh-lean/blob/72fead102f663d5c7a407964a13dc0cf9d977207/grh-compact/src/New/Filling/CrudeBoundDefs.lean\#L111-L111}{{\leanfont Filling.\allowbreak{}IsNear}}}{CrudeBoundDefs.lean}{111--111}
\crow{CharsAgree}{\href{https://github.com/lsernau/grh-lean/blob/72fead102f663d5c7a407964a13dc0cf9d977207/grh-compact/src/New/Filling/CrudeBoundDefs.lean\#L115-L116}{{\leanfont Filling.\allowbreak{}CharsAgree}}}{CrudeBoundDefs.lean}{115--116}
\crow{latMiddleGram}{\href{https://github.com/lsernau/grh-lean/blob/72fead102f663d5c7a407964a13dc0cf9d977207/grh-compact/src/New/Filling/EichlerFlag.lean\#L17-L19}{{\leanfont latMiddleGram}}}{EichlerFlag.lean}{17--19}
\crow{middleBaseChange}{\href{https://github.com/lsernau/grh-lean/blob/72fead102f663d5c7a407964a13dc0cf9d977207/grh-compact/src/New/Filling/EichlerFlag.lean\#L29-L29}{{\leanfont middleBaseChange}}}{EichlerFlag.lean}{29--29}
\crow{besselLatticeBasis}{\href{https://github.com/lsernau/grh-lean/blob/72fead102f663d5c7a407964a13dc0cf9d977207/grh-compact/src/New/Filling/EichlerFlag.lean\#L33-L37}{{\leanfont besselLatticeBasis}}}{EichlerFlag.lean}{33--37}
\crow{frameDisc}{\href{https://github.com/lsernau/grh-lean/blob/72fead102f663d5c7a407964a13dc0cf9d977207/grh-compact/src/New/Filling/EichlerFlag.lean\#L43-L43}{{\leanfont frameDisc}}}{EichlerFlag.lean}{43--43}
\crow{grambesselLatticeBasis}{\href{https://github.com/lsernau/grh-lean/blob/72fead102f663d5c7a407964a13dc0cf9d977207/grh-compact/src/New/Filling/EichlerFlag.lean\#L45-L48}{{\leanfont gram\_besselLatticeBasis}}}{EichlerFlag.lean}{45--48}
\crow{framePair}{\href{https://github.com/lsernau/grh-lean/blob/72fead102f663d5c7a407964a13dc0cf9d977207/grh-compact/src/New/Filling/EichlerFlag.lean\#L52-L52}{{\leanfont framePair}}}{EichlerFlag.lean}{52--52}
\crow{isotropicFramePair}{\href{https://github.com/lsernau/grh-lean/blob/72fead102f663d5c7a407964a13dc0cf9d977207/grh-compact/src/New/Filling/EichlerFlag.lean\#L55-L55}{{\leanfont isotropicFramePair}}}{EichlerFlag.lean}{55--55}
\crow{linkMatrix}{\href{https://github.com/lsernau/grh-lean/blob/72fead102f663d5c7a407964a13dc0cf9d977207/grh-compact/src/New/Filling/EichlerFlag.lean\#L59-L59}{{\leanfont linkMatrix}}}{EichlerFlag.lean}{59--59}
\crow{link}{\href{https://github.com/lsernau/grh-lean/blob/72fead102f663d5c7a407964a13dc0cf9d977207/grh-compact/src/New/Filling/EichlerFlag.lean\#L63-L63}{{\leanfont link}}}{EichlerFlag.lean}{63--63}
\crow{dampedPower}{\href{https://github.com/lsernau/grh-lean/blob/72fead102f663d5c7a407964a13dc0cf9d977207/grh-compact/src/New/Filling/EichlerFlag.lean\#L68-L68}{{\leanfont dampedPower}}}{EichlerFlag.lean}{68--68}
\crow{headWeight}{\href{https://github.com/lsernau/grh-lean/blob/72fead102f663d5c7a407964a13dc0cf9d977207/grh-compact/src/New/Filling/EichlerFlag.lean\#L72-L72}{{\leanfont headWeight}}}{EichlerFlag.lean}{72--72}
\crow{SieveSystem}{\href{https://github.com/lsernau/grh-lean/blob/72fead102f663d5c7a407964a13dc0cf9d977207/grh-compact/src/New/Filling/EichlerFlag.lean\#L81-L91}{{\leanfont SieveSystem}}}{EichlerFlag.lean}{81--91}
\crow{pi}{\href{https://github.com/lsernau/grh-lean/blob/72fead102f663d5c7a407964a13dc0cf9d977207/grh-compact/src/New/Filling/EichlerFlag.lean\#L96-L112}{{\leanfont SieveSystem.\allowbreak{}pi}}}{EichlerFlag.lean}{96--112}
\crow{stepsum}{\href{https://github.com/lsernau/grh-lean/blob/72fead102f663d5c7a407964a13dc0cf9d977207/grh-compact/src/New/Filling/EichlerFlag.lean\#L116-L120}{{\leanfont SieveSystem.\allowbreak{}step\_sum}}}{EichlerFlag.lean}{116--120}
\crow{nonnegandlowerSemicontinuous}{\href{https://github.com/lsernau/grh-lean/blob/72fead102f663d5c7a407964a13dc0cf9d977207/grh-compact/src/New/Filling/EichlerFlag.lean\#L125-L139}{{\leanfont SieveSystem.\allowbreak{}nonneg\_and\_lowerSemicontinuous}}}{EichlerFlag.lean}{125--139}
\crow{finHeadWeight}{\href{https://github.com/lsernau/grh-lean/blob/72fead102f663d5c7a407964a13dc0cf9d977207/grh-compact/src/New/Filling/EichlerFlag.lean\#L145-L146}{{\leanfont finHeadWeight}}}{EichlerFlag.lean}{145--146}
\crow{isUnimodularlatMiddleGram}{\href{https://github.com/lsernau/grh-lean/blob/72fead102f663d5c7a407964a13dc0cf9d977207/grh-compact/src/New/Filling/EichlerFlag.lean\#L148-L150}{{\leanfont isUnimodular\_latMiddleGram}}}{EichlerFlag.lean}{148--150}
\crow{frameImageonemulVec}{\href{https://github.com/lsernau/grh-lean/blob/72fead102f663d5c7a407964a13dc0cf9d977207/grh-compact/src/New/Filling/EichlerFlag.lean\#L152-L153}{{\leanfont frameImage\_one\_mulVec}}}{EichlerFlag.lean}{152--153}
\crow{frameImagemulVec}{\href{https://github.com/lsernau/grh-lean/blob/72fead102f663d5c7a407964a13dc0cf9d977207/grh-compact/src/New/Filling/EichlerFlag.lean\#L155-L156}{{\leanfont frameImage\_mulVec}}}{EichlerFlag.lean}{155--156}
\crow{existsstarvecMuleq}{\href{https://github.com/lsernau/grh-lean/blob/72fead102f663d5c7a407964a13dc0cf9d977207/grh-compact/src/New/Filling/EichlerFlag.lean\#L158-L159}{{\leanfont exists\_star\_vecMul\_eq}}}{EichlerFlag.lean}{158--159}
\crow{existsmembesselUeqeichlerMap}{\href{https://github.com/lsernau/grh-lean/blob/72fead102f663d5c7a407964a13dc0cf9d977207/grh-compact/src/New/Filling/EichlerFlag.lean\#L165-L175}{{\leanfont exists\_mem\_besselU\_eq\_eichlerMap}}}{EichlerFlag.lean}{165--175}
\crow{finAdeleMat}{\href{https://github.com/lsernau/grh-lean/blob/72fead102f663d5c7a407964a13dc0cf9d977207/grh-compact/src/New/Filling/GramLiouville.lean\#L18-L18}{{\leanfont Filling.\allowbreak{}finAdeleMat}}}{GramLiouville.lean}{18--18}
\crow{finAdeleMatconjTranspose}{\href{https://github.com/lsernau/grh-lean/blob/72fead102f663d5c7a407964a13dc0cf9d977207/grh-compact/src/New/Filling/GramLiouville.lean\#L20-L20}{{\leanfont Filling.\allowbreak{}finAdeleMat\_conjTranspose}}}{GramLiouville.lean}{20--20}
\crow{mapconjTranspose}{\href{https://github.com/lsernau/grh-lean/blob/72fead102f663d5c7a407964a13dc0cf9d977207/grh-compact/src/New/Filling/GramLiouville.lean\#L22-L22}{{\leanfont Filling.\allowbreak{}map\_conjTranspose}}}{GramLiouville.lean}{22--22}
\crow{finAdeleMatsmul}{\href{https://github.com/lsernau/grh-lean/blob/72fead102f663d5c7a407964a13dc0cf9d977207/grh-compact/src/New/Filling/GramLiouville.lean\#L25-L26}{{\leanfont Filling.\allowbreak{}finAdeleMat\_smul}}}{GramLiouville.lean}{25--26}
\crow{muleqzeroofrangeleker}{\href{https://github.com/lsernau/grh-lean/blob/72fead102f663d5c7a407964a13dc0cf9d977207/grh-compact/src/New/Filling/GramLiouville.lean\#L28-L29}{{\leanfont Filling.\allowbreak{}mul\_eq\_zero\_of\_range\_le\_ker}}}{GramLiouville.lean}{28--29}
\crow{existsbasisextendker}{\href{https://github.com/lsernau/grh-lean/blob/72fead102f663d5c7a407964a13dc0cf9d977207/grh-compact/src/New/Filling/GramLiouville.lean\#L33-L40}{{\leanfont Filling.\allowbreak{}exists\_basis\_extend\_ker}}}{GramLiouville.lean}{33--40}
\crow{sublatticeBasis}{\href{https://github.com/lsernau/grh-lean/blob/72fead102f663d5c7a407964a13dc0cf9d977207/grh-compact/src/New/Filling/GramLiouville.lean\#L49-L54}{{\leanfont Filling.\allowbreak{}sublatticeBasis}}}{GramLiouville.lean}{49--54}
\crow{sublatticeDisc}{\href{https://github.com/lsernau/grh-lean/blob/72fead102f663d5c7a407964a13dc0cf9d977207/grh-compact/src/New/Filling/GramLiouville.lean\#L58-L59}{{\leanfont Filling.\allowbreak{}sublatticeDisc}}}{GramLiouville.lean}{58--59}
\crow{sublatticeDiscmulindexsqle}{\href{https://github.com/lsernau/grh-lean/blob/72fead102f663d5c7a407964a13dc0cf9d977207/grh-compact/src/New/Filling/GramLiouville.lean\#L64-L79}{{\leanfont Filling.\allowbreak{}sublatticeDisc\_mul\_index\_sq\_le}}}{GramLiouville.lean}{64--79}
\crow{continuousbesselFn}{\href{https://github.com/lsernau/grh-lean/blob/72fead102f663d5c7a407964a13dc0cf9d977207/grh-compact/src/New/Filling/LatticeChart.lean\#L14-L14}{{\leanfont Filling.\allowbreak{}continuous\_besselFn}}}{LatticeChart.lean}{14--14}
\crow{continuousbesselKernel}{\href{https://github.com/lsernau/grh-lean/blob/72fead102f663d5c7a407964a13dc0cf9d977207/grh-compact/src/New/Filling/LatticeChart.lean\#L16-L17}{{\leanfont Filling.\allowbreak{}continuous\_besselKernel}}}{LatticeChart.lean}{16--17}
\crow{enormbesselKernel}{\href{https://github.com/lsernau/grh-lean/blob/72fead102f663d5c7a407964a13dc0cf9d977207/grh-compact/src/New/Filling/LatticeChart.lean\#L19-L22}{{\leanfont Filling.\allowbreak{}enorm\_besselKernel}}}{LatticeChart.lean}{19--22}
\crow{enormamplitudele}{\href{https://github.com/lsernau/grh-lean/blob/72fead102f663d5c7a407964a13dc0cf9d977207/grh-compact/src/New/Filling/LatticeChart.lean\#L24-L25}{{\leanfont Filling.\allowbreak{}enorm\_amplitude\_le}}}{LatticeChart.lean}{24--25}
\crow{isIntMatrixmiddleBaseChangemul}{\href{https://github.com/lsernau/grh-lean/blob/72fead102f663d5c7a407964a13dc0cf9d977207/grh-compact/src/New/Filling/LatticeChart.lean\#L34-L38}{{\leanfont Filling.\allowbreak{}isIntMatrix\_middleBaseChange\_mul}}}{LatticeChart.lean}{34--38}
\crow{forallinclHFinsuboneiff}{\href{https://github.com/lsernau/grh-lean/blob/72fead102f663d5c7a407964a13dc0cf9d977207/grh-compact/src/New/Filling/LatticeChart.lean\#L40-L43}{{\leanfont Filling.\allowbreak{}forall\_inclHFin\_sub\_one\_iff}}}{LatticeChart.lean}{40--43}
\crow{isIntMatrixscaleAtsubone}{\href{https://github.com/lsernau/grh-lean/blob/72fead102f663d5c7a407964a13dc0cf9d977207/grh-compact/src/New/Filling/LatticeChart.lean\#L49-L55}{{\leanfont Filling.\allowbreak{}isIntMatrix\_scaleAt\_sub\_one}}}{LatticeChart.lean}{49--55}
\crow{inclHFinmemlatStabilizer}{\href{https://github.com/lsernau/grh-lean/blob/72fead102f663d5c7a407964a13dc0cf9d977207/grh-compact/src/New/Filling/LatticeChart.lean\#L57-L65}{{\leanfont Filling.\allowbreak{}inclHFin\_mem\_latStabilizer}}}{LatticeChart.lean}{57--65}
\crow{reductioninclHFineqoneiff}{\href{https://github.com/lsernau/grh-lean/blob/72fead102f663d5c7a407964a13dc0cf9d977207/grh-compact/src/New/Filling/LatticeChart.lean\#L67-L72}{{\leanfont Filling.\allowbreak{}reduction\_inclHFin\_eq\_one\_iff}}}{LatticeChart.lean}{67--72}
\crow{classSumeqsum}{\href{https://github.com/lsernau/grh-lean/blob/72fead102f663d5c7a407964a13dc0cf9d977207/grh-compact/src/New/Filling/LatticeChart.lean\#L74-L75}{{\leanfont Filling.\allowbreak{}classSum\_eq\_sum}}}{LatticeChart.lean}{74--75}
\crow{cardresFieldeqlevelPrime}{\href{https://github.com/lsernau/grh-lean/blob/72fead102f663d5c7a407964a13dc0cf9d977207/grh-compact/src/New/Filling/LatticeChart.lean\#L77-L78}{{\leanfont Filling.\allowbreak{}card\_resField\_eq\_levelPrime}}}{LatticeChart.lean}{77--78}
\crow{mapsplitGramlatMiddle}{\href{https://github.com/lsernau/grh-lean/blob/72fead102f663d5c7a407964a13dc0cf9d977207/grh-compact/src/New/Filling/LevelBridge.lean\#L29-L30}{{\leanfont Filling.\allowbreak{}map\_splitGram\_latMiddle}}}{LevelBridge.lean}{29--30}
\crow{mapmemunitriangularOn}{\href{https://github.com/lsernau/grh-lean/blob/72fead102f663d5c7a407964a13dc0cf9d977207/grh-compact/src/New/Filling/LevelBridge.lean\#L32-L33}{{\leanfont Membership.\allowbreak{}mem.\allowbreak{}map\_mem\_unitriangularOn}}}{LevelBridge.lean}{32--33}
\crow{toLatBasismembesselU}{\href{https://github.com/lsernau/grh-lean/blob/72fead102f663d5c7a407964a13dc0cf9d977207/grh-compact/src/New/Filling/LevelBridge.lean\#L39-L41}{{\leanfont Filling.\allowbreak{}toLatBasis\_mem\_besselU}}}{LevelBridge.lean}{39--41}
\crow{ratToFinmembesselU}{\href{https://github.com/lsernau/grh-lean/blob/72fead102f663d5c7a407964a13dc0cf9d977207/grh-compact/src/New/Filling/LevelBridge.lean\#L43-L44}{{\leanfont Filling.\allowbreak{}ratToFin\_mem\_besselU}}}{LevelBridge.lean}{43--44}
\crow{inclHFinmembesselH}{\href{https://github.com/lsernau/grh-lean/blob/72fead102f663d5c7a407964a13dc0cf9d977207/grh-compact/src/New/Filling/LevelBridge.lean\#L46-L46}{{\leanfont Filling.\allowbreak{}inclHFin\_mem\_besselH}}}{LevelBridge.lean}{46--46}
\crow{rangeinclHFin}{\href{https://github.com/lsernau/grh-lean/blob/72fead102f663d5c7a407964a13dc0cf9d977207/grh-compact/src/New/Filling/LevelBridge.lean\#L48-L50}{{\leanfont Filling.\allowbreak{}range\_inclHFin}}}{LevelBridge.lean}{48--50}
\crow{ratToFintoBesselH}{\href{https://github.com/lsernau/grh-lean/blob/72fead102f663d5c7a407964a13dc0cf9d977207/grh-compact/src/New/Filling/LevelBridge.lean\#L52-L53}{{\leanfont Filling.\allowbreak{}ratToFin\_toBesselH}}}{LevelBridge.lean}{52--53}
\crow{classBasis}{\href{https://github.com/lsernau/grh-lean/blob/72fead102f663d5c7a407964a13dc0cf9d977207/grh-compact/src/New/Filling/LevelBridge.lean\#L57-L57}{{\leanfont Filling.\allowbreak{}classBasis}}}{LevelBridge.lean}{57--57}
\crow{coetermPosition}{\href{https://github.com/lsernau/grh-lean/blob/72fead102f663d5c7a407964a13dc0cf9d977207/grh-compact/src/New/Filling/LevelBridge.lean\#L59-L60}{{\leanfont Filling.\allowbreak{}coe\_termPosition}}}{LevelBridge.lean}{59--60}
\crow{termPositionmembesselU}{\href{https://github.com/lsernau/grh-lean/blob/72fead102f663d5c7a407964a13dc0cf9d977207/grh-compact/src/New/Filling/LevelBridge.lean\#L62-L64}{{\leanfont Membership.\allowbreak{}mem.\allowbreak{}termPosition\_mem\_besselU}}}{LevelBridge.lean}{62--64}
\crow{mul}{\href{https://github.com/lsernau/grh-lean/blob/72fead102f663d5c7a407964a13dc0cf9d977207/grh-compact/src/New/Filling/LevelBridge.lean\#L66-L67}{{\leanfont Filling.\allowbreak{}termsOver.\allowbreak{}mul}}}{LevelBridge.lean}{66--67}
\crow{inva}{\href{https://github.com/lsernau/grh-lean/blob/72fead102f663d5c7a407964a13dc0cf9d977207/grh-compact/src/New/Filling/LevelBridge.lean\#L69-L69}{{\leanfont Filling.\allowbreak{}termsOver.\allowbreak{}inv}}}{LevelBridge.lean}{69--69}
\crow{simpleRootSumadelicToFin}{\href{https://github.com/lsernau/grh-lean/blob/72fead102f663d5c7a407964a13dc0cf9d977207/grh-compact/src/New/Filling/LevelBridge.lean\#L71-L75}{{\leanfont Filling.\allowbreak{}simpleRootSum\_adelicToFin}}}{LevelBridge.lean}{71--75}
\crow{simpleRootSumratToAdelic}{\href{https://github.com/lsernau/grh-lean/blob/72fead102f663d5c7a407964a13dc0cf9d977207/grh-compact/src/New/Filling/LevelBridge.lean\#L77-L78}{{\leanfont Filling.\allowbreak{}simpleRootSum\_ratToAdelic}}}{LevelBridge.lean}{77--78}
\crow{normselfDualChar}{\href{https://github.com/lsernau/grh-lean/blob/72fead102f663d5c7a407964a13dc0cf9d977207/grh-compact/src/New/Filling/LevelBridge.lean\#L80-L80}{{\leanfont Filling.\allowbreak{}norm\_selfDualChar}}}{LevelBridge.lean}{80--80}
\crow{ratToInfmemunitriangularOn}{\href{https://github.com/lsernau/grh-lean/blob/72fead102f663d5c7a407964a13dc0cf9d977207/grh-compact/src/New/Filling/LevelBridge.lean\#L82-L83}{{\leanfont Membership.\allowbreak{}mem.\allowbreak{}ratToInf\_mem\_unitriangularOn}}}{LevelBridge.lean}{82--83}
\crow{besselFnadelicToInfmul}{\href{https://github.com/lsernau/grh-lean/blob/72fead102f663d5c7a407964a13dc0cf9d977207/grh-compact/src/New/Filling/LevelBridge.lean\#L88-L96}{{\leanfont Filling.\allowbreak{}besselFn\_adelicToInf\_mul}}}{LevelBridge.lean}{88--96}
\crow{besselFnratToInfmul}{\href{https://github.com/lsernau/grh-lean/blob/72fead102f663d5c7a407964a13dc0cf9d977207/grh-compact/src/New/Filling/LevelBridge.lean\#L98-L100}{{\leanfont Filling.\allowbreak{}besselFn\_ratToInf\_mul}}}{LevelBridge.lean}{98--100}
\crow{isUnimodularlatMiddle}{\href{https://github.com/lsernau/grh-lean/blob/72fead102f663d5c7a407964a13dc0cf9d977207/grh-compact/src/New/Filling/LevelBridge.lean\#L108-L109}{{\leanfont Filling.\allowbreak{}isUnimodular\_latMiddle}}}{LevelBridge.lean}{108--109}
\crow{transposesplitGram}{\href{https://github.com/lsernau/grh-lean/blob/72fead102f663d5c7a407964a13dc0cf9d977207/grh-compact/src/New/Filling/LevelBridge.lean\#L114-L115}{{\leanfont Filling.\allowbreak{}transpose\_splitGram}}}{LevelBridge.lean}{114--115}
\crow{idOrInvAdjoint}{\href{https://github.com/lsernau/grh-lean/blob/72fead102f663d5c7a407964a13dc0cf9d977207/grh-compact/src/New/Filling/LevelBridge.lean\#L119-L120}{{\leanfont Filling.\allowbreak{}idOrInvAdjoint}}}{LevelBridge.lean}{119--120}
\crow{reduceAboveG}{\href{https://github.com/lsernau/grh-lean/blob/72fead102f663d5c7a407964a13dc0cf9d977207/grh-compact/src/New/Filling/LevelBridge.lean\#L125-L126}{{\leanfont Filling.\allowbreak{}reduceAboveG}}}{LevelBridge.lean}{125--126}
\crow{intGLmapreduceAboveG}{\href{https://github.com/lsernau/grh-lean/blob/72fead102f663d5c7a407964a13dc0cf9d977207/grh-compact/src/New/Filling/LevelBridge.lean\#L130-L136}{{\leanfont Filling.\allowbreak{}intGL\_map\_reduceAboveG}}}{LevelBridge.lean}{130--136}
\crow{gIdeleval}{\href{https://github.com/lsernau/grh-lean/blob/72fead102f663d5c7a407964a13dc0cf9d977207/grh-compact/src/New/Filling/LevelBridge.lean\#L138-L138}{{\leanfont Filling.\allowbreak{}gIdele\_val}}}{LevelBridge.lean}{138--138}
\crow{gIdelemem}{\href{https://github.com/lsernau/grh-lean/blob/72fead102f663d5c7a407964a13dc0cf9d977207/grh-compact/src/New/Filling/LevelBridge.lean\#L140-L140}{{\leanfont Filling.\allowbreak{}gIdele\_mem}}}{LevelBridge.lean}{140--140}
\crow{forallselfDualCharmuleqoneiff}{\href{https://github.com/lsernau/grh-lean/blob/72fead102f663d5c7a407964a13dc0cf9d977207/grh-compact/src/New/Filling/LevelBridge.lean\#L142-L145}{{\leanfont Filling.\allowbreak{}forall\_selfDualChar\_mul\_eq\_one\_iff}}}{LevelBridge.lean}{142--145}
\crow{selfDualChareqoneofmem}{\href{https://github.com/lsernau/grh-lean/blob/72fead102f663d5c7a407964a13dc0cf9d977207/grh-compact/src/New/Filling/LevelBridge.lean\#L147-L148}{{\leanfont Filling.\allowbreak{}selfDualChar\_eq\_one\_of\_mem}}}{LevelBridge.lean}{147--148}
\crow{invgIdelesmulmemiff}{\href{https://github.com/lsernau/grh-lean/blob/72fead102f663d5c7a407964a13dc0cf9d977207/grh-compact/src/New/Filling/LevelBridge.lean\#L150-L152}{{\leanfont Filling.\allowbreak{}inv\_gIdele\_smul\_mem\_iff}}}{LevelBridge.lean}{150--152}
\crow{principalLevelequnitaryLevel}{\href{https://github.com/lsernau/grh-lean/blob/72fead102f663d5c7a407964a13dc0cf9d977207/grh-compact/src/New/Filling/LevelBridge.lean\#L154-L157}{{\leanfont Filling.\allowbreak{}principalLevel\_eq\_unitaryLevel}}}{LevelBridge.lean}{154--157}
\crow{siegellemma}{\href{https://github.com/lsernau/grh-lean/blob/72fead102f663d5c7a407964a13dc0cf9d977207/grh-compact/src/New/Filling/NearTerms.lean\#L27-L60}{{\leanfont Filling.\allowbreak{}siegel\_lemma}}}{NearTerms.lean}{27--60}
\crow{tsumsublatticeDiscinvle}{\href{https://github.com/lsernau/grh-lean/blob/72fead102f663d5c7a407964a13dc0cf9d977207/grh-compact/src/New/Filling/NearTerms.lean\#L65-L83}{{\leanfont Filling.\allowbreak{}tsum\_sublatticeDisc\_inv\_le}}}{NearTerms.lean}{65--83}
\crow{besselUdense}{\href{https://github.com/lsernau/grh-lean/blob/72fead102f663d5c7a407964a13dc0cf9d977207/grh-compact/src/New/Filling/NearTerms.lean\#L94-L110}{{\leanfont besselU\_dense}}}{NearTerms.lean}{94--110}
\crow{unipotentCharone}{\href{https://github.com/lsernau/grh-lean/blob/72fead102f663d5c7a407964a13dc0cf9d977207/grh-compact/src/New/Filling/NearTerms.lean\#L120-L121}{{\leanfont Filling.\allowbreak{}unipotentChar\_one}}}{NearTerms.lean}{120--121}
\crow{existsbesselUapprox}{\href{https://github.com/lsernau/grh-lean/blob/72fead102f663d5c7a407964a13dc0cf9d977207/grh-compact/src/New/Filling/NearTerms.lean\#L125-L131}{{\leanfont Filling.\allowbreak{}exists\_besselU\_approx}}}{NearTerms.lean}{125--131}
\crow{unipotentChareqofmemprincipalLevel}{\href{https://github.com/lsernau/grh-lean/blob/72fead102f663d5c7a407964a13dc0cf9d977207/grh-compact/src/New/Filling/NearTerms.lean\#L135-L141}{{\leanfont Filling.\allowbreak{}unipotentChar\_eq\_of\_mem\_principalLevel}}}{NearTerms.lean}{135--141}
\crow{eqandmemclassUnipofbesselS}{\href{https://github.com/lsernau/grh-lean/blob/72fead102f663d5c7a407964a13dc0cf9d977207/grh-compact/src/New/Filling/NearTerms.lean\#L145-L150}{{\leanfont Filling.\allowbreak{}eq\_and\_mem\_classUnip\_of\_besselS}}}{NearTerms.lean}{145--150}
\crow{existsratapproxunipotentChareq}{\href{https://github.com/lsernau/grh-lean/blob/72fead102f663d5c7a407964a13dc0cf9d977207/grh-compact/src/New/Filling/NearTerms.lean\#L152-L156}{{\leanfont Filling.\allowbreak{}exists\_rat\_approx\_unipotentChar\_eq}}}{NearTerms.lean}{152--156}
\crow{existsbesselSmemtermsOver}{\href{https://github.com/lsernau/grh-lean/blob/72fead102f663d5c7a407964a13dc0cf9d977207/grh-compact/src/New/Filling/NearTerms.lean\#L158-L165}{{\leanfont Filling.\allowbreak{}exists\_besselS\_mem\_termsOver}}}{NearTerms.lean}{158--165}
\crow{eqoneofreduceAteqone}{\href{https://github.com/lsernau/grh-lean/blob/72fead102f663d5c7a407964a13dc0cf9d977207/grh-compact/src/New/Filling/NearTerms.lean\#L170-L176}{{\leanfont Filling.\allowbreak{}eq\_one\_of\_reduceAt\_eq\_one}}}{NearTerms.lean}{170--176}
\crow{frameImagemul}{\href{https://github.com/lsernau/grh-lean/blob/72fead102f663d5c7a407964a13dc0cf9d977207/grh-compact/src/New/Filling/NearTerms.lean\#L187-L187}{{\leanfont Filling.\allowbreak{}frameImage\_mul}}}{NearTerms.lean}{187--187}
\crow{frameImagemulofmembesselU}{\href{https://github.com/lsernau/grh-lean/blob/72fead102f663d5c7a407964a13dc0cf9d977207/grh-compact/src/New/Filling/NearTerms.lean\#L189-L190}{{\leanfont Filling.\allowbreak{}frameImage\_mul\_of\_mem\_besselU0}}}{NearTerms.lean}{189--190}
\crow{frameImageinvmuleqiff}{\href{https://github.com/lsernau/grh-lean/blob/72fead102f663d5c7a407964a13dc0cf9d977207/grh-compact/src/New/Filling/NearTerms.lean\#L192-L193}{{\leanfont Filling.\allowbreak{}frameImage\_inv\_mul\_eq\_iff}}}{NearTerms.lean}{192--193}
\crow{starmullasteqofframeImageeq}{\href{https://github.com/lsernau/grh-lean/blob/72fead102f663d5c7a407964a13dc0cf9d977207/grh-compact/src/New/Filling/NearTerms.lean\#L201-L203}{{\leanfont Filling.\allowbreak{}star\_mul\_last\_eq\_of\_frameImage\_eq}}}{NearTerms.lean}{201--203}
\crow{frameProductsmulVecofframeImageeq}{\href{https://github.com/lsernau/grh-lean/blob/72fead102f663d5c7a407964a13dc0cf9d977207/grh-compact/src/New/Filling/NearTerms.lean\#L205-L207}{{\leanfont Filling.\allowbreak{}frameProducts\_mulVec\_of\_frameImage\_eq}}}{NearTerms.lean}{205--207}
\crow{eqzeroofratMiddleformeqzero}{\href{https://github.com/lsernau/grh-lean/blob/72fead102f663d5c7a407964a13dc0cf9d977207/grh-compact/src/New/Filling/NearTerms.lean\#L209-L210}{{\leanfont Filling.\allowbreak{}eq\_zero\_of\_ratMiddle\_form\_eq\_zero}}}{NearTerms.lean}{209--210}
\crow{normsqmiddleparteq}{\href{https://github.com/lsernau/grh-lean/blob/72fead102f663d5c7a407964a13dc0cf9d977207/grh-compact/src/New/Filling/NearTerms.lean\#L215-L223}{{\leanfont Filling.\allowbreak{}norm\_sq\_middle\_part\_eq}}}{NearTerms.lean}{215--223}
\crow{maptermDisceqnorm}{\href{https://github.com/lsernau/grh-lean/blob/72fead102f663d5c7a407964a13dc0cf9d977207/grh-compact/src/New/Filling/NearTerms.lean\#L225-L228}{{\leanfont Filling.\allowbreak{}map\_termDisc\_eq\_norm}}}{NearTerms.lean}{225--228}
\crow{hasLevellast}{\href{https://github.com/lsernau/grh-lean/blob/72fead102f663d5c7a407964a13dc0cf9d977207/grh-compact/src/New/Filling/NearTerms.lean\#L230-L231}{{\leanfont Filling.\allowbreak{}hasLevel\_last}}}{NearTerms.lean}{230--231}
\crow{gramDiscRealrealFrameGramratToInf}{\href{https://github.com/lsernau/grh-lean/blob/72fead102f663d5c7a407964a13dc0cf9d977207/grh-compact/src/New/Filling/NearTerms.lean\#L233-L236}{{\leanfont Filling.\allowbreak{}gramDiscReal\_realFrameGram\_ratToInf}}}{NearTerms.lean}{233--236}
\crow{ratToInfmembesselU}{\href{https://github.com/lsernau/grh-lean/blob/72fead102f663d5c7a407964a13dc0cf9d977207/grh-compact/src/New/Filling/NearTerms.lean\#L238-L239}{{\leanfont Filling.\allowbreak{}ratToInf\_mem\_besselU0}}}{NearTerms.lean}{238--239}
\crow{mulmemboxedPointsiff}{\href{https://github.com/lsernau/grh-lean/blob/72fead102f663d5c7a407964a13dc0cf9d977207/grh-compact/src/New/Filling/NearTerms.lean\#L241-L243}{{\leanfont Filling.\allowbreak{}mul\_mem\_boxedPoints\_iff}}}{NearTerms.lean}{241--243}
\crow{ratToInf}{\href{https://github.com/lsernau/grh-lean/blob/72fead102f663d5c7a407964a13dc0cf9d977207/grh-compact/src/New/Filling/NearTerms.lean\#L245-L250}{{\leanfont IsNormalForm.\allowbreak{}ratToInf}}}{NearTerms.lean}{245--250}
\crow{isIntMatrixtermPosition}{\href{https://github.com/lsernau/grh-lean/blob/72fead102f663d5c7a407964a13dc0cf9d977207/grh-compact/src/New/Filling/NearTerms.lean\#L252-L253}{{\leanfont Filling.\allowbreak{}termsOver.\allowbreak{}isIntMatrix\_termPosition}}}{NearTerms.lean}{252--253}
\crow{classCoordconjsplitGram}{\href{https://github.com/lsernau/grh-lean/blob/72fead102f663d5c7a407964a13dc0cf9d977207/grh-compact/src/New/Filling/NearTerms.lean\#L255-L257}{{\leanfont Filling.\allowbreak{}classCoord\_conj\_splitGram}}}{NearTerms.lean}{255--257}
\crow{frameImagetermPosition}{\href{https://github.com/lsernau/grh-lean/blob/72fead102f663d5c7a407964a13dc0cf9d977207/grh-compact/src/New/Filling/NearTerms.lean\#L259-L260}{{\leanfont Filling.\allowbreak{}frameImage\_termPosition}}}{NearTerms.lean}{259--260}
\crow{frameGramtermPosition}{\href{https://github.com/lsernau/grh-lean/blob/72fead102f663d5c7a407964a13dc0cf9d977207/grh-compact/src/New/Filling/NearTerms.lean\#L262-L266}{{\leanfont Filling.\allowbreak{}frameGram\_termPosition}}}{NearTerms.lean}{262--266}
\crow{mulofmembesselU}{\href{https://github.com/lsernau/grh-lean/blob/72fead102f663d5c7a407964a13dc0cf9d977207/grh-compact/src/New/Filling/NearTerms.lean\#L268-L270}{{\leanfont Filling.\allowbreak{}IsNear.\allowbreak{}mul\_of\_mem\_besselU0}}}{NearTerms.lean}{268--270}
\crow{membesselSiffframeImage}{\href{https://github.com/lsernau/grh-lean/blob/72fead102f663d5c7a407964a13dc0cf9d977207/grh-compact/src/New/Filling/NearTerms.lean\#L278-L279}{{\leanfont Filling.\allowbreak{}mem\_besselS\_iff\_frameImage}}}{NearTerms.lean}{278--279}
\crow{powletailScale}{\href{https://github.com/lsernau/grh-lean/blob/72fead102f663d5c7a407964a13dc0cf9d977207/grh-compact/src/New/Filling/NearTerms.lean\#L281-L283}{{\leanfont Filling.\allowbreak{}pow\_le\_tailScale}}}{NearTerms.lean}{281--283}
\crow{tailScalepowfour}{\href{https://github.com/lsernau/grh-lean/blob/72fead102f663d5c7a407964a13dc0cf9d977207/grh-compact/src/New/Filling/NearTerms.lean\#L285-L286}{{\leanfont Filling.\allowbreak{}tailScale\_pow\_four}}}{NearTerms.lean}{285--286}
\crow{ofsizebounds}{\href{https://github.com/lsernau/grh-lean/blob/72fead102f663d5c7a407964a13dc0cf9d977207/grh-compact/src/New/Filling/NearTerms.lean\#L288-L293}{{\leanfont Filling.\allowbreak{}of\_size\_bounds}}}{NearTerms.lean}{288--293}
\crow{normframeProductsle}{\href{https://github.com/lsernau/grh-lean/blob/72fead102f663d5c7a407964a13dc0cf9d977207/grh-compact/src/New/Filling/NearTerms.lean\#L295-L300}{{\leanfont Filling.\allowbreak{}IsNear.\allowbreak{}norm\_frameProducts\_le}}}{NearTerms.lean}{295--300}
\crow{normmapleadingMinorle}{\href{https://github.com/lsernau/grh-lean/blob/72fead102f663d5c7a407964a13dc0cf9d977207/grh-compact/src/New/Filling/NearTerms.lean\#L302-L305}{{\leanfont Filling.\allowbreak{}IsNear.\allowbreak{}norm\_map\_leadingMinor\_le}}}{NearTerms.lean}{302--305}
\crow{eqandmemclassUnip}{\href{https://github.com/lsernau/grh-lean/blob/72fead102f663d5c7a407964a13dc0cf9d977207/grh-compact/src/New/Filling/NearTerms.lean\#L309-L313}{{\leanfont Filling.\allowbreak{}termsOver.\allowbreak{}eq\_and\_mem\_classUnip₀}}}{NearTerms.lean}{309--313}
\crow{isIntegralframeProducts}{\href{https://github.com/lsernau/grh-lean/blob/72fead102f663d5c7a407964a13dc0cf9d977207/grh-compact/src/New/Filling/NearTerms.lean\#L315-L317}{{\leanfont Filling.\allowbreak{}isIntegral\_frameProducts}}}{NearTerms.lean}{315--317}
\crow{Filling}{\href{https://github.com/lsernau/grh-lean/blob/72fead102f663d5c7a407964a13dc0cf9d977207/grh-compact/src/New/Filling/Setup.lean\#L22-L58}{{\leanfont Filling}}}{Setup.lean}{22--58}
\crow{g}{\href{https://github.com/lsernau/grh-lean/blob/72fead102f663d5c7a407964a13dc0cf9d977207/grh-compact/src/New/Filling/Setup.lean\#L70-L70}{{\leanfont Filling.\allowbreak{}g}}}{Setup.lean}{70--70}
\crow{ratMiddle}{\href{https://github.com/lsernau/grh-lean/blob/72fead102f663d5c7a407964a13dc0cf9d977207/grh-compact/src/New/Filling/Setup.lean\#L74-L74}{{\leanfont Filling.\allowbreak{}ratMiddle}}}{Setup.lean}{74--74}
\crow{latMiddle}{\href{https://github.com/lsernau/grh-lean/blob/72fead102f663d5c7a407964a13dc0cf9d977207/grh-compact/src/New/Filling/Setup.lean\#L78-L78}{{\leanfont Filling.\allowbreak{}latMiddle}}}{Setup.lean}{78--78}
\crow{G}{\href{https://github.com/lsernau/grh-lean/blob/72fead102f663d5c7a407964a13dc0cf9d977207/grh-compact/src/New/Filling/Setup.lean\#L82-L82}{{\leanfont Filling.\allowbreak{}G}}}{Setup.lean}{82--82}
\crow{GFin}{\href{https://github.com/lsernau/grh-lean/blob/72fead102f663d5c7a407964a13dc0cf9d977207/grh-compact/src/New/Filling/Setup.lean\#L86-L86}{{\leanfont Filling.\allowbreak{}GFin}}}{Setup.lean}{86--86}
\crow{GRes}{\href{https://github.com/lsernau/grh-lean/blob/72fead102f663d5c7a407964a13dc0cf9d977207/grh-compact/src/New/Filling/Setup.lean\#L91-L91}{{\leanfont Filling.\allowbreak{}GRes}}}{Setup.lean}{91--91}
\crow{BesselCoord}{\href{https://github.com/lsernau/grh-lean/blob/72fead102f663d5c7a407964a13dc0cf9d977207/grh-compact/src/New/Filling/Setup.lean\#L95-L95}{{\leanfont Filling.\allowbreak{}BesselCoord}}}{Setup.lean}{95--95}
\crow{Classes}{\href{https://github.com/lsernau/grh-lean/blob/72fead102f663d5c7a407964a13dc0cf9d977207/grh-compact/src/New/Filling/Setup.lean\#L99-L99}{{\leanfont Filling.\allowbreak{}Classes}}}{Setup.lean}{99--99}
\crow{frameProducts}{\href{https://github.com/lsernau/grh-lean/blob/72fead102f663d5c7a407964a13dc0cf9d977207/grh-compact/src/New/Filling/Setup.lean\#L107-L107}{{\leanfont Filling.\allowbreak{}frameProducts}}}{Setup.lean}{107--107}
\crow{latStabilizer}{\href{https://github.com/lsernau/grh-lean/blob/72fead102f663d5c7a407964a13dc0cf9d977207/grh-compact/src/New/Filling/Setup.lean\#L111-L111}{{\leanfont Filling.\allowbreak{}latStabilizer}}}{Setup.lean}{111--111}
\crow{reduction}{\href{https://github.com/lsernau/grh-lean/blob/72fead102f663d5c7a407964a13dc0cf9d977207/grh-compact/src/New/Filling/Setup.lean\#L115-L116}{{\leanfont Filling.\allowbreak{}reduction}}}{Setup.lean}{115--116}
\crow{principalLevel}{\href{https://github.com/lsernau/grh-lean/blob/72fead102f663d5c7a407964a13dc0cf9d977207/grh-compact/src/New/Filling/Setup.lean\#L120-L120}{{\leanfont Filling.\allowbreak{}principalLevel}}}{Setup.lean}{120--120}
\crow{reducedImage}{\href{https://github.com/lsernau/grh-lean/blob/72fead102f663d5c7a407964a13dc0cf9d977207/grh-compact/src/New/Filling/Setup.lean\#L124-L124}{{\leanfont Filling.\allowbreak{}reducedImage}}}{Setup.lean}{124--124}
\crow{kappa}{\href{https://github.com/lsernau/grh-lean/blob/72fead102f663d5c7a407964a13dc0cf9d977207/grh-compact/src/New/Filling/Setup.lean\#L128-L128}{{\leanfont Filling.\allowbreak{}kappa}}}{Setup.lean}{128--128}
\crow{basisScale}{\href{https://github.com/lsernau/grh-lean/blob/72fead102f663d5c7a407964a13dc0cf9d977207/grh-compact/src/New/Filling/Setup.lean\#L132-L133}{{\leanfont Filling.\allowbreak{}basisScale}}}{Setup.lean}{132--133}
\crow{mapsplitGramratMiddle}{\href{https://github.com/lsernau/grh-lean/blob/72fead102f663d5c7a407964a13dc0cf9d977207/grh-compact/src/New/Filling/Setup.lean\#L138-L139}{{\leanfont Filling.\allowbreak{}map\_splitGram\_ratMiddle}}}{Setup.lean}{138--139}
\crow{selfDualChar}{\href{https://github.com/lsernau/grh-lean/blob/72fead102f663d5c7a407964a13dc0cf9d977207/grh-compact/src/New/Filling/Setup.lean\#L142-L142}{{\leanfont Filling.\allowbreak{}selfDualChar}}}{Setup.lean}{142--142}
\crow{varpinezero}{\href{https://github.com/lsernau/grh-lean/blob/72fead102f663d5c7a407964a13dc0cf9d977207/grh-compact/src/New/Filling/Setup.lean\#L144-L145}{{\leanfont Filling.\allowbreak{}varpi\_ne\_zero}}}{Setup.lean}{144--145}
\crow{frameScaleSq}{\href{https://github.com/lsernau/grh-lean/blob/72fead102f663d5c7a407964a13dc0cf9d977207/grh-compact/src/New/Filling/Setup.lean\#L148-L148}{{\leanfont Filling.\allowbreak{}frameScaleSq}}}{Setup.lean}{148--148}
\crow{frameScaleProd}{\href{https://github.com/lsernau/grh-lean/blob/72fead102f663d5c7a407964a13dc0cf9d977207/grh-compact/src/New/Filling/Setup.lean\#L152-L152}{{\leanfont Filling.\allowbreak{}frameScaleProd}}}{Setup.lean}{152--152}
\crow{unipCovolume}{\href{https://github.com/lsernau/grh-lean/blob/72fead102f663d5c7a407964a13dc0cf9d977207/grh-compact/src/New/Filling/Setup.lean\#L156-L157}{{\leanfont Filling.\allowbreak{}unipCovolume}}}{Setup.lean}{156--157}
\crow{siegelConst}{\href{https://github.com/lsernau/grh-lean/blob/72fead102f663d5c7a407964a13dc0cf9d977207/grh-compact/src/New/Filling/Setup.lean\#L161-L162}{{\leanfont Filling.\allowbreak{}siegelConst}}}{Setup.lean}{161--162}
\crow{termDisc}{\href{https://github.com/lsernau/grh-lean/blob/72fead102f663d5c7a407964a13dc0cf9d977207/grh-compact/src/New/Filling/Setup.lean\#L167-L167}{{\leanfont Filling.\allowbreak{}termDisc}}}{Setup.lean}{167--167}
\crow{basisScalebesselFrame}{\href{https://github.com/lsernau/grh-lean/blob/72fead102f663d5c7a407964a13dc0cf9d977207/grh-compact/src/New/Filling/Setup.lean\#L169-L170}{{\leanfont Filling.\allowbreak{}basisScale\_besselFrame}}}{Setup.lean}{169--170}
\crow{coekappa}{\href{https://github.com/lsernau/grh-lean/blob/72fead102f663d5c7a407964a13dc0cf9d977207/grh-compact/src/New/Filling/Setup.lean\#L172-L173}{{\leanfont Filling.\allowbreak{}coe\_kappa}}}{Setup.lean}{172--173}
\crow{kappanezero}{\href{https://github.com/lsernau/grh-lean/blob/72fead102f663d5c7a407964a13dc0cf9d977207/grh-compact/src/New/Filling/Setup.lean\#L175-L175}{{\leanfont Filling.\allowbreak{}kappa\_ne\_zero}}}{Setup.lean}{175--175}
\crow{frameScaleSqpos}{\href{https://github.com/lsernau/grh-lean/blob/72fead102f663d5c7a407964a13dc0cf9d977207/grh-compact/src/New/Filling/Setup.lean\#L177-L177}{{\leanfont Filling.\allowbreak{}frameScaleSq\_pos}}}{Setup.lean}{177--177}
\crow{frameScaleProdpos}{\href{https://github.com/lsernau/grh-lean/blob/72fead102f663d5c7a407964a13dc0cf9d977207/grh-compact/src/New/Filling/Setup.lean\#L179-L179}{{\leanfont Filling.\allowbreak{}frameScaleProd\_pos}}}{Setup.lean}{179--179}
\crow{basisScalenezero}{\href{https://github.com/lsernau/grh-lean/blob/72fead102f663d5c7a407964a13dc0cf9d977207/grh-compact/src/New/Filling/Setup.lean\#L181-L181}{{\leanfont Filling.\allowbreak{}basisScale\_ne\_zero}}}{Setup.lean}{181--181}
\crow{basisScaleGL}{\href{https://github.com/lsernau/grh-lean/blob/72fead102f663d5c7a407964a13dc0cf9d977207/grh-compact/src/New/Filling/Setup.lean\#L185-L185}{{\leanfont Filling.\allowbreak{}basisScaleGL}}}{Setup.lean}{185--185}
\crow{grambasisScaleGL}{\href{https://github.com/lsernau/grh-lean/blob/72fead102f663d5c7a407964a13dc0cf9d977207/grh-compact/src/New/Filling/Setup.lean\#L187-L188}{{\leanfont Filling.\allowbreak{}gram\_basisScaleGL}}}{Setup.lean}{187--188}
\crow{latBasis}{\href{https://github.com/lsernau/grh-lean/blob/72fead102f663d5c7a407964a13dc0cf9d977207/grh-compact/src/New/Filling/Setup.lean\#L192-L192}{{\leanfont Filling.\allowbreak{}latBasis}}}{Setup.lean}{192--192}
\crow{toLatBasis}{\href{https://github.com/lsernau/grh-lean/blob/72fead102f663d5c7a407964a13dc0cf9d977207/grh-compact/src/New/Filling/Setup.lean\#L196-L197}{{\leanfont Filling.\allowbreak{}toLatBasis}}}{Setup.lean}{196--197}
\crow{ratToAdelic}{\href{https://github.com/lsernau/grh-lean/blob/72fead102f663d5c7a407964a13dc0cf9d977207/grh-compact/src/New/Filling/Setup.lean\#L201-L201}{{\leanfont Filling.\allowbreak{}ratToAdelic}}}{Setup.lean}{201--201}
\crow{adelicToInf}{\href{https://github.com/lsernau/grh-lean/blob/72fead102f663d5c7a407964a13dc0cf9d977207/grh-compact/src/New/Filling/Setup.lean\#L205-L206}{{\leanfont Filling.\allowbreak{}adelicToInf}}}{Setup.lean}{205--206}
\crow{adelicToFin}{\href{https://github.com/lsernau/grh-lean/blob/72fead102f663d5c7a407964a13dc0cf9d977207/grh-compact/src/New/Filling/Setup.lean\#L210-L211}{{\leanfont Filling.\allowbreak{}adelicToFin}}}{Setup.lean}{210--211}
\crow{ratToInfa}{\href{https://github.com/lsernau/grh-lean/blob/72fead102f663d5c7a407964a13dc0cf9d977207/grh-compact/src/New/Filling/Setup.lean\#L215-L215}{{\leanfont Filling.\allowbreak{}ratToInf}}}{Setup.lean}{215--215}
\crow{ratToFin}{\href{https://github.com/lsernau/grh-lean/blob/72fead102f663d5c7a407964a13dc0cf9d977207/grh-compact/src/New/Filling/Setup.lean\#L219-L219}{{\leanfont Filling.\allowbreak{}ratToFin}}}{Setup.lean}{219--219}
\crow{inclHFin}{\href{https://github.com/lsernau/grh-lean/blob/72fead102f663d5c7a407964a13dc0cf9d977207/grh-compact/src/New/Filling/Setup.lean\#L223-L223}{{\leanfont Filling.\allowbreak{}inclHFin}}}{Setup.lean}{223--223}
\crow{h}{\href{https://github.com/lsernau/grh-lean/blob/72fead102f663d5c7a407964a13dc0cf9d977207/grh-compact/src/New/Filling/Setup.lean\#L229-L229}{{\leanfont Filling.\allowbreak{}h}}}{Setup.lean}{229--229}
\crow{termPosition}{\href{https://github.com/lsernau/grh-lean/blob/72fead102f663d5c7a407964a13dc0cf9d977207/grh-compact/src/New/Filling/Setup.lean\#L233-L233}{{\leanfont Filling.\allowbreak{}termPosition}}}{Setup.lean}{233--233}
\crow{termsOver}{\href{https://github.com/lsernau/grh-lean/blob/72fead102f663d5c7a407964a13dc0cf9d977207/grh-compact/src/New/Filling/Setup.lean\#L237-L237}{{\leanfont Filling.\allowbreak{}termsOver}}}{Setup.lean}{237--237}
\crow{classUnip}{\href{https://github.com/lsernau/grh-lean/blob/72fead102f663d5c7a407964a13dc0cf9d977207/grh-compact/src/New/Filling/Setup.lean\#L240-L241}{{\leanfont Filling.\allowbreak{}classUnip}}}{Setup.lean}{240--241}
\crow{classUnipa}{\href{https://github.com/lsernau/grh-lean/blob/72fead102f663d5c7a407964a13dc0cf9d977207/grh-compact/src/New/Filling/Setup.lean\#L244-L244}{{\leanfont Filling.\allowbreak{}classUnip₀}}}{Setup.lean}{244--244}
\crow{boxSize}{\href{https://github.com/lsernau/grh-lean/blob/72fead102f663d5c7a407964a13dc0cf9d977207/grh-compact/src/New/Filling/Setup.lean\#L247-L247}{{\leanfont Filling.\allowbreak{}boxSize}}}{Setup.lean}{247--247}
\crow{IsBoxPartition}{\href{https://github.com/lsernau/grh-lean/blob/72fead102f663d5c7a407964a13dc0cf9d977207/grh-compact/src/New/Filling/Setup.lean\#L252-L255}{{\leanfont Filling.\allowbreak{}IsBoxPartition}}}{Setup.lean}{252--255}
\crow{besselFn}{\href{https://github.com/lsernau/grh-lean/blob/72fead102f663d5c7a407964a13dc0cf9d977207/grh-compact/src/New/Filling/Setup.lean\#L259-L259}{{\leanfont Filling.\allowbreak{}besselFn}}}{Setup.lean}{259--259}
\crow{besselKernel}{\href{https://github.com/lsernau/grh-lean/blob/72fead102f663d5c7a407964a13dc0cf9d977207/grh-compact/src/New/Filling/Setup.lean\#L263-L264}{{\leanfont Filling.\allowbreak{}besselKernel}}}{Setup.lean}{263--264}
\crow{columnDim}{\href{https://github.com/lsernau/grh-lean/blob/72fead102f663d5c7a407964a13dc0cf9d977207/grh-compact/src/New/Filling/Setup.lean\#L268-L268}{{\leanfont Filling.\allowbreak{}columnDim}}}{Setup.lean}{268--268}
\crow{besselDegree}{\href{https://github.com/lsernau/grh-lean/blob/72fead102f663d5c7a407964a13dc0cf9d977207/grh-compact/src/New/Filling/Setup.lean\#L271-L271}{{\leanfont Filling.\allowbreak{}besselDegree}}}{Setup.lean}{271--271}
\crow{besselDegreepos}{\href{https://github.com/lsernau/grh-lean/blob/72fead102f663d5c7a407964a13dc0cf9d977207/grh-compact/src/New/Filling/Setup.lean\#L273-L273}{{\leanfont Filling.\allowbreak{}besselDegree\_pos}}}{Setup.lean}{273--273}
\crow{mainEnergy}{\href{https://github.com/lsernau/grh-lean/blob/72fead102f663d5c7a407964a13dc0cf9d977207/grh-compact/src/New/Filling/Setup.lean\#L276-L276}{{\leanfont Filling.\allowbreak{}mainEnergy}}}{Setup.lean}{276--276}
\crow{unipOrder}{\href{https://github.com/lsernau/grh-lean/blob/72fead102f663d5c7a407964a13dc0cf9d977207/grh-compact/src/New/Filling/Setup.lean\#L280-L280}{{\leanfont Filling.\allowbreak{}unipOrder}}}{Setup.lean}{280--280}
\crow{packetConst}{\href{https://github.com/lsernau/grh-lean/blob/72fead102f663d5c7a407964a13dc0cf9d977207/grh-compact/src/New/Filling/Setup.lean\#L284-L284}{{\leanfont Filling.\allowbreak{}packetConst}}}{Setup.lean}{284--284}
\crow{besselEnergy}{\href{https://github.com/lsernau/grh-lean/blob/72fead102f663d5c7a407964a13dc0cf9d977207/grh-compact/src/New/Filling/Setup.lean\#L289-L289}{{\leanfont Filling.\allowbreak{}besselEnergy}}}{Setup.lean}{289--289}
\crow{tailScale}{\href{https://github.com/lsernau/grh-lean/blob/72fead102f663d5c7a407964a13dc0cf9d977207/grh-compact/src/New/Filling/Setup.lean\#L292-L292}{{\leanfont Filling.\allowbreak{}tailScale}}}{Setup.lean}{292--292}
\crow{codimFactor}{\href{https://github.com/lsernau/grh-lean/blob/72fead102f663d5c7a407964a13dc0cf9d977207/grh-compact/src/New/Filling/Setup.lean\#L296-L296}{{\leanfont Filling.\allowbreak{}codimFactor}}}{Setup.lean}{296--296}
\crow{IsAdmissible}{\href{https://github.com/lsernau/grh-lean/blob/72fead102f663d5c7a407964a13dc0cf9d977207/grh-compact/src/New/Filling/Setup.lean\#L301-L318}{{\leanfont Filling.\allowbreak{}IsAdmissible}}}{Setup.lean}{301--318}
\crow{adelicToInfapply}{\href{https://github.com/lsernau/grh-lean/blob/72fead102f663d5c7a407964a13dc0cf9d977207/grh-compact/src/New/Filling/Setup.lean\#L320-L321}{{\leanfont Filling.\allowbreak{}adelicToInf\_apply}}}{Setup.lean}{320--321}
\crow{ratToInfapply}{\href{https://github.com/lsernau/grh-lean/blob/72fead102f663d5c7a407964a13dc0cf9d977207/grh-compact/src/New/Filling/Setup.lean\#L323-L324}{{\leanfont Filling.\allowbreak{}ratToInf\_apply}}}{Setup.lean}{323--324}
\crow{coeratToFin}{\href{https://github.com/lsernau/grh-lean/blob/72fead102f663d5c7a407964a13dc0cf9d977207/grh-compact/src/New/Filling/Setup.lean\#L326-L326}{{\leanfont Filling.\allowbreak{}coe\_ratToFin}}}{Setup.lean}{326--326}
\crow{finiteclassUnipintercompact}{\href{https://github.com/lsernau/grh-lean/blob/72fead102f663d5c7a407964a13dc0cf9d977207/grh-compact/src/New/Filling/Setup.lean\#L332-L339}{{\leanfont Filling.\allowbreak{}finite\_classUnip\_inter\_compact}}}{Setup.lean}{332--339}
\crow{existscutoffclassUnip}{\href{https://github.com/lsernau/grh-lean/blob/72fead102f663d5c7a407964a13dc0cf9d977207/grh-compact/src/New/Filling/Setup.lean\#L343-L352}{{\leanfont Filling.\allowbreak{}exists\_cutoff\_classUnip}}}{Setup.lean}{343--352}
\crow{classCutoff}{\href{https://github.com/lsernau/grh-lean/blob/72fead102f663d5c7a407964a13dc0cf9d977207/grh-compact/src/New/Filling/Setup.lean\#L356-L357}{{\leanfont Filling.\allowbreak{}classCutoff}}}{Setup.lean}{356--357}
\crow{amplitude}{\href{https://github.com/lsernau/grh-lean/blob/72fead102f663d5c7a407964a13dc0cf9d977207/grh-compact/src/New/Filling/Setup.lean\#L361-L362}{{\leanfont Filling.\allowbreak{}amplitude}}}{Setup.lean}{361--362}
\crow{amplitudeAbs}{\href{https://github.com/lsernau/grh-lean/blob/72fead102f663d5c7a407964a13dc0cf9d977207/grh-compact/src/New/Filling/Setup.lean\#L366-L367}{{\leanfont Filling.\allowbreak{}amplitudeAbs}}}{Setup.lean}{366--367}
\crow{amplitudeSum}{\href{https://github.com/lsernau/grh-lean/blob/72fead102f663d5c7a407964a13dc0cf9d977207/grh-compact/src/New/Filling/Setup.lean\#L370-L370}{{\leanfont Filling.\allowbreak{}amplitudeSum}}}{Setup.lean}{370--370}
\crow{abskappaeq}{\href{https://github.com/lsernau/grh-lean/blob/72fead102f663d5c7a407964a13dc0cf9d977207/grh-compact/src/New/Filling/Setup.lean\#L372-L373}{{\leanfont Filling.\allowbreak{}abs\_kappa\_eq}}}{Setup.lean}{372--373}
\crow{realFrameGramratToInf}{\href{https://github.com/lsernau/grh-lean/blob/72fead102f663d5c7a407964a13dc0cf9d977207/grh-compact/src/New/Filling/Setup.lean\#L375-L377}{{\leanfont Filling.\allowbreak{}realFrameGram\_ratToInf}}}{Setup.lean}{375--377}
\crow{detsubmatrixrealFrameGramratToInf}{\href{https://github.com/lsernau/grh-lean/blob/72fead102f663d5c7a407964a13dc0cf9d977207/grh-compact/src/New/Filling/Setup.lean\#L379-L382}{{\leanfont Filling.\allowbreak{}det\_submatrix\_realFrameGram\_ratToInf}}}{Setup.lean}{379--382}
\crow{unipCovolumepos}{\href{https://github.com/lsernau/grh-lean/blob/72fead102f663d5c7a407964a13dc0cf9d977207/grh-compact/src/New/Filling/Setup.lean\#L384-L384}{{\leanfont Filling.\allowbreak{}unipCovolume\_pos}}}{Setup.lean}{384--384}
\crow{gIdele}{\href{https://github.com/lsernau/grh-lean/blob/72fead102f663d5c7a407964a13dc0cf9d977207/grh-compact/src/New/Filling/Setup.lean\#L387-L387}{{\leanfont Filling.\allowbreak{}gIdele}}}{Setup.lean}{387--387}
\crow{unipotentChar}{\href{https://github.com/lsernau/grh-lean/blob/72fead102f663d5c7a407964a13dc0cf9d977207/grh-compact/src/New/Filling/Setup.lean\#L391-L391}{{\leanfont Filling.\allowbreak{}unipotentChar}}}{Setup.lean}{391--391}
\crow{Agrees}{\href{https://github.com/lsernau/grh-lean/blob/72fead102f663d5c7a407964a13dc0cf9d977207/grh-compact/src/New/Filling/Setup.lean\#L395-L396}{{\leanfont Filling.\allowbreak{}Agrees}}}{Setup.lean}{395--396}
\crow{IsStable}{\href{https://github.com/lsernau/grh-lean/blob/72fead102f663d5c7a407964a13dc0cf9d977207/grh-compact/src/New/Filling/Setup.lean\#L400-L400}{{\leanfont Filling.\allowbreak{}IsStable}}}{Setup.lean}{400--400}
\crow{termsOfLink}{\href{https://github.com/lsernau/grh-lean/blob/72fead102f663d5c7a407964a13dc0cf9d977207/grh-compact/src/New/Filling/Setup.lean\#L404-L404}{{\leanfont Filling.\allowbreak{}termsOfLink}}}{Setup.lean}{404--404}
\crow{splitGrammulVec}{\href{https://github.com/lsernau/grh-lean/blob/72fead102f663d5c7a407964a13dc0cf9d977207/grh-compact/src/New/Filling/StableCells.lean\#L15-L16}{{\leanfont splitGram\_mulVec}}}{StableCells.lean}{15--16}
\crow{linkMatrixmulVec}{\href{https://github.com/lsernau/grh-lean/blob/72fead102f663d5c7a407964a13dc0cf9d977207/grh-compact/src/New/Filling/StableCells.lean\#L18-L19}{{\leanfont linkMatrix\_mulVec}}}{StableCells.lean}{18--19}
\crow{isotropicFramePairmulVeceqzeroiff}{\href{https://github.com/lsernau/grh-lean/blob/72fead102f663d5c7a407964a13dc0cf9d977207/grh-compact/src/New/Filling/StableCells.lean\#L21-L24}{{\leanfont isotropicFramePair\_mulVec\_eq\_zero\_iff}}}{StableCells.lean}{21--24}
\crow{memlink}{\href{https://github.com/lsernau/grh-lean/blob/72fead102f663d5c7a407964a13dc0cf9d977207/grh-compact/src/New/Filling/StableCells.lean\#L28-L32}{{\leanfont mem\_link}}}{StableCells.lean}{28--32}
\crow{invb}{\href{https://github.com/lsernau/grh-lean/blob/72fead102f663d5c7a407964a13dc0cf9d977207/grh-compact/src/New/Filling/StableCells.lean\#L34-L36}{{\leanfont IsNormalForm.\allowbreak{}inv}}}{StableCells.lean}{34--36}
\crow{existsframeImageeqofisNormalForm}{\href{https://github.com/lsernau/grh-lean/blob/72fead102f663d5c7a407964a13dc0cf9d977207/grh-compact/src/New/Filling/StableCells.lean\#L38-L44}{{\leanfont Filling.\allowbreak{}exists\_frameImage\_eq\_of\_isNormalForm}}}{StableCells.lean}{38--44}
\crow{besselUmul}{\href{https://github.com/lsernau/grh-lean/blob/72fead102f663d5c7a407964a13dc0cf9d977207/grh-compact/src/New/Filling/StableCells.lean\#L46-L49}{{\leanfont IsNormalForm.\allowbreak{}besselU\_mul}}}{StableCells.lean}{46--49}
\crow{linkbesselUmul}{\href{https://github.com/lsernau/grh-lean/blob/72fead102f663d5c7a407964a13dc0cf9d977207/grh-compact/src/New/Filling/StableCells.lean\#L51-L55}{{\leanfont Filling.\allowbreak{}link\_besselU\_mul}}}{StableCells.lean}{51--55}
\crow{frameDiscframeGrambesselUmul}{\href{https://github.com/lsernau/grh-lean/blob/72fead102f663d5c7a407964a13dc0cf9d977207/grh-compact/src/New/Filling/StableCells.lean\#L57-L60}{{\leanfont Filling.\allowbreak{}frameDisc\_frameGram\_besselU\_mul}}}{StableCells.lean}{57--60}
\crow{isStablehasLevel}{\href{https://github.com/lsernau/grh-lean/blob/72fead102f663d5c7a407964a13dc0cf9d977207/grh-compact/src/New/Filling/StableCells.lean\#L64-L67}{{\leanfont Filling.\allowbreak{}isStable\_hasLevel}}}{StableCells.lean}{64--67}
\crow{memlinkifforthogonal}{\href{https://github.com/lsernau/grh-lean/blob/72fead102f663d5c7a407964a13dc0cf9d977207/grh-compact/src/New/Filling/StableCells.lean\#L71-L79}{{\leanfont mem\_link\_iff\_orthogonal}}}{StableCells.lean}{71--79}
\crow{frameDiscweylElemmul}{\href{https://github.com/lsernau/grh-lean/blob/72fead102f663d5c7a407964a13dc0cf9d977207/grh-compact/src/New/Filling/StableCells.lean\#L87-L91}{{\leanfont Filling.\allowbreak{}frameDisc\_weylElem\_mul}}}{StableCells.lean}{87--91}
\crow{linkbigcelleqiff}{\href{https://github.com/lsernau/grh-lean/blob/72fead102f663d5c7a407964a13dc0cf9d977207/grh-compact/src/New/Filling/StableCells.lean\#L96-L100}{{\leanfont Filling.\allowbreak{}link\_big\_cell\_eq\_iff}}}{StableCells.lean}{96--100}}
% the notations of the defined functions: one slot for each explicit argument of the Lean declaration
\newcommand{\occ}[1]{c_{G}(#1)}
\newcommand{\bandc}[1]{c_{\mathrm{band}}(#1)}
\newcommand{\bsize}[1]{\operatorname{bs}(#1)}
\newcommand{\InCut}[2]{#1\in\operatorname{Cut}(#2)}
\newcommand{\gdr}[1]{\mathsf{g}^{\mathbb{R}}(#1)}
\newcommand{\cellm}[4]{\operatorname{cm}(#1;#2,#3,#4)}
\newcommand{\ncorr}[3]{\mathrm{J}^{\circ}(#1;#2,#3)}
\newcommand{\BandL}[4]{\operatorname{BL}(#1,#2,#3,#4)}
\newcommand{\FixN}[2]{\operatorname{fix}(#1)\ge #2}
\newcommand{\zfp}[1]{\mathrm{L}_{z}(#1)}
\newcommand{\clink}[2]{\operatorname{link}^{\circ}(#1,#2)}
\newcommand{\HCFV}[4]{\operatorname{CF}_{#4}(#1;#2,#3)}
\newcommand{\ratG}[1]{J_{#1}}
\newcommand{\lform}[1]{Q^{\pm}_{#1}}
\newcommand{\clC}[2]{\mathrm{C}_{#1}(#2)}
\newcommand{\clCX}[2]{\mathrm{C}^{\circ}_{#1}(#2)}
\newcommand{\cw}[2]{\operatorname{cw}_{#2}(#1)}
\newcommand{\HasLev}[3]{#3\in\operatorname{Lv}_{#2}(#1)}
\newcommand{\Cong}[5]{#4\sim_{#1;#2,#3}#5}
\newcommand{\tdnorm}[2]{\|\mathsf{g}\|_{#1}(#2)}
\newcommand{\bxp}[1]{\operatorname{BxP}(#1)}
\newcommand{\IsNr}[2]{#2\in\operatorname{Near}(#1)}
\newcommand{\ChA}[2]{#2\in\operatorname{ChA}(#1)}
\newcommand{\Tlatg}[2]{T^{\mathrm{lat}}(#1;#2)}
\newcommand{\mbc}[2]{\Lambda^{\circ}(#1;#2)}
\newcommand{\blb}[2]{\Lambda(#1;#2)}
\newcommand{\fdisc}[2]{\mathsf{g}(#1;#2)}
\newcommand{\fpair}[1]{[\mathfrak{f}\mid #1\mathfrak{f}]}
\newcommand{\ifpair}[1]{[\mathfrak{e}\mid #1\mathfrak{e}]}
\newcommand{\lmat}[1]{\mathrm{L}(#1)}
\newcommand{\link}[1]{\operatorname{link}(#1)}
\newcommand{\damp}[2]{\operatorname{dp}_{#1}(#2)}
\newcommand{\hw}[3]{\Lambda^{\mathbb{C}}_{#1}(#2;#3)}
\newcommand{\Sieve}[8]{\operatorname{Sieve}(#1,#2,#3,#4,#5,#6,#7,#8)}
\newcommand{\fhw}[3]{\Lambda_{#1}(#2;#3)}
\newcommand{\fadm}[1]{(#1)_{\mathbb{A}_{f}}}
\newcommand{\sld}[2]{\operatorname{disc}_{#1}(#2)}
\newcommand{\ioia}[2]{\operatorname{adj}_{#1}^{#2}}
\newcommand{\rag}[2]{\mathrm{red}_{#1,#2}}
\newcommand{\Fill}[2]{\operatorname{Fill}(#1,#2)}
\newcommand{\gS}[1]{g(#1)}
\newcommand{\Trat}[2]{T^{\mathrm{rat}}_{#1}(#2)}
\newcommand{\Tlat}[2]{T^{\mathrm{lat}}_{#1}(#2)}
\newcommand{\GS}[2]{G_{#1}(#2)}
\newcommand{\GFin}[1]{G^{\mathrm{lat}}_{#1}}
\newcommand{\GRes}[1]{\mathsf{G}(#1)}
\newcommand{\Bc}[1]{\mathrm{Bc}(#1)}
\newcommand{\Cl}[1]{\mathrm{Cl}(#1)}
\newcommand{\Zf}[2]{Z_{#1}(#2)}
\newcommand{\KS}[1]{K(#1)}
\newcommand{\KgS}[1]{K_g(#1)}
\newcommand{\redim}[2]{\overline{#2}^{\,#1}}
\newcommand{\kapS}[2]{\kappa_{#2}(#1)}
\newcommand{\sdc}[1]{\psi_{#1}}
\newcommand{\fscale}[3]{\mathrm{fs}_{#1}(#2,#3)}
\newcommand{\siegC}[1]{c_{*}(#1)}
\newcommand{\tdisc}[2]{\mathsf{g}_{#1}(#2)}
\newcommand{\bsGL}[2]{\mathrm{Sc}_{#1,#2}}
\newcommand{\latB}[1]{\Lambda(#1)}
\newcommand{\tolb}[1]{\mathrm{lb}_{#1}}
\newcommand{\rinf}[1]{\mathrm{r}^{\infty}_{#1}}
\newcommand{\rfin}[1]{\mathrm{r}^{f}_{#1}}
\newcommand{\incH}[1]{\mathrm{inc}_{#1}}
\newcommand{\hc}[2]{h_{#1}(#2)}
\newcommand{\pos}[4]{\mathrm{pos}_{#1}(#2,#3;#4)}
\newcommand{\clU}[2]{\mathcal{U}_{#1,#2}}
\newcommand{\clUo}[2]{\mathcal{U}^{0}_{#1,#2}}
\newcommand{\bxs}[1]{\operatorname{bx}(#1)}
\newcommand{\cdim}[1]{\mathrm{cd}(#1)}
\newcommand{\mainE}[1]{\mathsf{d}(#1)}
\newcommand{\uord}[1]{\mathrm{uo}(#1)}
\newcommand{\packC}[1]{a_{*}(#1)}
\newcommand{\besE}[1]{\mathsf{d}_{1}(#1)}
\newcommand{\tailS}[1]{\mathrm{tl}(#1)}
\newcommand{\codimF}[1]{\mathrm{cf}(#1)}
\newcommand{\cutoff}[2]{\beta_{#1,#2}}
\newcommand{\gid}[1]{\hat{g}(#1)}
\newcommand{\uchar}[2]{\psi^{\mathcal{U}}_{#1}(#2)}
\newcommand{\Agr}[2]{#2\in\operatorname{Agr}(#1)}
\newcommand{\UnitMembers}{497}
\newcommand{\UnitInNote}{232}
\newcommand{\UnitProofs}{179}
\newcommand{\UnitSteps}{1371}
\newcommand{\UnitAllSteps}{1479}

% generated by notes/filling-setup/shared/mkspec.py from shared/notation-148ab.json; do not edit.
% \input after data.tex: a macro that the build made from the spec is kept, the others are defined here (same TeX).
\providecommand{\fs}{\mathfrak{s}}                      % the shape (rho, kappa, k_0)
\providecommand{\sT}{\mathsf{T}}                        % a finite set of finite places
\providecommand{\cL}{\mathcal{L}}                       % a level
\providecommand{\AEf}{\mathbb{A}_{E,f}}
\providecommand{\AdE}{\mathbb{A}_{E}}
\providecommand{\hO}{\hat{O}_E}
\providecommand{\GL}{\mathrm{GL}}
\providecommand{\full}{\mathrm{full}}
\providecommand{\cut}{\mathrm{cut}}
\providecommand{\rh}[1]{\varrho(#1)}                    % Shape.ϱ
\providecommand{\kk}[1]{k(#1)}                          % Shape.k = #κ
\providecommand{\nn}[1]{n(#1)}                          % Shape.n = ϱ + 1
\providecommand{\rr}[1]{r(#1)}                          % Shape.r = 2ϱ + k + 1
\providecommand{\Ix}[1]{\operatorname{Ind}(#1)}                          % Shape.BesselIndex
\providecommand{\Mx}[1]{\operatorname{Mid}(#1)}                          % Shape.MiddleIndex
\providecommand{\ninf}[1]{f(#1)}                        % numInfPlaces E
\providecommand{\ET}[1]{E_{#1}}                         % LocProd T
\providecommand{\bS}[1]{b(#1)}                          % S.b
\providecommand{\NS}[1]{N(#1)}                          % S.N
\providecommand{\qS}[2]{q_{#2}(#1)}                     % S.levelPrime i
\providecommand{\piS}[1]{\bar\pi(#1)}                   % S.piBar
\providecommand{\vpS}[1]{\varpi(#1)}                    % S.ϖ
\providecommand{\ChS}[2]{\mathrm{Ch}_{#2}(#1)}          % S.chosenPlaces i
\providecommand{\KcolS}[1]{\mathsf{K}^{\mathrm{col}}(#1)}        % S.columnLevel
\providecommand{\BsysS}[2]{\mathfrak{B}_{#2}(#1)}       % S.besselSys v
\providecommand{\mass}[2]{\operatorname{mass}(#1,#2)}   % mass E r
% Lib: forms, groups and frames (letters of the notes of unit 132)
\providecommand{\fr}[1]{\mathfrak{f}_{#1}}              % the frame vector of index a (e_a for a < rho, z for a = rho)
\providecommand{\Trg}[2]{T^{\mathrm{rat}}(#1;#2)}       % Shape.middleGram N: the form diag(1,...,1,-N)
\providecommand{\JT}[1]{J(#1)}                          % Shape.Form.splitGram
\providecommand{\GT}[1]{G(#1)}                          % Shape.Form.G
\providecommand{\Gone}[1]{G_{1}(#1)}                    % Shape.G₁
\providecommand{\UT}[1]{\mathcal{U}(#1)}                % besselU
\providecommand{\UoT}[1]{\mathcal{U}^{0}(#1)}           % besselU0
\providecommand{\HT}[1]{\mathcal{H}(#1)}                % besselH
\providecommand{\ST}[1]{\mathcal{S}(#1)}                % besselS
\providecommand{\Hl}[2]{\mathcal{H}_{#1}(#2)}           % GT l
\providecommand{\Ql}[2]{\mathcal{Q}_{#1}(#2)}           % unitriGroup l
\providecommand{\sd}[1]{\operatorname{sd}(#1)}          % simpleRootSum
\providecommand{\ZT}[2]{Z_{#1}(#2)}                     % Shape.Form.frameGram
\providecommand{\Zr}[1]{Z(#1)}                          % realFrameGram
\providecommand{\lmin}[2]{\mathrm{D}_{#1}(#2)}          % leadingMinor Z l
\providecommand{\weyl}[2]{w_{#1}(#2)}                   % weylElem l p
\providecommand{\gix}[1]{[#1]}                          % genIndex (the generalized index)
\providecommand{\den}[1]{\operatorname{den}(#1)}        % denomIndex
\providecommand{\Kf}[2]{\mathsf{K}_{f}(#1,#2)}                   % Kf E n
\providecommand{\Gf}[2]{G_{f}(#1,#2)}                   % Gf E n
\providecommand{\volO}[1]{\operatorname{vol}(\hat O_{#1})} % volÔ E
\providecommand{\rdE}[1]{\operatorname{rd}_{#1}}        % rootDiscr E
\providecommand{\dE}[1]{d_{#1}}                         % discr E
\providecommand{\CC}{\mathbb{C}}
\providecommand{\RR}{\mathbb{R}}
\providecommand{\NN}{\mathbb{N}}
\providecommand{\vol}{\operatorname{vol}}
% ----- the definitions of chunks a and b of unit 148 (and, for this note, its own additions)
\providecommand{\Tlatg}[2]{T^{\mathrm{lat}}(#1;#2)}                  % latMiddleGram (𝔰, N)
\providecommand{\mbc}[2]{\Lambda^{\circ}(#1;#2)}                     % middleBaseChange (𝔰, p)
\providecommand{\blb}[2]{\Lambda(#1;#2)}                             % besselLatticeBasis (𝔰, p)
\providecommand{\fdisc}[2]{\mathsf{g}(#1;#2)}                        % frameDisc (N, Z)
\providecommand{\fpair}[1]{[\mathfrak{f}\mid #1\mathfrak{f}]}        % framePair (γ)
\providecommand{\ifpair}[1]{[\mathfrak{e}\mid #1\mathfrak{e}]}       % isotropicFramePair (γ)
\providecommand{\lmat}[1]{\mathrm{L}(#1)}                            % linkMatrix (γ)
\providecommand{\link}[1]{\operatorname{link}(#1)}                   % link (γ)
\providecommand{\damp}[2]{\operatorname{dp}_{#1}(#2)}                % dampedPower (k, r)
\providecommand{\hw}[3]{\Lambda^{\mathbb{C}}_{#1}(#2;#3)}            % headWeight (K, i, p)
\providecommand{\Sieve}[8]{\operatorname{Sieve}(#1,#2,#3,#4,#5,#6,#7,#8)} % SieveSystem (u, ι, n, J, ψ, En, Λ, β)
\providecommand{\fhw}[3]{\Lambda_{#1}(#2;#3)}                        % finHeadWeight (K, i, p)
\providecommand{\occ}[1]{c_{G}(#1)}                                  % Shape.openCellConst (𝔰)
\providecommand{\bandc}[1]{c_{\mathrm{band}}(#1)}                    % Shape.bandConst (𝔰)
\providecommand{\bsize}[1]{\operatorname{bs}(#1)}                    % besselSize (y)
\providecommand{\InCut}[2]{#1\in\operatorname{Cut}(#2)}              % InCut (y, b)
\providecommand{\frbox}[1]{\operatorname{FBox}(#1)}                  % frameBox (ρ)
\providecommand{\gdr}[1]{\mathsf{g}^{\mathbb{R}}(#1)}                % gramDiscReal (Z)
\providecommand{\cellm}[4]{\operatorname{cm}(#1;#2,#3,#4)}           % cellMass (y, deg, ρ, d)
\providecommand{\bcorr}[3]{\mathrm{J}(#1,#2;#3)}                     % besselCorrelation (μ, W, γ)
\providecommand{\SchurLe}[2]{\operatorname{Sch}(#1)\le #2}           % SchurLe (W, D)
\providecommand{\HasSB}[4]{\operatorname{SchB}(#1,#2;#3,#4)}         % HasSchurBound (W, rep, B, D)
\providecommand{\ncorr}[3]{\mathrm{J}^{\circ}(#1;#2,#3)}             % normCorrelation (y, W, p)
\providecommand{\BandL}[4]{\operatorname{BL}(#1,#2,#3,#4)}           % HasBandLimits (b, W, deg, J)
\providecommand{\CutPair}[4]{\operatorname{CP}(#1,#2,#3,#4)}         % IsCutPair (b, f, rep, W)
\providecommand{\BSys}[3]{\operatorname{BSys}(#1,#2,#3)}             % BesselSystem (𝔰, b, f)
\providecommand{\Fill}[2]{\operatorname{Fill}(#1,#2)}                % Filling (E, 𝔰)
\providecommand{\gS}[1]{g(#1)}                                       % Filling.g (S)
\providecommand{\Trat}[2]{T^{\mathrm{rat}}_{#1}(#2)}                 % Filling.ratMiddle (S, R)
\providecommand{\Tlat}[2]{T^{\mathrm{lat}}_{#1}(#2)}                 % Filling.latMiddle (S, R)
\providecommand{\GS}[2]{G_{#1}(#2)}                                  % Filling.G (S, R)
\providecommand{\GFin}[1]{G^{\mathrm{lat}}_{#1}}                     % Filling.GFin (S)
\providecommand{\GRes}[1]{\mathsf{G}(#1)}                            % Filling.GRes (S)
\providecommand{\Bc}[1]{\mathrm{Bc}(#1)}                             % Filling.BesselCoord (S)
\providecommand{\Cl}[1]{\mathrm{Cl}(#1)}                             % Filling.Classes (S)
\providecommand{\Zf}[2]{Z_{#1}(#2)}                                  % Filling.frameProducts (S, γ)
\providecommand{\KS}[1]{\mathsf{K}(#1)}                              % Filling.latStabilizer (S)
\providecommand{\redS}[1]{\mathrm{red}_{#1}}                         % Filling.reduction (S)
\providecommand{\KgS}[1]{\mathsf{K}_g(#1)}                           % Filling.principalLevel (S)
\providecommand{\redim}[2]{\overline{#2}^{\,#1}}                     % Filling.reducedImage (S, A)
\providecommand{\kapS}[2]{\kappa_{#2}(#1)}                           % Filling.kappa (S, v)
\providecommand{\bscale}[3]{\mathrm{sc}_{#1,#2}(#3)}                 % Filling.basisScale (S, v, i)
\providecommand{\sdc}[1]{\psi_{#1}}                                  % Filling.selfDualChar (S)
\providecommand{\fscale}[3]{\mathrm{fs}_{#1}(#2,#3)}                 % Filling.frameScaleSq (S, v, a)
\providecommand{\fsprod}[1]{\mathrm{fs}(#1)}                         % Filling.frameScaleProd (S)
\providecommand{\ucov}[1]{\mathrm{cov}(#1)}                          % Filling.unipCovolume (S)
\providecommand{\siegC}[1]{c_{*}(#1)}                                % Filling.siegelConst (S)
\providecommand{\tdisc}[2]{\mathsf{g}_{#1}(#2)}                      % Filling.termDisc (S, γ)
\providecommand{\bsGL}[2]{\mathrm{Sc}_{#1,#2}}                       % Filling.basisScaleGL (S, v)
\providecommand{\latB}[1]{\Lambda(#1)}                               % Filling.latBasis (S)
\providecommand{\tolb}[1]{\mathrm{lb}_{#1}}                          % Filling.toLatBasis (S)
\providecommand{\radl}[1]{\mathrm{r}^{\mathbb{A}}_{#1}}              % Filling.ratToAdelic (S)
\providecommand{\ainf}[1]{\mathrm{a}^{\infty}_{#1}}                  % Filling.adelicToInf (S)
\providecommand{\afin}[1]{\mathrm{a}^{f}_{#1}}                       % Filling.adelicToFin (S)
\providecommand{\rinf}[1]{\mathrm{r}^{\infty}_{#1}}                  % Filling.ratToInf (S)
\providecommand{\rfin}[1]{\mathrm{r}^{f}_{#1}}                       % Filling.ratToFin (S)
\providecommand{\incH}[1]{\mathrm{inc}_{#1}}                         % Filling.inclHFin (S)
\providecommand{\hc}[2]{h_{#1}(#2)}                                  % Filling.h (S, c)
\providecommand{\pos}[4]{\mathrm{pos}_{#1}(#2,#3;#4)}                % Filling.termPosition (S, c, c', γ)
\providecommand{\terms}[4]{\mathrm{Tm}_{#1}(#2,#3,#4)}               % Filling.termsOver (S, c, s, c')
\providecommand{\clU}[2]{\mathcal{U}_{#1,#2}}                        % Filling.classUnip (S, c)
\providecommand{\clUo}[2]{\mathcal{U}^{0}_{#1,#2}}                   % Filling.classUnip₀ (S, c)
\providecommand{\bxs}[1]{\operatorname{bx}(#1)}                      % Filling.boxSize (S)
\providecommand{\IsBP}[3]{#3\in\operatorname{BP}(#1,#2)}             % Filling.IsBoxPartition (S, c, ω)
\providecommand{\besfn}[4]{W^{#1}_{#2}(#3)_{#4}}                     % Filling.besselFn (σ, S, y, i)
\providecommand{\besker}[4]{\mathcal{W}^{#1}_{#2}(#3;#4)}            % Filling.besselKernel (σ, S, γ, y)
\providecommand{\cdim}[1]{\mathrm{cd}(#1)}                           % Filling.columnDim (S)
\providecommand{\bdeg}[1]{\deg(#1)}                                  % Filling.besselDegree (S)
\providecommand{\mainE}[1]{\mathsf{d}(#1)}                           % Filling.mainEnergy (S)
\providecommand{\uord}[1]{\mathrm{uo}(#1)}                           % Filling.unipOrder (S)
\providecommand{\packC}[1]{a_{*}(#1)}                                % Filling.packetConst (S)
\providecommand{\besE}[1]{\mathsf{d}_{1}(#1)}                        % Filling.besselEnergy (S)
\providecommand{\tailS}[1]{\mathrm{tl}(#1)}                          % Filling.tailScale (S)
\providecommand{\codimF}[1]{\mathrm{cf}(#1)}                         % Filling.codimFactor (S)
\providecommand{\IsAdm}[1]{#1\text{ is admissible}}                  % Filling.IsAdmissible (S)
\providecommand{\cutoff}[2]{\beta_{#1,#2}}                           % Filling.classCutoff (S, c)
\providecommand{\amp}[5]{a^{#1}_{#2}(#3,#4;#5)}                      % Filling.amplitude (σ, S, c, c', γ)
\providecommand{\ampabs}[5]{|a|^{#1}_{#2}(#3,#4;#5)}                 % Filling.amplitudeAbs (σ, S, c, c', γ)
\providecommand{\ampsum}[6]{A^{#1}_{#2}(#3;#4,#5,#6)}                % Filling.amplitudeSum (σ, S, T, c, s, c')
\providecommand{\gid}[1]{\hat{g}(#1)}                                % Filling.gIdele (S)
\providecommand{\uchar}[2]{\psi^{\mathcal{U}}_{#1}(#2)}              % Filling.unipotentChar (S, u)
\providecommand{\Agr}[2]{#2\in\operatorname{Agr}(#1)}                % Filling.Agrees (S, s)
\providecommand{\IsStab}[2]{#2\in\operatorname{Stab}(#1)}            % Filling.IsStable (S, T)
\providecommand{\tlink}[2]{\mathrm{Tm}^{\mathrm{lk}}_{#1}(#2)}       % Filling.termsOfLink (S, Λ)
\providecommand{\clB}[2]{\mathrm{B}_{#1}(#2)}                        % Filling.classBasis (S, c)
\providecommand{\ioia}[2]{\operatorname{adj}_{#1}^{#2}}              % Filling.idOrInvAdjoint (S, b)
\providecommand{\rag}[2]{\mathrm{red}_{#1,#2}}                       % Filling.reduceAboveG (S, x)
\providecommand{\vfix}[3]{\mathrm{Fix}_{#1,#2}(#3)}                  % Filling.vectorFixer (S, w, d)
\providecommand{\FixN}[2]{\operatorname{fix}(#1)\ge #2}              % Filling.FixesNormAtLeast (s, x)
\providecommand{\fixset}[3]{\mathrm{Fix}_{#1}(#2,#3)}                % Filling.fixingSet (S, Sg, d)
\providecommand{\zfp}[1]{\mathrm{L}_{z}(#1)}                         % Filling.zFormPair (T)
\providecommand{\clink}[2]{\operatorname{link}^{\circ}(#1,#2)}       % commonLink (γ, γ')
\providecommand{\HCFV}[4]{\operatorname{CF}_{#4}(#1;#2,#3)}          % HasCommonFlagVector (γ, c, c', a)
\providecommand{\IsNF}[2]{#2\in\operatorname{NF}_{#1}}               % IsNormalForm (j, γ)
\providecommand{\ratG}[1]{J_{#1}}                                    % Filling.ratGram (S)
\providecommand{\lform}[1]{Q^{\pm}_{#1}}                             % Filling.linkForm (S)
\providecommand{\clC}[2]{\mathrm{C}_{#1}(#2)}                        % Filling.classCoord (S, c)
\providecommand{\clCX}[2]{\mathrm{C}^{\circ}_{#1}(#2)}               % Filling.classCoordX₀ (S, c)
\providecommand{\cw}[2]{\operatorname{cw}_{#2}(#1)}                  % Filling.countWeight (S, i)
\providecommand{\HasLev}[3]{#3\in\operatorname{Lv}_{#2}(#1)}         % Filling.HasLevel (S, j, γ)
\providecommand{\Cong}[5]{#4\sim_{#1;#2,#3}#5}                       % Filling.CongruentTerms (S, c, c', γ, γ')
\providecommand{\tdnorm}[2]{\|\mathsf{g}\|_{#1}(#2)}                 % Filling.termDiscNorm (S, γ)
\providecommand{\bxp}[1]{\operatorname{BxP}(#1)}                     % Filling.boxedPoints (S)
\providecommand{\IsNr}[2]{#2\in\operatorname{Near}(#1)}              % Filling.IsNear (S, γ)
\providecommand{\ChA}[2]{#2\in\operatorname{ChA}(#1)}                % Filling.CharsAgree (S, γ)
\providecommand{\fadm}[1]{(#1)_{\mathbb{A}_{f}}}                     % Filling.finAdeleMat (M)
\providecommand{\slb}[2]{\mathrm{m}_{#1}(#2)}                        % Filling.sublatticeBasis (H, V)
\providecommand{\sld}[2]{\operatorname{disc}_{#1}(#2)}               % Filling.sublatticeDisc (u, V)
\providecommand{\lmet}[1]{Q^{\mathrm{link}}_{#1}}                    % Filling.linkMetric (S)
\providecommand{\ldsq}[3]{\mathrm{ld}_{#1}(#2,#3)}                   % Filling.linkDiscSq (S, v, e)
\providecommand{\HIL}[3]{#3\in\operatorname{IL}(#1;#2)}              % Filling.HasIntegralLifts (c, γ, m)
\providecommand{\IdxB}[1]{\mathrm{Ib}(#1)}                           % Filling.IndexBelow (j)
\providecommand{\xrows}[1]{#1^{\mathrm{x}}}                          % Filling.xRows (Z)
\providecommand{\padl}[2]{\operatorname{pad}_{#1}(#2)}               % Filling.paddedLeading (j, Z)
\providecommand{\fmid}[1]{(#1\mathfrak{f})^{\mathrm{mid}}}           % Filling.frameMiddle (γ)
\providecommand{\latd}[2]{\operatorname{disc}_{#1,#2}}               % Filling.latticeDisc (S, c)
\providecommand{\fvi}[5]{\mathrm{o}_{#1,#2}(#3;#4,#5)}               % Filling.frameValueIndex (S, c, k, a, V)
\providecommand{\xri}[3]{\mathrm{i}_{#1,#2}(#3)}                     % Filling.xRowsIndex (S, j, k)
\providecommand{\sati}[2]{\mathfrak{d}_{#1}(#2)}                     % Filling.saturationIndex (S, k)
\providecommand{\pinv}[5]{\mathrm{pi}_{#1}(#2,#3;#4,#5)}             % Filling.pairInvariant (S, c, c', γ, γ')
\providecommand{\telem}[3]{\mathrm{el}_{#1}(#2;#3)}                  % Filling.termElem (S, c, p)
\providecommand{\mnorm}[3]{\mathrm{mn}_{#1}(#2;#3)}                  % Filling.minorNorm (S, j, γ₁)
\providecommand{\cgauge}[5]{\mathrm{cg}_{#1}(#2;#3,#4;#5)}           % Filling.classGauge (S, j, c, c', γ)
\providecommand{\IsNRep}[4]{#4\in\operatorname{NRep}_{#1}(#2,#3)}    % Filling.IsNearRep (S, c, c', γ)
\providecommand{\ipw}[4]{\mathrm{ipw}_{#1}(#2;#3,#4)}                % Filling.invPowWeight (S, c, n, p)
\providecommand{\termw}[7]{\mathrm{w}_{#1}(#2,#3;#4;#5,#6;#7)}       % Filling.termWeight (S, j, i, c, p, p', V)
\providecommand{\pairw}[6]{\mathrm{pw}_{#1}(#2;#3,#4;#5,#6)}         % Filling.pairWeight (S, j, c, c', γ, γ')
\providecommand{\IsLab}[5]{#5\in\operatorname{Lab}_{#1}(#2;#3,#4)}   % Filling.IsLabelledTerm (S, j, c, t, p)
\providecommand{\energy}[3]{\mathrm{En}_{#1}(#2;#3)}                 % Filling.energy (S, A, Ω)
\providecommand{\linkE}[5]{\mathrm{LE}^{#1}_{#2}(#3;#4,#5)}          % Filling.linkEnergy (σ, S, T, Ω', R)
\providecommand{\csize}[3]{\operatorname{cs}_{#1}(#2;#3)}            % Filling.classSize (S, j, x)
\providecommand{\bcut}[2]{\beta^{0}_{#1,#2}}                         % Filling.boxCutoff (S, c)
\providecommand{\schw}[4]{\mathrm{sw}_{#1}(#2;#3,#4)}                % Filling.schurWeight (S, c, γ, y)
\providecommand{\schi}[4]{\mathrm{SI}_{#1}(#2,#3;#4)}                % Filling.schurIntegral (S, c, c', γ')
\providecommand{\charT}[2]{\psi_{#1,#2}}                             % Filling.charOnPlaces (S, T)
\providecommand{\fgain}[3]{\beta_{#1,#2}(#3)}                        % Filling.frequencyGain (S, T, ξ)
\providecommand{\ratpt}[3]{\mathrm{pt}_{#1}(#2;#3)}                  % Filling.ratPoint (T, l, t)
\providecommand{\Fits}[6]{\operatorname{Fits}(#1;#2,#3,#4,#5,#6)}    % Filling.Fits (S, V, l, kw, Y, η)
\providecommand{\stdK}[1]{\mathsf{K}^{\mathrm{std}}_{g}(#1)}         % Filling.stdPrincipalLevel (S)
\providecommand{\tailv}[3]{\operatorname{tv}(#1,#2,#3)}              % tailVector (π, t, x)


\linespread{1.08}                       % 11 pt on 14.7 pt
\setlength{\parindent}{0pt}
\setlength{\parskip}{5.5pt plus 1.5pt minus 0.5pt}
\setlength{\abovedisplayskip}{7.5pt plus 2pt minus 3pt}
\setlength{\belowdisplayskip}{7.5pt plus 2pt minus 3pt}
\setlength{\abovedisplayshortskip}{3pt plus 2pt}
\setlength{\belowdisplayshortskip}{6pt plus 2pt minus 2pt}
\emergencystretch=1.5em
\clubpenalty=10000 \widowpenalty=10000 \displaywidowpenalty=10000
\relpenalty=7000 \binoppenalty=9000
\raggedbottom
\pagestyle{fancy}\fancyhf{}\renewcommand{\headrulewidth}{0pt}
\fancyfoot[C]{\sffamily\fontsize{7.6}{9}\selectfont\color{inkfaint}\thepage\,/\,\pageref*{LastPage}}
\setlist{topsep=2pt,partopsep=0pt,parsep=0pt,itemsep=4pt,leftmargin=2.1em,labelsep=0.55em}

\newcommand{\Q}{\mathbb{Q}}
\newcommand{\Z}{\mathbb{Z}}
\newcommand{\F}{\mathbb{F}}
\newcommand{\len}{\operatorname{len}}
\newcommand{\cok}{\operatorname{cok}}
\newcommand{\End}{\operatorname{End}}
\newcommand{\Hom}{\operatorname{Hom}}
\newcommand{\Gal}{\operatorname{Gal}}
\newcommand{\diag}{\operatorname{diag}}
\newcommand{\kw}[1]{\emph{#1}}
\newcommand{\m}[1]{\mbox{$#1$}}        % a short formula that must not be broken across lines
\newcommand{\head}[1]{\par\addvspace{17pt plus 4pt minus 2pt}\Needspace{6\baselineskip}%
  \noindent{\fontsize{13.2}{16}\selectfont\bfseries\color{ink}#1}\par\nopagebreak\vspace{-1pt}}
\newcommand{\sect}[2]{\refstepcounter{section}\label{#2}\par\addvspace{19pt plus 4pt minus 2pt}\Needspace{7\baselineskip}%
  \noindent\llap{{\fontsize{13.2}{16}\selectfont\bfseries\color{inkfaint}\thesection}\hspace{9pt}}%
  {\fontsize{13.2}{16}\selectfont\bfseries\color{ink}#1}\par\nopagebreak\vspace{-1pt}}
\newcommand{\subhead}[1]{\par\addvspace{9pt plus 2pt}\Needspace{4\baselineskip}\noindent\textbf{#1}\enspace\ignorespaces}
% where a statement lives in the Lean archive
\newcommand{\leantag}{{\sffamily\bfseries\fontsize{6.3}{8}\selectfont\color{leancol}\addfontfeatures{LetterSpace=7}LEAN}}
\newcommand{\sep}{\unskip\nobreak\hspace{5pt}{\color{inkfaint}\textperiodcentered}\hspace{5pt plus 2pt}\ignorespaces}
\newcommand{\leanline}[1]{\par\nopagebreak\vspace{-2.5pt}{\fontsize{8.9}{12.4}\selectfont\color{inksoft}\raggedright
  \noindent\llap{\leantag\hspace{9pt}}#1\par}\vspace{1pt}}
% a displayed statement in a tinted box: #1 colour, #2 text
\newcommand{\statement}[2]{\par\addvspace{7pt}\begingroup\setlength{\fboxsep}{8.5pt}\setlength{\fboxrule}{0.6pt}%
  \noindent\fcolorbox{#1!30!white}{#1!3!white}{\parbox{\dimexpr\linewidth-2\fboxsep-2\fboxrule\relax}{#2}}%
  \endgroup\par\addvspace{7pt}}
\newcommand{\nt}[2]{\par\noindent\hangindent=1.1em\hangafter=1 #1\enspace{\color{inksoft}#2}\par\vspace{1pt}}
\newcommand{\legdot}[1]{\tikz[baseline=-0.62ex]\fill[#1] (0,0) circle (2.5pt);\hspace{3pt}}
\newcommand{\crow}[4]{\csname CS@#1\endcsname & #2 & \textcolor{inksoft}{#3} & \textcolor{inksoft}{#4}\\}
\newcommand{\leanquoteline}[2]{\par\noindent\hangindent=\dimexpr#1\fontdimen2\font+2.2em\relax\hangafter=1
  \hspace*{#1\fontdimen2\font}#2}
\newcommand{\quoted}[1]{\par\vspace{3pt}{\ttfamily\fontsize{7.5}{10.4}\selectfont\color{leancol}\raggedright
  \setlength{\parskip}{0pt}\leanquote{#1}\par}\vspace{4pt}}
\newcommand{\bx}[3]{{\sffamily\bfseries\fontsize{7.3}{9.2}\selectfont\color{#1!82!black}#2}\\[1.5pt]#3}
\newcommand{\cs}[2]{\expandafter\newcommand\csname CS@#1\endcsname{#2}}

% theorems
\newtheoremstyle{plainx}{7pt plus 2pt}{7pt plus 2pt}{\itshape}{}{\bfseries}{.}{0.5em}{}
\newtheoremstyle{uprightx}{7pt plus 2pt}{7pt plus 2pt}{}{}{\bfseries}{.}{0.5em}{}
\theoremstyle{plainx}
\newtheorem{theorem}{Theorem}[section]
\newtheorem{lemma}[theorem]{Lemma}
\newtheorem{proposition}[theorem]{Proposition}
\newtheorem{corollary}[theorem]{Corollary}
\theoremstyle{uprightx}
\newtheorem{definition}[theorem]{Definition}
\newtheorem{setting}[theorem]{Setting}          % hypotheses shared by several statements; each of them cites it by number
\newtheorem{remark}[theorem]{Remark}
\newtheorem{example}[theorem]{Example}
\renewcommand{\qedsymbol}{{\setlength{\fboxsep}{0pt}\setlength{\fboxrule}{0.5pt}\fbox{\rule{0pt}{0.46em}\hspace{0.46em}}}}
% what holds throughout the note: the ambient objects only (implicit and instance arguments), never a hypothesis
\newenvironment{conventions}{\par\addvspace{5pt}\noindent\textbf{Conventions.}\enspace\ignorespaces}{\par\addvspace{3pt}}
\makeatletter
\renewenvironment{proof}[1][Proof]{\par\pushQED{\qed}\normalfont\topsep0pt\partopsep0pt
  \trivlist\item[\hskip\labelsep\itshape #1.]\ignorespaces}{\popQED\endtrivlist\@endpefalse}
\makeatother

% ------------------------------------------------------------------------------------- local macros of this note
\newcommand{\cj}[1]{\overline{#1}}                  % the involution of a ring; complex conjugation
\newcommand{\adj}[1]{#1^{*}}                        % conjugate transpose of a matrix
\newcommand{\tr}[1]{#1^{\mathsf{t}}}                % transpose
\newcommand{\at}[2]{#1_{#2}}                        % the image of an element of E at an infinite place
\newcommand{\Wsys}[3]{W^{#1}_{#2,#3}}               % the Bessel function of side #1 of the system of S at the place #3
\newcommand{\Isys}[2]{I_{#2}(#1)}                  % the index set of the Bessel system of S at the place
\newcommand{\degsys}[2]{\deg_{#2}(#1)}             % its constant deg
\newcommand{\repsys}[2]{\mathrm{rep}_{#1,#2}}      % its representation
\newcommand{\Ginf}[2]{G_{\infty}(#1,#2)}           % the product of the groups G_1(s) over the infinite places of E
\newcommand{\fl}{\operatorname{fl}}
\newcommand{\lev}{\operatorname{lev}}
\newcommand{\idx}[1]{[#1]}
\newcommand{\Sp}{\operatorname{span}}
\newcommand{\rk}{\operatorname{rk}}
\newcommand{\id}{\mathrm{id}}
\newcommand{\Uni}{\mathrm{U}}
\newcommand{\Res}[1]{\mathbb{F}_{#1}}              % the residue field at a finite place
\newcommand{\IntE}{O_E}
\newcommand{\abs}[1]{\lvert #1\rvert}
\newcommand{\nrm}[1]{\lVert #1\rVert}

% where each declaration of the concordance is treated: \cs{key}{place}, the place naming a \label or a lettered item
\cs{openCellConst}{Definition~\ref{d:bconst}}
\cs{openCellConstpos}{Definition~\ref{d:bconst}}
\cs{bandConst}{Definition~\ref{d:bconst}}
\cs{besselSize}{Definition~\ref{d:bconst}}
\cs{InCut}{Definition~\ref{d:bconst}}
\cs{frameBox}{Definition~\ref{d:bconst}}
\cs{gramDiscReal}{Definition~\ref{d:bconst}}
\cs{cellMass}{Definition~\ref{d:bconst}}
\cs{besselCorrelation}{Definition~\ref{d:schur}}
\cs{SchurLe}{Definition~\ref{d:schur}}
\cs{HasSchurBound}{Definition~\ref{d:schur}}
\cs{normCorrelation}{Definition~\ref{d:schur}}
\cs{HasBandLimits}{Definition~\ref{d:schur}}
\cs{Side}{Definition~\ref{d:bsys}}
\cs{IsCutPair}{Definition~\ref{d:bsys}}
\cs{BesselSystem}{Definition~\ref{d:bsys}}
\cs{reductionapplycongr}{Definition~\ref{d:fixers}}
\cs{existsreduceAboveGeq}{Lemma~\ref{l:char}}
\cs{vectorFixer}{Definition~\ref{d:fixers}}
\cs{FixesNormAtLeast}{Definition~\ref{d:fixers}}
\cs{fixingSet}{Definition~\ref{d:fixers}}
\cs{zFormPair}{Definition~\ref{d:nf}}
\cs{commonLink}{Definition~\ref{d:nf}}
\cs{HasCommonFlagVector}{Definition~\ref{d:nf}}
\cs{IsNormalForm}{Definition~\ref{d:nf}}
\cs{ratGram}{Definition~\ref{d:levels}}
\cs{linkForm}{Definition~\ref{d:levels}}
\cs{classCoord}{Definition~\ref{d:classbasis}}
\cs{classCoordX}{Definition~\ref{d:classbasis}}
\cs{countWeight}{Definition~\ref{d:levels}}
\cs{HasLevel}{Definition~\ref{d:levels}}
\cs{CongruentTerms}{Definition~\ref{d:levels}}
\cs{termDiscNorm}{Definition~\ref{d:levels}}
\cs{boxedPoints}{Definition~\ref{d:levels}}
\cs{IsNear}{Definition~\ref{d:levels}}
\cs{CharsAgree}{Definition~\ref{d:levels}}
\cs{latMiddleGram}{Definition~\ref{d:latform}}
\cs{middleBaseChange}{Definition~\ref{d:latform}}
\cs{besselLatticeBasis}{Definition~\ref{d:latform}}
\cs{frameDisc}{Definition~\ref{d:fdisc}}
\cs{grambesselLatticeBasis}{Lemma~\ref{l:latgram}}
\cs{framePair}{Definition~\ref{d:link}}
\cs{isotropicFramePair}{Definition~\ref{d:link}}
\cs{linkMatrix}{Definition~\ref{d:link}}
\cs{link}{Definition~\ref{d:link}}
\cs{dampedPower}{Definition~\ref{d:head}}
\cs{headWeight}{Definition~\ref{d:head}}
\cs{SieveSystem}{Definition~\ref{d:sieve}}
\cs{pi}{Proposition~\ref{p:pi}}
\cs{stepsum}{Lemma~\ref{l:stepsum}}
\cs{nonnegandlowerSemicontinuous}{Proposition~\ref{p:lsc}}
\cs{finHeadWeight}{Definition~\ref{d:head}}
\cs{isUnimodularlatMiddleGram}{Lemma~\ref{l:latgram}}
\cs{frameImageonemulVec}{Definition~\ref{d:link}}
\cs{frameImagemulVec}{Definition~\ref{d:link}}
\cs{existsstarvecMuleq}{Definition~\ref{d:link}}
\cs{existsmembesselUeqeichlerMap}{Proposition~\ref{p:eichler}}
\cs{finAdeleMat}{Definition~\ref{d:sublattice}}
\cs{finAdeleMatconjTranspose}{Definition~\ref{d:sublattice}}
\cs{mapconjTranspose}{Definition~\ref{d:sublattice}}
\cs{finAdeleMatsmul}{Definition~\ref{d:sublattice}}
\cs{muleqzeroofrangeleker}{Lemma~\ref{l:extend}}
\cs{existsbasisextendker}{Lemma~\ref{l:extend}}
\cs{sublatticeBasis}{Definition~\ref{d:sublattice}}
\cs{sublatticeDisc}{Definition~\ref{d:sublattice}}
\cs{sublatticeDiscmulindexsqle}{Proposition~\ref{p:discindex}}
\cs{continuousbesselFn}{Lemma~\ref{l:kernel}}
\cs{continuousbesselKernel}{Lemma~\ref{l:kernel}}
\cs{enormbesselKernel}{Lemma~\ref{l:kernel}}
\cs{enormamplitudele}{Lemma~\ref{l:kernel}}
\cs{isIntMatrixmiddleBaseChangemul}{Lemma~\ref{l:mbc}}
\cs{forallinclHFinsuboneiff}{Lemma~\ref{l:mbc}}
\cs{isIntMatrixscaleAtsubone}{Lemma~\ref{l:mbc}}
\cs{inclHFinmemlatStabilizer}{Theorem E}
\cs{reductioninclHFineqoneiff}{Lemma~\ref{l:redinc}}
\cs{classSumeqsum}{Lemma~\ref{l:redinc}}
\cs{cardresFieldeqlevelPrime}{Lemma~\ref{l:redinc}}
\cs{mapsplitGramlatMiddle}{Lemma~\ref{l:scales}}
\cs{mapmemunitriangularOn}{Lemma~\ref{l:scales}}
\cs{toLatBasismembesselU}{Lemma~\ref{l:dictUH}}
\cs{ratToFinmembesselU}{Lemma~\ref{l:dictUH}}
\cs{inclHFinmembesselH}{Lemma~\ref{l:dictUH}}
\cs{rangeinclHFin}{Lemma~\ref{l:dictUH}}
\cs{ratToFintoBesselH}{Lemma~\ref{l:dictUH}}
\cs{classBasis}{Definition~\ref{d:classbasis}}
\cs{coetermPosition}{Lemma~\ref{l:dictUH}}
\cs{termPositionmembesselU}{Lemma~\ref{l:dictUH}}
\cs{mul}{Lemma~\ref{l:dictUH}}
\cs{inva}{Lemma~\ref{l:dictUH}}
\cs{simpleRootSumadelicToFin}{Lemma~\ref{l:coords}}
\cs{simpleRootSumratToAdelic}{Lemma~\ref{l:coords}}
\cs{normselfDualChar}{Lemma~\ref{l:char}}
\cs{ratToInfmemunitriangularOn}{Lemma~\ref{l:scales}}
\cs{besselFnadelicToInfmul}{Theorem E}
\cs{besselFnratToInfmul}{Corollary~\ref{c:besselrat}}
\cs{isUnimodularlatMiddle}{Lemma~\ref{l:redab}}
\cs{transposesplitGram}{Lemma~\ref{l:redab}}
\cs{idOrInvAdjoint}{Definition~\ref{d:fixers}}
\cs{reduceAboveG}{Definition~\ref{d:fixers}}
\cs{intGLmapreduceAboveG}{Lemma~\ref{l:redab}}
\cs{gIdeleval}{Lemma~\ref{l:char}}
\cs{gIdelemem}{Lemma~\ref{l:char}}
\cs{forallselfDualCharmuleqoneiff}{Lemma~\ref{l:char}}
\cs{selfDualChareqoneofmem}{Lemma~\ref{l:char}}
\cs{invgIdelesmulmemiff}{Lemma~\ref{l:char}}
\cs{principalLevelequnitaryLevel}{Theorem E}
\cs{siegellemma}{Theorem C}
\cs{tsumsublatticeDiscinvle}{Theorem C}
\cs{besselUdense}{Theorem A}
\cs{unipotentCharone}{Lemma~\ref{l:char}}
\cs{existsbesselUapprox}{Theorem A}
\cs{unipotentChareqofmemprincipalLevel}{Lemma~\ref{l:charlevel}}
\cs{eqandmemclassUnipofbesselS}{Theorem D}
\cs{existsratapproxunipotentChareq}{Lemma~\ref{l:charlevel}}
\cs{existsbesselSmemtermsOver}{Theorem D}
\cs{eqoneofreduceAteqone}{Proposition~\ref{p:free}}
\cs{frameImagemul}{Lemma~\ref{l:common}}
\cs{frameImagemulofmembesselU}{Lemma~\ref{l:common}}
\cs{frameImageinvmuleqiff}{Lemma~\ref{l:common}}
\cs{starmullasteqofframeImageeq}{Lemma~\ref{l:common}}
\cs{frameProductsmulVecofframeImageeq}{Lemma~\ref{l:common}}
\cs{eqzeroofratMiddleformeqzero}{Lemma~\ref{l:common}}
\cs{normsqmiddleparteq}{Proposition~\ref{p:normid}}
\cs{maptermDisceqnorm}{Corollary~\ref{c:tdisc}}
\cs{hasLevellast}{Corollary~\ref{c:tdisc}}
\cs{gramDiscRealrealFrameGramratToInf}{Corollary~\ref{c:tdisc}}
\cs{ratToInfmembesselU}{Lemma~\ref{l:near}}
\cs{mulmemboxedPointsiff}{Lemma~\ref{l:near}}
\cs{ratToInf}{Lemma~\ref{l:near}}
\cs{isIntMatrixtermPosition}{Lemma~\ref{l:intterm}}
\cs{classCoordconjsplitGram}{Lemma~\ref{l:intterm}}
\cs{frameImagetermPosition}{Lemma~\ref{l:intterm}}
\cs{frameGramtermPosition}{Lemma~\ref{l:intterm}}
\cs{mulofmembesselU}{Lemma~\ref{l:near}}
\cs{membesselSiffframeImage}{Lemma~\ref{l:common}}
\cs{powletailScale}{Lemma~\ref{l:sizes}}
\cs{tailScalepowfour}{Lemma~\ref{l:sizes}}
\cs{ofsizebounds}{Lemma~\ref{l:sizes}}
\cs{normframeProductsle}{Proposition~\ref{p:nearsize}}
\cs{normmapleadingMinorle}{Proposition~\ref{p:nearsize}}
\cs{eqandmemclassUnip}{Theorem D}
\cs{isIntegralframeProducts}{Lemma~\ref{l:intterm}}
\cs{Filling}{Definition~\ref{d:filling}}
\cs{g}{Definition~\ref{d:groups}}
\cs{ratMiddle}{Definition~\ref{d:groups}}
\cs{latMiddle}{Definition~\ref{d:groups}}
\cs{G}{Definition~\ref{d:groups}}
\cs{GFin}{Definition~\ref{d:groups}}
\cs{GRes}{Definition~\ref{d:groups}}
\cs{BesselCoord}{Definition~\ref{d:groups}}
\cs{Classes}{Definition~\ref{d:groups}}
\cs{frameProducts}{Definition~\ref{d:groups}}
\cs{latStabilizer}{Definition~\ref{d:level}}
\cs{reduction}{Definition~\ref{d:level}}
\cs{principalLevel}{Definition~\ref{d:level}}
\cs{reducedImage}{Definition~\ref{d:level}}
\cs{kappa}{Definition~\ref{d:scales}}
\cs{basisScale}{Definition~\ref{d:scales}}
\cs{mapsplitGramratMiddle}{Lemma~\ref{l:scales}}
\cs{selfDualChar}{Definition~\ref{d:maps}}
\cs{varpinezero}{Lemma~\ref{l:scales}}
\cs{frameScaleSq}{Definition~\ref{d:scales}}
\cs{frameScaleProd}{Definition~\ref{d:scales}}
\cs{unipCovolume}{Definition~\ref{d:const}}
\cs{siegelConst}{Definition~\ref{d:const}}
\cs{termDisc}{Definition~\ref{d:groups}}
\cs{basisScalebesselFrame}{Lemma~\ref{l:scales}}
\cs{coekappa}{Lemma~\ref{l:scales}}
\cs{kappanezero}{Lemma~\ref{l:scales}}
\cs{frameScaleSqpos}{Lemma~\ref{l:scales}}
\cs{frameScaleProdpos}{Lemma~\ref{l:scales}}
\cs{basisScalenezero}{Lemma~\ref{l:scales}}
\cs{basisScaleGL}{Definition~\ref{d:scales}}
\cs{grambasisScaleGL}{Lemma~\ref{l:scales}}
\cs{latBasis}{Definition~\ref{d:maps}}
\cs{toLatBasis}{Definition~\ref{d:maps}}
\cs{ratToAdelic}{Definition~\ref{d:maps}}
\cs{adelicToInf}{Definition~\ref{d:maps}}
\cs{adelicToFin}{Definition~\ref{d:maps}}
\cs{ratToInfa}{Definition~\ref{d:maps}}
\cs{ratToFin}{Definition~\ref{d:maps}}
\cs{inclHFin}{Definition~\ref{d:maps}}
\cs{h}{Definition~\ref{d:maps}}
\cs{termPosition}{Definition~\ref{d:terms}}
\cs{termsOver}{Definition~\ref{d:terms}}
\cs{classUnip}{Definition~\ref{d:terms}}
\cs{classUnipa}{Definition~\ref{d:terms}}
\cs{boxSize}{Definition~\ref{d:amp}}
\cs{IsBoxPartition}{Definition~\ref{d:amp}}
\cs{besselFn}{Definition~\ref{d:amp}}
\cs{besselKernel}{Definition~\ref{d:amp}}
\cs{columnDim}{Definition~\ref{d:const}}
\cs{besselDegree}{Definition~\ref{d:const}}
\cs{besselDegreepos}{Lemma~\ref{l:scales}}
\cs{mainEnergy}{Definition~\ref{d:const}}
\cs{unipOrder}{Definition~\ref{d:const}}
\cs{packetConst}{Definition~\ref{d:const}}
\cs{besselEnergy}{Definition~\ref{d:const}}
\cs{tailScale}{Definition~\ref{d:const}}
\cs{codimFactor}{Definition~\ref{d:const}}
\cs{IsAdmissible}{Definition~\ref{d:adm}}
\cs{adelicToInfapply}{Lemma~\ref{l:coords}}
\cs{ratToInfapply}{Lemma~\ref{l:coords}}
\cs{coeratToFin}{Lemma~\ref{l:coords}}
\cs{finiteclassUnipintercompact}{Lemma~\ref{l:cutoff}}
\cs{existscutoffclassUnip}{Lemma~\ref{l:cutoff}}
\cs{classCutoff}{Definition~\ref{d:amp}}
\cs{amplitude}{Definition~\ref{d:amp}}
\cs{amplitudeAbs}{Definition~\ref{d:amp}}
\cs{amplitudeSum}{Definition~\ref{d:amp}}
\cs{abskappaeq}{Lemma~\ref{l:scales}}
\cs{realFrameGramratToInf}{Lemma~\ref{l:coords}}
\cs{detsubmatrixrealFrameGramratToInf}{Lemma~\ref{l:coords}}
\cs{unipCovolumepos}{Lemma~\ref{l:scales}}
\cs{gIdele}{Definition~\ref{d:level}}
\cs{unipotentChar}{Definition~\ref{d:maps}}
\cs{Agrees}{Definition~\ref{d:terms}}
\cs{IsStable}{Definition~\ref{d:terms}}
\cs{termsOfLink}{Definition~\ref{d:terms}}
\cs{splitGrammulVec}{Lemma~\ref{l:memlink}}
\cs{linkMatrixmulVec}{Lemma~\ref{l:memlink}}
\cs{isotropicFramePairmulVeceqzeroiff}{Lemma~\ref{l:memlink}}
\cs{memlink}{Lemma~\ref{l:memlink}}
\cs{invb}{Lemma~\ref{l:nf}}
\cs{existsframeImageeqofisNormalForm}{Lemma~\ref{l:nf}}
\cs{besselUmul}{Lemma~\ref{l:nf}}
\cs{linkbesselUmul}{Theorem B}
\cs{frameDiscframeGrambesselUmul}{Lemma~\ref{l:nf}}
\cs{isStablehasLevel}{Corollary~\ref{c:stablelevel}}
\cs{memlinkifforthogonal}{Theorem B}
\cs{frameDiscweylElemmul}{Proposition~\ref{p:bigcell}}
\cs{linkbigcelleqiff}{Proposition~\ref{p:bigcell}}

\begin{document}
\color{ink}

{\raggedright
{\sffamily\bfseries\fontsize{7.2}{9}\selectfont\color{leancol}\addfontfeatures{LetterSpace=8}READBACK NOTE 41\par}
\vspace{5pt}
{\fontsize{21.5}{25.5}\selectfont\bfseries The filling data: Bessel subgroups, links of frames, Siegel's lemma and strong
approximation\par}
\vspace{5pt}
{\fontsize{12.4}{16}\selectfont\itshape\color{inksoft} One unit of the Lean artifact \ArtifactName, read back into
mathematics: Note S of the unit on the filling theorem\par}}
\vspace{8pt}
{\color{rulegrey}\hrule height 0.6pt}
\vspace{7pt}
\thispagestyle{empty}

\textbf{Summary.} The filling theorem behind Theorem~B of the write-up (\emph{enough forms}) is a second moment of
Poincar\'e series on a unitary group of signature $(\varrho+k,\varrho+1)$ at every infinite place of a CM field $E$,
built from Bessel functions along the subgroup $\mathcal S=\mathcal U\mathcal H$ that is unitriangular on a frame
of $\varrho$ isotropic vectors and one negative vector. This note sets the stage. It defines the record of the
\emph{filling data} with its standing hypotheses, the groups, double cosets, terms, cutoffs and amplitudes of the
second moment, and the abstract inputs from the real places (Bessel systems, sieve systems). It then proves the
tools of the later notes: the dictionary between rational, finite adelic and real coordinates; that the \emph{link}
of two frames is a Lagrangian subspace and an invariant of the double coset of $\mathcal U$; Siegel's lemma for
hermitian lattices, with the convergence of a height zeta function of subspaces; strong approximation for
$\mathcal U$; and that a term in the Bessel subgroup with trivial reduction is unipotent, by the freeness of the
classes.

\statement{leancol}{\fontsize{9.6}{13.2}\selectfont\textbf{What this note is.} The artifact \ArtifactName{} is a
body of Lean code whose final theorems state the Riemann hypothesis and its analogue for Dirichlet $L$-functions.
That is a claim: its authors report that Lean's kernel accepts the code, a replay of that check by the owner of
this write-up is recorded elsewhere, and nothing was compiled for this note; no mathematician has read the proof
through. This note treats the first part of one unit (the folder of the filling theorem): \UnitInNote{} of its
\UnitMembers{} declarations, the content of nine files. Three further notes of the same writer treat the rest of the
first half of the unit: Note~B (\emph{Box partitions of arithmetic unipotent groups}), Note~E (\emph{The energy of
amplitude sums}) and Note~U (\emph{Poincar\'e series of cut Bessel functions}); the second half has four notes of
another writer: Note~F (the filling theorem), Note~P (the pairing of links), Note~T (the top floor at a finite
place) and Note~G (from amplitude sums to local correlations). This note is Note~S; the eight notes cite each other
by these letters. Every statement was compared, argument by argument, with the Lean text. Proofs are complete at the
level of a research paper, with these exceptions, each said where it occurs: results of other units and of the
background library are quoted as \emph{Inputs}; identities between explicit matrices that the formal text checks
entry by entry are stated and not recomputed; and the argument lists that the census of the artifact reads from the
text differ from Lean's for some declarations (Section~\ref{s:notes}). The Lean text read is \PinText; Appendix~A is
the concordance of \ConcDecls{} declarations.}

\sect{The statements}{s:intro}
The notions are defined in Section~\ref{s:setting}; here is what they are in a few words. A \emph{shape}
$\fs=(\varrho,\kappa,k_0)$ fixes a hermitian space with basis $e_0,\dots,e_{\varrho-1}$; $x_i$ $(i\in\kappa)$; $z$;
$e^*_{\varrho-1},\dots,e^*_0$, in which $e_a$ and $e^*_a$ are dual isotropic vectors and the middle part, spanned by
the $x_i$ and $z$, carries a form $T$. The \emph{frame} is $(\fr{0},\dots,\fr{\varrho})=(e_0,\dots,e_{\varrho-1},z)$.
The group $\UT{T}$ is the unipotent radical of the stabilizer of the flag of the $e_a$, the group $\HT{T}$ fixes all
$e_a$, $e_a^*$ and $z$, and $\ST{T}=\UT{T}\HT{T}$ is the Bessel subgroup. Filling data $S$ over a CM field $E$ single
out the form $\Trat{S}{E}=\operatorname{diag}(1,\dots,1,-\NS{S})$, a lattice, two primes $\qS{S}{0},\qS{S}{1}$ with
product $\gS{S}$ and the principal congruence subgroup $\KgS{S}$ of level $\gS{S}$ in the finite adelic group
$\GFin{S}$.

\statement{arith}{\phantomsection\label{t:A}\textbf{Theorem A (strong approximation for the unipotent group).}\enspace
\emph{(a)} Let $E$ be a number field that is a CM field, $\fs$ a shape, $T$ an $\fs$-form over $E$ that is hermitian
and invertible, and $T'$ an $\fs$-form over the finite adeles $\AEf$ such that $\JT{T'}$ is the image of $\JT{T}$
under $E\to\AEf$. Then $\UT{T'}$ is contained in the closure, in $\GT{T'}$, of the image of $\UT{T}$ under the
homomorphism $\GT{T}\to\GT{T'}$ that embeds every entry.

\emph{(b)} Let $E$ be a CM field with a CM type, $\fs$ a shape and $S\in\Fill{E}{\fs}$. For every
$x\in\GFin{S}$ with $x\in\UT{\Tlat{S}{\AEf}}$ and every $y\in\GFin{S}$ there is a rational element
$u\in\UT{\Trat{S}{E}}$ such that $y^{-1}\,\bigl(x^{-1}\,\rfin{S}(u)\bigr)\,y\in\KgS{S}$.}

\statement{arith}{\phantomsection\label{t:B}\textbf{Theorem B (the link of two frames).}\enspace
\emph{(a)} Let $F$ be a commutative ring with involution, $\fs$ a shape, $T$ an $\fs$-form over $F$ and
$\gamma,u,u'\in\GT{T}$ with $u\in\UT{T}$ and $u'\in\UT{T}$. Then $\link{u^{-1}\gamma u'}=\link{\gamma}$.

\emph{(b)} Let $E$ be a CM field with a CM type, $\fs$ a shape, $S\in\Fill{E}{\fs}$, $\gamma\in\GS{S}{E}$ and
$p\in E^{\Mx{\fs}}\times E^{\Mx{\fs}}$. Then $p\in\link{\gamma}$ if and only if $\adj{q}\,\lform{S}\,p=0$ for every
$q\in\link{\gamma}$: the link is its own orthogonal complement for the form $\lform{S}$ of signature zero.}

\statement{arith}{\phantomsection\label{t:C}\textbf{Theorem C (Siegel's lemma for hermitian lattices; a height zeta
function).}\enspace Let $E$ be a CM field with a CM type that is stably unramified, $\kappa$ a finite set, and
$u\in\Gf{E}{\kappa}$ an element of the finite adelic unitary group of the standard form. Write
$\nrm{X}=\bigl(\sum_{\tau}\sum_{i\in\kappa}\abs{\at{(X_i)}{\tau}}^2\bigr)^{1/2}$ for $X\in E^\kappa$, the sum over
the infinite places $\tau$ of $E$, and $\sld{u^{-1}}{V}$ for the discriminant of the lattice of the adelic points of a
subspace $V\subseteq E^{\kappa}$ that lie in $u\,\hO^{\kappa}$ (Definition~\ref{d:sublattice}).

\emph{(a)} Let $V\subseteq E^\kappa$ be a subspace and $s$ a natural number with $\dim_EV+s=\#\kappa$. Then there
are vectors $X_1,\dots,X_s\in E^\kappa$, all non-zero and in $u\,\hO^\kappa$, such that $V$ is the set of the
$v\in E^\kappa$ with $\sum_{i\in\kappa}\cj{(X_j)_i}\,v_i=0$ for $j=1,\dots,s$, and
\[\Bigl(\prod_{j=1}^{s}\nrm{X_j}\Bigr)^{2\ninf{E}}\ \le\ \bigl(8\,\ninf{E}\,s\,\rdE{E}\bigr)^{\ninf{E}s}\cdot\sld{u^{-1}}{V}.\]

\emph{(b)} For every natural number $s$,
\[\sum_{\substack{V\subseteq E^\kappa\\ \dim_EV+s=\#\kappa}}\sld{u^{-1}}{V}^{-(\#\kappa+1)}\ \le\
\bigl(200\,\ninf{E}\,\#\kappa\,\rdE{E}\bigr)^{\ninf{E}(\#\kappa+1)s}.\]}

\statement{arith}{\phantomsection\label{t:D}\textbf{Theorem D (terms in the Bessel subgroup).}\enspace Let $E$ be a CM
field with a CM type, $\fs$ a shape and $S\in\Fill{E}{\fs}$ admissible filling data.

\emph{(a)} Let $c,d\in\Cl{S}$ and $\gamma\in\GS{S}{E}$ with $\gamma\in\ST{\Trat{S}{E}}$ and
$\pos{S}{c}{d}{\gamma}\in\KgS{S}$. Then $c=d$ and $\gamma\in\clU{S}{c}$.

\emph{(b)} Let $c\in\Cl{S}$ and $s\in\redim{S}{\ST{\Tlat{S}{\AEf}}}$. Then there are a class $d\in\Cl{S}$ and an
element $\gamma\in\ST{\Trat{S}{E}}$ with $\gamma\in\terms{S}{c}{s}{d}$.

\emph{(c)} Let $c,c',d\in\Cl{S}$, $s\in\GRes{S}$ and $\gamma,u\in\GS{S}{E}$ with $\gamma\in\terms{S}{c}{s}{c'}$,
$u\in\terms{S}{c}{s}{d}$ and $\gamma\fr{a}=u\fr{a}$ for all frame vectors. Then $c'=d$ and
$\gamma^{-1}u\in\clUo{S}{c'}$.}

\statement{arith}{\phantomsection\label{t:E}\textbf{Theorem E (the dictionary).}\enspace Let $E$ be a CM field with a
CM type, $\fs$ a shape and $S\in\Fill{E}{\fs}$.

\emph{(a)} $\KgS{S}$ is the group of the $k\in\GFin{S}$ such that $k$ and $k^{-1}$ have entries in $\hO$ and all
entries of $\gid{S}^{-1}(k-1)$ lie in $\hO$.

\emph{(b)} Let $k\in\Gf{E}{\kappa}$, where $\kappa$ is the index set of $\fs$. Then $\incH{S}(k)\in\KS{S}$ if and only
if $k$ and $k^{-1}$ have entries in $\hO$ and the row $k_0$ of $k-1$ has entries in $\piS{S}\,\hO$.

\emph{(c)} Let $\sigma$ be a side, $x\in\GS{S}{\AdE}$ with $x\in\UT{\Trat{S}{\AdE}}$, and $y\in\Ginf{E}{\fs}$. Then, as
vectors indexed by $\Bc{S}$,
\[\besfn{\sigma}{S}{\ainf{S}(x)\,y}{\,\cdot\,}=
\frac{\psi_E\bigl(-(\gS{S}\,\vpS{S})^{-1}\sd{x}\bigr)}{\uchar{S}{\afin{S}(x)}}\ \besfn{\sigma}{S}{y}{\,\cdot\,},\]
where $\psi_E$ is the standard additive character of $\AdE$.}

Theorem~A is the classical strong approximation theorem for a unipotent group, in the form needed here (the
artifact names Platonov and Rapinchuk~\cite{pr} for it); part (b) feeds Theorem~D and the construction of
Poincar\'e series in Note~U. Theorem~B is used in Note~E to sort the terms of
the second moment by their links. Theorem~C counts subspaces by the discriminants of their lattices; the sum in (b)
is what makes the contribution of degenerate terms small. Theorem~D identifies the terms that lie in the Bessel
subgroup, and Theorem~E translates between the three systems of coordinates in which the data are written.

\sect{The setting}{s:setting}
\begin{conventions}\phantomsection\label{c:conventions}
Rings are commutative with $1$. A \emph{ring with involution} is a commutative ring with a ring automorphism
$x\mapsto\cj{x}$ of order at most two; letters $R$, $R'$, $F$, $A$, $C$ denote such rings (each statement says which
structure it uses), and $\adj{M}$ is the conjugate transpose of a matrix $M$, $\tr{M}$ its transpose. All index sets
$\kappa$, $\Omega$, $P$, $m$, $n$, $I$ of matrices are finite. A \emph{shape} is denoted $\fs$, an $\fs$-form $T$ or
$T'$, and $\gamma,\gamma',\delta,u,u',s,s',x,y,\sigma_0,\sigma_1$ are elements of the groups named in the statements;
$p,q$ are sequences of units; $j$, $l$, $a$, $a'$, $i$ are indices and $c,c',c'',d$ classes or ring elements, as each
statement says. $E$ is a number field that is a CM field: a totally complex quadratic extension of a totally real
field, with complex conjugation $x\mapsto\cj{x}$. A \emph{CM type} of $E$ is the choice, for every infinite place
$\tau$ of $E$, of an embedding $E\to\CC$ inducing $\tau$; the image of $x\in E$ is written $\at{x}{\tau}$. When $E$
has a CM type we say \emph{CM field with a CM type}. $S$ always denotes filling data over $E$ for $\fs$
(Definition~\ref{d:filling}); in statements about $S$ the field $E$ and the shape $\fs$ are those of $S$. A
\emph{side} is one of the two words $\full$, $\cut$, and $\sigma$ denotes a side. Measures on locally compact groups
are Haar measures, fixed by the formal text; integrals of functions with values in $[0,\infty]$ are Lebesgue
integrals with values in $[0,\infty]$. The letters $S$, $s$, $\sigma$, $q$, $\psi$, $\zeta$ have in the notes of this
unit the meanings given here, not those of the overview of the write-up (there: a set of columns, the complex
variable, a field automorphism, a rational number, a column, Riemann's function). The letters $D$, $N$, $O$, $P$, $d$, $f$, $r$, $w$, $x$, $\mu$, $\nu$, $\varsigma$ of
the handout-wide Conventions~0.1 to~0.3, 0.6 and~0.7 (the data, the level, the auxiliary pair, the parameters, the prime, the
number of infinite places, the rank, a finite place, the scale, the exponent, a character of the ideles, the chosen root of
unity) are used here only for the local objects that the notation list names, and each use says which: $\mu$ is a measure, $\nu$ the
Heisenberg element $\nu_l(v,\beta)$, $\varsigma$ the scalar of an Eichler map, $w$ is a finite place and, in
Proposition~\ref{p:discindex}, an invertible matrix. For the filling data of admissible
data (Note~G), the primes $\qS{S}{0}$, $\qS{S}{1}$ are the primes that the overview calls $g_0$, $g_1$, and $\gS{S}$ is
their product.
\end{conventions}

\subhead{Adeles.} For a CM field $E$ with ring of integers $\IntE$: $\AdE$ is the ring of adeles, $\AEf$ the ring of
finite adeles, $\hO\subset\AEf$ the closure of $\IntE$, a compact open subring. Complex conjugation acts on $\AEf$ and
on the set of finite places $w\mapsto\cj{w}$; $\Res{w}$ is the residue field at $w$, $q_w$ its number of elements, and
$\hO\to\Res{w}$ the reduction modulo $w$. The number of infinite places is $\ninf{E}$, the discriminant $\dE{E}$, the
root discriminant $\rdE{E}=\abs{\dE{E}}^{1/(2\ninf{E})}$, and $\volO{E}=2^{\ninf{E}}/\sqrt{\abs{\dE{E}}}$. The standard
additive character $\psi_f$ of $\AEf$ is $\psi_f(x)=e^{2\pi i\,\mathrm{Tr}_{E/\Q}(a)}$ for any $a\in E$ with $x-a\in\hO$, and
$\psi_E(x)=e^{-2\pi i\sum_\tau2\operatorname{Re} x_\tau}\,\psi_f(x_f)$ is the standard additive character of $\AdE$. For a unit
$a$ of $\AEf$ (or of a finite product of local fields), $\gix{a}$ is the \emph{generalized index}: the inverse of
the factor by which multiplication by $a$ scales Haar measure, equal to the index $[\hO:a\hO]$ when $a\in\hO$; it is
$0$ at non-units. $E$ is \emph{stably unramified} if it contains a primitive root of unity of order $5040$.

\begin{setting}[shapes, forms, groups; background]\label{set:shape}
A \emph{shape} $\fs=(\varrho,\kappa,k_0)$ is a natural number $\varrho$, a finite set $\kappa$ and an element
$k_0\in\kappa$. Put $\kk{\fs}=\#\kappa$, $\nn{\fs}=\varrho+1$, $\Mx{\fs}=\kappa\sqcup\{z\}$ (the \emph{middle
indices}), and let $\Ix{\fs}$ be the disjoint union of $\{0,\dots,\varrho-1\}$, $\Mx{\fs}$ and a second copy of
$\{0,\dots,\varrho-1\}$; its $\rr{\fs}=2\varrho+\kk{\fs}+1$ elements index a basis written
\[e_0,\dots,e_{\varrho-1};\qquad x_i\ (i\in\kappa),\ z;\qquad e^*_{\varrho-1},\dots,e^*_0 .\]
The \emph{frame} is $\fr{a}=e_a$ for $a<\varrho$ and $\fr{\varrho}=z$. The \emph{level} of a basis vector is
$\lev(e_a)=a$, $\lev(x_i)=\lev(z)=\varrho$, $\lev(e_a^*)=2\varrho-a$, and its \emph{floor} is
$\fl(e_a)=\fl(e_a^*)=a$, $\fl(z)=\varrho$, $\fl(x_i)=\varrho+1$.

Let $R$ be a ring. An \emph{$\fs$-form} over $R$ is a square matrix $T$ over $R$ indexed by $\Mx{\fs}$; for $N\in R$,
$\Trg{\fs}{N}$ is the diagonal $\fs$-form with entries $1$ on $\kappa$ and $-N$ at $z$. The \emph{split Gram matrix}
$\JT{T}$, indexed by $\Ix{\fs}$, has the entries $1$ at $(e_a,e_a^*)$ and $(e_a^*,e_a)$, the block $T$ on $\Mx{\fs}$,
and $0$ elsewhere. If $R$ has an involution, $\GT{T}$ is the group of the invertible matrices $g$ over $R$ indexed by
$\Ix{\fs}$ with $\adj{g}\,\JT{T}\,g=\JT{T}$, and $h_T(v,w)=\adj{v}\,\JT{T}\,w$. For the form $\Trg{\fs}{1}$ over $\CC$ the
group is $\Gone{\fs}$, a unitary group of signature $(\varrho+\kk{\fs},\varrho+1)$. Subgroups of $\GT{T}$:
\begin{itemize}
\item $\UT{T}$: the $g$ with $g_{ij}=\delta_{ij}$ whenever $\lev(j)\le\lev(i)$ (unitriangular for the blocks of equal
level);
\item $\UoT{T}$: the $g\in\UT{T}$ with $g\fr{a}=\fr{a}$ for all $a\le\varrho$;
\item $\Hl{l}{T}$, for a natural number $l$: the $g$ whose column $j$ is that of $1$ whenever $\fl(j)<l$; and
$\HT{T}=\Hl{\varrho+1}{T}$, which fixes all $e_a$, $e_a^*$ and $z$;
\item $\ST{T}$, the \emph{Bessel subgroup}: the $g$ with $g_{i\,\fr{a}}=\delta_{i\,\fr{a}}$ for all $a\le\varrho$ and
all $i$ with $\lev(i)\ge\lev(\fr{a})$, that is, $g\fr{a}-\fr{a}\in\Sp(e_{a'}:a'<a)$.
\end{itemize}
For a matrix $\gamma$ indexed by $\Ix{\fs}$, the \emph{frame Gram matrix} is the square matrix of size $\varrho+1$
\[\ZT{T}{\gamma}=\bigl((\JT{T}\gamma)_{\fr{a},\fr{a'}}\bigr)_{a,a'\le\varrho}=\bigl(h_T(\fr{a},\gamma\fr{a'})\bigr)_{a,a'},\]
and for $y\in\Gone{\fs}$ we write $\Zr{y}$ for the frame Gram matrix of $y$ for $\Trg{\fs}{1}$. For a square matrix
$Z$ and $l$ at most its size, $\lmin{l}{Z}$ is the determinant of its upper left $l\times l$ block. For
$s\in\GT{T}$, $\sd{s}=\sum_{a<\varrho}s_{\fr{a},\fr{a+1}}$. For $v\in R^{\varrho+1}$ we write $\mathfrak f v=\sum_av_a\fr{a}$
and $(\gamma\mathfrak f)v=\gamma(\mathfrak fv)$. For a sequence $p$ of units of $R$ and $l\le\varrho$, the \emph{Weyl
element} $\weyl{l}{p}\in\GT{T}$ maps $e_a\mapsto p_a e_a^*$ and $e_a^*\mapsto\cj{p_a}^{-1}e_a$ for $a<l$ and fixes the
other basis vectors.
\end{setting}

The following facts about these groups are used throughout; the formal text proves them in its background library.

\statement{ana}{\textbf{Inputs (groups).}\enspace In Setting~\ref{set:shape}, let $R$ be a ring with involution and $T$
an $\fs$-form over $R$.
(G1), unit~112: if $d$ is an invertible matrix that is block upper triangular for the levels and $u$ is
unitriangular for the levels, then $d^{-1}ud$ is unitriangular for the levels.
(G2), unit~1: $\ST{T}$ normalizes $\UT{T}$; $\sd{hu}=\sd{h}+\sd{u}$ for $h,u\in\ST{T}$; $\sd{h}=0$ for $h\in\HT{T}$.
(G3), units~119 and~1: $\UT{T}\subseteq\ST{T}$ and $\HT{T}\subseteq\ST{T}$; if $\JT{T}$ is hermitian and invertible, then
$s\in\ST{T}$ if and only if $s=uh$ with $u\in\UT{T}$ and $h\in\HT{T}$. For $T=\Trg{\fs}{N}$ and $g$ in the unitary
group $\Uni_\kappa(R)$ of the identity matrix on $\kappa$, let $\iota_N(g)\in\GT{T}$ act as $g$ on the $x_i$ and fix the
other basis vectors; then $\iota_N(g)\in\HT{T}$, and if $N$ is a unit every $h\in\HT{T}$ is $\iota_N(g)$ for some $g$; if
moreover $\JT{T}$ is hermitian and invertible and $R$ has an element $t$ with $t+\cj{t}=1$, every $s\in\ST{T}$ is
$u\,\iota_N(g)$ with $u\in\UT{T}$ and $g\in\Uni_\kappa(R)$.
(G4), unit~119: for $u,u'\in\UT{T}$ and $\gamma\in\GT{T}$, every leading minor satisfies
$\lmin{l}{\ZT{T}{u^{-1}\gamma u'}}=\lmin{l}{\ZT{T}{\gamma}}$; and $\ZT{T}{\gamma}=\adj{\mathfrak f}\,\JT{T}\,\gamma\mathfrak f$.
(G5), unit~120: for $u\in\UT{T}$, $u\fr{a'}-\fr{a'}\in\Sp(\fr{a}:a<a')$ for all $a'\le\varrho$ (the matrix of $u$ on the frame is
upper unitriangular), and the entries of $u$ in the rows of $\Mx{\fs}$ and of the $e_a^*$ and the columns of the $e_a$
and of $\Mx{\fs}$ are those of $1$: $u$ preserves $X=\Sp(e_a)$ and acts trivially on $X^{\perp}/X$.
(G6), unit~29: for $g\in\GT{T}$, $\JT{T}\,g^{-1}=\adj{g}\,\JT{T}$ and $h_T(gv,gw)=h_T(v,w)$.}

\subhead{The lattice form and the lattice basis.}
\begin{definition}[lattice form, lattice basis]\label{d:latform}
Let $\fs=(\varrho,\kappa,k_0)$ be a shape and $R$ a ring.

\emph{(i)} For $N\in R$, $\Tlatg{\fs}{N}$ is the $\fs$-form with diagonal entries $1$ at the $i\in\kappa$, $i\ne k_0$, $0$
at $k_0$ and $-N$ at $z$, with entries $-1$ at $(k_0,z)$ and $(z,k_0)$, and $0$ elsewhere.

\emph{(ii)} Let $R$ have an involution and let $p$ be a unit of $R$. The matrix $\mbc{\fs}{p}$, with rows indexed by
$\Mx{\fs}$ and columns by $\kappa$, has the rows $i\in\kappa$ of the diagonal matrix with $\cj{p}$ at $k_0$ and $1$
elsewhere, and its row $z$ has the entry $-p^{-1}$ in column $k_0$ and $0$ elsewhere.

\emph{(iii)} With $R$ and $p$ as in (ii), $\blb{\fs}{p}$ is the invertible matrix indexed by $\Ix{\fs}$ that is the
identity on the $e_a$ and $e_a^*$ and whose middle block has the columns $x_i$ ($i\ne k_0$),
$v=\cj{p}^{-1}x_{k_0}+(p\cj{p})^{-1}z$ in the place of $x_{k_0}$, and $z$. Its inverse has the middle block
$(\mbc{\fs}{p}\mid z)$: in the basis $x_i$ $(i\neq k_0)$, $v$, $z$ one has $x_{k_0}=\cj{p}\,v-p^{-1}z$.
\end{definition}
\leanline{\lean{latMiddleGram}\sep \lean{middleBaseChange}\sep \lean{besselLatticeBasis}}

So $\Tlatg{\fs}{N}$ is the Gram matrix of the form $\Trg{\fs}{N}$ in the basis $x_i$ $(i\ne k_0)$, $v$, $z$ when
$N=p\cj{p}$: the vector $v$ is isotropic, $h(v,z)=-1$ and $h(z,z)=-N$.

\begin{lemma}\label{l:latgram}
Let $\fs$ be a shape.
\emph{(a)} Let $R$ be a ring with involution, $p$ a unit of $R$ and $N\in R$ with $p\,\cj{p}=N$. Then
\[\adj{\blb{\fs}{p}}\ \JT{\Trg{\fs}{N}}\ \blb{\fs}{p}=\JT{\Tlatg{\fs}{N}}.\]
\emph{(b)} Let $A$ be a ring with involution, $R\subseteq A$ a subring and $N\in A$ with $N\in R$ and $\cj{N}=N$. Then
$\Tlatg{\fs}{N}$ is \emph{unimodular over $R$}: it is hermitian, and it and its inverse have entries in~$R$.
\end{lemma}
\begin{proof}
(a) Both sides are the identity pairing on the $e_a$, $e_a^*$, so only the middle block matters. There the columns
of $\blb{\fs}{p}$ are $x_i$ $(i\neq k_0)$, $v$, $z$, and $h(x_i,x_j)=\delta_{ij}$, $h(x_i,v)=0$ for $i\ne k_0$,
$h(x_i,z)=0$, $h(v,v)=(p\cj{p})^{-1}-N(p\cj{p})^{-2}=0$, $h(v,z)=-N(p\cj{p})^{-1}=-1$, $h(z,z)=-N$. (b) The form is
symmetric with entries $0,\pm1,-N$ fixed by the involution. Its inverse is the matrix with $1$ at the diagonal
places $i\ne k_0$, with $N$ at $(k_0,k_0)$, $-1$ at $(k_0,z)$ and $(z,k_0)$, and $0$ at $(z,z)$, that is
$\Tlatg{\fs}{0}$ plus $N$ at $(k_0,k_0)$: on the plane of $k_0$ and $z$ the product of
$\bigl(\begin{smallmatrix}0&-1\\-1&-N\end{smallmatrix}\bigr)$ and
$\bigl(\begin{smallmatrix}N&-1\\-1&0\end{smallmatrix}\bigr)$ is $1$. Both have entries in $R$.
\end{proof}
\leanline{\lean{grambesselLatticeBasis}\sep \lean{isUnimodularlatMiddleGram}}

\subhead{The frame discriminant, frame pairs and the link.}
\begin{definition}[frame discriminant]\label{d:fdisc}
Let $R$ be a ring with involution, $\fs$ a shape, $N\in R$ and $Z$ a square matrix over $R$ of size $\varrho+1$. Then
\[\fdisc{N}{Z}=\det Z\cdot\cj{\det Z}-N^2\,\lmin{\varrho}{Z}\cdot\cj{\lmin{\varrho}{Z}} .\]
\end{definition}
\leanline{\lean{frameDisc}}

\begin{definition}[frame pairs, link]\label{d:link}
Let $R$ be a ring with involution, $\fs$ a shape, $T$ an $\fs$-form over $R$ and $\gamma\in\GT{T}$.

\emph{(i)} $\fpair{\gamma}$ is the matrix with $2(\varrho+1)$ columns $\fr{0},\dots,\fr{\varrho},\gamma\fr{0},\dots,\gamma\fr{\varrho}$,
and $\ifpair{\gamma}$ the matrix with the $2\varrho$ columns $e_0,\dots,e_{\varrho-1},\gamma e_0,\dots,\gamma e_{\varrho-1}$.

\emph{(ii)} $\lmat{\gamma}$ is the matrix of the linear map $R^{\Ix{\fs}}\to R^{\Mx{\fs}}\times R^{\Mx{\fs}}$,
$t\mapsto(t_{\mathrm{mid}},(\gamma^{-1}t)_{\mathrm{mid}})$, where $t_{\mathrm{mid}}$ denotes the coordinates of $t$ at
the middle indices.

\emph{(iii)} The \emph{link} of $\gamma$ is the submodule
\[\link{\gamma}=\bigl\{\lmat{\gamma}\,t\ :\ t\in R^{\Ix{\fs}},\ \adj{\ifpair{\gamma}}\,\JT{T}\,t=0\bigr\}\ \subseteq\ R^{\Mx{\fs}}\times R^{\Mx{\fs}}.\]
\end{definition}
\leanline{\lean{framePair}\sep \lean{isotropicFramePair}\sep \lean{linkMatrix}\sep \lean{link}}

With $X=\Sp(e_0,\dots,e_{\varrho-1})$, the condition on $t$ says $t\perp X$ and $t\perp\gamma X$, and for such $t$ both
$t$ and $\gamma^{-1}t$ lie in $X^\perp$, whose quotient by $X$ is the middle space. So $\link{\gamma}$ is the set of the
pairs (class of $t$, class of $\gamma^{-1}t$) in $(X^\perp/X)^2$, for $t\perp X+\gamma X$. Three unfoldings are used
below without mention: $\mathfrak f v$ has the $e$-coordinates $v_0,\dots,v_{\varrho-1}$, the $z$-coordinate
$v_\varrho$ and no other; $(\gamma\mathfrak f)v=\gamma(\mathfrak fv)$; and for a hermitian invertible matrix $J$ every
row vector $l$ is $\adj{v}J$ for $v=\adj{(lJ^{-1})}$, so every linear form is $h(v,\cdot)$.
\leanline{\lean{frameImageonemulVec}\sep \lean{frameImagemulVec}\sep \lean{existsstarvecMuleq}}

\subhead{Sieve systems.} The local theory at the real places and at the finite places produces, floor by floor, a
family of positive kernels. Their common shape is the following; it is used in Note~U and in the notes of the
second half.

\begin{definition}[damped powers, head weights]\label{d:head}
\emph{(i)} For a natural number $k$ and a real number $r$, $\damp{k}{r}=r^k\cdot\min(1,r)^6$.

\emph{(ii)} For natural numbers $K$, $i$ and a sequence $p=(p_0,p_1,\dots)$ of non-zero complex numbers,
$\hw{K}{i}{p}=\prod_{l<i}\damp{K}{\abs{p_0p_1\cdots p_l}^{-1}}$.

\emph{(iii)} Let $A$ be a finite product of non-archimedean local fields with its ring of integers $O_A$ (in the formal
text: a ring satisfying the axioms of such a product). For natural numbers $K$, $i$ and a sequence $p$ of units of
$A$, with $D_l=p_0p_1\cdots p_l$,
\[\fhw{K}{i}{p}=\prod_{l<i}\begin{cases}\sqrt{\gix{D_l}}^{\,K}&\text{if }D_l\in O_A,\\0&\text{otherwise.}\end{cases}\]
\end{definition}
\leanline{\lean{dampedPower}\sep \lean{headWeight}\sep \lean{finHeadWeight}}

For a sequence $p$, an index $i$ and an element $x$ we write $p[i\mapsto x]$ for the sequence $p$ with its entry
of index $i$ replaced by $x$.

\begin{definition}[sieve system]\label{d:sieve}
Let $C$ be a ring with a measure and $U$ a topological space. Let $u\colon C\to U$ and $\iota\colon U\to C$ be maps,
$n$ a natural number, $J=(J_i)_{i\in\NN}$ a family of functions $J_i\colon U^{\NN}\times U^{\NN}\to\CC$,
$\psi\colon C\to\CC$ a function, $\mathit{En}$ a real number, $\Lambda=(\Lambda_i)_{i\in\NN}$ a family of functions
$\Lambda_i\colon U^{\NN}\to\RR$ and $\beta\colon C\to\RR$ a function. We write
$\Sieve{u}{\iota}{n}{J}{\psi}{\mathit{En}}{\Lambda}{\beta}$, and say that these data are a \emph{sieve system} on
$C$, if:
\begin{enumerate}[label=(S\arabic*),leftmargin=3em]
\item $\psi(0)=1$;
\item for all $i\le n$, all $p,q\in U^{\NN}$ and all $\xi\in C$, the function
$f(c)=\psi(\xi c)\,J_{i+1}\bigl(p[i\mapsto u(c)],q[i\mapsto u(c)]\bigr)$ on $C$ is integrable; if $\xi=0$ then
$\int_Cf=J_i(p,q)$; and if $\xi$ is a unit of $C$ then $\abs{\int_Cf}\le\beta(\xi)\,\Lambda_i(p)\,\Lambda_i(q)$;
\item for every finite set $\Omega$, every family $(a_k)_{k\in\Omega}$ in $U^{\NN}$ and every $t\in U$ with
$(a_k)_n=t$ for all $k$, the matrix $\bigl(J_{n+1}(a_k,a_l)\bigr)_{k,l\in\Omega}$ is positive semidefinite;
\item $J_{n+1}$ is continuous on $U^{\NN}\times U^{\NN}$;
\item $\abs{J_0(p,q)}\le\mathit{En}$ for all $p,q$;
\item for all $i<n$ and all $p,q$ such that $\xi=\iota(p_i)-\iota(q_i)$ is a unit of $C$,
$\abs{J_{i+1}(p,q)}\le\beta(\xi)\,\Lambda_{i+1}(p)\,\Lambda_{i+1}(q)$.
\end{enumerate}
\end{definition}
\leanline{\lean{SieveSystem}}

So $J_i$ is obtained from $J_{i+1}$ by integrating over a common entry of index $i$ (S2, $\xi=0$), the top form is a
positive kernel in the remaining variables (S3), and non-trivial frequencies and unequal entries are small (S2, S6).
On $\CC$ the standard choices are $u(c)=c$ for $c\ne0$ and $u(0)=1$, with values in $\CC^\times$, which we write
$u_{\CC}$, and $\psi(c)=e^{c-\cj{c}}$, which we write $\mathbf e$.

\begin{proposition}[products of sieve systems]\label{p:pi}
Let $C$ be a ring with an $s$-finite measure, $U$ a topological space, $\Omega$ a finite set, and let
$K$ be a sieve system $\Sieve{u}{\iota}{n}{J}{\psi}{\mathit{En}}{\Lambda}{\beta}$ on $C$. Let
$\iota'\colon\CC^\times\to\CC$ be a map, $\Lambda'=(\Lambda'_i)_i$ a family of functions
$(\CC^{\times})^{\NN}\to\RR$, and for every $v\in\Omega$ let $J'_v$, $\mathit{En}'_v$ and
$\beta'_v\colon\CC\to\RR$ be such that $H_v$ is a sieve system
$\Sieve{u_{\CC}}{\iota'}{n}{J'_v}{\mathbf e}{\mathit{En}'_v}{\Lambda'}{\beta'_v}$ on $\CC$. Then
\[\Sieve{u^{+}}{\iota^{+}}{n}{J^{+}}{\psi^{+}}{\mathit{En}^{+}}{\Lambda^{+}}{\beta^{+}}\]
holds on the ring $\CC^{\Omega}\times C$ with the product measure, where, for $x=((x_v)_v,x_C)$ and for sequences
$p$ in $(\CC^\times)^\Omega\times U$ with components $p^{(v)}$ and $p^{(C)}$,
\begin{align*}
u^{+}(x)&=\bigl((u_{\CC}(x_v))_v,u(x_C)\bigr), & \iota^{+}(y)&=\bigl((\iota'(y_v))_v,\iota(y_U)\bigr),\\
J^{+}_i(p,q)&=\prod_vJ'_{v,i}\bigl(p^{(v)},q^{(v)}\bigr)\cdot J_i\bigl(p^{(C)},q^{(C)}\bigr), & \psi^{+}(x)&=\prod_v\mathbf e(x_v)\cdot\psi(x_C),\\
\Lambda^{+}_i(p)&=\prod_v\Lambda'_i\bigl(p^{(v)}\bigr)\cdot\Lambda_i\bigl(p^{(C)}\bigr), & \beta^{+}(x)&=\prod_v\beta'_v(x_v)\cdot\beta(x_C),
\end{align*}
and $\mathit{En}^{+}=\prod_v\mathit{En}'_v\cdot\mathit{En}$.
\end{proposition}
\begin{proof}
(S1) holds since $\mathbf e(0)=1$ and $\psi(0)=1$. A unit of $\CC^\Omega\times C$ is a tuple of units. For (S2), the
function attached to $(i,p,q,\xi)$ is the product, in independent variables, of the functions attached to
$(i,p^{(v)},q^{(v)},\xi_v)$ for $H_v$ and to $(i,p^{(C)},q^{(C)},\xi_C)$ for $K$. Each factor is integrable, hence so
is the product, and by Fubini's theorem (the measures are $s$-finite) its integral is the product of the integrals
of the factors. For $\xi=0$ this is $J^{+}_i(p,q)$; for a unit $\xi$ every factor is bounded by its own (S2), and
the product of the bounds is $\beta^{+}(\xi)\Lambda^{+}_i(p)\Lambda^{+}_i(q)$. For (S3), the matrix
$(J^{+}_{n+1}(a_k,a_l))_{k,l}$ is the entrywise product of the matrices of the factors, which are positive
semidefinite because the components of the $a_k$ share their entries of index $n$; by Schur's product theorem,
applied $\#\Omega$ times, it is positive semidefinite. (S4) holds because products of continuous functions composed
with projections are continuous. (S5) and (S6) follow by multiplying the bounds of the factors, all of them
non-negative; in (S6), $\iota^{+}(p_i)-\iota^{+}(q_i)$ is a unit exactly when all its components are.
\end{proof}
\leanline{\lean{pi}}

\begin{lemma}[one step, summed]\label{l:stepsum}
Let $C$ be a ring with a measure, $U$ a topological space, $\Omega$ a finite set, and let $K$ be a sieve system
$\Sieve{u}{\iota}{n}{J}{\psi}{\mathit{En}}{\Lambda}{\beta}$ on $C$. Let $i\le n$ and let $(a_p)_{p\in\Omega}$ be a family
in $U^{\NN}$. Put $F(c)=\sum_{p,q\in\Omega}J_{i+1}\bigl(a_p[i\mapsto u(c)],a_q[i\mapsto u(c)]\bigr)$. Then $F$ is
integrable, $\int_CF=\sum_{p,q}J_i(a_p,a_q)$, and for every unit $\xi$ of $C$,
$\abs{\int_C\psi(\xi c)F(c)\,dc}\le\beta(\xi)\bigl(\sum_{p}\Lambda_i(a_p)\bigr)^2$.
\end{lemma}
\begin{proof}
Sum (S2) over the pairs $(p,q)$, once with $\xi=0$, where $\psi(0)=1$ by (S1), and once with the given unit; note
$\sum_{p,q}\Lambda_i(a_p)\Lambda_i(a_q)=(\sum_p\Lambda_i(a_p))^2$.
\end{proof}
\leanline{\lean{stepsum}}

\begin{proposition}[positivity and semicontinuity on every floor]\label{p:lsc}
Let $C$ be a ring with a measure and a first countable topology, $U$ a topological space, $\Omega$ a finite set,
and let $K$ be a sieve system $\Sieve{u}{\iota}{n}{J}{\psi}{\mathit{En}}{\Lambda}{\beta}$ on $C$. Let $i\le n$, let
$b\colon C\to(U^{\NN})^{\Omega}$ and $c\colon C\to U$ be maps such that $(b(x)_p)_i=c(x)$ for all $x\in C$ and
$p\in\Omega$, and put $F(x)=\sum_{p,q\in\Omega}J_{i+1}\bigl(b(x)_p,b(x)_q\bigr)$. Then $F(x)$ is a non-negative real
number for every $x$, and $F$ is lower semicontinuous at every point $x$ at which all the maps
$x'\mapsto b(x')_p$ are continuous.
\end{proposition}
\begin{proof}
Downward induction on $i$, for all $b$ and $c$ at once. For $i=n$, $F(x)$ is the value at the vector
$(1,\dots,1)$ of the hermitian form with matrix $(J_{n+1}(b(x)_p,b(x)_q))_{p,q}$, which is positive semidefinite by
(S3), the sequences sharing the entry $c(x)$ of index $n$; and $F$ is continuous at $x$ by (S4), hence lower
semicontinuous there. Let $i<n$ and suppose the claim for $i+1$ and all families. For $z\in C$ put
$b^{z}(x)_p=b(x)_p[i+1\mapsto u(z)]$. These sequences share their entry $u(z)$ of index $i+1$, and
$x'\mapsto b^{z}(x')_p$ is continuous where $x'\mapsto b(x')_p$ is. By the induction hypothesis
$F^{z}(x)=\sum_{p,q}J_{i+2}\bigl(b^{z}(x)_p,b^{z}(x)_q\bigr)$ is non-negative for all $x$, and lower semicontinuous in
$x$ at the points in question. By Lemma~\ref{l:stepsum} with the index $i+1\le n$ and the family
$(b(x)_p)_p$, the function $z\mapsto F^{z}(x)$ is integrable and $\int_CF^{z}(x)\,dz=F(x)$. So $F(x)\ge0$. If
$x_k\to x$, then by Fatou's lemma and the semicontinuity of the $F^{z}$,
\begin{gather*}F(x)=\int_CF^{z}(x)\,dz\le\int_C\liminf_kF^{z}(x_k)\,dz\\
\le\liminf_k\int_CF^{z}(x_k)\,dz=\liminf_kF(x_k);\end{gather*}
since the topology of $C$ is first countable, this is lower semicontinuity at $x$. The hypothesis on the common
entry of index $i$ is used only for $i=n$.
\end{proof}
\leanline{\lean{nonnegandlowerSemicontinuous}}

\subhead{Eichler maps in the unipotent group.}
For a hermitian matrix $J$ with form $h(v,w)=\adj{v}Jw$, vectors $x$, $v$ and a scalar $\varsigma$, the \emph{Eichler
map} is $w\mapsto w+x\,h(v,w)-(v+\varsigma x)\,h(x,w)$.

\statement{ana}{\textbf{Input (E1)}, unit~119.\enspace If $J$ is hermitian, $h(x,x)=0$, $h(v,x)=0$ and
$\varsigma+\cj{\varsigma}=h(v,v)$, then the Eichler map of $(x,v,\varsigma)$ is an isometry of $h$.}

\begin{proposition}[Eichler maps along the frame]\label{p:eichler}
Let $A$ be a ring with involution that has an element $t$ with $t+\cj{t}=1$, $\fs$ a shape and $T$ an $\fs$-form
over $A$ such that $\JT{T}$ is hermitian and invertible. Let $a\le\varrho$ be a frame index, $j\in A^{\varrho+1}$ with
$j_{a'}=0$ for all $a'\ge a$, and $r\in A$. Let $v\in A^{\Ix{\fs}}$ and let $m=\adj{v}\,\JT{T}$ be the linear form
$h_T(v,\cdot)$, and suppose $m(\fr{a'})=0$ for all $a'<a$. Then there is $u\in\UT{T}$ with
\[\sd{u}=r\sum_{a'=1}^{\varrho}m(\fr{a'})\,j_{a'-1},\]
whose matrix is the Eichler map of $\bigl(r\,\mathfrak fj,\ v,\ t\,h_T(v,v)\bigr)$.
\end{proposition}
\begin{proof}
Put $x=r\,\mathfrak fj=r\sum_{a'<a}j_{a'}\fr{a'}$ and $\varsigma=t\,h_T(v,v)$. Since $a\le\varrho$, $x$ lies in
$\Sp(e_{a'}:a'<a)$, which is isotropic, and $h_T(v,x)=m(x)=r\sum_{a'<a}j_{a'}\,m(\fr{a'})=0$. As $\JT{T}$ is hermitian,
$h_T(v,v)$ is fixed by the involution, so $\varsigma+\cj{\varsigma}=(t+\cj{t})h_T(v,v)=h_T(v,v)$. By (E1) the Eichler
map $u$ lies in $\GT{T}$. The form $h_T(x,\cdot)$ vanishes on all basis vectors except the $e_{a'}^*$ with $a'<a$, and
the $e_{a''}^*$-coordinate of $v$ is $h_T(e_{a''},v)=\cj{m(\fr{a''})}=0$ for $a''<a$. So the matrix
$(v+\varsigma x)\,h_T(x,\cdot)$ has non-zero entries only in columns $e_{a'}^*$, $a'<a$, and in rows other than the
$e_{a''}^*$ with $a''<a$: these rows have smaller level than the column. The matrix $x\,m$ has non-zero entries only
in rows $e_c$, $c<a$, and its entry at a column of level $\le c$, a vector $e_{c'}$ with $c'\le c$, is
$x_c\,m(\fr{c'})=0$. Hence $u\in\UT{T}$. Finally $h_T(x,\fr{a'+1})=0$ for all $a'<\varrho$, so
$\sd{u}=\sum_{a'<\varrho}x_{\fr{a'}}\,m(\fr{a'+1})=r\sum_{a'<\varrho}j_{a'}\,m(\fr{a'+1})$.
\end{proof}
\leanline{\lean{existsmembesselUeqeichlerMap}}

\subhead{Bessel systems.} The inputs from the real places are packaged in a record. In this paragraph $\fs$ is a
shape, $n=\nn{\fs}$, $k=\kk{\fs}$, $r=\rr{\fs}$, and $\Gone{\fs}$ carries its Haar measure. The group $\Uni_\kappa(\CC)$ is embedded
in $\Gone{\fs}$ by $\iota_1$ (Input (G3)). For a finite set $I$ and a homomorphism $\mathrm{rep}$ from
$\Uni_\kappa(\CC)$ to the unitary group of $\CC^I$, a \emph{Bessel function for $\mathrm{rep}$} is a bounded measurable
function $\Phi\colon\Gone{\fs}\to\CC^I$ with $\Phi(uy)=\mathbf e(\sd{u})\,\Phi(y)$ for $u\in\UT{\Trg{\fs}{1}}$ and
$\Phi(\iota_1(h)y)=\mathrm{rep}(h)\Phi(y)$ for $h\in\Uni_\kappa(\CC)$ (background, unit~117). The measure
$\mu_{\mathcal U\backslash G}$ on $\Gone{\fs}$ is Haar measure times a chosen cutoff for the left action of
$\UT{\Trg{\fs}{1}}$ with its Haar measure, so that it integrates left $\mathcal U$-invariant functions over the
quotient (background, unit~118). The numbers $\vol\Uni(j)=\prod_{i<j}2\pi^{i+1}/i!$ are the volumes of the compact
unitary groups.

\begin{definition}[constants, sizes, boxes]\label{d:bconst}
Let $\fs$ be a shape, with $n$, $k$, $r$ as above.

\emph{(i)} $\occ{\fs}=\dfrac{\vol\Uni(k)}{\vol\Uni(r)}\cdot\dfrac{2\pi^{k}}{(k-1)!}$ and
$\bandc{\fs}=\dfrac{\vol\Uni(k)}{\vol\Uni(r)}\cdot\pi^{(2\varrho+1)k}$.

\emph{(ii)} For $y\in\Gone{\fs}$, $\bsize{y}=\inf_{s}\bigl(\nrm{sy}+1\bigr)$, the infimum over $s\in\ST{\Trg{\fs}{1}}$ of the
operator norms of the matrices $sy$. For a real number $b$, $\InCut{y}{b}$ means $\bsize{y}\le b^{6n+2}$.

\emph{(iii)} For a real number $\rho$, $\frbox{\rho}$ is the set of the $y\in\Gone{\fs}$ with
$\nrm{y^{-1}\fr{a}}\le\rho$ for all $a\le\varrho$ (Euclidean norms in $\CC^{\Ix{\fs}}$).

\emph{(iv)} For a complex square matrix $Z$ of size $\varrho+1$, $\gdr{Z}=\operatorname{Re}\fdisc{1}{Z}=\abs{\det Z}^2-\abs{\lmin{\varrho}{Z}}^2$.

\emph{(v)} For $y\in\Gone{\fs}$ and real numbers $\deg$, $\rho$, $d$, with $\mathsf g=\gdr{\Zr{y}}$,
\[\cellm{y}{\deg}{\rho}{d}=\Bigl(\frac{\rho^{9n}}{\mathsf g}\Bigr)^{10^5n^4}\cdot\frac{d^2}{\deg\cdot\occ{\fs}\cdot(\mathsf g/2)^{k-1}}.\]
\end{definition}
\leanline{\lean{openCellConst}\sep \lean{bandConst}\sep \lean{besselSize}\sep \lean{InCut}\sep \lean{frameBox}\sep
\lean{gramDiscReal}\sep \lean{cellMass}\sep \lean{openCellConstpos}}

The constant $\occ{\fs}$ is positive, being a quotient of positive numbers. For $k\ge1$ its second factor is
$\vol\Uni(k)/\vol\Uni(k-1)$.

\begin{definition}[correlations, Schur bounds, band limits]\label{d:schur}
Let $\fs$ be a shape, $I$ a finite set and $W\colon\Gone{\fs}\to\CC^I$ a function.

\emph{(i)} For a group $G$ with a measure $\mu$, a function $W\colon G\to\CC^I$ and $\gamma\in G$,
$\bcorr{\mu}{W}{\gamma}\in\CC^{I\times I}$ is the vector with the coordinates
$\int_GW(y)_i\,\cj{W(\gamma^{-1}y)_{i'}}\,d\mu(y)$, $(i,i')\in I\times I$.

\emph{(ii)} For $D\in[0,\infty]$, $\SchurLe{W}{D}$ means: for every integrable $\varphi\colon\Gone{\fs}\to\CC$,
\[\int\Bigl\lVert\int_{\Gone{\fs}}\varphi(y)\,W(\gamma y)\,dy\Bigr\rVert^2d\mu_{\mathcal U\backslash G}(\gamma)\ \le\
D\int_{\Gone{\fs}}\abs{\varphi(y)}^2\,dy .\]

\emph{(iii)} For a homomorphism $\mathrm{rep}$ from $\Uni_\kappa(\CC)$ to the unitary group of $\CC^I$, a real number
$B$ and $D\in[0,\infty]$, $\HasSB{W}{\mathrm{rep}}{B}{D}$ means: there are continuous functions $V_w\colon\Gone{\fs}\to\CC^I$,
one for each finite list $w$ of pairs of elements of $\Ix{\fs}$, with $V_{\emptyset}=W$, such that for every $w$:
$V_w$ is a Bessel function for $\mathrm{rep}$; for every curve $g\colon\RR\to\Gone{\fs}$, every matrix $X$ and every
real $t$ with $g'(t)=g(t)X$, the function $V_w\circ g$ has at $t$ the derivative
$\sum_{a,a'\in\Ix{\fs}}X_{aa'}\,V_{(a,a')w}(g(t))$; and if the length $\abs{w}$ is at most $2n^2$, then
$\SchurLe{V_w}{B^{2\abs{w}}D}$.

\emph{(iv)} For $y\in\Gone{\fs}$ and a sequence $p$ of non-zero complex numbers,
\[\ncorr{y}{W}{p}=\frac{\sqrt{\occ{\fs}}\,\sqrt{\abs{p_\varrho}^2-1}^{\,k-1}}{\#I}\ \prod_{a<\varrho}\abs{p_a}^{k+1+2(\varrho-a-1)}\cdot
\bcorr{\vol}{W}{\weyl{\varrho}{p}\,y},\]
with $\vol$ the Haar measure of $\Gone{\fs}$.

\emph{(v)} For real numbers $b$, $\deg$ and a family $J=(J_i)_i$ of functions on pairs of sequences of non-zero complex
numbers, $\BandL{b}{W}{\deg}{J}$ means:
$\Sieve{u_{\CC}}{\,x\mapsto x^{-1}\,}{\varrho}{J}{\mathbf e}{(1+b^{-10})/\deg}{\Lambda}{\beta}$ holds on $\CC$ with
$\Lambda_i(p)=\hw{n^6}{i}{p}$ and $\beta(\xi)=\bandc{\fs}\,b^{2000n^9}\,\damp{k-n^7}{\abs{\xi}^{-1}}$; and for every
$\sigma_0\in\Hl{\varrho}{\Trg{\fs}{1}}$ and all $p$, $q$ with $p_\varrho=\Zr{\sigma_0}_{\varrho\varrho}$ and $q_\varrho=p_\varrho$,
\[J_{n}(p,q)=\bigl\langle\ncorr{\sigma_0}{W}{p},\ \ncorr{\sigma_0}{W}{q}\bigr\rangle .\]
\end{definition}
\leanline{\lean{besselCorrelation}\sep \lean{SchurLe}\sep \lean{HasSchurBound}\sep \lean{normCorrelation}\sep \lean{HasBandLimits}}

In (iii) the function $V_{(a,a')w}$ is the derivative of $V_w$ along the elementary matrix $E_{aa'}$ acting on the
right; in (v), $\Zr{\sigma_0}_{\varrho\varrho}=h(z,\sigma_0z)$, and the inner product is that of $\CC^{I\times I}$,
conjugate linear in the first variable.

\begin{definition}[sides, cut pairs, Bessel systems]\label{d:bsys}
Let $\fs$ be a shape.

\emph{(i)} A \emph{side} is one of the two values $\full$ and $\cut$. In the formal text a side is an ambient
argument, with $\full$ as its default value.

\emph{(ii)} Let $I$ be a finite set, $b$ a real number, $f$ a natural number, $\mathrm{rep}$ a homomorphism from
$\Uni_\kappa(\CC)$ to the unitary group of $\CC^I$, and $W=(W^{\full},W^{\cut})$ a pair of continuous functions
$\Gone{\fs}\to\CC^I$. Then $\CutPair{b}{f}{\mathrm{rep}}{W}$ is the conjunction of: $W^{\full}$ and $W^{\cut}$ are Bessel functions
for $\mathrm{rep}$; $W^{\full}$ is integrable for $\mu_{\mathcal U\backslash G}$;
$\int\nrm{W^{\full}}^2d\mu_{\mathcal U\backslash G}=\#I$; $\int\nrm{W^{\cut}}^2d\mu_{\mathcal U\backslash G}\le\#I$;
$W^{\cut}(y)\ne0$ implies $\InCut{y}{b}$;
$\int\nrm{W^{\cut}-W^{\full}}^2d\mu_{\mathcal U\backslash G}\le(b^{-5r^2})^2\,\#I$; and there are a measurable
function $A\colon\Gone{\fs}\to[0,\infty]$, invariant under left translation by $\UT{\Trg{\fs}{1}}$, and a finite measure
$\mathsf m$ on $\Gone{\fs}$ with
\begin{align*}
\epsilon(y)\,\nrm{W^\sigma(y)}&\le A(y)\le\mathsf m\{g:\InCut{yg}{b}\}\quad\text{for all }y\text{ and both sides }\sigma,\\
\mathsf m(\Gone{\fs})\int A\,d\mu_{\mathcal U\backslash G}&\le b^{2r^3}\,\#I,
\end{align*}
where $\epsilon(y)=1$ if $W^{\full}(y)=W^{\cut}(y)$ and $\epsilon(y)=b^{3r^3f}$ otherwise.

\emph{(iii)} For a real number $b$ and a natural number $f$, a \emph{Bessel system} $\mathfrak B\in\BSys{\fs}{b}{f}$ is a
record of: a non-empty finite set $I$; a pair $W=(W^{\full},W^{\cut})$ of continuous functions $\Gone{\fs}\to\CC^I$; a real
number $\deg>0$; a homomorphism $\mathrm{rep}$ from $\Uni_\kappa(\CC)$ to the unitary group of $\CC^I$; and a family
$J$ as in Definition~\ref{d:schur}(v); such that $\CutPair{b}{f}{\mathrm{rep}}{W}$ holds; every matrix that commutes with
all $\mathrm{rep}(h)$ is scalar; $\HasSB{W^{\full}}{\mathrm{rep}}{P}{\deg^{-1}\#I}$ for every real $P\ge b^5$;
$8\,r\le b$; and $\BandL{b}{W^{\cut}}{\deg}{J}$.
\end{definition}
\leanline{\lean{Side}\sep \lean{IsCutPair}\sep \lean{BesselSystem}}

Bessel systems are constructed, from discrete series representations, in the notes on the real place (unit~132)
and used here only through this record.

\subhead{The filling data.}
For a CM field $E$ and a finite set $\kappa$: $\Gf{E}{\kappa}$ is the unitary group of the identity matrix on $\kappa$
over $\AEf$ (the matrices $g$ with $\adj{g}g=1$), $\Kf{E}{\kappa}$ its subgroup of the matrices that have entries in
$\hO$ together with their inverses, and $\Gamma_f(E,\kappa)$ the image of $\Uni_\kappa(E)$ in $\Gf{E}{\kappa}$. A
\emph{level} is an open subgroup of $\Gf{E}{\kappa}$ contained in $\Kf{E}{\kappa}$. For a shape $\fs$, $\Ginf{E}{\fs}$ is the
product of copies of $\Gone{\fs}$ over the infinite places of $E$, with the product of the Haar measures.

\begin{definition}[filling data]\label{d:filling}
Let $E$ be a CM field with a CM type and $\fs=(\varrho,\kappa,k_0)$ a shape; write $k=\kk{\fs}$, $n=\nn{\fs}$,
$f=\ninf{E}$. \emph{Filling data} $S\in\Fill{E}{\fs}$ are a record of
\begin{itemize}
\item a natural number $\NS{S}$; two natural numbers $\qS{S}{0}$, $\qS{S}{1}$; a unit $\piS{S}$ of $\AEf$; an integer
$\vpS{S}\in\IntE$; a real number $\bS{S}$;
\item two finite sets $\ChS{S}{0}$, $\ChS{S}{1}$ of finite places of $E$, the \emph{chosen places};
\item a level $\KcolS{S}\subseteq\Gf{E}{\kappa}$, the \emph{column level};
\item for every infinite place $\tau$ of $E$ a Bessel system $\BsysS{S}{\tau}\in\BSys{\fs}{\bS{S}}{f}$, whose components
we write $\Isys{S}{\tau}$, $\Wsys{\sigma}{S}{\tau}$, $\degsys{S}{\tau}$, $\repsys{S}{\tau}$;
\end{itemize}
subject to the following conditions, in which $b=\bS{S}$, $N=\NS{S}$, $q_i=\qS{S}{i}$, $\bar\pi=\piS{S}$:
\begin{enumerate}[label=(F\arabic*),leftmargin=3em]
\item $\Gamma_f(E,\kappa)\cap g\,\KcolS{S}\,g^{-1}=1$ for every $g\in\Gf{E}{\kappa}$;
\item $\bar\pi\cdot\cj{\bar\pi}=N$ in $\AEf$; $N$ is a prime number; $\bar\pi\in\hO$; and for every finite place
$w$, either $\bar\pi$ is a unit at $w$, or $q_w=N$ and $\bar\pi$ is a unit at $\cj{w}$;
\item $q_0$ and $q_1$ are prime numbers different from $N$, with $b^{10}\le q_i\le b^{11}$;
\item every $w\in\ChS{S}{i}$ has $q_w=q_i$; the sets $\ChS{S}{0}$ and $\ChS{S}{1}$ are disjoint; if $w$ is chosen then $\cj{w}$ is
not; and each $\ChS{S}{i}$ has $f$ elements;
\item for every $x\in\hO$: $q_0q_1$ divides $x$ in $\hO$ if and only if $x$ and $\cj{x}$ reduce to $0$ modulo every
chosen place;
\item $\cj{\vpS{S}}=-\vpS{S}$, and the different of $E$ is the ideal generated by $\vpS{S}$;
\item $10^5\le\varrho$; $2^{-\varrho}\le b^{-6}$; $(\varrho+1)[E:\Q]\le2^{2\varrho}$; $1000(2\varrho+2)^9\le k$;
$1000(2\varrho+2)\le k$; $k\le f$; $10^{10}n^4\le n^6$, $3n^6+12\le n^7$ and $n^7\le k$;
\item $(200\,f\,k\,\rdE{E})^{1000}\le b$; $f\le b/1000$; $k\le b$; $\rdE{E}\le b$; $\varrho+1\le b$; $1000\le b$;
$1\le b$; $b^{19}\le q_0q_1\le b^{23}$; $5\le N$.
\end{enumerate}
\end{definition}
\leanline{\lean{Filling}}

\begin{definition}[forms and groups of the data]\label{d:groups}
Let $S\in\Fill{E}{\fs}$ and let $R$ be a ring.

\emph{(i)} $\gS{S}=\qS{S}{0}\cdot\qS{S}{1}$.

\emph{(ii)} $\Trat{S}{R}=\Trg{\fs}{\NS{S}}$ and $\Tlat{S}{R}=\Tlatg{\fs}{\NS{S}}$, with $\NS{S}$ read in $R$.

\emph{(iii)} If $R$ has an involution, $\GS{S}{R}=\GT{\Trat{S}{R}}$. Further $\GFin{S}=\GT{\Tlat{S}{\AEf}}$, the finite adelic
group \emph{in the lattice basis}.

\emph{(iv)} $\GRes{S}=\prod_{i\in\{0,1\}}\prod_{w\in\ChS{S}{i}}\GL_{\Ix{\fs}}(\Res{w})$, a finite group.

\emph{(v)} $\Bc{S}=\prod_\tau\Isys{S}{\tau}$, the product over the infinite places.

\emph{(vi)} $\Cl{S}=\Gamma_f(E,\kappa)\backslash\Gf{E}{\kappa}/\KcolS{S}$, the set of \emph{classes}. It is finite and
non-empty; every class $c$ has a chosen representative $\dot c\in\Gf{E}{\kappa}$.

\emph{(vii)} For $\gamma\in\GS{S}{E}$, $\Zf{S}{\gamma}=\ZT{\Trat{S}{E}}{\gamma}$ and $\tdisc{S}{\gamma}=\fdisc{\NS{S}}{\Zf{S}{\gamma}}$.
\end{definition}
\leanline{\lean{g}\sep \lean{ratMiddle}\sep \lean{latMiddle}\sep \lean{G}\sep \lean{GFin}\sep
\lean{GRes}\sep \lean{BesselCoord}\sep \lean{Classes}\sep \lean{frameProducts}\sep \lean{termDisc}}

So $\Zf{S}{\gamma}=(h(\fr{a},\gamma\fr{a'}))_{a,a'}$ for the form $h$ of signature $(\varrho+k,\varrho+1)$ at every
infinite place, and $\tdisc{S}{\gamma}=\abs{\lmin{\varrho+1}{Z}}^2-\NS{S}^2\abs{\lmin{\varrho}{Z}}^2$ with $Z=\Zf{S}{\gamma}$ and
$\abs{x}^2=x\cj{x}$.

\begin{definition}[the integral subgroup and its reduction]\label{d:level}
Let $S\in\Fill{E}{\fs}$.

\emph{(i)} $\KS{S}$ is the subgroup of the $k\in\GFin{S}$ that have entries in $\hO$ together with their inverses.

\emph{(ii)} $\redS{S}\colon\KS{S}\to\GRes{S}$ is the homomorphism that reduces every entry modulo each chosen place.

\emph{(iii)} $\KgS{S}\subseteq\GFin{S}$ is the kernel of $\redS{S}$, the \emph{principal level}.

\emph{(iv)} For a subgroup $A\subseteq\GFin{S}$, $\redim{S}{A}=\redS{S}(A\cap\KS{S})$.

\emph{(v)} $\gid{S}$ is the integer $\gS{S}$ as a unit of $\AEf$.
\end{definition}
\leanline{\lean{latStabilizer}\sep \lean{reduction}\sep \lean{principalLevel}\sep \lean{reducedImage}\sep
\lean{gIdele}}

\begin{definition}[scales at the infinite places]\label{d:scales}
Let $S\in\Fill{E}{\fs}$, $\tau$ an infinite place of $E$, $a$ a natural number and $i\in\Ix{\fs}$.

\emph{(i)} $\kapS{S}{\tau}=\operatorname{Re}\bigl(\gS{S}\,\at{\vpS{S}}{\tau}/(2\pi i)\bigr)$.

\emph{(ii)} With $s_a=\kapS{S}{\tau}^{\varrho-a}/\sqrt{\NS{S}}$: $\bscale{S}{\tau}{i}$ is $s_a$ for $i=e_a$, $1$ for $i=x_j$,
$s_\varrho=1/\sqrt{\NS{S}}$ for $i=z$, and $s_a^{-1}$ for $i=e_a^*$; and $\bsGL{S}{\tau}$ is the invertible diagonal complex
matrix with the entries $\bscale{S}{\tau}{i}$.

\emph{(iii)} $\fscale{S}{\tau}{a}=\kapS{S}{\tau}^{2(\varrho-a)}/\NS{S}$, and
$\fsprod{S}=\prod_{\tau}\prod_{a=0}^{\varrho}\fscale{S}{\tau}{a}^{\,k+2(\varrho-a)}$, with $k=\kk{\fs}$.
\end{definition}
\leanline{\lean{kappa}\sep \lean{basisScale}\sep \lean{basisScaleGL}\sep \lean{frameScaleSq}\sep
\lean{frameScaleProd}}

\begin{definition}[coordinates]\label{d:maps}
Let $S\in\Fill{E}{\fs}$.

\emph{(i)} $\latB{S}=\blb{\fs}{\piS{S}}$, an invertible matrix over $\AEf$, and
$\tolb{S}\colon\GS{S}{\AEf}\to\GFin{S}$ is the isomorphism $x\mapsto\latB{S}^{-1}x\,\latB{S}$.

\emph{(ii)} $\radl{S}\colon\GS{S}{E}\to\GS{S}{\AdE}$ embeds every entry diagonally. For $x\in\GS{S}{\AdE}$,
$\ainf{S}(x)\in\Ginf{E}{\fs}$ has the component $\bsGL{S}{\tau}^{-1}\,x_\tau\,\bsGL{S}{\tau}$ at $\tau$, where $x_\tau$ is the
component of $x$ at $\tau$; and $\afin{S}(x)=\tolb{S}(x_f)$ for the finite part $x_f$ of $x$. All three are
homomorphisms.

\emph{(iii)} $\rinf{S}=\ainf{S}\circ\radl{S}$ and $\rfin{S}=\afin{S}\circ\radl{S}$.

\emph{(iv)} $\incH{S}\colon\Gf{E}{\kappa}\to\GFin{S}$ is $\tolb{S}\circ\iota_{\NS{S}}$, with $\iota_{\NS{S}}$ the embedding of
Input~(G3) over $\AEf$; and $\hc{S}{c}=\incH{S}(\dot c)$ for a class $c\in\Cl{S}$.

\emph{(v)} $\sdc{S}$ is the additive character $x\mapsto\psi_f(-x/\vpS{S})$ of $\AEf$, with values in $\CC$, and for
$u\in\GFin{S}$, $\uchar{S}{u}=\sdc{S}\bigl(\gid{S}^{-1}\sd{u}\bigr)$.
\end{definition}
\leanline{\lean{latBasis}\sep \lean{toLatBasis}\sep \lean{ratToAdelic}\sep \lean{adelicToInf}\sep
\lean{adelicToFin}\sep \lean{ratToInfa}\sep \lean{ratToFin}\sep \lean{inclHFin}\sep \lean{h}\sep
\lean{selfDualChar}\sep \lean{unipotentChar}}

\begin{definition}[terms and arithmetic unipotent groups]\label{d:terms}
Let $S\in\Fill{E}{\fs}$, $c,c'\in\Cl{S}$, $s\in\GRes{S}$ and $\gamma\in\GS{S}{E}$.

\emph{(i)} $\pos{S}{c}{c'}{\gamma}=\hc{S}{c}^{-1}\,\rfin{S}(\gamma)\,\hc{S}{c'}\in\GFin{S}$, the \emph{position} of $\gamma$.

\emph{(ii)} $\terms{S}{c}{s}{c'}$ is the set of the $\gamma\in\GS{S}{E}$ whose position $\pos{S}{c}{c'}{\gamma}$ lies in $\KS{S}$
and has reduction $s$; its elements are the \emph{terms over $(c,s,c')$}.

\emph{(iii)} $\clU{S}{c}$ is the group of the $u\in\UT{\Trat{S}{E}}$ with $\pos{S}{c}{c}{u}\in\KgS{S}$, and
$\clUo{S}{c}=\clU{S}{c}\cap\UoT{\Trat{S}{E}}$.

\emph{(iv)} $\Agr{S}{s}$ means: for all $u,u'\in\KS{S}$ that lie in $\UT{\Tlat{S}{\AEf}}$ and satisfy
$\redS{S}(u)\,s=s\,\redS{S}(u')$, one has $\uchar{S}{u}=\uchar{S}{u'}$.

\emph{(v)} A set $\Xi\subseteq\GS{S}{E}$ is \emph{stable}, $\IsStab{S}{\Xi}$, if for all $u,u'\in\UT{\Trat{S}{E}}$ and all
$\gamma\in\GS{S}{E}$: $u^{-1}\gamma u'\in\Xi$ if and only if $\gamma\in\Xi$.

\emph{(vi)} For a submodule $\Lambda\subseteq E^{\Mx{\fs}}\times E^{\Mx{\fs}}$, $\tlink{S}{\Lambda}$ is the set of the
$\gamma\in\GS{S}{E}$ with $\tdisc{S}{\gamma}\ne0$ and $\link{\gamma}=\Lambda$.
\end{definition}
\leanline{\lean{termPosition}\sep \lean{termsOver}\sep \lean{classUnip}\sep \lean{classUnipa}\sep
\lean{Agrees}\sep \lean{IsStable}\sep \lean{termsOfLink}}

\begin{definition}[Bessel functions, boxes, cutoffs, amplitudes]\label{d:amp}
Let $S\in\Fill{E}{\fs}$, $\sigma$ a side, $c,c'\in\Cl{S}$, $s\in\GRes{S}$, $\gamma\in\GS{S}{E}$, $y\in\Ginf{E}{\fs}$ and
$i\in\Bc{S}$; write $n=\nn{\fs}$.

\emph{(i)} $\besfn{\sigma}{S}{y}{i}=\prod_\tau\Wsys{\sigma}{S}{\tau}(y_\tau)_{i_\tau}$, and $\besker{\sigma}{S}{\gamma}{y}\in\CC^{\Bc{S}\times\Bc{S}}$
is the vector with the coordinates $\besfn{\sigma}{S}{y}{i}\cdot\cj{\besfn{\sigma}{S}{\rinf{S}(\gamma)^{-1}y}{i'}}$.

\emph{(ii)} $\bxs{S}=\bS{S}^{2n^2}$.

\emph{(iii)} For a function $\omega\colon\Ginf{E}{\fs}\to[0,\infty]$, $\IsBP{S}{c}{\omega}$ means: $\omega$ is measurable;
$\omega\le1$; $\omega(ty)=\omega(y)$ for all $y$ and all $t\in\Ginf{E}{\fs}$ with
$t_\tau\in\UoT{\Trg{\fs}{1}}\cdot\HT{\Trg{\fs}{1}}$ for every $\tau$ (more exactly, in the subgroup generated by the two); for
all $y$ and $\tau$, if $\omega(y)\ne0$ and $\InCut{y_\tau}{\bS{S}}$ then $y_\tau\in\frbox{\bxs{S}}$; and for all $y$
and every function $\alpha\colon\clU{S}{c}\to[0,\infty]$ with $\sum_{u\in\clUo{S}{c}}\alpha(u\gamma_0)=1$ for all
$\gamma_0\in\clU{S}{c}$, one has $\sum_{\gamma_0\in\clU{S}{c}}\alpha(\gamma_0)\,\omega(\rinf{S}(\gamma_0)y)=1$.

\emph{(iv)} $\cutoff{S}{c}\colon\Ginf{E}{\fs}\to[0,\infty]$ is the product $\beta^0\cdot\omega$ of a chosen measurable
function $\beta^0$ with $\sum_{u\in\clUo{S}{c}}\beta^0(\rinf{S}(u)y)=1$ for all $y$ (it exists by
Lemma~\ref{l:cutoff}) and a chosen function $\omega$ with $\IsBP{S}{c}{\omega}$ if there is one, an arbitrary function
otherwise.

\emph{(v)} The \emph{amplitude} and its majorant are
\begin{align*}
\amp{\sigma}{S}{c}{c'}{\gamma}&=\int_{\Ginf{E}{\fs}}\cutoff{S}{c}(y)\,\cutoff{S}{c'}\bigl(\rinf{S}(\gamma)^{-1}y\bigr)\ \besker{\sigma}{S}{\gamma}{y}\,dy\ \in\CC^{\Bc{S}\times\Bc{S}},\\
\ampabs{\sigma}{S}{c}{c'}{\gamma}&=\int_{\Ginf{E}{\fs}}\cutoff{S}{c}(y)\,\cutoff{S}{c'}\bigl(\rinf{S}(\gamma)^{-1}y\bigr)\ \nrm{\besker{\sigma}{S}{\gamma}{y}}\,dy\ \in[0,\infty],
\end{align*}
where in the first integral the product of the two cutoffs is read as a real number ($0$ where it is infinite) and
the integral is $0$ if the integrand is not integrable.

\emph{(vi)} For a set $\Xi\subseteq\GS{S}{E}$, the \emph{amplitude sum} is
\[\ampsum{\sigma}{S}{\Xi}{c}{s}{c'}=\sum_{\gamma\in\Xi\cap\terms{S}{c}{s}{c'}}\amp{\sigma}{S}{c}{c'}{\gamma},\]
read as $0$ if the family is not summable.
\end{definition}
\leanline{\lean{besselFn}\sep \lean{besselKernel}\sep \lean{boxSize}\sep \lean{IsBoxPartition}\sep
\lean{classCutoff}\sep \lean{amplitude}\sep \lean{amplitudeAbs}\sep \lean{amplitudeSum}}

Note~B shows that under the standing hypotheses, for $E$ stably unramified, a function with $\IsBP{S}{c}{\omega}$ exists, so that $\cutoff{S}{c}$ is a
cutoff for $\clU{S}{c}$ acting through $\rinf{S}$, supported, where the cut Bessel function is non-zero, in a box.

\begin{definition}[constants]\label{d:const}
Let $S\in\Fill{E}{\fs}$; write $k=\kk{\fs}$, $n=\nn{\fs}$, $r=\rr{\fs}$, $f=\ninf{E}$, $b=\bS{S}$, $g=\gS{S}$, $N=\NS{S}$.

\emph{(i)} $\cdim{S}=\#\Cl{S}\cdot\#\Bc{S}$ and $\bdeg{S}=\prod_\tau\degsys{S}{\tau}$.

\emph{(ii)} $\mainE{S}=\cdim{S}^2/\bigl(\bdeg{S}\cdot\mass{E}{r}\bigr)$, where $\mass{E}{r}$ is the mass of the genus of the
standard hermitian lattice of rank $r$ over $E$ (background).

\emph{(iii)} $\uord{S}=g^{\,f\varrho(2k+2\varrho+1)}$ and $\tailS{S}=b^{kf/4}$.

\emph{(iv)} $\siegC{S}=\bigl(\pi^k/((k-1)!\,N^k)\bigr)^{f}\cdot\abs{\dE{E}}^{1/4}\cdot\volO{E}^{k}$.

\emph{(v)} $\packC{S}=\#\redim{S}{\ST{\Tlat{S}{\AEf}}}\cdot\uord{S}\cdot\#\Cl{S}\cdot\siegC{S}$ and
$\besE{S}=\#\redim{S}{\ST{\Tlat{S}{\AEf}}}\cdot\cdim{S}$.

\emph{(vi)} $\codimF{S}=(b/g)^{kf}\cdot2\,\tailS{S}^2$.

\emph{(vii)} With $d=\sum_{a<\varrho}(2a+k+1)=\varrho(\varrho+k)$,
\[\ucov{S}=\Bigl(\frac{g^{2d+\varrho}}{N^k}\Bigr)^{f}\cdot\volO{E}^{-d}\cdot\abs{\dE{E}}^{\varrho/4}\cdot\fsprod{S}^{-1}.\]
\end{definition}
\leanline{\lean{columnDim}\sep \lean{besselDegree}\sep \lean{mainEnergy}\sep \lean{unipOrder}\sep
\lean{tailScale}\sep \lean{siegelConst}\sep \lean{packetConst}\sep \lean{besselEnergy}\sep
\lean{codimFactor}\sep \lean{unipCovolume}}

Since $\volO{E}=2^f\abs{\dE{E}}^{-1/2}$ and $\abs{\dE{E}}=\rdE{E}^{2f}$, one has
$\siegC{S}=\bigl((2\pi)^k\rdE{E}^{1/2-k}/((k-1)!N^k)\bigr)^f$. Note~B proves that $\ucov{S}$ is the covolume of the
image of $\clU{S}{c}$ in the product of the groups $\UT{\Trg{\fs}{1}}$ over the infinite places.

\begin{definition}[standing hypotheses]\label{d:adm}
Let $S\in\Fill{E}{\fs}$, with the abbreviations of Definition~\ref{d:const}. We say that $\IsAdm{S}$ if:
\begin{enumerate}[label=(A\arabic*),leftmargin=3em]
\item $\KcolS{S}=\incH{S}^{-1}(\KgS{S})$, and $\redS{S}\colon\KS{S}\to\GRes{S}$ is surjective;
\item for every infinite place $\tau$: $\abs{\at{\vpS{S}}{\tau}}\le b$ and $1\le\abs{\kapS{S}{\tau}}\le b^{23}$;
\item $N\le b^{60n}$;
\item $\prod_\tau\ \bandc{\fs}\,b^{7k}\Bigl(\dfrac{\abs{\kapS{S}{\tau}}^{2\varrho+1}}{g\,N}\Bigr)^{k-n^7}\le\bdeg{S}^{-1}$;
\item $\tailS{S}^{44}\,\besE{S}\le\mainE{S}$, $\ \tailS{S}^{18}\,\codimF{S}\,\mainE{S}\le\besE{S}$ and $\codimF{S}\le1$;
\item $\#\Bc{S}\le b^{2kf}$;
\item $(N^k/b)^f\le\bigl[\Kf{E}{\kappa}:\Kf{E}{\kappa}\cap\incH{S}^{-1}(\KS{S})\bigr]$;
\item $\packC{S}\le(1+b^{-9})\,\mainE{S}\,\bdeg{S}\,\occ{\fs}^{\,f}\,\volO{E}^{-n}\,\bigl(\ucov{S}\,\fsprod{S}/\#\Bc{S}\bigr)^2$.
\end{enumerate}
\end{definition}
\leanline{\lean{IsAdmissible}}

\subhead{Lattices of classes, levels of terms, near terms.}
\begin{definition}[class bases]\label{d:classbasis}
Let $S\in\Fill{E}{\fs}$ and $c\in\Cl{S}$.

\emph{(i)} $\clB{S}{c}=\iota_{\NS{S}}(\dot c)\cdot\latB{S}$, an invertible matrix over $\AEf$ indexed by $\Ix{\fs}$. Its
columns are a basis of the lattice of the class $c$ in rational coordinates.

\emph{(ii)} $\clC{S}{c}=\clB{S}{c}^{-1}$, as a matrix.

\emph{(iii)} $\clCX{S}{c}=\mbc{\fs}{\piS{S}}\cdot\dot c^{-1}$, a matrix with rows indexed by $\Mx{\fs}$ and columns by $\kappa$.
\end{definition}
\leanline{\lean{classBasis}\sep \lean{classCoord}\sep \lean{classCoordX}}

\begin{definition}[reduction above $g$, fixers]\label{d:fixers}
Let $S\in\Fill{E}{\fs}$.

\emph{(i)} For a ring $F$ and a truth value $\epsilon$, $\ioia{S}{\epsilon}$ is the automorphism of the group of invertible
matrices over $F$ indexed by $\Ix{\fs}$ that is the identity if $\epsilon$ is false, and $s\mapsto J^{-1}\,\tr{(s^{-1})}\,J$
with $J=\JT{\Tlat{S}{F}}$ if $\epsilon$ is true.

\emph{(ii)} For $x=((i,w),\epsilon)$, with $i\in\{0,1\}$, $w\in\ChS{S}{i}$ and a truth value $\epsilon$,
$\rag{S}{x}\colon\hO\to\Res{w}$ is the reduction modulo $w$, preceded by complex conjugation if $\epsilon$ is true.

\emph{(iii)} For $w\in\ChS{S}{i}$ and a pair $d=(\epsilon,v)$ of a truth value and a non-zero vector
$v\in\Res{w}^{\varrho+1}$, $\vfix{S}{w}{d}$ is the subgroup of the $s\in\GRes{S}$ such that $\ioia{S}{\epsilon}(s_{i,w})$ fixes
the vector $\mathfrak fv\in\Res{w}^{\Ix{\fs}}$.

\emph{(iv)} For $s\in\GRes{S}$ and a natural number $x$, $\FixN{s}{x}$ means: there is a finite set $P$ of pairs
$((i,w),\epsilon)$ as in (ii) with $x\le\prod_{((i,w),\epsilon)\in P}\qS{S}{i}$, such that for each of them $s\in\vfix{S}{w}{(\epsilon,v)}$ for
some non-zero $v$.

\emph{(v)} For a finite set $\Sigma\subseteq\ChS{S}{0}$ and a family $d=(d_w)_w$ of pairs as in (iii), $\fixset{S}{\Sigma}{d}$
is the set of the $s\in\GRes{S}$ with $s\in\vfix{S}{w}{d_w}$ for all $w\in\Sigma$ and such that either $\Sigma=\ChS{S}{0}$,
or there is $w\in\ChS{S}{0}$ with $\lmin{l}{\ZT{\Tlat{S}{\Res{w}}}{s_{0,w}}}\ne0$ for all $l\le\varrho+1$.
\end{definition}
\leanline{\lean{idOrInvAdjoint}\sep \lean{reduceAboveG}\sep \lean{vectorFixer}\sep
\lean{FixesNormAtLeast}\sep \lean{fixingSet}}

The truth value records whether one looks at $w$ or at $\cj{w}$: by Lemma~\ref{l:redab} the reduction modulo $\cj{w}$
of an element of $\KS{S}$ is the inverse adjoint of its reduction modulo $w$. An entry of $\redS{S}(k)$ depends only
on the corresponding entry of $k$.
\leanline{\lean{reductionapplycongr}}

\begin{definition}[common links, flag vectors, normal forms]\label{d:nf}
Let $R$ be a ring with involution, $\fs$ a shape, $T$ an $\fs$-form over $R$ and $\gamma,\gamma'\in\GT{T}$.

\emph{(i)} $\zfp{T}$ is the matrix with two rows and columns indexed by $\Mx{\fs}\sqcup\Mx{\fs}$ that is block diagonal
with both blocks equal to the row of $T$ of index $z$ (for this $R$ need not have an involution).

\emph{(ii)} $\clink{\gamma}{\gamma'}\subseteq R^{\kappa}$ is the set of the restrictions $p_1|_\kappa$ of the pairs
$(p_1,p_2)\in\link{\gamma}\cap\link{\gamma'}$ with $\zfp{T}\,(p_1,p_2)=0$, that is, with $(Tp_1)_z=(Tp_2)_z=0$.

\emph{(iii)} For $c,c'\in R$ and a frame index $a\le\varrho$, $\HCFV{\gamma}{c}{c'}{a}$ means: there are
$x,y\in R^{\varrho+1}$ with $\mathfrak fx=(\gamma\mathfrak f)y$, $x_a=c$, $y_a=c'$, and $x_{a'}=y_{a'}=0$ for all $a'>a$.

\emph{(iv)} For a natural number $j$, $\IsNF{j}{\gamma}$ means: there is $c\in R$ with $c\,\cj{c}=1$ such that
$\HCFV{\gamma}{c}{1}{a}$ for every frame index $a\ge j$. We then say that $\gamma$ is in \emph{normal form from $j$}.
\end{definition}
\leanline{\lean{zFormPair}\sep \lean{commonLink}\sep \lean{HasCommonFlagVector}\sep \lean{IsNormalForm}}

In words: $\HCFV{\gamma}{c}{c'}{a}$ says that the flags of the frames $\mathfrak f$ and $\gamma\mathfrak f$ share a vector
of the $a$-th steps with leading coefficients $c$ and $c'$.

\begin{definition}[levels, congruent terms, near terms]\label{d:levels}
Let $S\in\Fill{E}{\fs}$, $c,c'\in\Cl{S}$, $\gamma,\gamma'\in\GS{S}{E}$, $j\in\{0,\dots,\varrho+1\}$ and $i$ a natural number;
write $k=\kk{\fs}$, $f=\ninf{E}$.

\emph{(i)} $\ratG{S}=\JT{\Trat{S}{E}}$, and $\lform{S}$ is the block diagonal matrix with the blocks $\Trat{S}{E}$ and
$-\Trat{S}{E}$, indexed by $\Mx{\fs}\sqcup\Mx{\fs}$.

\emph{(ii)} $\cw{S}{i}=\#\redim{S}{\ST{\Tlat{S}{\AEf}}}\cdot\#\Cl{S}\cdot\NS{S}^{-kf}\cdot\bigl(2\,\tailS{S}^2\bigr)^i$.

\emph{(iii)} $\HasLev{S}{j}{\gamma}$, in words \emph{$\gamma$ has level $j$}, means: $\IsNF{j}{\gamma}$,
$\lmin{j}{\Zf{S}{\gamma}}\ne0$, and $\tdisc{S}{\gamma}\ne0$ if and only if $j=\varrho+1$.

\emph{(iv)} $\Cong{S}{c}{c'}{\gamma}{\gamma'}$, in words \emph{$\gamma$ and $\gamma'$ are congruent terms}, means: there is
$s\in\GRes{S}$ with $\gamma\in\terms{S}{c}{s}{c'}$ and $\gamma'\in\terms{S}{c}{s}{c'}$.

\emph{(v)} $\tdnorm{S}{\gamma}=\prod_\tau\abs{\at{\tdisc{S}{\gamma}}{\tau}}$.

\emph{(vi)} $\bxp{S}$ is the set of the $y\in\Ginf{E}{\fs}$ with $y_\tau\in\frbox{\bxs{S}}$ and $\InCut{y_\tau}{\bS{S}}$ for
all $\tau$; and $\IsNr{S}{\gamma}$, in words \emph{$\gamma$ is near}, means that there is $y\in\bxp{S}$ with
$\rinf{S}(\gamma)^{-1}y\in\bxp{S}$.

\emph{(vii)} $\ChA{S}{\gamma}$, in words \emph{the characters agree at $\gamma$}, means: for all $u,u'\in\UT{\Trat{S}{E}}$
with $u\gamma=\gamma u'$, $\uchar{S}{\rfin{S}(u)}=\uchar{S}{\rfin{S}(u')}$.
\end{definition}
\leanline{\lean{ratGram}\sep \lean{linkForm}\sep \lean{countWeight}\sep \lean{HasLevel}\sep
\lean{CongruentTerms}\sep \lean{termDiscNorm}\sep \lean{boxedPoints}\sep \lean{IsNear}\sep
\lean{CharsAgree}}

\subhead{Lattices in subspaces.}
For a ring $A$ with a subring $O$ and a matrix $M$ over $A$: $M$ is \emph{integral} if its entries lie in $O$,
\emph{primitive} if it is integral and has an integral left inverse; the \emph{column lattice} of $M$ is the
$O$-module spanned by its columns, its \emph{column space} the $A$-module they span; the \emph{index} $[L':L]$ of two
$O$-submodules of $A^m$ is the index of $L\cap L'$ in $L'$, and $0$ if that is infinite. Over $\AEf$ the subring
is $\hO$.

\begin{definition}[lattice bases and discriminants of subspaces]\label{d:sublattice}
Let $E$ be a CM field (a number field for (i)), $\Omega$ and $P$ finite sets.

\emph{(i)} For a matrix $M$ over $E$, $\fadm{M}$ is the matrix over $\AEf$ with the same entries.

\emph{(ii)} Let $E$ be stably unramified, let $H$ be a $P\times\Omega$ matrix over $\AEf$ with a left inverse $L$,
$LH=1$, and let $V\subseteq E^\Omega$ be a subspace of dimension $s$. Then $\slb{H}{V}$ is a chosen
$\Omega\times s$ matrix $\mathrm m$ over $\AEf$ such that $H\mathrm m$ is primitive over $\hO$ and
$\mathrm m=\fadm{y}\,a$ for a matrix $y$ over $E$ with a left inverse whose columns span $V$, and an invertible
$s\times s$ matrix $a$ over $\AEf$. (It exists by the Inputs (La1) and (La2) of Section~\ref{s:siegel}.)

\emph{(iii)} In the situation of (ii), $\sld{H}{V}=\sqrt{\gix{\det(\adj{\mathrm m}\mathrm m)}}$ for $\mathrm m=\slb{H}{V}$.
\end{definition}
\leanline{\lean{finAdeleMat}\sep \lean{sublatticeBasis}\sep \lean{sublatticeDisc}}

The columns of $\slb{H}{V}$ are a basis of the lattice of the adelic points $v$ of $V$ with $Hv$ integral; for
$H=u^{-1}$ with $u\in\Gf{E}{\Omega}$ this is $V(\AEf)\cap u\,\hO^{\Omega}$, the lattice of Theorem~C. Two elementary
compatibilities are used without mention: $\fadm{\adj{M}}=\adj{\fadm{M}}$; $\fadm{aM}=a\fadm{M}$ for $a\in E$; and for an
embedding $\varphi\colon E\to\CC$, which commutes with the conjugations because $E$ is a CM field, the matrix of the
images of $\adj{M}$ is the conjugate transpose of that of $M$.
\leanline{\lean{finAdeleMatconjTranspose}\sep \lean{finAdeleMatsmul}\sep \lean{mapconjTranspose}}

\sect{First properties and the dictionary}{s:dict}
Throughout this section $E$ is a CM field with a CM type, $\fs=(\varrho,\kappa,k_0)$ a shape and
$S\in\Fill{E}{\fs}$; we abbreviate $N=\NS{S}$, $g=\gS{S}$, $b=\bS{S}$, $k=\kk{\fs}$, $n=\nn{\fs}$, $f=\ninf{E}$ in proofs.

\begin{lemma}[scales]\label{l:scales}
Let $S\in\Fill{E}{\fs}$, $\tau$ an infinite place of $E$, $a$ a natural number, $a'\le\varrho$ a frame index and
$i\in\Ix{\fs}$.

\emph{(a)} $\vpS{S}\ne0$; $\kapS{S}{\tau}=\gS{S}\,\at{\vpS{S}}{\tau}/(2\pi i)$, so this complex number is real;
$\abs{\kapS{S}{\tau}}=\gS{S}\,\abs{\at{\vpS{S}}{\tau}}/(2\pi)$; and $\kapS{S}{\tau}\ne0$.

\emph{(b)} $\bscale{S}{\tau}{\fr{a'}}=\kapS{S}{\tau}^{\varrho-a'}/\sqrt{\NS{S}}$; $\bscale{S}{\tau}{i}\ne0$; $\fscale{S}{\tau}{a}>0$;
$\fsprod{S}>0$; $\bdeg{S}>0$; and $\ucov{S}>0$.

\emph{(c)} $\adj{\bsGL{S}{\tau}}\ \JT{\Trat{S}{\CC}}\ \bsGL{S}{\tau}=\JT{\Trg{\fs}{1}}$.
\end{lemma}
\begin{proof}
(a) The different is a non-zero ideal, so its generator $\vpS{S}$ is not zero (F6). The embedding at $\tau$ commutes
with complex conjugation, so $\cj{\at{\vpS{S}}{\tau}}=-\at{\vpS{S}}{\tau}$ by (F6), and $\at{\vpS{S}}{\tau}/(2\pi i)$ is
real; hence the real part in the definition of $\kapS{S}{\tau}$ is the number itself. The rest follows, $\gS{S}$ being a
product of two primes. (b) The first formula is the definition, and all scales are powers of the non-zero real
number $\kapS{S}{\tau}$ divided or multiplied by $\sqrt{\NS{S}}>0$. Further $\fscale{S}{\tau}{a}$ is the square of a
non-zero real number divided by $\NS{S}$, the degrees of Bessel systems are positive by definition, and $\ucov{S}$ is
a quotient of positive numbers. (c) The scales are real, with product $1$ on each pair $e_a,e_a^*$, equal to $1$
on the $x_j$ and to $1/\sqrt{\NS{S}}$ on $z$, where the form has the value $-\NS{S}$.
\end{proof}
\leanline{\lean{varpinezero}\sep \lean{coekappa}\sep \lean{abskappaeq}\sep \lean{kappanezero}\sep
\lean{basisScalebesselFrame}\sep \lean{basisScalenezero}\sep \lean{frameScaleSqpos}\sep
\lean{frameScaleProdpos}\sep \lean{besselDegreepos}\sep \lean{unipCovolumepos}\sep
\lean{grambasisScaleGL}}

Part (c) is what makes $\ainf{S}$ take values in $\Ginf{E}{\fs}$. That $\radl{S}$, $\afin{S}$ and the reduction maps are
defined uses that a ring homomorphism $R\to R'$ carries $\JT{\Trat{S}{R}}$ to $\JT{\Trat{S}{R'}}$ and
$\JT{\Tlat{S}{R}}$ to $\JT{\Tlat{S}{R'}}$, and that it maps a matrix that is the identity at a set of positions to a
matrix with the same property; also $\Trat{S}{E}$ is hermitian and invertible, and $\GS{S}{E}$ is countable.
\leanline{\lean{mapsplitGramratMiddle}\sep \lean{mapsplitGramlatMiddle}\sep
\lean{mapmemunitriangularOn}\sep \lean{ratToInfmemunitriangularOn}}

\begin{lemma}[coordinates]\label{l:coords}
Let $S\in\Fill{E}{\fs}$, $\tau$ an infinite place, $i,j\in\Ix{\fs}$, and $a,a'\le\varrho$ frame indices.

\emph{(a)} For $x\in\GS{S}{\AdE}$: $\bigl(\ainf{S}(x)_\tau\bigr)_{ij}=(x_{ij})_\tau\cdot\bscale{S}{\tau}{j}/\bscale{S}{\tau}{i}$. For
$\gamma\in\GS{S}{E}$: $\bigl(\rinf{S}(\gamma)_\tau\bigr)_{ij}=\at{(\gamma_{ij})}{\tau}\cdot\bscale{S}{\tau}{j}/\bscale{S}{\tau}{i}$ and
$\rfin{S}(\gamma)=\latB{S}^{-1}\,\gamma\,\latB{S}$, with $\gamma$ read over $\AEf$.

\emph{(b)} For $\gamma\in\GS{S}{E}$:
$\Zr{\rinf{S}(\gamma)_\tau}_{aa'}=\dfrac{\kapS{S}{\tau}^{\varrho-a}\,\kapS{S}{\tau}^{\varrho-a'}}{\NS{S}}\cdot\at{\bigl(\Zf{S}{\gamma}_{aa'}\bigr)}{\tau}$;
and for a finite set $m$ and a map $e$ from $m$ to the frame indices,
\[\det\bigl(\Zr{\rinf{S}(\gamma)_\tau}_{e(\mathfrak g),e(\mathfrak k)}\bigr)_{\mathfrak g,\mathfrak k\in m}=\prod_{\mathfrak g\in m}\fscale{S}{\tau}{e(\mathfrak g)}\cdot
\at{\bigl(\det(\Zf{S}{\gamma}_{e(\mathfrak g),e(\mathfrak k)})_{\mathfrak g,\mathfrak k\in m}\bigr)}{\tau}.\]

\emph{(c)} For $u\in\GS{S}{\AdE}$: $\sd{\afin{S}(u)}=\sd{u}_f$, the finite part; for $u\in\GS{S}{E}$:
$\sd{\radl{S}(u)}=\sd{u}$, embedded diagonally.
\end{lemma}
\begin{proof}
(a) is the definition of $\ainf{S}$ as conjugation by the diagonal matrix $\bsGL{S}{\tau}$, and of $\tolb{S}$. (b) By
Lemma~\ref{l:scales}(c), $h_1(v,\bsGL{S}{\tau}^{-1}w)=h_N(\bsGL{S}{\tau} v,w)$ for the forms $h_1$, $h_N$ of $\Trg{\fs}{1}$ and
$\Trat{S}{\CC}$. So the $(a,a')$ entry is
$h_N(\bsGL{S}{\tau}\fr{a},\gamma\,\bsGL{S}{\tau}\fr{a'})=\bscale{S}{\tau}{\fr{a}}\,\bscale{S}{\tau}{\fr{a'}}\,\at{h_N(\fr{a},\gamma\fr{a'})}{\tau}$,
the scales being real, and Lemma~\ref{l:scales}(b) gives the factor. In a minor, the row and the column of index
$e(\mathfrak g)$ each carry the factor $\kapS{S}{\tau}^{\varrho-e(\mathfrak g)}/\sqrt{\NS{S}}$, whose square is $\fscale{S}{\tau}{e(\mathfrak g)}$.
(c) The entries of $\sd{\cdot}$ are at the positions $(\fr{a},\fr{a+1})$; the matrix $\latB{S}$ is the identity on the
$e_a$ and its column $z$ is $z$, so conjugation by it does not change these entries.
\end{proof}
\leanline{\lean{adelicToInfapply}\sep \lean{ratToInfapply}\sep \lean{coeratToFin}\sep
\lean{realFrameGramratToInf}\sep \lean{detsubmatrixrealFrameGramratToInf}\sep
\lean{simpleRootSumadelicToFin}\sep \lean{simpleRootSumratToAdelic}}

\statement{ana}{\textbf{Inputs (adeles).}\enspace Let $E$ be a CM field.
(Ad1), unit~123: $E$, embedded diagonally, is a lattice of covolume one in $\AdE$; in particular a compact subset
of $\AdE$ contains only finitely many elements of $E$. $\hO$ is compact and open in $\AEf$ and stable under
conjugation, and $E$ is dense in $\AEf$.
(Ad2), unit~122: the group of the invertible matrices over $\hO$ is compact and open in the group of invertible
matrices over $\AEf$.
(Ad3), unit~114: a countable discrete subgroup $\Gamma$ of a locally compact second countable group has a
\emph{cutoff}: a measurable function $\beta$ with values in $[0,\infty]$ and $\sum_{\gamma\in\Gamma}\beta(\gamma y)=1$
for all $y$.
(Ad4), unit~123: for $x\in\AEf$, $\psi_f(ux)=1$ for all $u\in\hO$ if and only if $dx\in\hO$ for every $d$ in the
different of $E$; and $\psi_E$ is trivial on $E$.}

\begin{lemma}[discreteness and cutoffs]\label{l:cutoff}
Let $S\in\Fill{E}{\fs}$ and $c\in\Cl{S}$.

\emph{(a)} Let $m$ be a finite set, $\phi$ a map from $m$ to $\Ix{\fs}$, and $C$ a compact set of families
$(M_\tau)_\tau$ of complex $m\times m$ matrices indexed by the infinite places $\tau$. Then only finitely many
points of $C$ are of the form $\bigl((\rinf{S}(u)_\tau)_{\phi(i),\phi(j)}\bigr)_{\tau;\,i,j\in m}$ with $u\in\clU{S}{c}$.

\emph{(b)} Let $\mathsf{K}$ be a subgroup of $\clU{S}{c}$ and $\Gamma=\rinf{S}(\mathsf{K})\subseteq\Ginf{E}{\fs}$. Then $\Gamma$ is discrete,
and there is a measurable function $\beta\colon\Ginf{E}{\fs}\to[0,\infty]$ with $\sum_{u\in \mathsf{K}}\beta(\rinf{S}(u)y)=1$ for all
$y$; it is a cutoff for $\Gamma$ with the counting measure.
\end{lemma}
\begin{proof}
(a) Let $u\in\clU{S}{c}$ and $k=\pos{S}{c}{c}{u}\in\KgS{S}$, a matrix with entries in $\hO$ together with its inverse. By
Lemma~\ref{l:coords}(a), $u=(\latB{S}\hc{S}{c})\,k\,(\latB{S}\hc{S}{c})^{-1}$ over $\AEf$. So the finite adele of an entry
$x=u_{\phi(i)\phi(j)}\in E$ lies in a compact set that does not depend on $u$, the image of the compact group of
(Ad2) under a continuous map. If the point attached to $u$ lies in $C$, then by Lemma~\ref{l:coords}(a) the
archimedean components of $x$ are those of a point of $C$ times fixed scales, and lie in a compact set too. By
(Ad1) only finitely many $x\in E$ are possible for each pair $(i,j)$, and the point is determined by these entries.
(b) The map $u\mapsto\rinf{S}(u)$ is injective, the embeddings $E\to\CC$ being injective. By (a) with $\phi$ the
identity, every compact subset of $\Ginf{E}{\fs}$ contains only finitely many elements of $\Gamma$, so $\Gamma$ is
discrete. It is countable, and (Ad3) gives a cutoff $\beta$ for $\Gamma$; the sum over $\mathsf{K}$ is the sum over
$\Gamma$.
\end{proof}
\leanline{\lean{finiteclassUnipintercompact}\sep \lean{existscutoffclassUnip}}

\begin{lemma}[the groups $\mathcal U$ and $\mathcal H$ in the three systems of coordinates]\label{l:dictUH}
Let $S\in\Fill{E}{\fs}$.

\emph{(a)} For $x\in\GS{S}{\AEf}$: $\tolb{S}(x)\in\UT{\Tlat{S}{\AEf}}$ if and only if $x\in\UT{\Trat{S}{\AEf}}$. For
$\gamma\in\GS{S}{E}$: $\rfin{S}(\gamma)\in\UT{\Tlat{S}{\AEf}}$ if and only if $\gamma\in\UT{\Trat{S}{E}}$.

\emph{(b)} For every $k\in\Gf{E}{\kappa}$, $\incH{S}(k)\in\HT{\Tlat{S}{\AEf}}$; and the image of $\incH{S}$ is
$\HT{\Tlat{S}{\AEf}}$. For $\gamma_0\in\Uni_\kappa(E)$,
$\rfin{S}(\iota_{\NS{S}}(\gamma_0))=\incH{S}(\gamma_0)$, where on the right $\gamma_0$ is read in $\Gf{E}{\kappa}$.

\emph{(c)} For $c,c'\in\Cl{S}$ and $\gamma\in\GS{S}{E}$: $\pos{S}{c}{c'}{\gamma}=\clB{S}{c}^{-1}\,\gamma\,\clB{S}{c'}$ as matrices.
For $c\in\Cl{S}$ and $\gamma\in\GS{S}{E}$ with $\gamma\in\UT{\Trat{S}{E}}$: $\pos{S}{c}{c}{\gamma}\in\UT{\Tlat{S}{\AEf}}$.

\emph{(d)} Let $c,c',c''\in\Cl{S}$, $s,s'\in\GRes{S}$ and $\gamma,\delta\in\GS{S}{E}$. If $\gamma\in\terms{S}{c}{s}{c'}$ and
$\delta\in\terms{S}{c'}{s'}{c''}$, then $\gamma\delta\in\terms{S}{c}{ss'}{c''}$. If $\gamma\in\terms{S}{c}{s}{c'}$, then
$\gamma^{-1}\in\terms{S}{c'}{s^{-1}}{c}$.
\end{lemma}
\begin{proof}
(a) The matrix $\latB{S}$ and its inverse are block diagonal for the levels: they are the identity on the $e_a$
and $e_a^*$ and mix only the middle block, of level $\varrho$. By (G1), conjugation by them preserves
unitriangularity for the levels, in both directions. The second claim follows, since an entry of $\gamma$ is $0$
or $1$ in $E$ exactly when it is so in $\AEf$. (b) $\iota_N(k)$ lies in $\HT{\Trat{S}{\AEf}}$ by (G3), and
conjugation by $\latB{S}$ keeps the columns of the $e_a$, $e_a^*$ and $z$, since $\latB{S}$ fixes these vectors.
Conversely let $h\in\HT{\Tlat{S}{\AEf}}$. Then $\latB{S}\,h\,\latB{S}^{-1}\in\HT{\Trat{S}{\AEf}}$, and it is $\iota_N(k)$
for some $k$ by (G3), $N$ being a unit of $\AEf$ by (F2); so $h=\incH{S}(k)$. The last formula is the definition.
(c) By definition $\hc{S}{c}=\latB{S}^{-1}\iota_N(\dot c)\latB{S}$ and $\rfin{S}(\gamma)=\latB{S}^{-1}\gamma\latB{S}$, which
gives the formula. For the second claim, $\rfin{S}(\gamma)\in\UT{\Tlat{S}{\AEf}}$ by (a), and $\hc{S}{c}^{-1}$ lies in
$\HT{\Tlat{S}{\AEf}}\subseteq\ST{\Tlat{S}{\AEf}}$, which normalizes $\UT{\Tlat{S}{\AEf}}$ by (G2). (d) Positions multiply,
$\pos{S}{c}{c''}{\gamma\delta}=\pos{S}{c}{c'}{\gamma}\cdot\pos{S}{c'}{c''}{\delta}$, and $\pos{S}{c'}{c}{\gamma^{-1}}=\pos{S}{c}{c'}{\gamma}^{-1}$;
$\KS{S}$ is a group and $\redS{S}$ a homomorphism.
\end{proof}
\leanline{\lean{toLatBasismembesselU}\sep \lean{ratToFinmembesselU}\sep \lean{inclHFinmembesselH}\sep
\lean{rangeinclHFin}\sep \lean{ratToFintoBesselH}\sep \lean{coetermPosition}\sep
\lean{termPositionmembesselU}\sep \lean{mul}\sep \lean{inva}}

\begin{lemma}[the character]\label{l:char}
Let $S\in\Fill{E}{\fs}$.

\emph{(a)} $\abs{\sdc{S}(x)}=1$ for all $x\in\AEf$. For $a\in\AEf$: $\sdc{S}(ra)=1$ for all $r\in\hO$ if and only if
$a\in\hO$. In particular $\sdc{S}(x)=1$ for every $x\in\AEf$ with $x\in\hO$.

\emph{(b)} $\gid{S}$ has the value $\gS{S}$ in $\AEf$ and lies in $\hO$; and $\uchar{S}{1}=1$.

\emph{(c)} For $r\in\hO$: $\gid{S}^{-1}r\in\hO$ if and only if $\rag{S}{x}(r)=0$ for all $x=((i,w),\epsilon)$ as in
Definition~\ref{d:fixers}(ii).

\emph{(d)} For every family $(t_x)_x$ of elements $t_x\in\Res{w}$, indexed by the $x=((i,w),\epsilon)$, there is
$y\in\hO$ with $\rag{S}{x}(y)=t_x$ for all $x$.
\end{lemma}
\begin{proof}
(a) $\sdc{S}(ra)=\psi_f(r\cdot(-a/\vpS{S}))$. By (Ad4) this is $1$ for all $r\in\hO$ if and only if $-da/\vpS{S}\in\hO$
for all $d$ in the different, which is generated by $\vpS{S}$ (F6); that is, if and only if $a\in\hO$. Take $a=1$
for the last claim. (b) is clear; $\sd{1}=0$. (c) is (F5), with the two reductions of $r$ and of $\cj{r}$ written as
the four maps $\rag{S}{x}$. (d) Lift $t_x$ to $\tilde t_x\in\IntE$ and let $\varepsilon_w\in\hO$ be the idempotent that
is $1$ at $w$ and $0$ at the other places. Put $y=\sum_{x}\tilde t_x\varepsilon_w$ over the $x$ with $\epsilon$
false, plus $\sum_x\cj{\tilde t_x\varepsilon_w}$ over the $x$ with $\epsilon$ true. Since
$\cj{\varepsilon_w}=\varepsilon_{\cj{w}}$, the chosen places are distinct and no conjugate of a chosen place is
chosen (F4), exactly one term survives each $\rag{S}{x}$, and it gives $t_x$.
\end{proof}
\leanline{\lean{normselfDualChar}\sep \lean{forallselfDualCharmuleqoneiff}\sep
\lean{selfDualChareqoneofmem}\sep \lean{gIdeleval}\sep \lean{gIdelemem}\sep
\lean{unipotentCharone}\sep \lean{invgIdelesmulmemiff}\sep \lean{existsreduceAboveGeq}}

\begin{lemma}[reduction modulo a conjugate place]\label{l:redab}
Let $S\in\Fill{E}{\fs}$ and let $F$ be a ring.

\emph{(a)} $\JT{\Tlat{S}{F}}$ is a symmetric matrix. If $F$ has an involution, then $\Tlat{S}{F}$ is unimodular over
every subring $R$ of $F$; in particular $\JT{\Tlat{S}{F}}$ is hermitian and invertible.

\emph{(b)} Let $k\in\KS{S}$, $i\in\{0,1\}$, $w\in\ChS{S}{i}$ and $\epsilon$ a truth value. Applying $\rag{S}{((i,w),\epsilon)}$
to every entry of $k$ gives the matrix $\ioia{S}{\epsilon}\bigl(\redS{S}(k)_{i,w}\bigr)$.
\end{lemma}
\begin{proof}
(a) The matrix has the entries $0$, $\pm1$, $-\NS{S}$ and is symmetric; the image of the natural number $\NS{S}$
lies in every subring and is fixed by every involution, so Lemma~\ref{l:latgram}(b) applies. (b) For $\epsilon$
false this is the definition of $\redS{S}$. Let $\epsilon$ be true. The matrix $k$ is an isometry of
$J=\JT{\Tlat{S}{\hO}}$ over $\hO$, so $Jk^{-1}=\adj{k}J$ by (G6), and $\cj{k}=\tr{(\adj{k})}=\tr{(Jk^{-1}J^{-1})}=J^{-1}\,\tr{(k^{-1})}\,J$,
$J$ being symmetric. Reducing modulo $w$ gives the claim.
\end{proof}
\leanline{\lean{isUnimodularlatMiddle}\sep \lean{transposesplitGram}\sep \lean{intGLmapreduceAboveG}}

\begin{proof}[Proof of Theorem E(a)]
Let $k\in\KgS{S}$, that is $k\in\KS{S}$ with $\redS{S}(k)=1$. By Lemma~\ref{l:redab}(b), for every $x=((i,w),\epsilon)$
the map $\rag{S}{x}$ sends $k$ to $\ioia{S}{\epsilon}(1)=1$, so it kills every entry of $k-1$; by Lemma~\ref{l:char}(c) the
entries of $\gid{S}^{-1}(k-1)$ lie in $\hO$. Conversely, if $k\in\GFin{S}$ has entries in $\hO$ with its inverse and
$\gid{S}^{-1}(k-1)$ is integral, then by Lemma~\ref{l:char}(c) with $\epsilon$ false every entry of $k-1$ vanishes
modulo every chosen place, so $\redS{S}(k)=1$.
\end{proof}
\leanline{\lean{principalLevelequnitaryLevel}}

\begin{proof}[Proof of Theorem E(c)]
Put $\mathsf z=-(g\,\vpS{S})^{-1}\sd{x}\in\AdE$. By Lemma~\ref{l:coords}(c),
$\uchar{S}{\afin{S}(x)}=\psi_f(-\gid{S}^{-1}\sd{x}_f/\vpS{S})=\psi_f(\mathsf z_f)$. Let $\tau$ be an infinite place and
$x'_\tau=\ainf{S}(x)_\tau$. It lies in $\UT{\Trg{\fs}{1}}$, a conjugate of a unitriangular matrix by a diagonal one
(G1), and by Lemma~\ref{l:coords}(a) and Lemma~\ref{l:scales},
\[\sd{x'_\tau}=\sum_{a<\varrho}(x_{\fr{a}\fr{a+1}})_\tau\,\frac{\bscale{S}{\tau}{\fr{a+1}}}{\bscale{S}{\tau}{\fr{a}}}
=\kapS{S}{\tau}^{-1}\,\sd{x}_\tau=\frac{2\pi i}{g\,\at{\vpS{S}}{\tau}}\,\sd{x}_\tau=-2\pi i\,\mathsf z_\tau .\]
Since $\Wsys{\sigma}{S}{\tau}$ is a Bessel function,
$\Wsys{\sigma}{S}{\tau}(x'_\tau y_\tau)=\mathbf e(-2\pi i\mathsf z_\tau)\,\Wsys{\sigma}{S}{\tau}(y_\tau)$ and
$\mathbf e(-2\pi i\mathsf z_\tau)=e^{-2\pi i\cdot2\operatorname{Re}\mathsf z_\tau}$. The product over $\tau$ of these factors is the
archimedean part of $\psi_E(\mathsf z)$, that is $\psi_E(\mathsf z)/\psi_f(\mathsf z_f)$.
\end{proof}
\leanline{\lean{besselFnadelicToInfmul}}

\begin{corollary}\label{c:besselrat}
Let $S\in\Fill{E}{\fs}$, $\sigma$ a side, $\gamma\in\GS{S}{E}$ with $\gamma\in\UT{\Trat{S}{E}}$, and $y\in\Ginf{E}{\fs}$. Then
$\besfn{\sigma}{S}{\rinf{S}(\gamma)y}{\,\cdot\,}=\uchar{S}{\rfin{S}(\gamma)}^{-1}\,\besfn{\sigma}{S}{y}{\,\cdot\,}$.
\end{corollary}
\begin{proof}
Theorem~E(c) for $x=\radl{S}(\gamma)$: here $\mathsf z\in E$, and $\psi_E$ is trivial on $E$ (Ad4).
\end{proof}
\leanline{\lean{besselFnratToInfmul}}

\begin{lemma}[norms of kernels and amplitudes]\label{l:kernel}
Let $S\in\Fill{E}{\fs}$, $\sigma$ a side, $\gamma\in\GS{S}{E}$, $c,c'\in\Cl{S}$, $y\in\Ginf{E}{\fs}$ and $i\in\Bc{S}$. Then
$y'\mapsto\besfn{\sigma}{S}{y'}{i}$ and $y'\mapsto\besker{\sigma}{S}{\gamma}{y'}$ are continuous,
\[\nrm{\besker{\sigma}{S}{\gamma}{y}}=\prod_\tau\nrm{\Wsys{\sigma}{S}{\tau}(y_\tau)}\cdot\prod_\tau\nrm{\Wsys{\sigma}{S}{\tau}\bigl((\rinf{S}(\gamma)^{-1}y)_\tau\bigr)},\]
and $\nrm{\amp{\sigma}{S}{c}{c'}{\gamma}}\le\ampabs{\sigma}{S}{c}{c'}{\gamma}$.
\end{lemma}
\begin{proof}
The functions of a Bessel system are continuous. The squared norm of $\besker{\sigma}{S}{\gamma}{y}$ is
$\sum_{i,i'}\abs{\besfn{\sigma}{S}{y}{i}}^2\abs{\besfn{\sigma}{S}{\rinf{S}(\gamma)^{-1}y}{i'}}^2$, and a sum over tuples of a product
is the product of the sums. The last inequality is the triangle inequality for the integral.
\end{proof}
\leanline{\lean{continuousbesselFn}\sep \lean{continuousbesselKernel}\sep \lean{enormbesselKernel}\sep
\lean{enormamplitudele}}

\subhead{The group $\mathcal H$ and the integral subgroup.}
\begin{lemma}\label{l:mbc}
Let $S\in\Fill{E}{\fs}$ and write $\bar\pi=\piS{S}$; for a unit $p$ of $\AEf$ let $D_p$ be the diagonal matrix indexed
by $\kappa$ with $p$ at $k_0$ and $1$ elsewhere.

\emph{(a)} Let $m$ be a finite set and $l$ a matrix over $\AEf$ with rows indexed by $\kappa$ and columns by $m$. If
$D_{\bar\pi^{-1}}\,l$ is integral, then $\mbc{\fs}{\bar\pi}\,l$ is integral. And $\mbc{\fs}{\bar\pi}\,l$ is integral if and
only if $l$ is integral and its row $k_0$ has entries in $\bar\pi\hO$.

\emph{(b)} Let $\kappa'$ be a finite set, $p$ a unit of $\AEf$ such that at every finite place $w$ either $p$ is
a unit at $w$ or $p$ is a unit at $\cj{w}$, $o\in\kappa'$, and $k\in\Gf{E}{\kappa'}$ such that $k$ and $k^{-1}$ are
integral and the row $o$ of $k-1$ has entries in $p\hO$. Then $D_{p^{-1}}\,(k-1)\,D_{\cj{p}^{-1}}$ is integral, with
$D$ now indexed by $\kappa'$ and $o$ in the place of $k_0$.

\emph{(c)} Let $k\in\Gf{E}{\kappa}$ and let $\Pi$ be a property of elements of $\AEf$ that holds for $0$. Then all
entries of $\incH{S}(k)-1$ have $\Pi$ if and only if all entries of $\mbc{\fs}{\bar\pi}\,(k-1)\,D_{\cj{\bar\pi}^{-1}}$
have $\Pi$.
\end{lemma}
\begin{proof}
(a) The rows of $\mbc{\fs}{\bar\pi}\,l$ are the rows $l_i$ for $i\ne k_0$, the row $\cj{\bar\pi}\,l_{k_0}$ and the row
$-\bar\pi^{-1}l_{k_0}$. Both $\bar\pi$ and $\cj{\bar\pi}$ lie in $\hO$ by (F2) and (Ad1), and
$\cj{\bar\pi}\,l_{k_0}=N\cdot\bar\pi^{-1}l_{k_0}$; both claims follow. (b) The row $o$ of $k-1$ lies in $p\hO$. So does the row $o$ of $k^{-1}-1=-(k-1)\,k^{-1}$, because $k^{-1}$ is
integral. As $k^{-1}=\adj{k}$, this row is the conjugate of the column $o$ of $k-1$, which therefore lies in
$\cj{p}\hO$. The entries of the scaled matrix are $(k-1)_{ij}$, $p^{-1}(k-1)_{oj}$, $(k-1)_{io}\cj{p}^{-1}$ and
$p^{-1}(k-1)_{oo}\cj{p}^{-1}$. At a place $w$ where $p$ is a unit, division by $p$ is harmless and the column
condition handles $\cj{p}$; at a place $w$ where $p$ is a unit at $\cj{w}$, $\cj{p}$ is a unit at $w$ and the row
condition handles $p$. (c) $\incH{S}(k)-1=\latB{S}^{-1}(\iota_N(k)-1)\latB{S}$. The matrix $\iota_N(k)-1$ is zero
outside the block of $\kappa$, where it is $k-1$; the columns of $\latB{S}^{-1}$ at $\kappa$ form the matrix
$\mbc{\fs}{\bar\pi}$ in the middle rows, and the rows of $\latB{S}$ at $\kappa$ form $D_{\cj{\bar\pi}^{-1}}$ in the
columns of $\kappa$ and vanish elsewhere. So the entries of $\incH{S}(k)-1$ are those of the product and zeros.
\end{proof}
\leanline{\lean{isIntMatrixmiddleBaseChangemul}\sep \lean{isIntMatrixscaleAtsubone}\sep
\lean{forallinclHFinsuboneiff}}

\begin{proof}[Proof of Theorem E(b)]
An isometry $x$ of the unimodular matrix $J=\JT{\Tlat{S}{\AEf}}$ lies in $\KS{S}$ if and only if $x-1$ is integral,
since then $x^{-1}=J^{-1}\adj{x}J$ is integral too (Lemma~\ref{l:redab}(a)). By Lemma~\ref{l:mbc}(c), with $\Pi$
integrality, $\incH{S}(k)\in\KS{S}$ if and only if $M=\mbc{\fs}{\bar\pi}\,(k-1)\,D_{\cj{\bar\pi}^{-1}}$ is integral. If $M$ is
integral, so is $M\,D_{\cj{\bar\pi}}=\mbc{\fs}{\bar\pi}(k-1)$, and Lemma~\ref{l:mbc}(a) shows that $k-1$ is integral with
row $k_0$ in $\bar\pi\hO$; then $k$ is integral, and so is $k^{-1}=\adj{k}$. Conversely, if $k$ satisfies these
conditions, Lemma~\ref{l:mbc}(b) with $p=\bar\pi$ and $o=k_0$, whose hypothesis is (F2), shows that
$D_{\bar\pi^{-1}}(k-1)D_{\cj{\bar\pi}^{-1}}$ is integral, and Lemma~\ref{l:mbc}(a) with $l=(k-1)D_{\cj{\bar\pi}^{-1}}$ that
$M$ is.
\end{proof}
\leanline{\lean{inclHFinmemlatStabilizer}}

\begin{lemma}\label{l:redinc}
Let $S\in\Fill{E}{\fs}$.

\emph{(a)} Let $k\in\Gf{E}{\kappa}$ with $\incH{S}(k)\in\KS{S}$, $i\in\{0,1\}$ and $w\in\ChS{S}{i}$. Then
$\redS{S}(\incH{S}(k))_{i,w}=1$ if and only if $k$, which is integral by Theorem~E(b), reduces to $1$ modulo $w$.

\emph{(b)} For every real function $\phi$ on $\Gf{E}{\kappa}$, the class sum
$\sum_{c}\phi(\dot c)/\#(\Gamma_f(E,\kappa)\cap\dot c\,\KcolS{S}\,\dot c^{-1})$ over $c\in\Cl{S}$ equals
$\sum_{c\in\Cl{S}}\phi(\dot c)$.

\emph{(c)} For $i\in\{0,1\}$ and $w\in\ChS{S}{i}$, the residue field $\Res{w}$ has $\qS{S}{i}$ elements.
\end{lemma}
\begin{proof}
(a) Both conditions say that every entry of the matrix minus $1$ lies in the maximal ideal at~$w$. By
Lemma~\ref{l:mbc}(c) with this property, the condition for $\incH{S}(k)$ is the condition for
$\mbc{\fs}{\bar\pi}(k-1)D_{\cj{\bar\pi}^{-1}}$. A chosen place has $q_w=\qS{S}{i}\ne N$ by (F3) and (F4), so $\bar\pi$ is a
unit at $w$ by (F2), and so is $\cj{\bar\pi}$, as $\bar\pi\cj{\bar\pi}=N$ is a unit at $w$. Multiplying rows and
columns by units at $w$ does not change the condition, and the row $z$ of $\mbc{\fs}{\bar\pi}(k-1)$ is a unit multiple
of the row $k_0$ of $k-1$. (b) The stabilizers are trivial by (F1). (c) is (F4).
\end{proof}
\leanline{\lean{reductioninclHFineqoneiff}\sep \lean{classSumeqsum}\sep
\lean{cardresFieldeqlevelPrime}}

\sect{Links and normal forms}{s:links}
In this section $F$ is a ring with involution, $\fs$ a shape and $T$ an $\fs$-form over $F$; vectors of $F^{\Ix{\fs}}$
are written $(a,p,a')$ with $e$-coordinates $a$, middle part $p$ and $e^*$-coordinates $a'$. By the shape of
$\JT{T}$, $\JT{T}\,(a,p,a')=(a',Tp,a)$; hence a vector is orthogonal to $X=\Sp(e_0,\dots,e_{\varrho-1})$ if and
only if its $e^*$-part vanishes, and $\adj{\ifpair{\gamma}}\JT{T}\,t=0$ if and only if the $e^*$-parts of $t$ and of
$\gamma^{-1}t$ vanish; and $\lmat{\gamma} t=(t_{\mathrm{mid}},(\gamma^{-1}t)_{\mathrm{mid}})$.
\leanline{\lean{splitGrammulVec}\sep \lean{isotropicFramePairmulVeceqzeroiff}\sep \lean{linkMatrixmulVec}}

\begin{lemma}[the link in coordinates]\label{l:memlink}
Let $F$ be a ring with involution, $\fs$ a shape, $T$ an $\fs$-form over $F$, $\gamma\in\GT{T}$ and
$p=(p_1,p_2)\in F^{\Mx{\fs}}\times F^{\Mx{\fs}}$. Then $p\in\link{\gamma}$ if and only if there is $a\in F^{\varrho}$ such
that $\gamma\,(a,p_2,0)$ has middle part $p_1$ and vanishing $e^*$-part.
\end{lemma}
\begin{proof}
Every $t$ is $\gamma w$ for a unique $w$. The conditions on $t$ in Definition~\ref{d:link} say that $\gamma w$ and
$w$ have vanishing $e^*$-parts, and then $\lmat{\gamma} t=((\gamma w)_{\mathrm{mid}},w_{\mathrm{mid}})$.
\end{proof}
\leanline{\lean{memlink}}

\begin{proof}[Proof of Theorem B(a)]
Let $u,w\in\UT{T}$. We show $\link{\gamma}\subseteq\link{u\gamma w^{-1}}$. Let $t\perp X+\gamma X$. By (G5), $u$ and
$w$ preserve $X$, hence $X^\perp$, and act trivially on $X^\perp/X$. So $ut\perp X$,
$(u\gamma w^{-1})^{-1}(ut)=w\gamma^{-1}t\perp X$, and the middle parts of $ut$ and of $w\gamma^{-1}t$ are those of
$t$ and of $\gamma^{-1}t$. Thus $\lmat{\gamma} t\in\link{u\gamma w^{-1}}$. Applying the inclusion to
$(u^{-1},u'^{-1},\gamma)$ and to $(u,u',u^{-1}\gamma u')$ gives the two inclusions of the theorem.
\end{proof}
\leanline{\lean{linkbesselUmul}}

\begin{lemma}[normal forms]\label{l:nf}
Let $F$ be a ring with involution, $\fs$ a shape, $T$ an $\fs$-form over $F$, $\gamma,u,u'\in\GT{T}$, and $j$ a natural
number.

\emph{(a)} If $\IsNF{j}{\gamma}$, then $\IsNF{j}{\gamma^{-1}}$.

\emph{(b)} If $u\in\UT{T}$ and $\IsNF{j}{\gamma}$, then $\IsNF{j}{u\gamma}$.

\emph{(c)} Let $j\le\varrho+1$ and suppose $\IsNF{j}{\gamma}$. Then for every $y\in F^{\varrho+1}$ there are
$x,y'\in F^{\varrho+1}$ with $y'_a=y_a$ for all frame indices $a\ge j$ and $\mathfrak fx=(\gamma\mathfrak f)y'$.

\emph{(d)} If $u\in\UT{T}$ and $u'\in\UT{T}$, then for every $N\in F$,
$\fdisc{N}{\ZT{T}{u^{-1}\gamma u'}}=\fdisc{N}{\ZT{T}{\gamma}}$.
\end{lemma}
\begin{proof}
(a) Let $c$ with $c\cj{c}=1$ be as in Definition~\ref{d:nf}, and for $a\ge j$ let $x,y$ be the vectors with
$\mathfrak fx=(\gamma\mathfrak f)y$, $x_a=c$, $y_a=1$. Then $\mathfrak f(\cj{c}y)=(\gamma^{-1}\mathfrak f)(\cj{c}x)$, the
coefficients at $a$ are $\cj{c}$ and $\cj{c}c=1$, and those above $a$ vanish; so $\cj{c}$ serves for $\gamma^{-1}$.
(b) Keep $c$ and $y$, and replace $x$ by the vector $x'$ with $\mathfrak fx'=u\,\mathfrak fx$. By (G5),
$u\fr{a'}-\fr{a'}\in\Sp(\fr{a''}:a''<a')$, so $x'_a=x_a$ and $x'_{a'}=0$ for $a'>a$; and
$\mathfrak fx'=u(\gamma\mathfrak f)y=((u\gamma)\mathfrak f)y$. (c) Downward induction on $j$. For $j=\varrho+1$ take
$x=y'=0$. Let $j\le\varrho$ and let $x,y''$ solve the problem for $j+1$, which is possible because normal form from
$j$ implies normal form from $j+1$. With $X_j$, $Y_j$ the vectors of Definition~\ref{d:nf} for $a=j$, so that
$(Y_j)_j=1$ and both vanish above $j$, put $x'=x+(y_j-y''_j)X_j$ and $y'=y''+(y_j-y''_j)Y_j$. (d) By (G4) the
determinant of $\ZT{T}{\gamma}$, which is its leading minor of order $\varrho+1$, and its leading minor of order
$\varrho$ do not change.
\end{proof}
\leanline{\lean{invb}\sep \lean{besselUmul}\sep \lean{existsframeImageeqofisNormalForm}\sep
\lean{frameDiscframeGrambesselUmul}}

\begin{corollary}\label{c:stablelevel}
Let $E$ be a CM field with a CM type, $\fs$ a shape, $S\in\Fill{E}{\fs}$ and $j\in\{0,\dots,\varrho+1\}$. The set of the
$\gamma\in\GS{S}{E}$ with $\HasLev{S}{j}{\gamma}$ is stable: $\IsStab{S}{\{\gamma:\HasLev{S}{j}{\gamma}\}}$.
\end{corollary}
\begin{proof}
Let $u,u'\in\UT{\Trat{S}{E}}$. From $\IsNF{j}{\gamma}$ one gets normal form from $j$ of $u^{-1}\gamma$, of
$\gamma^{-1}u$, of $u'^{-1}\gamma^{-1}u$ and of $u^{-1}\gamma u'$, by Lemma~\ref{l:nf}(b), (a), (b), (a); the converse is
the same with $u^{-1}$, $u'^{-1}$. The minor $\lmin{j}{\Zf{S}{\gamma}}$ and $\tdisc{S}{\gamma}$ are unchanged by (G4) and
Lemma~\ref{l:nf}(d).
\end{proof}
\leanline{\lean{isStablehasLevel}}

\statement{ana}{\textbf{Input (L1)}, unit~112.\enspace Let $P$, $P'$ be matrices over a field with the same finite index
sets, $s$, $s'$ sets of column indices, and $x^0$, $y^0$ column vectors. If every row vector $l$ with $(lP)_b=0$ for
$b\in s$ and $(lP')_b=0$ for $b\in s'$ satisfies $lPx^0=lP'y^0$, then there are $x$, $y$ with $Px=P'y$, $x_b=x^0_b$ for
$b\notin s$ and $y_b=y^0_b$ for $b\notin s'$.}

\begin{proof}[Proof of Theorem B(b)]
Write $T=\Trat{S}{E}$, $h=h_T$, and $p=(p_1,p_2)$. For vectors $s$, $t$ with vanishing $e^*$-parts,
$h(s,t)=\adj{s_{\mathrm{mid}}}\,T\,t_{\mathrm{mid}}$. Let $p=\lmat{\gamma} t$ and $q=\lmat{\gamma} s$ be in the link. Then
by (G6)
\[\adj{q}\,\lform{S}\,p=h(s,t)-h(\gamma^{-1}s,\gamma^{-1}t)=0 .\]
Conversely let $p$ be orthogonal to the link. By Lemma~\ref{l:memlink} we need $x=(a',p_1,0)$ and $y=(a,p_2,0)$ with
$x=\gamma y$. Apply (L1) to $P=1$, $P'=\gamma$, the sets $s=s'$ of the indices of the $e_a$, and $x^0=(0,p_1,0)$,
$y^0=(0,p_2,0)$. A row vector $l$ is $h(t,\cdot)$ for a vector $t$, since $\ratG{S}$ is hermitian and invertible. The
conditions on $l$ say $h(t,e_a)=0$ and $h(t,\gamma e_a)=0$ for all $a$, that is, $t\perp X+\gamma X$; then
$q=\lmat{\gamma} t$ lies in the link, and
\[lx^0-l\gamma y^0=h(t,x^0)-h(\gamma^{-1}t,y^0)=\adj{t_{\mathrm{mid}}}Tp_1-\adj{(\gamma^{-1}t)_{\mathrm{mid}}}Tp_2=\adj{q}\,\lform{S}\,p=0.\]
So (L1) applies.
\end{proof}
\leanline{\lean{memlinkifforthogonal}}

\statement{ana}{\textbf{Input (L2)}, unit~119.\enspace Let $A$ be a field with involution, $T$ an $\fs$-form over $A$,
$p$ a sequence of units of $A$ and $\sigma_0\in\Hl{\varrho}{T}$. The matrix of $\weyl{\varrho}{p}\,\sigma_0$ maps
$(a,m,a')\mapsto\bigl((\cj{p_b}^{-1}a'_b)_b,\ \sigma_0^{\mathrm{mid}}m,\ (p_ba_b)_b\bigr)$, where $\sigma_0^{\mathrm{mid}}$ is
the middle block of $\sigma_0$; and $\ZT{T}{\weyl{\varrho}{p}\,\sigma_0}$ is the diagonal matrix with the entries
$p_0,\dots,p_{\varrho-1}$ and $\ZT{T}{\sigma_0}_{\varrho\varrho}=h_T(z,\sigma_0z)$.}

\begin{proposition}[the big cell]\label{p:bigcell}
Let $A$ be a field with involution, $\fs$ a shape, $T$ an $\fs$-form over $A$, $p$, $q$ sequences of units of $A$ and
$\sigma_0,\sigma_1\in\GT{T}$ with $\sigma_0\in\Hl{\varrho}{T}$ and $\sigma_1\in\Hl{\varrho}{T}$.

\emph{(a)} For $N\in A$, with $t=\ZT{T}{\sigma_0}_{\varrho\varrho}$ and $d=\prod_{a<\varrho}p_a$:
$\fdisc{N}{\ZT{T}{\weyl{\varrho}{p}\,\sigma_0}}=d\,\cj{d}\,(t\,\cj{t}-N^2)$.

\emph{(b)} $\link{\weyl{\varrho}{p}\,\sigma_0}=\link{\weyl{\varrho}{q}\,\sigma_1}$ if and only if $\sigma_0=\sigma_1$.
\end{proposition}
\begin{proof}
(a) By (L2) the frame Gram matrix is diagonal with determinant $dt$ and leading minor $d$ of order $\varrho$.
(b) By Lemma~\ref{l:memlink} and (L2), $(p_1,p_2)\in\link{\weyl{\varrho}{p}\sigma_0}$ if and only if there is $a$ such
that $\weyl{\varrho}{p}\sigma_0(a,p_2,0)=(0,\sigma_0^{\mathrm{mid}}p_2,(p_ba_b)_b)$ has middle part $p_1$ and no
$e^*$-part: this forces $a=0$ and $p_1=\sigma_0^{\mathrm{mid}}p_2$. So the link is the graph of the middle block of
$\sigma_0$, whatever $p$ is, and two elements of $\Hl{\varrho}{T}$ with the same middle block are equal.
\end{proof}
\leanline{\lean{frameDiscweylElemmul}\sep \lean{linkbigcelleqiff}}

\sect{Lattices in subspaces and Siegel's lemma}{s:siegel}
\begin{lemma}[a basis adapted to a kernel]\label{l:extend}
Let $K_0$ be a field, $P$ and $\Omega$ finite sets, $A$ an $\Omega\times P$ matrix and $x$ a $P\times\Omega$ matrix
over $K_0$ with $\det(Ax)\ne0$, and $L\subseteq K_0^{P}$ a subspace that contains the columns of $x$. Then there are
a natural number $s$ and a matrix $y$ with $s$ columns such that: $s=\dim(L\cap\ker A)$; $(x\mid y)$ has a left
inverse; the columns of $(x\mid y)$ span $L$; the columns of $y$ lie in $L$; $Ay=0$; and if $yt$ lies in the span
of the columns of $x$ for a vector $t$, then $t=0$.
\end{lemma}
\begin{proof}
Let $y$ be a matrix whose columns are a basis of $L\cap\ker A$, with a left inverse $l$, and put
$C=(Ax)^{-1}A$. Then $Cx=1$ and $Cy=0$, since $Ay=0$. The matrix with the rows $C$ and $l-lxC$ is a left inverse of
$(x\mid y)$. Every $v\in L$ is $x(Cv)+(v-xCv)$ with $v-xCv\in L\cap\ker A$. If $yt=xw$, then $w=Cxw=Cyt=0$, so
$yt=0$ and $t=l(yt)=0$.
\end{proof}
\leanline{\lean{existsbasisextendker}\sep \lean{muleqzeroofrangeleker}}

\statement{ana}{\textbf{Inputs (lattices).}\enspace Let $E$ be a stably unramified CM field.
(La1), unit~120: every subspace of $K_0^{P}$, $K_0$ a field, is the column space of a matrix with a left inverse.
(La2), unit~57: for a matrix $M$ over $\AEf$ with a left inverse there is an invertible matrix $a$ over $\AEf$
such that $Ma$ is primitive over $\hO$ and the columns of $a$ span the lattice of the $v$ with $Mv$ integral. With
(La1) this gives the matrix $\slb{H}{V}$ of Definition~\ref{d:sublattice}: take a matrix $y$ over $E$ with a left
inverse $l$ whose columns span $V$, and apply (La2) to $M=H\fadm{y}$, which has the left inverse $\fadm{l}\,L$.
(La3), unit~57: $\gix{\det(\adj{J}MJ)}=\gix{\det M}\,\gix{\det J}^2$ for square matrices over $\AEf$; for an integral
invertible matrix $M$, the index of its column lattice in $\hO^{s}$ is $\gix{\det M}$; indices of lattices are
multiplicative in towers, and the index of the images of two lattices under a linear map divides their index; and
(units~57 and~123) $\gix{x}=\abs{N_{E/\Q}(x)}=\prod_\tau\abs{\at{x}{\tau}}^2$ for $x\in E^\times$.
(La4), unit~57, \emph{Jacobi's identity for complements}: let $J$ be a hermitian matrix that is unimodular over
$\hO$, $F$ an integral matrix and $\mathsf n$ a primitive matrix with $\adj{F}J\mathsf n=0$ whose numbers of columns
add up to the size of $J$, and suppose $F$ has a left inverse. Then
$\gix{\det(\adj{\mathsf n}J\mathsf n)}\cdot i^2=\gix{\det(\adj{F}JF)}$, where $i$ is the index of
$\hO^{s}\cap\{v:Fv\text{ integral}\}$ in $\{v:Fv\text{ integral}\}$.}

\begin{proposition}[discriminants and indices]\label{p:discindex}
Let $E$ be a stably unramified CM field, $P$, $\Omega$, $\Omega'$ finite sets, $H$ a $P\times\Omega$ matrix over $\AEf$
with a left inverse $L$, $LH=1$, $V\subseteq E^{\Omega}$ a subspace, $s$ a natural number with $s=\dim_EV$, $x$ an
$\Omega\times s$ matrix over $E$ whose columns lie in $V$, $w$ an invertible $s\times s$ matrix over $\AEf$, and $D$ a
real number with
\[\gix{\det\fadm{\adj{x}x}}\cdot\gix{\det w}^2=D\qquad\text{and}\qquad D\ne0 .\]
Let $Y$ be an $\Omega'\times P$ matrix over $\AEf$ and $K$ an $\hO$-submodule of $\AEf^{\Omega'}$. Put
$\mathsf n=H\,\fadm{x}\,w$. If $\mathsf n$ is integral and every column of $Y\mathsf n$ lies in $K$, then, with $i$ the
index of $K\cap\Lambda$ in the column lattice $\Lambda$ of $Y\,H\,\slb{H}{V}$,
\[\bigl(\sld{H}{V}\cdot i\bigr)^2\le D,\qquad\sld{H}{V}>0,\qquad i\ge1 .\]
\end{proposition}
\begin{proof}
Let $\mathrm m=\slb{H}{V}=\fadm{y}\,a$ with $y$, $a$ as in Definition~\ref{d:sublattice}, and let $l$ be an integral left
inverse of $H\mathrm m$. The columns of $x$ lie in $V$, the column space of $y$, so $x=yB$ for a matrix $B$ over
$E$, and $\fadm{x}\,w=\mathrm m\,A$ with $A=a^{-1}\fadm{B}\,w$. By (La3),
\[D=\gix{\det(\adj{w}\,\fadm{\adj{x}x}\,w)}=\gix{\det(\adj{A}\,\adj{\mathrm m}\mathrm m\,A)}=\sld{H}{V}^2\cdot\gix{\det A}^2 .\]
As $D\ne0$, $\gix{\det A}\ne0$ and $\sld{H}{V}\neq0$, so $A$ is invertible and $\sld{H}{V}>0$. Moreover
$A=l\,(H\mathrm mA)=l\,\mathsf n$ is integral, so by (La3) $\gix{\det A}$ is the index of the column lattice of $A$
in $\hO^{s}$. The map $v\mapsto YH\mathrm mv$ carries $\hO^s$ onto $\Lambda$ and the column lattice of $A$ into
$K\cap\Lambda$, the columns of $Y\mathsf n=YH\mathrm mA$ lying in $K$. Hence $i$ divides the index of the image of
the column lattice of $A$ in $\Lambda$, which divides $\gix{\det A}$; so $1\le i\le\gix{\det A}$ and
$(\sld{H}{V}\,i)^2\le D$.
\end{proof}
\leanline{\lean{sublatticeDiscmulindexsqle}}

\statement{ana}{\textbf{Inputs (geometry of numbers).}\enspace Let $E$ be a CM field with $f=\ninf{E}$.
(M1), unit~121, \emph{Minkowski's second theorem in adelic form}: let $s$ be a natural number, $r_\tau$ an invertible
complex $s\times s$ matrix for each infinite place $\tau$ and $a$ an invertible $s\times s$ matrix over $\AEf$. Then
there are a basis $B_1,\dots,B_s$ of $E^s$ and integers $z_1,\dots,z_s$ such that $a^{-1}B_j\in\hO^s$,
$\sum_\tau\nrm{r_\tau^{-1}\at{(B_j)}{\tau}}^2\le4^{z_j}$ for all $j$, and
$\prod_\tau\abs{\det r_\tau}^{2}\cdot\gix{\det a}^{-1}\cdot4^{f\sum_jz_j}\le(8fs\,\rdE{E})^{fs}$. The second reader of
this write-up compared it with Bombieri and Vaaler~\cite{bv}, Theorem~3, of which it is a weaker form.
(M2), unit~120, Gram--Schmidt: a complex matrix $X$ with a left inverse has $Xr$ with orthonormal columns for some
invertible $r$.
(M3), unit~121: for $u\in\Gf{E}{\kappa}$, every non-zero $x\in E^\kappa\cap u\hO^\kappa$ has $\nrm{x}\ge1$.
(M4), unit~114: if points of a real normed space of dimension $d$ have norms and mutual distances at least $1$, then
for every integer $\mathfrak k\ge d+2$ the sum of their norms to the power $-\mathfrak k$ is at most $5^{\mathfrak k}$.}

\begin{proof}[Proof of Theorem C(a)]
Let $\mathrm m=\slb{u^{-1}}{V}=\fadm{x}\,b$, with $x$ a matrix over $E$ with a left inverse $l_0$ whose columns span
$V$, and $b$ invertible over $\AEf$. The orthogonal complement $V^\perp=\ker\adj{x}$ has dimension $s$, because
$\adj{x}$ is surjective, with right inverse $\adj{l_0}$. Let $\mathrm m'=\slb{u^{-1}}{V^\perp}=\fadm{y}\,a$, with $y$ a
matrix over $E$ with a left inverse $l$ whose columns span $V^\perp$, and $a$ invertible. Then $\adj{y}x=0$, and
since $\ker\adj{x}$ lies in the column space of $y$ there are matrices $f_0$, $t_0$ with $yf_0+t_0\adj{x}=1$; hence a
vector $v$ with $\adj{y}v=0$ is $v=x(\adj{t_0}v)\in V$, and $V=\ker\adj{y}$.

\emph{The two discriminants agree.} Apply (La4) with $J=1$, $F=u^{-1}\mathrm m$ and $\mathsf n=u^{-1}\mathrm m'$: both
are primitive, their numbers of columns add up to $\#\kappa$, and
$\adj{F}\mathsf n=\adj{\mathrm m}\mathrm m'=\adj{b}\,\fadm{\adj{x}y}\,a=0$ because $u$ is unitary. Since $F$ is
primitive, $Fv$ is integral only for integral $v$, so the index in (La4) is $1$ and
\begin{equation}\label{e:jacobi}
\sld{u^{-1}}{V^\perp}^2=\gix{\det(\adj{\mathrm m'}\mathrm m')}=\gix{\det(\adj{\mathrm m}\mathrm m)}=\sld{u^{-1}}{V}^2 .
\end{equation}

\emph{Short vectors in $V^\perp$.} By (M2) choose for each $\tau$ an invertible $r_\tau$ such that $\at{y}{\tau} r_\tau$
has orthonormal columns; then $\abs{\det r_\tau}^{-2}=\det\at{(\adj{y}y)}{\tau}$. Apply (M1) to $(r_\tau)_\tau$ and $a$,
and put $X_j=yB_j$. These vectors span $V^\perp$, so $V=\ker\adj{y}$ is the set of the $v$ with $\adj{X_j}v=0$ for all
$j$. They are non-zero, as $y$ has a left inverse, and $u^{-1}X_j=(u^{-1}\mathrm m')(a^{-1}B_j)$ is integral. Since
$\at{y}{\tau} r_\tau$ is an isometry,
\[\nrm{X_j}^2=\sum_\tau\nrm{\at{y}{\tau} r_\tau\cdot r_\tau^{-1}\at{(B_j)}{\tau}}^2=\sum_\tau\nrm{r_\tau^{-1}\at{(B_j)}{\tau}}^2\le4^{z_j}.\]
Finally, by (La3) and \eqref{e:jacobi},
\[\sld{u^{-1}}{V}^2=\gix{\det(\adj{a}\,\fadm{\adj{y}y}\,a)}=\prod_\tau\abs{\det\at{(\adj{y}y)}{\tau}}^2\cdot\gix{\det a}^2
=\Bigl(\prod_\tau\abs{\det r_\tau}^{-2}\,\gix{\det a}\Bigr)^2 .\]
So
\[\Bigl(\prod_j\nrm{X_j}\Bigr)^{2f}\le4^{f\sum_jz_j}\le(8fs\,\rdE{E})^{fs}\,\prod_\tau\abs{\det r_\tau}^{-2}\gix{\det a}
=(8fs\,\rdE{E})^{fs}\,\sld{u^{-1}}{V}.\]
\end{proof}
\leanline{\lean{siegellemma}}

\begin{proof}[Proof of Theorem C(b)]
Write $k=\#\kappa$ and $\Lambda_u=E^\kappa\cap u\hO^\kappa$. If $s>k$, no subspace $V$ has $\dim V+s=k$, the sum is
empty, and the bound holds because its right side is at least~$1$. Let $s\le k$. For each $V$ with $\dim V+s=k$ choose
$X(V)=(X_1,\dots,X_s)$ as in (a). The subspace $V$ is determined by $X(V)$, and $\prod_j\nrm{X_j}\ge1$ by (M3). Raising
the bound of (a) to the power $k+1$ and using $(8fs\,\rdE{E})^{fs(k+1)}\le(8fk\,\rdE{E})^{f(k+1)s}$, which holds as
$s\le k$,
\[\sld{u^{-1}}{V}^{-(k+1)}\le(8fk\,\rdE{E})^{f(k+1)s}\prod_{j=1}^s\nrm{X_j}^{-2f(k+1)} .\]
Summing over $V$, and then over all $s$-tuples of non-zero vectors of $\Lambda_u$, the right side is at most
$(8fk\,\rdE{E})^{f(k+1)s}\bigl(\sum_{x}\nrm{x}^{-2f(k+1)}\bigr)^s$, the sum over the non-zero $x\in\Lambda_u$. These
points of the real vector space $\CC^{f\kappa}$ of dimension $2fk$ have norms and mutual distances at least $1$ by
(M3), and $2f(k+1)\ge2fk+2$; by (M4) the inner sum is at most $5^{2f(k+1)}=25^{f(k+1)}$. As $25\cdot8=200$, the
claim follows.
\end{proof}
\leanline{\lean{tsumsublatticeDiscinvle}}

\sect{Strong approximation and terms in the Bessel subgroup}{s:approx}
\statement{ana}{\textbf{Inputs (the unipotent group, floor by floor).}\enspace Let $A$ be a ring with involution with an
element $t$, $t+\cj{t}=1$, and $T$ an $\fs$-form over $A$ with $\JT{T}$ hermitian and invertible. For $l<\varrho$ and a
pair $(v,\beta)$ of a vector $v\in A^{\Ix{\fs}}$ and an element $\beta\in A$ with $\cj{\beta}=-\beta$, the
\emph{Heisenberg element} $\nu_l(v,\beta)\in\GT{T}$ is an Eichler map along $e_l$ whose entries are polynomials in the
coordinates of $v$, their conjugates, $\beta$ and $t$.
(U1), unit~119: $\nu_l(v,\beta)\in\UT{T}$, and $\UT{T}$ is generated by these elements: if a subset of $\GT{T}$ contains
$1$ and is stable under left multiplication by all $\nu_l(v,\beta)$, $l<\varrho$, then it contains $\UT{T}$.
(U2), unit~29: if $\JT{T}$ is invertible over a Hausdorff topological ring with continuous involution, the topology of
$\GT{T}$ is induced by the entries.
(U3), unit~1: for a subring $Q\subseteq A$, a unit $g$ of $A$, and $k\in\GT{T}$ with entries in $Q$ together with its
inverse and $g^{-1}(k-1)$ with entries in $Q$: $g^{-1}\sd{k}\in Q$; and if $k=uh$ with $u\in\UT{T}$ and $h\in\HT{T}$,
then $h$ has the same two properties.
(U4), unit~112: let $\Gamma\to G$ be a homomorphism, $\mathsf{K}\subseteq G$ a subgroup and $(r_c)_c$ representatives of the
double cosets $\Gamma\backslash G/\mathsf{K}$. Every $g\in G$ is $\gamma r_ck$ with $\gamma\in\Gamma$, a class $c$ and $k\in \mathsf{K}$.
If $r_d^{-1}\gamma r_d\in \mathsf{K}$ only for $\gamma=1$, then $r_c^{-1}\gamma r_d\in \mathsf{K}$ implies $c=d$ and $\gamma=1$.}

\begin{proof}[Proof of Theorem A(a)]
$\JT{T'}$ is hermitian and invertible, as the image of $\JT{T}$. Let $\mathcal C$ be the closure in $\GT{T'}$ of the
image of $\UT{T}$; it is a subgroup. By (U1) for $T'$, with $t=\tfrac12$, it suffices to show
$\nu_l(v,\beta)\in\mathcal C$ for all $l<\varrho$, all $v$ over $\AEf$ and all skew $\beta\in\AEf$. A skew $\beta$ is
$\beta'-\cj{\beta'}$ with $\beta'=\tfrac12\beta$. The map $(v,\beta')\mapsto\nu_l(v,\beta'-\cj{\beta'})$ is continuous
by (U2), its entries being polynomials. The pairs with entries in $E$ are dense by (Ad1), and for such a pair
$\nu_l(v,\beta'-\cj{\beta'})$ is the image of the Heisenberg element over $E$ (an identity of matrices that the formal
text checks entry by entry), which lies in $\UT{T}$ by (U1). Hence all $\nu_l(v,\beta)$ lie in $\mathcal C$.
\end{proof}
\leanline{\lean{besselUdense}}

\begin{proof}[Proof of Theorem A(b)]
By Theorem~E(a) and (Ad1), (Ad2), $\KgS{S}$ is open in $\GFin{S}$: integrality of a matrix and of $\gid{S}^{-1}(k-1)$
are open conditions. The map $\zeta\mapsto y^{-1}x^{-1}\,\tolb{S}(\zeta)\,y$ from $\GS{S}{\AEf}$ to $\GFin{S}$ is
continuous, so the preimage of $\KgS{S}$ is open; it contains $\zeta_0=\tolb{S}^{-1}(x)$, which lies in
$\UT{\Trat{S}{\AEf}}$ by Lemma~\ref{l:dictUH}(a). By Theorem~A(a) for $T=\Trat{S}{E}$ and $T'=\Trat{S}{\AEf}$, the point
$\zeta_0$ is a limit of images of elements of $\UT{\Trat{S}{E}}$; one of them, $u$, has its image in the open set, and
$\tolb{S}$ of this image is $\rfin{S}(u)$.
\end{proof}
\leanline{\lean{existsbesselUapprox}}

\begin{lemma}[characters and the principal level]\label{l:charlevel}
Let $S\in\Fill{E}{\fs}$.

\emph{(a)} Let $u,v\in\GFin{S}$ with $u\in\UT{\Tlat{S}{\AEf}}$ and $v\in\UT{\Tlat{S}{\AEf}}$, and $g_0\in\Gf{E}{\kappa}$ with
$u^{-1}\bigl(\incH{S}(g_0)^{-1}\,v\,\incH{S}(g_0)\bigr)\in\KgS{S}$. Then $\uchar{S}{v}=\uchar{S}{u}$.

\emph{(b)} Let $c\in\Cl{S}$ and $u\in\GFin{S}$ with $u\in\UT{\Tlat{S}{\AEf}}$. Then there is $u_0\in\UT{\Trat{S}{E}}$ with
$u^{-1}\,\pos{S}{c}{c}{u_0}\in\KgS{S}$ and $\uchar{S}{\rfin{S}(u_0)}=\uchar{S}{u}$.
\end{lemma}
\begin{proof}
(a) Put $h=\incH{S}(g_0)\in\HT{\Tlat{S}{\AEf}}$ and $k=u^{-1}h^{-1}vh\in\KgS{S}$. By Theorem~E(a) and (U3),
$\gid{S}^{-1}\sd{k}\in\hO$, so $\sdc{S}(\gid{S}^{-1}\sd{k})=1$ by Lemma~\ref{l:char}(a). All elements involved lie in
$\ST{\Tlat{S}{\AEf}}$, where $\sd{\cdot}$ is additive and vanishes on $\HT{\Tlat{S}{\AEf}}$ by (G2). Hence
$\sd{v}=\sd{h^{-1}vh}=\sd{uk}=\sd{u}+\sd{k}$ and
$\uchar{S}{v}=\sdc{S}(\gid{S}^{-1}\sd{u})\,\sdc{S}(\gid{S}^{-1}\sd{k})=\uchar{S}{u}$.
(b) The element $x=\hc{S}{c}\,u\,\hc{S}{c}^{-1}$ lies in $\UT{\Tlat{S}{\AEf}}$ by (G2). Theorem~A(b), with $y=\hc{S}{c}$, gives
$u_0\in\UT{\Trat{S}{E}}$ with $\hc{S}{c}^{-1}x^{-1}\rfin{S}(u_0)\hc{S}{c}=u^{-1}\pos{S}{c}{c}{u_0}\in\KgS{S}$. Now (a) applies to $u$,
$v=\rfin{S}(u_0)$ and $g_0=\dot c$.
\end{proof}
\leanline{\lean{unipotentChareqofmemprincipalLevel}\sep \lean{existsratapproxunipotentChareq}}

\begin{proof}[Proof of Theorem D(a)]
By (G3) over $E$, with $t=\tfrac12$, write $\gamma=u\,\iota_N(\eta)$ with $u\in\UT{\Trat{S}{E}}$ and
$\eta\in\Uni_\kappa(E)$. By Lemma~\ref{l:dictUH}(b),
\[\pos{S}{c}{d}{\gamma}=\bigl(\hc{S}{c}^{-1}\rfin{S}(u)\hc{S}{c}\bigr)\cdot\incH{S}\bigl(\dot c^{-1}\eta\,\dot d\bigr),\]
the first factor in $\UT{\Tlat{S}{\AEf}}$ by Lemma~\ref{l:dictUH}(c), the second in $\HT{\Tlat{S}{\AEf}}$. The product lies
in $\KgS{S}$; by Theorem~E(a) and (U3) so does the second factor. By (A1), $\dot c^{-1}\eta\,\dot d\in\KcolS{S}$. By (F1)
and the injectivity of $\Uni_\kappa(E)\to\Gf{E}{\kappa}$, the hypothesis of the second part of (U4) holds for
$\Gamma=\Uni_\kappa(E)$, $\mathsf{K}=\KcolS{S}$; hence $c=d$ and $\eta=1$. So $\gamma=u$, and $\pos{S}{c}{c}{\gamma}\in\KgS{S}$ says
$\gamma\in\clU{S}{c}$.
\end{proof}
\leanline{\lean{eqandmemclassUnipofbesselS}}

\begin{proof}[Proof of Theorem D(b)]
Let $s=\redS{S}(k)$ with $k\in\KS{S}$, $k\in\ST{\Tlat{S}{\AEf}}$. Then $\hc{S}{c}k\in\ST{\Tlat{S}{\AEf}}$, so by (G3) and
Lemma~\ref{l:dictUH}(b), $\hc{S}{c}k=u\,\incH{S}(g_0)$ with $u\in\UT{\Tlat{S}{\AEf}}$ and $g_0\in\Gf{E}{\kappa}$. Theorem~A(b),
with $x=u$ and $y=\incH{S}(g_0)$, gives $x_0\in\UT{\Trat{S}{E}}$ with
$k'=\incH{S}(g_0)^{-1}u^{-1}\rfin{S}(x_0)\incH{S}(g_0)\in\KgS{S}$. By the first part of (U4), $g_0=\eta\,\dot d\,a$ with
$\eta\in\Uni_\kappa(E)$, a class $d$ and $a\in\KcolS{S}$. Put $\gamma=x_0\,\iota_N(\eta)$; it lies in $\ST{\Trat{S}{E}}$.
Using $\rfin{S}(x_0)=u\,\incH{S}(g_0)\,k'\,\incH{S}(g_0)^{-1}$, $\hc{S}{c}^{-1}u\,\incH{S}(g_0)=k$ and
$\incH{S}(g_0)^{-1}\incH{S}(\eta)\hc{S}{d}=\incH{S}(a)^{-1}$,
\[\pos{S}{c}{d}{\gamma}=\hc{S}{c}^{-1}\rfin{S}(x_0)\,\incH{S}(\eta)\,\hc{S}{d}=k\,k'\,\incH{S}(a)^{-1}.\]
Here $k'\in\KgS{S}$ and $\incH{S}(a)\in\KgS{S}$ by (A1). So the position lies in $\KS{S}$ and has the reduction
$\redS{S}(k)=s$.
\end{proof}
\leanline{\lean{existsbesselSmemtermsOver}}

\statement{ana}{\textbf{Input (N1)}, unit~50, neatness; proved below.\enspace Let $E$ be a CM field with a CM type and
at least two infinite places, $\Omega$ a finite set, $w_1,\dots,w_{\mathfrak k}$ distinct finite places of $E$ with
$(2\#\Omega)^{[E:\Q]}<\prod_jq_{w_j}$, and $\mathsf{K}\subseteq\Kf{E}{\Omega}$ a subgroup of elements that reduce to $1$ modulo
every $w_j$. Then $\Gamma_f(E,\Omega)\cap gKg^{-1}=1$ for every $g\in\Gf{E}{\Omega}$. This is the statement of the
background library. The catalogue of the write-up names Serre~\cite{serre}, appendix, for the method (place not
confirmed); by the second reader's verdict Serre's statement, on automorphisms of finite order that are congruent
to $1$ modulo an integer at least~$3$, is a different one, and the relation is unresolved. No reference is needed:
the statement has the short proof that follows, which is the argument of the formal proof (a trace, its norm, the
eigenvalues at the infinite places).}

\begin{proof}[Proof of (N1)]
Let $\gamma\in\Gamma_f(E,\Omega)\cap gKg^{-1}$ be the image of $\gamma_0\in\Uni_\Omega(E)$, and put
$t=\operatorname{tr}\gamma_0-\#\Omega\in E$. The matrix $g^{-1}\gamma g$ lies in $\mathsf{K}\subseteq\Kf{E}{\Omega}$ and has the
trace of~$\gamma$, so $t\in E\cap\hO=\IntE$; and $t=\operatorname{tr}(g^{-1}\gamma g-1)$ lies in the prime ideal of every
$w_j$, since $g^{-1}\gamma g$ reduces to $1$ modulo~$w_j$. As the $w_j$ are distinct, $t$ lies in the product of
these prime ideals, so $\prod_jq_{w_j}\le\abs{N_{E/\Q}(t)}$ if $t\ne0$. At an infinite place $\tau$ the complex matrix
$\at{(\gamma_0)}{\tau}$ is unitary, because the embedding of the CM type commutes with the complex conjugations; so
its eigenvalues $\lambda_1,\dots,\lambda_{\#\Omega}$ have absolute value~$1$, and
$\abs{\at{t}{\tau}}=\abs{\sum_i\lambda_i-\#\Omega}\le2\#\Omega$. Hence
$\abs{N_{E/\Q}(t)}=\prod_\tau\abs{\at{t}{\tau}}^2\le(2\#\Omega)^{[E:\Q]}$, and the hypothesis forces $t=0$. Then
$\sum_i\operatorname{Re}\lambda_i=\#\Omega$ with $\operatorname{Re}\lambda_i\le1$, so every $\lambda_i$ is~$1$; a unitary
matrix all of whose eigenvalues are $1$ is the identity, so $\at{(\gamma_0)}{\tau}=1$, and $\gamma_0=1$. The second
infinite place was not used; the formal proof needs it because it passes through a lemma on forms that are definite
at the places other than a fixed one.
\end{proof}

\begin{proposition}[freeness]\label{p:free}
Let $E$ be a stably unramified CM field with a CM type, $\fs$ a shape, $S\in\Fill{E}{\fs}$, $\Omega$ a finite set with
$\#\Omega\le\kk{\fs}$, and $i\in\{0,1\}$. Let $\gamma\in\Uni_\Omega(E)$ and $g_0\in\Gf{E}{\Omega}$ be such that
$g_0^{-1}\gamma g_0\in\Kf{E}{\Omega}$ and that $g_0^{-1}\gamma g_0$ reduces to $1$ modulo every $w\in\ChS{S}{i}$. Then
$\gamma=1$.
\end{proposition}
\begin{proof}
Apply (N1) to the $f=\ninf{E}$ places of $\ChS{S}{i}$, all of norm $\qS{S}{i}$ by (F4), and the group $\mathsf{K}$ of all elements of
$\Kf{E}{\Omega}$ that reduce to $1$ modulo them. Since $[E:\Q]=2f$, the condition is $(2\#\Omega)^{2}<\qS{S}{i}$, and
$(2\#\Omega)^2\le4k^2\le4b^2<b^{10}\le\qS{S}{i}$ by (F3), (F8). There are at least two infinite places, as
$f\ge k\ge1000$ by (F7); the formal statement of (N1) asks for this. The image of $\gamma$ lies in $\Gamma_f(E,\Omega)\cap g_0Kg_0^{-1}$, so it is $1$, and
$\Uni_\Omega(E)\to\Gf{E}{\Omega}$ is injective.
\end{proof}
\leanline{\lean{eqoneofreduceAteqone}}

\sect{Frames of terms}{s:frames}
In this section $S\in\Fill{E}{\fs}$, and $N$, $b$, $n$ abbreviate $\NS{S}$, $\bS{S}$, $\nn{\fs}$. Three remarks on frames
are used below. For $s,s'\in\GT{T}$ over any ring with involution: $(ss')\mathfrak f=s\,(s'\mathfrak f)$; for
$u\in\UoT{T}$, $(\gamma u)\mathfrak f=\gamma\mathfrak f$; and $(s^{-1}s')\mathfrak f=\mathfrak f$ if and only if
$s\mathfrak f=s'\mathfrak f$. Also $s\in\ST{T}$ if and only if $(s\mathfrak f)_{i,a}=\mathfrak f_{i,a}$ for all frame
indices $a$ and all $i\in\Ix{\fs}$ with $\lev(i)\ge a$, which is the definition in Setting~\ref{set:shape}.
\leanline{\lean{frameImagemul}\sep \lean{frameImagemulofmembesselU}\sep
\lean{frameImageinvmuleqiff}\sep \lean{membesselSiffframeImage}}

\begin{lemma}[a common vector of two frames]\label{l:common}
Let $S\in\Fill{E}{\fs}$, $\gamma\in\GS{S}{E}$ and $x,y\in E^{\varrho+1}$ with $\mathfrak fx=(\gamma\mathfrak f)y$. Then
$\cj{x_\varrho}\,x_\varrho=\cj{y_\varrho}\,y_\varrho$, and $\Zf{S}{\gamma}\,y$ is the vector with the entry
$-\NS{S}\,x_\varrho$ at the index $\varrho$ and $0$ elsewhere. Moreover, if $v\in E^{\Mx{\fs}}$ has $v_z=0$ and
$\adj{v}\,\Trat{S}{E}\,v=0$, then $v=0$.
\end{lemma}
\begin{proof}
For the form $h$ of $\Trat{S}{E}$, $h(\mathfrak fx,\mathfrak fx)=-N\cj{x_\varrho}x_\varrho$, the $e_a$ being isotropic
and orthogonal to $z$; and $h((\gamma\mathfrak f)y,(\gamma\mathfrak f)y)=h(\mathfrak fy,\mathfrak fy)$ by (G6). As
$N\ne0$, the first claim follows. The entry of index $a$ of $\Zf{S}{\gamma}\,y$ is
$h(\fr{a},(\gamma\mathfrak f)y)=h(\fr{a},\mathfrak fx)$, which is $0$ for $a<\varrho$ and $-Nx_\varrho$ for
$a=\varrho$. In the last claim $\sum_{i\in\kappa}\abs{\at{(v_i)}{\tau}}^2=0$ at any infinite place.
\end{proof}
\leanline{\lean{starmullasteqofframeImageeq}\sep \lean{frameProductsmulVecofframeImageeq}\sep
\lean{eqzeroofratMiddleformeqzero}}

\begin{proposition}[the norm identity]\label{p:normid}
Let $S\in\Fill{E}{\fs}$, $\gamma\in\GS{S}{E}$, $y\in E^{\varrho+1}$, $D\in E$ and $t\in E^{\Ix{\fs}}$ such that
$\Zf{S}{\gamma}\,y$ is the vector with the entry $D$ at the index $\varrho$ and $0$ elsewhere, and
$(\gamma\mathfrak f)y=t$. Then the $e^*$-coordinates of $t$ vanish, $-\NS{S}\,t_z=D$, and for every infinite place
$\tau$,
\[\NS{S}\sum_{i\in\kappa}\abs{\at{(t_{x_i})}{\tau}}^2=\abs{\at{D}{\tau}}^2-\bigl(\NS{S}\,\abs{\at{(y_\varrho)}{\tau}}\bigr)^2;\]
if this number is $\le0$ for some $\tau$, then $t=\mathfrak fx$ for some $x\in E^{\varrho+1}$.
\end{proposition}
\begin{proof}
The $e_a^*$-coordinate of $t$ is $h(e_a,t)=(\Zf{S}{\gamma} y)_a=0$, and $h(z,t)=-Nt_z$ is the entry of index
$\varrho$, which is $D$. As $t$ has no $e^*$-part, $h(t,t)=\sum_i\cj{t_{x_i}}t_{x_i}-N\cj{t_z}t_z$; on the other hand
$h(t,t)=h(\mathfrak fy,\mathfrak fy)=-N\cj{y_\varrho}y_\varrho$ by (G6). Substituting $t_z=-D/N$, multiplying by $N$
and applying the embedding at $\tau$ gives the identity. If the left side is $\le0$ at one place, it is $0$ there,
so all $t_{x_i}$ vanish and $t$ is a combination of the $e_a$ and $z$.
\end{proof}
\leanline{\lean{normsqmiddleparteq}}

\begin{corollary}\label{c:tdisc}
Let $S\in\Fill{E}{\fs}$, $\gamma\in\GS{S}{E}$ and $\tau$ an infinite place of $E$.

\emph{(a)} $\at{\tdisc{S}{\gamma}}{\tau}$ is a non-negative real number: $\at{\tdisc{S}{\gamma}}{\tau}=\abs{\at{\tdisc{S}{\gamma}}{\tau}}$.

\emph{(b)} If $\tdisc{S}{\gamma}\ne0$, then $\HasLev{S}{\varrho+1}{\gamma}$.

\emph{(c)} $\gdr{\Zr{\rinf{S}(\gamma)_\tau}}=\Bigl(\prod_{a\le\varrho}\fscale{S}{\tau}{a}\Bigr)^2\cdot\abs{\at{\tdisc{S}{\gamma}}{\tau}}$.
\end{corollary}
\begin{proof}
(a) Let $Z=\Zf{S}{\gamma}$ and let $y$ be the last column of the adjugate of $Z$, so that $Zy$ has the entry $\det Z$
at $\varrho$ and $0$ elsewhere, and $y_\varrho=\lmin{\varrho}{Z}$. Proposition~\ref{p:normid} with $D=\det Z$ shows that
$\at{\tdisc{S}{\gamma}}{\tau}=\abs{\at{\det Z}{\tau}}^2-N^2\abs{\at{\lmin{\varrho}{Z}}{\tau}}^2$ is $N$ times a sum of squares.
(b) Normal form from $\varrho+1$ is an empty condition. If $\det Z=0$, then $\at{\tdisc{S}{\gamma}}{\tau}\le0$, hence
$=0$ by (a), against the hypothesis; so $\lmin{\varrho+1}{Z}\ne0$. (c) By Lemma~\ref{l:coords}(b) the determinant of
$\Zr{\rinf{S}(\gamma)_\tau}$ is $\prod_{a\le\varrho}\fscale{S}{\tau}{a}\cdot\at{\det Z}{\tau}$ and its leading minor of order
$\varrho$ is $\prod_{a<\varrho}\fscale{S}{\tau}{a}\cdot\at{\lmin{\varrho}{Z}}{\tau}$; and $\fscale{S}{\tau}{\varrho}=1/N$. With (a)
this gives the formula.
\end{proof}
\leanline{\lean{maptermDisceqnorm}\sep \lean{hasLevellast}\sep \lean{gramDiscRealrealFrameGramratToInf}}

\statement{ana}{\textbf{Inputs (the real place)}, unit~132.\enspace For $x,y\in\Gone{\fs}$, frame indices $a,a'$, and
$t\in\Gone{\fs}$: (R1) $\abs{\Zr{x}_{aa'}}\le\nrm{y^{-1}\fr{a}}\cdot\nrm{(x^{-1}y)^{-1}\fr{a'}}$; (R2) $\bsize{ty}=\bsize{y}$ if
$t\in\ST{\Trg{\fs}{1}}$; (R3) $(ty)^{-1}\fr{a}=y^{-1}\fr{a}$ if $t$ lies in the group generated by $\UoT{\Trg{\fs}{1}}$ and
$\HT{\Trg{\fs}{1}}$.}

\begin{lemma}[near terms and the archimedean coordinates]\label{l:near}
Let $S\in\Fill{E}{\fs}$, $\gamma,u,u'\in\GS{S}{E}$ and $\tau$ an infinite place of $E$.

\emph{(a)} If $u\in\UoT{\Trat{S}{E}}$, then $\rinf{S}(u)_\tau\in\UoT{\Trg{\fs}{1}}$.

\emph{(b)} Let $y,t\in\Ginf{E}{\fs}$ with $t_{\tau'}\in\UoT{\Trg{\fs}{1}}$ for every infinite place $\tau'$. Then
$ty\in\bxp{S}$ if and only if $y\in\bxp{S}$.

\emph{(c)} If $\IsNr{S}{\gamma}$, $u\in\UoT{\Trat{S}{E}}$ and $u'\in\UoT{\Trat{S}{E}}$, then $\IsNr{S}{u\gamma u'}$.

\emph{(d)} For a natural number $j$: if $\IsNF{j}{\gamma}$, then $\IsNF{j}{\rinf{S}(\gamma)_\tau}$.
\end{lemma}
\begin{proof}
(a) Conjugation by the diagonal matrix $\bsGL{S}{\tau}$ and the embedding at $\tau$ preserve the entries equal to $0$
or $1$ that define $\UoT{\cdot}$. (b) By (R2), (R3), since $\UoT{\Trg{\fs}{1}}\subseteq\ST{\Trg{\fs}{1}}$. (c) If $y$ and
$\rinf{S}(\gamma)^{-1}y$ lie in $\bxp{S}$, then so do $\rinf{S}(u)y$ and
$\rinf{S}(u\gamma u')^{-1}\rinf{S}(u)y=\rinf{S}(u')^{-1}\rinf{S}(\gamma)^{-1}y$, by (a) and (b). (d) Let $c$ be the
constant of the normal form; then $\abs{\at{c}{\tau}}^2=1$. For a frame index $a\ge j$ and $x,y$ with
$\mathfrak fx=(\gamma\mathfrak f)y$, put $d_{a'}=\bscale{S}{\tau}{\fr{a}}/\bscale{S}{\tau}{\fr{a'}}$ and
$x'_{a'}=d_{a'}\,\at{(x_{a'})}{\tau}$, $y'_{a'}=d_{a'}\,\at{(y_{a'})}{\tau}$. Applying the embedding and the diagonal matrix
$\bscale{S}{\tau}{\fr{a}}\,\bsGL{S}{\tau}^{-1}$ to the equation gives
$\mathfrak fx'=(\rinf{S}(\gamma)_\tau\mathfrak f)y'$; the supports are unchanged, and $x'_a=\at{c}{\tau}$, $y'_a=1$.
\end{proof}
\leanline{\lean{ratToInfmembesselU}\sep \lean{mulmemboxedPointsiff}\sep
\lean{mulofmembesselU}\sep \lean{ratToInf}}

\begin{lemma}[terms in the lattice basis]\label{l:intterm}
Let $S\in\Fill{E}{\fs}$, $c,c'\in\Cl{S}$, $s\in\GRes{S}$ and $\gamma\in\GS{S}{E}$.

\emph{(a)} $\adj{\clC{S}{c}}\ \JT{\Tlat{S}{\AEf}}\ \clC{S}{c}=\JT{\Trat{S}{\AEf}}$.

\emph{(b)} $\pos{S}{c}{c'}{\gamma}\,\mathfrak f=\clC{S}{c}\cdot\fadm{\gamma\mathfrak f}$, and
$\ZT{\Tlat{S}{\AEf}}{\pos{S}{c}{c'}{\gamma}}=\fadm{\Zf{S}{\gamma}}$.

\emph{(c)} If $\gamma\in\terms{S}{c}{s}{c'}$, then $\pos{S}{c}{c'}{\gamma}$ has entries in $\hO$, and for all frame indices $a$,
$a'$ the entry $\Zf{S}{\gamma}_{aa'}$ is an algebraic integer.
\end{lemma}
\begin{proof}
(a) By Lemma~\ref{l:latgram}(a) and (F2), $\adj{\latB{S}}\,\JT{\Trat{S}{\AEf}}\,\latB{S}=\JT{\Tlat{S}{\AEf}}$, and
$\iota_N(\dot c)$ is an isometry of $\Trat{S}{\AEf}$; so $\adj{\clB{S}{c}}\JT{\Trat{S}{\AEf}}\clB{S}{c}=\JT{\Tlat{S}{\AEf}}$, which is
(a) for the inverse matrix. (b) By Lemma~\ref{l:dictUH}(c), $\pos{S}{c}{c'}{\gamma}=\clC{S}{c}\,\gamma\,\clB{S}{c'}$, and
$\clB{S}{c'}$ fixes the frame vectors: $\iota_N(\dot c')$ acts on the $x_i$ only and $\latB{S}$ fixes the $e_a$ and
$z$. This is the first formula. By (G4) and (a),
\begin{align*}
\ZT{\Tlat{S}{\AEf}}{\pos{S}{c}{c'}{\gamma}}&=\adj{\mathfrak f}\,\JT{\Tlat{S}{\AEf}}\,\clC{S}{c}\,\gamma\mathfrak f\\
&=\adj{(\clB{S}{c}\mathfrak f)}\,\JT{\Trat{S}{\AEf}}\,\gamma\mathfrak f=\adj{\mathfrak f}\JT{\Trat{S}{\AEf}}\gamma\mathfrak f .
\end{align*}
(c) The position lies in $\KS{S}$. By (b) the matrix $\fadm{\Zf{S}{\gamma}}$ is a block of the product of the integral
matrices $\JT{\Tlat{S}{\AEf}}$ and $\pos{S}{c}{c'}{\gamma}$; an element of $E$ that is integral at all finite places is an
algebraic integer.
\end{proof}
\leanline{\lean{classCoordconjsplitGram}\sep \lean{frameImagetermPosition}\sep
\lean{frameGramtermPosition}\sep \lean{isIntMatrixtermPosition}\sep
\lean{isIntegralframeProducts}}

\begin{proof}[Proof of Theorem D(c)]
The element $\delta=\gamma^{-1}u$ fixes the frame, so it lies in $\ST{\Trat{S}{E}}$. By Lemma~\ref{l:dictUH}(d) it is a
term over $(c',s^{-1}s,d)=(c',1,d)$, that is, $\pos{S}{c'}{d}{\delta}$ lies in $\KgS{S}$. Theorem~D(a) gives $c'=d$ and
$\delta\in\clU{S}{c'}$; as $\delta$ fixes the frame, $\delta\in\clUo{S}{c'}$.
\end{proof}
\leanline{\lean{eqandmemclassUnip}}

\begin{lemma}[sizes]\label{l:sizes}
Let $S\in\Fill{E}{\fs}$ and $e$ a natural number.

\emph{(a)} If $e\le10^5\,\nn{\fs}^9\,\ninf{E}$, then $\bS{S}^{e}\le\tailS{S}$.

\emph{(b)} $\tailS{S}^4=\bS{S}^{\kk{\fs}\,\ninf{E}}$.

\emph{(c)} $2\le\gS{S}$; $2\le\tailS{S}$; $1000\le\bS{S}$; $0\le\besE{S}$; $0\le\mainE{S}$; $0<\NS{S}$; $0<\ucov{S}$; and
$\redim{S}{\ST{\Tlat{S}{\AEf}}}$ is not empty. (In the formal text: a statement $p$ holds as soon as it follows from
these facts together with $2\le\bS{S}$, $0<\bS{S}$ and $1\le\bS{S}$.)
\end{lemma}
\begin{proof}
(a) $\tailS{S}=b^{kf/4}$ with $b\ge1000>1$, and $e\le10^5n^9f\le1000(2n)^9f/4\le kf/4$ by (F7). (b) is the
definition. (c) $g$ is a product of two primes; $\tailS{S}\ge b\ge2$ by (a) with $e=1$; $\mainE{S}$ is a square divided
by the product of the positive numbers $\bdeg{S}$ and $\mass{E}{r}$ (the mass is positive: background, unit~99);
$N$ is a prime; $\ucov{S}>0$ by Lemma~\ref{l:scales}; and the reduced image is a group.
\end{proof}
\leanline{\lean{powletailScale}\sep \lean{tailScalepowfour}\sep \lean{ofsizebounds}}

\begin{proposition}[frame products of near terms are small]\label{p:nearsize}
Let $S\in\Fill{E}{\fs}$ be admissible filling data, $\gamma\in\GS{S}{E}$ with $\IsNr{S}{\gamma}$, $\tau$ an infinite place of
$E$, $a,a'\le\varrho$ frame indices and $l\le\varrho+1$. Then, with $n=\nn{\fs}$,
\[\abs{\at{(\Zf{S}{\gamma}_{aa'})}{\tau}}\le\bS{S}^{5n^2}\qquad\text{and}\qquad\abs{\at{\lmin{l}{\Zf{S}{\gamma}}}{\tau}}^2\le\bS{S}^{11n^3}.\]
\end{proposition}
\begin{proof}
Let $y\in\bxp{S}$ with $y'=\rinf{S}(\gamma)^{-1}y\in\bxp{S}$. By (R1) for $x=\rinf{S}(\gamma)_\tau$,
$\abs{\Zr{\rinf{S}(\gamma)_\tau}_{aa'}}\le\nrm{y_\tau^{-1}\fr{a}}\,\nrm{y'^{-1}_\tau\fr{a'}}\le\bxs{S}^2=b^{4n^2}$. By
Lemma~\ref{l:coords}(b) the left side is $\abs{\kapS{S}{\tau}}^{2\varrho-a-a'}N^{-1}\abs{\at{(\Zf{S}{\gamma}_{aa'})}{\tau}}$, and
$\abs{\kapS{S}{\tau}}\ge1$, $N\le b^{60n}$ by (A2), (A3). So $\abs{\at{(\Zf{S}{\gamma}_{aa'})}{\tau}}\le b^{60n+4n^2}\le b^{5n^2}$,
since $n\ge10^5$ by (F7). A determinant of size $l$ with entries of absolute value at most $B=b^{5n^2}$ is at most
$l!\,B^l\le(lB)^l$ in absolute value, and $l\le n\le b$ by (F8); so its square is at most
$(b^{1+5n^2})^{2n}=b^{2n+10n^3}\le b^{11n^3}$.
\end{proof}
\leanline{\lean{normframeProductsle}\sep \lean{normmapleadingMinorle}}

\sect{Notes on the formal text}{s:notes}
\subhead{Arguments as the census reads them.} The census of the artifact reads the arguments of a declaration from
its text, without running Lean, and for some declarations of this note the reading is evidently not Lean's
signature: a variable of the file is counted because the declaration binds a variable of the same name, or because
a name such as that of $\bS{S}$ contains its letter; or the reading stops at a definition inside a statement. The
reader of the unit recorded these cases, with Lean's presumable signature and the evidence for it; nothing was
compiled. The definitions concerned are written here with Lean's arguments: the reduction $\redS{S}$ and the maps
$\radl{S}$, $\ainf{S}$, $\afin{S}$ (the census reads the source of the homomorphism as an argument); the scales
$\bscale{S}{\tau}{i}$, the numbers $\fsprod{S}$, $\ucov{S}$, $\bdeg{S}$; the sets $\terms{S}{c}{s}{c'}$, $\tlink{S}{\Lambda}$,
$\frbox{\rho}$ and the properties $\IsBP{S}{c}{\omega}$, $\IsStab{S}{\Xi}$, $\IsNF{j}{\gamma}$, $\IsAdm{S}$ (a bound variable is
counted); the notions $\bcorr{\mu}{W}{\gamma}$, $\SchurLe{W}{D}$, $\HasSB{W}{\mathrm{rep}}{B}{D}$, $\CutPair{b}{f}{\mathrm{rep}}{W}$,
$\BSys{\fs}{b}{f}$ (variables of the file are counted, for a Bessel system the components $I$, $W$, $\deg$ of the record
itself); the matrix $\slb{H}{V}$ (the hypothesis $LH=1$ is missed). The Bessel function $\besfn{\sigma}{S}{y}{i}$, the kernel
$\besker{\sigma}{S}{\gamma}{y}$, the amplitudes and the amplitude sum take the side as an ambient argument in Lean, which
the census does not list; the notation shows it. The notations $\clB{S}{c}$, $\vfix{S}{w}{d}$, $\fixset{S}{\Sigma}{d}$ show the
data $S$, which Lean takes implicitly. For theorems the same happens in the statements of
Propositions~\ref{p:eichler}, \ref{p:discindex}, \ref{p:normid}, \ref{p:nearsize}, of Lemmas~\ref{l:coords}(b),
\ref{l:cutoff}, \ref{l:extend}, \ref{l:near}(b), \ref{l:sizes} and of Corollary~\ref{c:stablelevel}; in each case
the statement here has Lean's arguments, as far as the text shows them, and the account of the census's extra
arguments points to the words of the statement that name the bound variable. In Proposition~\ref{p:bigcell}(a) the
census misses the sequence $p$, which the statement uses. For Proposition~\ref{p:nearsize} the census reads the
proof of the first bound as part of its statement, and lists as an argument a hypothesis ($\gamma$ is a term) that
the statement does not have; the reader gave the argument in a note, and the proof here follows the Lean text.

\subhead{Where the formal text differs from the statements above.}
(1) Theorem~C(b) is, in the formal text, a sum over all subspaces of a term that is $0$ unless the dimension
condition holds, in $[0,\infty]$. (2) The conclusion of Proposition~\ref{p:lsc} is stated there for the function
$x\mapsto\max(\operatorname{Re} F(x),0)$ with values in $[0,\infty]$; it equals $F$ by the first assertion. Its hypothesis on the
common entry of index $i$ is used only for $i=n$. (3) Lemma~\ref{l:cutoff}(b) gives the cutoff in both forms, as a
sum over $\mathsf{K}$ and as an integral for the counting measure of $\Gamma$. (4) The cutoff $\cutoff{S}{c}$ is defined by two
choices; the second is a choice among the functions $\omega$ with $\IsBP{S}{c}{\omega}$ and is arbitrary when there is
none. Nothing in this note depends on it; Note~B proves that such functions exist under the standing hypotheses when $E$ is stably unramified.
(5) The record of filling data lists some consequences of its own conditions: $1\le b$ beside $1000\le b$, and
$1000(2\varrho+2)\le k$ beside $1000(2\varrho+2)^9\le k$. (6) The comment of the formal text on $\siegC{S}$ gives the
second formula of the remark after Definition~\ref{d:const}; the definition is the first. (7) Lemma~\ref{l:sizes}(c)
is stated in the formal text as a rule: a proposition holds if it follows from the listed facts. (8) In
Lemma~\ref{l:redab}(a) and in the facts that $\Trat{S}{E}$ and $\Tlat{S}{F}$ are hermitian and invertible, that
$\tfrac12\in\AEf$ has trace $1$, that $\Cl{S}$ is finite and $\GS{S}{E}$ countable, the formal text has instances
without names; they are not in the concordance. (9) Theorem~B(b) is the precise form of ``the link is Lagrangian'';
the dimension of the link is computed in Note~B. (10) The Inputs (U1), (La2), (La4) and (M1) are stated here in the
special cases that are used: (U1) combines an induction principle over the floors with the description of the
elements of one floor; (La4) is the case of the formal lemma in which its additional block is absent.

\subhead{What this note takes for granted.} From the background library: the groups of
Setting~\ref{set:shape} and the Inputs (G1)--(G6), (E1), (L1), (L2), (U1)--(U4) (units~1, 29, 112, 119, 120);
adeles, the standard characters and cutoffs (Ad1)--(Ad4) (units~114, 122, 123); lattices and indices
(La1)--(La4) (units~57, 120); Minkowski's second theorem, Gram--Schmidt, the first minimum and a packing bound
(M1)--(M4) (units~114, 120, 121); the positivity of the mass (unit~99); and neatness (N1) (unit~50), which is not taken for granted any more: since the revision after the referee's report a proof follows its statement. From the notes on
the real place (unit~132): (R1)--(R3). The notions of Bessel function and of the measure
$\mu_{\mathcal U\backslash G}$ are background (units~117, 118).

\subhead{How this note was written.} From the dossier of the unit: the texts of the members of the nine files
treated here and their structured proofs. I read the nine Lean files completely, and compared every statement of
this note with the Lean statement, argument by argument; for the statements that contain a local definition I read
the full Lean text and not the census's span. For Proposition~\ref{p:nearsize}, which has no structured proof, I
read the proof in the Lean text. The proofs follow the formal ones in their structure. Written out here where the
formal text calls a decision procedure or a simplifier: the verification that the Eichler map of
Proposition~\ref{p:eichler} is unitriangular; the inverse of the lattice form in Lemma~\ref{l:latgram}; the
bookkeeping of Lemma~\ref{l:mbc}; the exponents of Proposition~\ref{p:nearsize} and of Proposition~\ref{p:free}. In
the proof of Theorem~C(a) I use that the index in Jacobi's identity is $1$, so that the discriminants of a subspace
and of its orthogonal complement agree; the formal text uses only that it is at least $1$. Identities between
explicit matrices that the formal text checks entry by entry (the image of a rational Heisenberg element, the
matrix of the lattice basis and its inverse) are stated and not recomputed. Of the Inputs I read the Lean
statements of (La2), (La4), (M1), (M4), (N1), (L1), (R1)--(R3), (G3) and (U3) and the dossier's readings of the
others; I did not open their proofs. The references are those of the second reader's verdicts and of the catalogue
of the write-up; I did not open the cited works. The program check of the unit's dossier covered \UnitSteps{} leaf
steps of \UnitProofs{} proofs.

\head{Appendix A. Concordance with the Lean artifact}
The formal names of the objects are these. The fields \texttt{S.N}, \texttt{S.levelPrime i}, \texttt{S.piBar},
\texttt{S.ϖ}, \texttt{S.b}, \texttt{S.chosenPlaces i}, \texttt{S.columnLevel}, \texttt{S.besselSys v} are $\NS{S}$, $\qS{S}{i}$, $\piS{S}$, $\vpS{S}$,
$\bS{S}$, $\ChS{S}{i}$, $\KcolS{S}$, $\BsysS{S}{\tau}$; the Lean variable \texttt{v} for an infinite place is $\tau$ and
\texttt{Φ v x} is $\at{x}{\tau}$; \texttt{besselU}, \texttt{besselU0}, \texttt{besselH}, \texttt{besselS}, \texttt{GT l} are
$\UT{T}$, $\UoT{T}$, $\HT{T}$, $\ST{T}$, $\Hl{l}{T}$; \texttt{frameGram}, \texttt{realFrameGram}, \texttt{leadingMinor},
\texttt{simpleRootSum}, \texttt{frameImage} are $\ZT{T}{\gamma}$, $\Zr{y}$, $\lmin{l}{Z}$, $\sd{u}$, $\gamma\mathfrak f$;
\texttt{genIndex} is $\gix{\,\cdot\,}$; \texttt{toBesselH} is $\iota_N$; a set \texttt{T} of rational elements is
$\Xi$, the subspace \texttt{T} of Siegel's lemma is $V$, and the finite index sets \texttt{τ}, \texttt{σ} are
$\Omega$. Four statements as the artifact has them:
\quoted{siegellemma}
\quoted{memlinkifforthogonal}
\quoted{existsbesselUapprox}
\quoted{eqandmemclassUnipofbesselS}
{\fontsize{8.6}{11.6}\selectfont\renewcommand{\leanfont}{\ttfamily\fontsize{7.7}{10}\selectfont\color{leancol}}
\setlength{\LTleft}{0pt}\setlength{\LTright}{\fill}\setlength{\tabcolsep}{0pt}
\begin{longtable}{@{}p{27mm}p{62mm}p{40mm}r@{}}
\textbf{In this note} & \textbf{Lean declaration} & \textbf{File} & \textbf{Lines}\\\midrule\endhead
\ConcordanceRows
\end{longtable}}

\head{Notation}
\begingroup\fontsize{8.6}{10.9}\selectfont\raggedright
\setlength{\columnsep}{9mm}\setlength{\parskip}{0pt}
\begin{multicols}{2}
\nt{$E$, $\IntE$, $\cj{x}$}{a CM field, its integers, complex conjugation or the involution of a ring}
\nt{$\tau$, $\at{x}{\tau}$}{an infinite place of $E$; the image of $x\in E$ under the embedding of the CM type}
\nt{$w$, $\cj{w}$, $\Res{w}$, $q_w$}{a finite place, its conjugate, the residue field, its number of elements; $w$ also an invertible matrix over $\AEf$ in Proposition~\ref{p:discindex}}
\nt{$\AdE$, $\AEf$, $\hO$}{adeles, finite adeles, integral finite adeles}
\nt{$\ninf{E}$, $\dE{E}$, $\rdE{E}$, $\volO{E}$}{number of infinite places, discriminant, root discriminant, $2^{f}/\sqrt{|d_E|}$}
\nt{$\psi_f$, $\psi_E$, $\mathbf e$}{standard additive characters; $\mathbf e(c)=e^{c-\cj{c}}$}
\nt{$\pi$, $\varepsilon_w$, $\epsilon$}{the number $\pi$; the idempotent of $\hO$ at $w$; a truth value; $\epsilon$ also the factor of Definition~\ref{d:bsys}}
\nt{$\gix{a}$}{generalized index of a unit $a$}
\nt{$\fs=(\varrho,\kappa,k_0)$}{a shape; $\kk{\fs}$, $\nn{\fs}$, $\rr{\fs}$, $\Mx{\fs}$, $\Ix{\fs}$: Setting~\ref{set:shape}}
\nt{$e_a$, $x_i$, $z$, $e_a^*$; $\fr{a}$, $\mathfrak f$}{the basis; the frame}
\nt{$\lev$, $\fl$}{level and floor of a basis vector}
\nt{$R$, $R'$, $F$, $A$, $C$, $K_0$}{rings; $K_0$ a field}
\nt{$T$, $T'$, $\Trg{\fs}{N}$, $\JT{T}$, $h_T$, $h$}{forms, the diagonal form, the split Gram matrix, the hermitian form}
\nt{$\GT{T}$, $\Gone{\fs}$, $\Ginf{E}{\fs}$}{unitary groups; the product over the infinite places}
\nt{$\UT{T}$, $\UoT{T}$, $\HT{T}$, $\Hl{l}{T}$, $\ST{T}$}{unipotent group, its frame-fixing part, the groups $\mathcal H$, the Bessel subgroup}
\nt{$\ZT{T}{\gamma}$, $\Zr{y}$, $Z$, $\lmin{l}{Z}$}{frame Gram matrices; leading principal minors}
\nt{$\sd{s}$, $\weyl{l}{p}$, $\iota_N$, $\nu_l(v,\beta)$}{sum of the superdiagonal frame entries, Weyl element, embedding of $\Uni_\kappa$, Heisenberg element}
\nt{$\adj{M}$, $\tr{M}$, $\fadm{M}$}{conjugate transpose, transpose, the same matrix over $\AEf$}
\nt{$\Tlatg{\fs}{N}$, $\mbc{\fs}{p}$, $\blb{\fs}{p}$}{lattice form, base change, lattice basis: Definition~\ref{d:latform}}
\nt{$\fdisc{N}{Z}$, $\gdr{Z}$}{frame discriminants: Definitions~\ref{d:fdisc}, \ref{d:bconst}}
\nt{$\fpair{\gamma}$, $\ifpair{\gamma}$, $\lmat{\gamma}$, $\link{\gamma}$}{frame pairs, link matrix, link: Definition~\ref{d:link}}
\nt{$X$, $t_{\mathrm{mid}}$}{the span of the $e_a$; middle part of a vector}
\nt{$\damp{k}{r}$, $\hw{K}{i}{p}$, $\fhw{K}{i}{p}$}{damped power, head weights: Definition~\ref{d:head}}
\nt{$\Sieve{u}{\iota}{n}{J}{\psi}{\mathit{En}}{\Lambda}{\beta}$}{sieve system: Definition~\ref{d:sieve}; $U$, $\xi$, $\Omega$, $u_{\CC}$}
\nt{$\occ{\fs}$, $\bandc{\fs}$, $\bsize{y}$, $\InCut{y}{b}$, $\frbox{\rho}$, $\cellm{y}{\deg}{\rho}{d}$}{Definition~\ref{d:bconst}}
\nt{$\bcorr{\mu}{W}{\gamma}$, $\SchurLe{W}{D}$, $\HasSB{W}{\mathrm{rep}}{B}{D}$, $\ncorr{y}{W}{p}$, $\BandL{b}{W}{\deg}{J}$}{Definition~\ref{d:schur}; $\varphi$, $V_w$, $X$, $B$, $D$}
\nt{$\sigma$, $\full$, $\cut$; $\CutPair{b}{f}{\mathrm{rep}}{W}$, $\BSys{\fs}{b}{f}$}{sides; cut pairs, Bessel systems: Definition~\ref{d:bsys}; $\mathsf m$, $\epsilon$, $P$, $\mathfrak B$, $\Phi$, $g$}
\nt{$\mu_{\mathcal U\backslash G}$, $\vol$}{the measure for $\mathcal U\backslash G_1$; volumes}
\nt{$S\in\Fill{E}{\fs}$}{filling data: Definition~\ref{d:filling}; $\NS{S}$, $\qS{S}{i}$, $\piS{S}$, $\vpS{S}$, $\bS{S}$, $\ChS{S}{i}$, $\KcolS{S}$, $\BsysS{S}{\tau}$}
\nt{$\Isys{S}{\tau}$, $\Wsys{\sigma}{S}{\tau}$, $\degsys{S}{\tau}$, $\repsys{S}{\tau}$}{the components of the Bessel system at $\tau$}
\nt{$\Gf{E}{\kappa}$, $\Kf{E}{\kappa}$, $\Gamma_f(E,\kappa)$, $\Uni_\kappa$}{finite adelic unitary group, its integral subgroup, rational points}
\nt{$\gS{S}$, $\Trat{S}{R}$, $\Tlat{S}{R}$, $\GS{S}{R}$, $\GFin{S}$, $\GRes{S}$, $\Bc{S}$, $\Cl{S}$, $\dot c$, $\Zf{S}{\gamma}$, $\tdisc{S}{\gamma}$}{Definition~\ref{d:groups}}
\nt{$\KS{S}$, $\redS{S}$, $\KgS{S}$, $\redim{S}{A}$, $\gid{S}$}{Definition~\ref{d:level}}
\nt{$\kapS{S}{\tau}$, $\bscale{S}{\tau}{i}$, $\bsGL{S}{\tau}$, $\fscale{S}{\tau}{a}$, $\fsprod{S}$}{Definition~\ref{d:scales}}
\nt{$\latB{S}$, $\tolb{S}$, $\radl{S}$, $\ainf{S}$, $\afin{S}$, $\rinf{S}$, $\rfin{S}$, $\incH{S}$, $\hc{S}{c}$, $\sdc{S}$, $\uchar{S}{u}$}{Definition~\ref{d:maps}}
\nt{$\pos{S}{c}{c'}{\gamma}$, $\terms{S}{c}{s}{c'}$, $\clU{S}{c}$, $\clUo{S}{c}$, $\Agr{S}{s}$, $\IsStab{S}{\Xi}$, $\tlink{S}{\Lambda}$}{Definition~\ref{d:terms}}
\nt{$\besfn{\sigma}{S}{y}{i}$, $\besker{\sigma}{S}{\gamma}{y}$, $\bxs{S}$, $\IsBP{S}{c}{\omega}$, $\cutoff{S}{c}$, $\amp{\sigma}{S}{c}{c'}{\gamma}$, $\ampabs{\sigma}{S}{c}{c'}{\gamma}$, $\ampsum{\sigma}{S}{\Xi}{c}{s}{c'}$}{Definition~\ref{d:amp}; $\alpha$, $\beta^0$}
\nt{$\cdim{S}$, $\bdeg{S}$, $\mainE{S}$, $\uord{S}$, $\tailS{S}$, $\siegC{S}$, $\packC{S}$, $\besE{S}$, $\codimF{S}$, $\ucov{S}$, $\mass{E}{r}$}{Definition~\ref{d:const}}
\nt{$\IsAdm{S}$}{the standing hypotheses (A1)--(A8): Definition~\ref{d:adm}}
\nt{$\clB{S}{c}$, $\clC{S}{c}$, $\clCX{S}{c}$}{Definition~\ref{d:classbasis}}
\nt{$\ioia{S}{\epsilon}$, $\rag{S}{x}$, $\vfix{S}{w}{d}$, $\FixN{s}{x}$, $\fixset{S}{\Sigma}{d}$}{Definition~\ref{d:fixers}}
\nt{$\zfp{T}$, $\clink{\gamma}{\gamma'}$, $\HCFV{\gamma}{c}{c'}{a}$, $\IsNF{j}{\gamma}$}{Definition~\ref{d:nf}}
\nt{$\ratG{S}$, $\lform{S}$, $\cw{S}{i}$, $\HasLev{S}{j}{\gamma}$, $\Cong{S}{c}{c'}{\gamma}{\gamma'}$, $\tdnorm{S}{\gamma}$, $\bxp{S}$, $\IsNr{S}{\gamma}$, $\ChA{S}{\gamma}$}{Definition~\ref{d:levels}}
\nt{$\slb{H}{V}$, $\sld{H}{V}$, $H$, $L$, $V$, $Y$, $K$, $\mathsf n$, $\mathrm m$}{Definition~\ref{d:sublattice}, Proposition~\ref{p:discindex}}
\nt{$\nrm{X}$, $\abs{x}$}{Euclidean norms over all infinite places; absolute values}
\nt{$c,c',c''$}{classes; also ring elements in Definition~\ref{d:nf}}
\nt{$\gamma,\gamma',\delta,\gamma_0$; $u,u',u_0$; $t$}{rational elements; unipotent elements; elements of $\GRes{S}$ or of groups, natural numbers; $\gamma_0$ also a rational unitary matrix in the proof of (N1); $t$ also the trace of $\gamma_0$ minus $\#\Omega$ in the proof of (N1)}
\nt{$y,y',x'$}{points of $\Ginf{E}{\fs}$, vectors, group elements}
\nt{$a,a',a'',j,l,i,i'$; $m,n,k$; $e$}{indices; sizes ($n=\nn{\fs}$, $k=\kk{\fs}$, $b=\bS{S}$, $g=\gS{S}$ in proofs); an exponent or a map of indices}
\nt{$p,q$; $\xi$; $\rho$; $\zeta,\zeta_0$; $\eta$; $\phi$; $\mathsf z$}{sequences of units; a frequency; a radius; points of $\GS{S}{\AEf}$; an element of $\Uni_\kappa(E)$; a map or a function; an adele}
\nt{$\sigma_0,\sigma_1$; $o$; $v$; $z_j$, $B_j$, $r_\tau$; $x^0$, $y^0$}{elements of $\Hl{\varrho}{T}$; an index; a vector or an index; the data of Minkowski's theorem; vectors of Input (L1)}
\nt{$\lambda_i$, $w_j$, $N_{E/\Q}$}{in the proof of (N1): its eigenvalues at an infinite place, the places of the level, the norm}
\nt{$M$, $J$, $O_A$, $Q$, $G$, $I$, $U$, $W$, $\Pi$, $\Gamma$, $\Lambda$, $\Sigma$, $\Xi$, $\mathcal C$}{matrices; a ring of integers; index sets; a subring; groups; a finite set; a space; Bessel functions; a property; discrete groups; a submodule or head bounds; a set of places; a set of rational elements; a closure}
\nt{$D$, $N$, $O$, $P$, $d$, $f$, $r$, $s$, $x$, $\mu$, $\nu$, $\varsigma$}{in the Conventions of this note also local objects, not those of
Conventions~0.1 to~0.3, 0.6 and~0.7: $D$ a number (Definition~\ref{d:schur}, Proposition~\ref{p:discindex}) and $N$ an element of a ring
(Definition~\ref{d:fdisc}), $O$ and $P$ matrices or index sets, $d$ a class, a ring element, a product of units or a real number, $f$ and $r$
sizes ($f=\ninf{E}$, $r=\rr{\fs}$ in proofs), $s$ an element of $\GRes{S}$ or of a group, $x$ a point, a vector or a group element, $\mu$ a measure,
$\nu$ the Heisenberg element $\nu_l(v,\beta)$, $\varsigma$ the scalar of an Eichler map}
\end{multicols}
\endgroup

\head{References}
\begingroup\fontsize{9.4}{12.6}\selectfont
\begin{thebibliography}{9}
\bibitem{bv} E. Bombieri, J. Vaaler, On Siegel's lemma, Invent. Math. 73 (1983). Theorem~3 (adelic successive
minima), as recorded in the second reader's comparison; the formal statement (M1) is a weaker form. The artifact
names the same paper for Theorem~C(a). Not opened.
\bibitem{pr} V. Platonov, A. Rapinchuk, \emph{Algebraic groups and number theory}, 1994, \S7.1. As cited by the
artifact for strong approximation in unipotent groups (Theorem~A); place not confirmed, not opened.
\bibitem{serre} J.-P. Serre, Rigidit\'e du foncteur de Jacobi d'\'echelon $n\ge3$, S\'eminaire Henri Cartan 13
(1960/61), appendix. As recorded in the catalogue of the write-up for (N1); place not confirmed, relation
unresolved by the second reader. Not used: (N1) is proved in Section~\ref{s:approx}.
\end{thebibliography}
\endgroup

\end{document}
